\documentclass{article}
\usepackage[english]{babel}
\usepackage{geometry,amsmath,amssymb,stmaryrd,hyperref,latexsym,alltt,theorem}
\geometry{letterpaper}

%%%%%%%%%% Start TeXmacs macros
\catcode`\<=\active \def<{
\fontencoding{T1}\selectfont\symbol{60}\fontencoding{\encodingdefault}}
\catcode`\>=\active \def>{
\fontencoding{T1}\selectfont\symbol{62}\fontencoding{\encodingdefault}}
\newcommand{\assign}{:=}
\newcommand{\cdummy}{\cdot}
\newcommand{\colons}{\,:\,}
\newcommand{\nin}{\not\in}
\newcommand{\nobracket}{}
\newcommand{\nosymbol}{}
\newcommand{\tmSep}{; }
\newcommand{\tmaffiliation}[1]{\thanks{\textit{Affiliation:} #1}}
\newcommand{\tmem}[1]{{\em #1\/}}
\newcommand{\tmmathbf}[1]{\ensuremath{\boldsymbol{#1}}}
\newcommand{\tmmisc}[1]{\thanks{\textit{Misc:} #1}}
\newcommand{\tmname}[1]{\textsc{#1}}
\newcommand{\tmop}[1]{\ensuremath{\operatorname{#1}}}
\newcommand{\tmrsub}[1]{\ensuremath{_{\textrm{#1}}}}
\newcommand{\tmstrong}[1]{\textbf{#1}}
\newcommand{\tmtextit}[1]{{\itshape{#1}}}
\newcommand{\tmverbatim}[1]{{\ttfamily{#1}}}
\newenvironment{proof}{\noindent\textbf{Proof\ }}{\hspace*{\fill}$\Box$\medskip}
\newenvironment{tmcode}[1][]{\begin{alltt}}{\end{alltt}}
\newtheorem{axiom}{Axiom}
\newtheorem{corollary}{Corollary}
\newtheorem{definition}{Definition}
{\theorembodyfont{\rmfamily}\newtheorem{example}{Example}}
\newtheorem{lemma}{Lemma}
\newtheorem{proposition}{Proposition}
\newtheorem{theorem}{Theorem}
%%%%%%%%%% End TeXmacs macros

%

\newcommand{\bytext}{ }
\newcommand{\renderabstract}[1]{{\noindent}\chapter*{{\abstracttext}}

#1{\hflush}}

\begin{document}

\title{
  GADTs for Reconstruction of Invariants and Postconditions
  \tmmisc{{\tmstrong{PhD Thesis}}{\tmSep}
  {\tmstrong{Doctoral advisor: prof. dr hab. Leszek Pacholski}}}
}

\author{
  {\L}ukasz Stafiniak
  \tmaffiliation{University of Wroc{\l}aw{\tmSep}
  Department of Mathematics and Computer Science{\tmSep}
  Institute of Computer Science}
}

\date{Wroc{\l}aw 2015}

\maketitle

\title{
  Typy GADT dla rekonstrukcji niezmiennik{\'o}w i warunk{\'o}w
  ko{\'n}{\nobreak}cowych
  \tmmisc{{\tmstrong{Rozprawa doktorska napisana pod kierunkiem}}{\tmSep}
  {\tmstrong{prof. dra hab. Leszka Pacholskiego}}}
}

\author{
  {\L}ukasz Stafiniak
  \tmaffiliation{Uniwersytet Wroc{\l}awski{\tmSep}
  Wydzia{\l} Matematyki i Informatyki{\tmSep}
  Instytut Informatyki}
}

\date{Wroc{\l}aw 2015}

\maketitle

\begin{abstract}
  Type systems for programming languages are both a first line of defense
  against programmer mistakes, and an aid in structuring programs around data
  structures and functions that manipulate them. Type systems are a natural
  formalism for language extensions that express those aspects of program
  specifications which can be automatically, statically checked. Type
  inference enhances programmer efficiency and code adaptability by providing
  the types of expressions to the programmer, rather than asking the
  programmer for the types. There is a long line of research on automated
  program analysis and verification by generating specifications that a
  program does obey, for example generating loop invariants. This thesis
  extends a familiar type system for functional programming languages and
  provides a type inference algorithm that generates invariants and
  postconditions of recursive functions.
  
  Generalized Algebraic Data Types (GADTs) extend polymorphic type systems by
  introducing reasoning by cases into type-checking of definitions by cases.
  We present a GADTs type system $\tmop{MMG} (X)$ based on Fran{\c c}ois
  Pottier and Vincent Simonet's $\tmop{HMG} (X)$ but without type annotations.
  We extend it to a language with existential types represented as implicitly
  defined and used GADTs. We show that the type inference problem reduces to
  satisfaction of second order constraints.
  
  We solve the constraints by iterating: (1) constraint generalization, which
  finds most specific common consequences; and (2) joint constraint abduction,
  which finds most general common explanations. Abduction is used to generate
  invariants, infer and check types. Generalization is used to generate
  existential types, which serve as postconditions. The constraints include
  linear arithmetic (equations and inequalities).
  
  We called the system implementing these techniques {\tmname{InvarGenT}}. It
  solves a vast majority of inference tasks we attempted, without type
  annotations. {\tmname{InvarGenT}} solves nearly all meaningful GADTs
  inference tasks we tried, clearly more than all earlier type inference
  implementations for GADTs. It also solves clearly more invariant and
  postcondition inference tasks than the {\tmem{Liquid Types}} approach,
  except for some programs with higher-order functions. In addition, we
  present a selection of new inference tasks, for programs manipulating lists
  with length, binary numbers and AVL trees of imbalance 2. These programs can
  serve as baseline tests for future research on invariant and postcondition
  inference.
\end{abstract}

\chapter*{Streszczenie}

Systemy typ{\'o}w dla j{\k e}zyk{\'o}w programowania s{\k a} zar{\'o}wno
pierwsz{\k a} lini{\k a} obrony przed b{\l}{\k e}dami programistycznymi, jak i
pomoc{\k a} przy strukturowaniu program{\'o}w wok{\'o}{\l} struktur danych
oraz funkcji kt{\'o}re nimi manipuluj{\k a}. Systemy typ{\'o}w to dobry
formalizm do wyra{\.z}ania tych aspekt{\'o}w specyfikacji program{\'o}w,
kt{\'o}re mog{\k a} by{\'c} automatycznie, statycznie sprawdzone. Inferencja
typ{\'o}w zwi{\k e}ksza wydajno{\'s}{\'c} programisty i adaptowalno{\'s}{\'c}
kodu dostarczaj{\k a}c programi{\'s}cie typy wyra{\.z}e{\'n}, zamiast
wymaga{\'c} podawania typ{\'o}w. Badania nad automatyczn{\k a} analiz{\k a} i
weryfikacj{\k a} program{\'o}w od dawna obejmowa{\l}y mi{\k e}dzy innymi
generowanie specyfikacji spe{\l}nianych przez dane programy, na przyk{\l}ad
generowanie niezmiennik{\'o}w p{\k e}tli. Ta praca rozszerza znany system
typ{\'o}w dla j{\k e}zyk{\'o}w funkcyjnych i podaje algorytm inferencji
typ{\'o}w, generuj{\k a}cy niezmienniki i warunki ko{\'n}cowe funkcji
rekurencyjnych.

Generalized Algebraic Data Types (GADTs) rozszerzaj{\k a} systemy typ{\'o}w
polimorficznych o wnioskowanie przez przypadki podczas sprawdzania typu
definicji przez przypadki. Prezentuj{\k e} system typ{\'o}w $\tmop{MMG} (X)$ z
GADTs, oparty o system typ{\'o}w $\tmop{HMG} (X)$ Fran{\c c}ois Pottiera i
Vincenta Simoneta ale bez annotacji typami. Rozszerzam go do j{\k e}zyka z
typami egzystencjalnymi reprezentowanymi jako domy{\'s}lnie definiowane i
u{\.z}ywane struktury GADTs. Pokazuj{\k e} redukcj{\k e} problemu inferencji
typ{\'o}w do spe{\l}nialno{\'s}ci wi{\k e}z{\'o}w drugiego rz{\k
e}du.{\hspace*{\fill}}

Wi{\k e}zy drugiego rz{\k e}du rozwi{\k a}zuj{\k e} iteruj{\k a}c dwa
algorytmy: (1) generalizacj{\k e} wi{\k e}z{\'o}w, znajduj{\k a}c{\k a}
najbardziej specyficzn{\k a} wsp{\'o}ln{\k a} konsekwencj{\k e} wi{\k
e}z{\'o}w; (2) {\l}{\k a}czn{\k a} abdukcj{\k e} wi{\k e}z{\'o}w,
zna{\nobreak}jduj{\k a}ca najog{\'o}lniejsze wsp{\'o}lne obja{\'s}nienie,
przes{\l}ank{\k e} implikuj{\k a}c{\k a} wi{\k e}zy. Abdukcji u{\.z}ywamy
g{\l}{\'o}wnie do generowania niezmiennik{\'o}w, inferencji i sprawdzania
typ{\'o}w. Generalizacji u{\.z}ywamy do generowania typ{\'o}w
egzystencjalnych, s{\l}u{\.z}{\k a}cych jako warunki ko{\'n}cowe. Wi{\k e}zy
obejmuj{\k a} arytmetyk{\k e} liniow{\k a} (r{\'o}wnania i
nier{\'o}wno{\'s}ci).

System implementuj{\k a}cy te techniki nazwa{\l}em {\tmname{InvarGenT}}.
Rozwi{\k a}zuje on zdecydowan{\k a} wi{\k e}kszo{\'s}{\'c} zada{\'n}
inferencji kt{\'o}re opracowa{\l}em lub zaadaptowa{\l}em, bez annotacji
typami. {\tmname{InvarGenT}} rozwi{\k a}zuje zdecydowanie wi{\k e}cej
zada{\'n} inferencji typ{\'o}w dla GADTs ni{\.z} dotychczasowe systemy.
Rozwi{\k a}zuje te{\.z} wi{\k e}cej problem{\'o}w inferencji niezmiennik{\'o}w
i warunk{\'o}w ko{\'n}cowych ni{\.z} podej{\'s}cie {\tmem{Liquid Types}},
je{\'s}li ograniczymy si{\k e} do zada{\'n} nie potrzebuj{\k a}cych inferencji
niezmiennik{\'o}w dla argument{\'o}w funkcji wy{\.z}szego rz{\k e}du.
Dodatkowo, prezentuj{\k e} kilka program{\'o}w operuj{\k a}cych na listach z
d{\l}ugo{\'s}ci{\k a}, liczbach binarnych i drzewach AVL o niezbalansowaniu
nie przekraczaj{\k a}cym 2. Te programy mog{\k a} s{\l}u{\.z}y{\'c} jako cz{\k
e}{\'s}{\'c} test{\'o}w dla przysz{\l}ych prac nad wszechstronnymi systemami
inferencji niezmiennik{\'o}w i warunk{\'o}w ko{\'n}cowych.

\chapter*{Acknowledgements}

I thank my advisor Leszek Pacholski for supporting the great opportunity to
pursue this work despite my excursions to other topics and projects. I thank
{\L}ukasz Kaiser for feedback and support through the years. My office
colleague Jan Otop offered discussions early in my work. Dariusz Biernacki
reviewed and helped improve my conference paper. Filip Sieczkowski, Jakub
Michaliszyn and Marek Materzok remarks helped improve the thesis presentation.
All the technical errors are my sole responsibility. Most of all, I dedicate
this thesis to my parents Maria and Piotr, without whose support this work
would not have been possible.

\

\

\

\

\

\begin{flushright}
  {\tmem{Every problem that is interesting is also soluble.}}
  
  David Deutsch's Law
\end{flushright}

{\tableofcontents}

\chapter{Introduction}\label{Intro}

Type systems are established natural deduction-style means to specify
programs. Dependent types can represent arbitrarily complex properties as they
use the same language for both types and programs. The type of the value
returned by a function can itself be a function of the argument. Generalized
Algebraic Data Types (GADTs) bring some of that expressiveness to type systems
with data-types and parametric polymorphism, by introducing the ability to
reason about the return type by case analysis on the input value. Work over
the past decade (Pottier and R{\'e}gis-Gianas {\cite{StratifiedGADTs}},
Schrijvers, Peyton Jones, Sulzmann  and  Vytiniotis
{\cite{ComplDecidableGADTs}}, Lin  and  Sheard {\cite{PointwiseGADTs}}) shows
that type systems with GADTs can remain tractable from a compiler
implementer's perspective if appropriately constrained, i.e. have principal
typings and efficient type inference. Our work is based on Simonet and Pottier
{\cite{simonet-pottier-hmg-toplas}} instead, which is the most general
presentation of GADTs. Our methods are computationally intensive and involve
search with backtracking. We expect to be able to infer types for more
programs than all of the efficiency-oriented approaches. Our work could bear
the title {\tmem{Type Inference for GADTs and Existentials}}, but we stress
that we do not compete with the work of {\cite{StratifiedGADTs}},
{\cite{ComplDecidableGADTs}} in particular. We are concerned with type
inference for recursive definitions which are not given their type beforehand.
This polymorphic recursion problem has been tackled by {\cite{PointwiseGADTs}}
with GADTs, by Schrijvers and Bruynooghe {\cite{ADTReconstruct}} without
GADTs. The line of work following Unno and Kobayashi
{\cite{InterpolationInfer}} also tackles reconstruction of types for recursive
definitions, but without ADTs and without polymorphic recursion. Our intended
usage is that the generated types of recursive definitions be integrated into
the source code, via integration with an IDE.

Existential types hide some information conveyed in a type. We may opt to use
them to expose a more abstract interface. However, sometimes we are forced to
use existential types to hide what cannot be expressed in the type system.
GADTs provide existential types by using local type variables for the hidden
parts of the type encapsulated in a GADT. With no other way to express
existential types, programmers need to introduce by hand spurious data-types
for them. For example, if we want to hide the length of a list in OCaml, we
need to define \tmverbatim{type \_ elist = List : ('a, 'b) llist -> 'a elist},
where \tmverbatim{('a, 'b) llist} is the type of lists of length
\tmverbatim{'b} with elements of type \tmverbatim{'a}. The type system in Xi 
and  Pfenning {\cite{DML99}} for a language called Dependent ML, a precursor
for GADTs, has a more light-weight approach: explicit existential types. But
{\cite{DML99}} has limited type inference, and we find its type system harder
to grasp than the more recent presentations of GADTs. Knowles  and  Flanagan
{\cite{KnowlesF07}} and Unno  and  Kobayashi {\cite{InterpolationInfer}} can
be seen as performing type inference for forms of existential types. We
provide full reconstruction of existential types.

The feedback provided by type inference makes strongly typed programming more
convenient in several ways. It softens the learning curve by enabling learning
of types by experimentation. It facilitates rapid prototyping and code
refactoring by reporting the types of functions when the programmer
experiments with invariants of data-types. Besides providing certain
correctness guarantees, types also serve as documentation -- it is good to
keep them around. If type inference results are incorporated in the source
code, type checking might suffice during slow evolution of a code base, while
the rich types document the code.

Our type system for GADTs differs from others in that we do not require any
type annotations on expressions, even on recursive functions. Inference in our
implementation sometimes requires guidance by \tmverbatim{assert} clauses. Our
implementation {\tmname{InvarGenT}} includes linear arithmetic in the language
used to express invariants and postconditions, with the possibility to
introduce more domains in the future. The arithmetic properties are
conjunctions of equations and inequalities, where a side can be either a
linear combination, or a maximum or minimum of (at most two) linear
combinations.

\begin{example}
  We can easily specify the data-type of AVL trees with height imbalance of at
  most 2:
  \begin{tmcode}
  datatype Avl : type * num
datacons Empty : a. Avl (a, 0)
datacons Node : a,k,m,n
  [k=max(m,n)  0m  0n  nm+2  mn+2].
     Avl (a, m) * a * Avl (a, n) * Num (k+1)
      Avl (a, k+1)
  \end{tmcode}
  A similar definition can be given in the DML and ATS languages by Hongwei Xi
  (see {\cite{DML99}}). Given this type definition and AVL tree algorithms,
  {\tmname{InvarGenT}} automatically finds out that the height of a tree with
  added element can be the same as original tree or bigger by 1, and that
  removing an element can decrease the height of a tree by at most 1,
  returning, among others, the types:
  \begin{tmcode}
  add :
  a,n.aAvl(a, n)k[kn+1  1k  nk].Avl(a, k)
remove :
  a,n.aAvl(a, n)k[nk+1  0k  kn].Avl(a, k)
  \end{tmcode}
  There are no places requiring type annotations or informative assertions in
  the source code of AVL tree algorithms including the above functions and
  their helper functions.
\end{example}

We reduce the type inference task to {\tmem{second order}} constraint
satisfaction, which in addition to finding substitutions for first order
variables, finds solutions to the predicate unknowns. As the predicate-finding
techniques that we present in Chapter \ref{ChapSolv} (especially abduction)
improve, type inference will work for more programs. The downside is that we
do not specify declaratively the limitations of type inference. That is, while
we provide some guarantees that our approach to inference does not overlook
solutions (completeness-like result at the end of Section
\ref{PredVarSolvDetails}), there is no concise formulation of when type
inference succeeds.

Constraints appear in three contexts in programs and inferred signatures:
\begin{itemize}
  \item In specifications of value constructors (symbolically $K \colons
  \forall \bar{\beta} [D] . \tau_1 \times \ldots \times \tau_n \longrightarrow
  \varepsilon (\bar{\tau})$). We call the constraint (i.e. $D$) the invariant
  (of $K$).
  
  \item In type schemes, which usually are signatures of values (symbolically
  $x : \forall \bar{a} [D] . \tau$). We call the constraint (i.e. $D$) the
  invariant or the precondition (of $x$), interchangeably.
  
  \item In existential types (symbolically $\exists \bar{\alpha} [D] . \tau$);
  existential types usually occur in result positions of function types in
  type schemes. We call the constraint (i.e. $D$) the postcondition (of the
  corresponding function or computation).
\end{itemize}

\section{Contributions}

In order to implement {\tmname{InvarGenT}}, we needed to resolve several
issues, leading to our contributions:
\begin{itemize}
  \item Design a type system that captures invariants (Section \ref{TypSys})
  and postconditions (Section \ref{ExTypes}).
  
  \item Reduce type inference to satisfaction of second order constraints
  (Sections \ref{TypInfCns} and \ref{ExTypesCns}).
  
  \item Employ invariant-finding abduction and postcondition-finding
  generalization in an algorithm that reconstructs correct types (Section
  \ref{MainAlg}).
  
  \item Develop constraint abduction algorithms that handle universally
  quantified variables (Section \ref{AbdAlg}).
  
  \item Develop constraint generalization algorithms. Constraint
  generalization computes anti-unification in case of free terms and extended
  convex hull in case of linear inequalities (Section \ref{DisjElim}).
\end{itemize}
To summarize, the novelty of our work lies in:
\begin{enumerate}
  \item performing type inference for GADTs with polymorphic recursion for
  more programs than in prior work ({\cite{PointwiseGADTs}}),
  
  \item introducing existentials into the GADT type system with minimal added
  complexity, by using GADT encodings, and performing type inference for them,
  
  \item implementing the type inference over a constraint domain with linear
  arithmetics, thus performing numerical precondition and postcondition
  generation for recursive functions (in a different manner than in Rondon,
  Kawaguchi and Jhala {\cite{DSolvePaper}}).
\end{enumerate}
We discuss related work in Chapter \ref{ChapBackground} and Section
\ref{SecRelWork}. There remains work to be done, as we discuss in Section
\ref{FailureCases} and Section \ref{SecFutWork}.

{\tmname{InvarGenT}} can be found at
\tmverbatim{\href{https://github.com/lukstafi/invargent}{https://github.com/lukstafi/invargent}}.
It solves a vast majority of inference tasks we attempted, without type
annotations. {\tmname{InvarGenT}} solves all but 3 tasks from Lin
{\cite{PracticalGADTsInfer}} (2 unsolved are practical, one of them is solved
after a slight meaning-preserving modification); the system from
{\cite{PracticalGADTsInfer}} does not solve 8 of those tasks (7 unsolved are
practical). We translated into {\tmname{InvarGenT}} all but 3 tasks from
Rondon, Kawaguchi and Jhala {\cite{DSolvePaper}} -- all tasks that
{\cite{DSolvePaper}} uses for comparison with Xi and Pfenning's DML system
{\cite{DML99}}, {\cite{DMLArrays}}. {\tmname{InvarGenT}} solves all but 2 of
these inference tasks without any type annotation, and solves the remaining 2
tasks after a slight meaning-preserving change to the programs (still without
any type annotation). The system {\tmname{DSolve}} from {\cite{DSolvePaper}}
implementing the {\tmem{Liquid Types}} approach needs a type annotation for 3
of these inference tasks. The inference times of {\tmname{InvarGenT}} and
{\tmname{DSolve}} are of the same order. Overall, the {\tmname{InvarGenT}}
approach solves clearly more inference tasks than the {\tmem{Liquid Types}}
approach, when the programs do not have higher-order functions that require
universally quantified invariants for arguments, including existentially
quantified invariants for arguments of arguments. In addition, we present a
selection of new inference tasks, for programs manipulating lists with length,
binary numbers and AVL trees of imbalance 2. They can serve as a baseline for
future research. We expect that some of them are beyond the capabilities of
any current type inference system, except {\tmname{InvarGenT}}.
{\tmname{DSolve}} does not introduce new linear combinations into generated
constraints, therefore will not solve some of these inference tasks. The
recent system {\tmname{SpecLearn}}, see He Zhu, Aditya Nori and Suresh
Jagannathan {\cite{SpecLearnArrays}}, cannot solve tasks like the AVL trees
example, without any type annotations or assertions.

\section{Examples}

To motivate the language constructs and type system rules, we start with some
examples. The concrete syntax of {\tmname{InvarGenT}} is similar to that of
OCaml. The sort of a type variable is identified by the first letter of the
variable. \tmverbatim{a}, \tmverbatim{b}, \tmverbatim{c}, \tmverbatim{r},
\tmverbatim{s}, \tmverbatim{t}, \tmverbatim{a1},... are in the sort of
``proper'' types. \tmverbatim{i}, \tmverbatim{j}, \tmverbatim{k},
\tmverbatim{l}, \tmverbatim{m}, \tmverbatim{n}, \tmverbatim{i1},... are in the
sort of linear arithmetic over rational numbers. Type constructors and value
constructors have the same syntax: capitalized name followed by a tuple of
arguments. They are introduced by the keywords \tmverbatim{datatype} and
\tmverbatim{datacons} respectively. The result sort of \tmverbatim{datatype}
is always \tmverbatim{type} and is omitted. By {\tmem{toplevel}} of a source
file we mean the environment of names available for potential code appended to
the end of the file. Values assumed into the toplevel without
{\tmname{InvarGenT}} definition are introduced by the keyword
\tmverbatim{external}.

We can introduce existential types directly in type declarations. To have an
existential type inferred, we have to use \tmverbatim{efunction},
\tmverbatim{ematch} or \tmverbatim{eif} expressions, which differ from
\tmverbatim{function}, \tmverbatim{match} and \tmverbatim{if} correspondingly
only in that the (return) type is an existential type. To use, i.e. unpack, a
value of an existential type, we have to bind it with a
\tmverbatim{let}..\tmverbatim{in} expression. An existential type will be
automatically unpacked before being ``repackaged'' as another existential
type. In a future version, it might be possible to use existential types
without explicit introduction (the \tmverbatim{efunction}, \tmverbatim{ematch}
or \tmverbatim{eif} syntax) and explicit elimination (the need of
\tmverbatim{let}..\tmverbatim{in} expressions) at a cost of longer inference
times.

\begin{example}
  Table \ref{CodeEx2} defines an artificial example using phantom types (i.e.
  types used for type information, without values), resembling binary numbers.
  The function \tmverbatim{erase\_zeros} erases the ``zeros''
  \tmverbatim{BinaryO} and adds one \tmverbatim{BinaryI}.\begin{table}[t]
    \begin{tmcode}
    datatype Z
datatype I : type
datatype O : type
datatype Binary : type
datacons BinaryZ : Binary Z
datacons BinaryO : a. Binary a  Binary (O a)
datacons BinaryI : a. Binary a  Binary (I a)

let rec erase_zeros = efunction
  | BinaryZ -> BinaryI BinaryZ
  | BinaryO x -> erase_zeros x
  | BinaryI x -> let y = erase_zeros x in BinaryI y
    \end{tmcode}
    
    \caption{Use of existential types}
  \end{table}
  
  We get \tmverbatim{erase\_zeros:$\forall$a.Binary a$\rightarrow
  \exists$a.Binary (I a)}. The resulting type forgets the shape of the
  content, other than the fact that it starts with a \tmverbatim{BinaryI}. The
  first branch of \tmverbatim{erase\_zeros} returns a value directly. The
  second branch directly performs the recursive call -- the existential type
  of the result is unpacked and ``repackaged''. The third branch unpacks the
  result of the recusrive call explicitly. By design, the inlined variant
  \tmverbatim{\textbar  BinaryI x -> BinaryI (erase\_zeros x)} would not
  type-check. The type system does not allow passing values of existential
  types as arguments.
\end{example}

{\tmname{InvarGenT}} commits to a type of a toplevel definition before
proceeding to the next one, so sometimes we need to provide more information
in the program. Besides type annotations, there are three means to enrich the
generated constraints: \tmverbatim{assert false} indicates an unreachable code
location and provides a negative constraint, \tmverbatim{assert num
e$_1$<=e$_2$} and \tmverbatim{assert type e$_1$=e$_2$} provide positive
constraints, and the \tmverbatim{test} syntax includes constraints of the use
cases appearing after \tmverbatim{test} with the constraint of a toplevel
definition.

\begin{example}
  Table \ref{CodeEx3} defines a function \tmverbatim{equal} comparing values
  provided representation of their types. The value constructors
  \tmverbatim{TInt}, \tmverbatim{TPair}, \tmverbatim{TList} are used to
  represent the types of values compared. The function \tmverbatim{equal}
  checks whether the second and third argument are of the same type, and if
  so, whether they are equal.\begin{table}[t]
    \begin{tmcode}
    datatype List : type
datacons Nil : a. List a
datacons Cons : a. a * List a  List a
datatype TypeRepr : type
datacons TInt : TypeRepr Int
datacons TPair : a, b. TypeRepr a * TypeRepr b  TypeRepr (a, b)
datacons TList : a. TypeRepr a  TypeRepr (List a)
external let eq_int : Int  Int  Bool = "(=)"
external let b_and : Bool  Bool  Bool = "fun a b -> a && b"
external let b_not : Bool  Bool = "fun b -> not b"
external let forall2 :
  a, b. (abBool)  List a  List b  Bool =
  "fun f a b -> List.for_all2 f a b"
external let zero : Int = "0"

let rec equal = function
  | TInt, TInt -> fun x y -> eq_int x y
  | TPair (t1, t2), TPair (u1, u2) ->  
    (fun (x1, x2) (y1, y2) ->
        b_and (equal (t1, u1) x1 y1)
              (equal (t2, u2) x2 y2))
  | TList t, TList u -> forall2 (equal (t, u))
  | _ -> fun _ _ -> False
  | TInt, TList l ->
    (function Nil -> assert false)
  | TList l, TInt ->
    (function _ -> function Nil -> assert false)
test b_not (equal (TInt, TList TInt) zero Nil)
    \end{tmcode}
    
    \caption{Two ways of constraining types}
  \end{table}
  
  We get \tmverbatim{equal$: \forall$a,b.(TypeRepr a, TypeRepr
  b)$\rightarrow$a$\rightarrow$b$\rightarrow$Bool}. To ensure only one
  maximally general type for \tmverbatim{equal}, it is sufficient to provide
  either the two \tmverbatim{assert false} clauses, or the \tmverbatim{test}
  clause; both are illustrated above. The first assertion excludes
  independence of the first encoded type and the second argument. The second
  assertion excludes independence of the second encoded type and the third
  argument. The test ensures that arguments of distinct types can be~given.
\end{example}

Besides displaying types of toplevel definitions, {\tmname{InvarGenT}} also
exports an OCaml source file with all the required GADT definitions and type
annotations.

\chapter{Background and Related Work}\label{ChapBackground}

In this chapter, we provide background knowledge and context by taking an
in-depth look at a selection of related work. Sections \ref{HMGSection} and
\ref{HerbSection} constitute a proper background for later chapters. The
remaining sections describe alternative approaches to type inferece for GADTs,
and more broadly to the task of automatic generation and verification of
program specifications for functional programming languages like OCaml.
Readers familiar with the works cited below may skip directly to the
subsections {\tmem{Relevance}}, where we contrast the corresponding work with
{\tmname{InvarGenT}}. The final chapter's Section \ref{SecRelWork} provides a
broader perspective on related work.

The theoretical underpinnings of type systems for program specification trace
back to the work on dependent types by Per Martin-L{\"o}f, for example
{\cite{MartinLof}}. An early example of a practical programming language based
on dependent types is Cayenne by Lennart Augustsson, {\cite{Cayenne}}. A
currently popular example of such language is Idris by Edwin Brady
{\cite{IdrisCurrent}}, see also Brady, Herrmann and Hammond
{\cite{IdrisEarly}}. But the family of approaches our thesis belongs to,
maintains the separation of types and values characteristic of languages like
Pascal, C, OCaml and Haskell. This separation allows to employ domain-specific
decision procedures to reason about types, and also to automatically infer
types.

We start with Hongwei Xi and Frank Pfenning {\cite{DML99}}, which not only is
a precursor of GADTs research, but also stresses the importance of existential
types and type inference. We present the $\tmop{HMG} (X)$ system from Simonet
and Pottier {\cite{simonet-pottier-hmg-toplas}}, to motivate and gain
familiarity with the formalism of our constraint-based type system. We then
present Schrijvers, Peyton Jones, Sulzmann and Vytiniotis {\cite{OutsideIn}},
and Lin and Sheard {\cite{PointwiseGADTs}}, to illustrate one thread of
approaches to type inference for type systems with GADTs. Then, we switch
gears to discuss inference of invariants, known as {\tmem{liquid types}} from
Rondon, Kawaguchi and Jhala {\cite{DSolvePaper}}, expanded in
{\cite{DSolveTechRep}} and {\cite{DSolveThesis}}. Finally, we describe Maher
and Huang {\cite{AbductionSolvMaher}}, which provides the basis for our
constraint abduction algorithm for terms.

We take liberties with the presentation of the systems described, except
$\tmop{HMG} (X)$ where we stay close to the letter of
{\cite{simonet-pottier-hmg-toplas}}. We apologize for any unfortunate
resulting errors and misunderstandings, and encourage interested readers to
consult the original publications.

\section{The DML System}\label{DMLSection}

The work of Hongwei Xi and Frank Pfenning {\cite{DML99}}, see also
{\cite{DMLThesis}}, was at the forefront of research introducing type system
features modeled on dependent types into practical programming languages of
the ML family. These features mediate the dependence of types on terms by the
use of strongly constrained types, including singleton types, i.e. types
inhabited by single terms. They were later simplified and further developed to
fit within the polymorphic (non-dependent) type systems. Together with an
independent work at the time by Christoph Zenger {\cite{IndexedTypes}}, and an
earlier work by Konstantin L{\"a}ufer and Martin Odersky {\cite{ExTyADTs}},
these efforts are the origins of Generalized Algebraic Data Types.

Hongwei Xi's $\tmop{DML} (X)$ is parameterized by a constraint domain $X$,
and, like {\tmname{InvarGenT}}, the implementation includes linear
arithmetics. We will discuss a small selection of type system rules and then
the approach to type inference taken by the DML system. The reader should not
expect to gain understanding of DML from this short exposition, in part
because it is a complex system.

\subsection{The Type System}

Type judgments in $\tmop{DML} (X)$ have naturally a different form for
patterns than for expressions. For pattern $p$, type $\tau$, formula
$\varphi$, and environment assigning expression variables to types $\Gamma$,
we write $p \downarrow \tau \vartriangleright (\varphi ; \Gamma)$ to mean that
pattern $p$, when matched against an expression of type $\tau$, allows us to
type-check its corresponding branch in the context of an environment enriched
by $\Gamma$ and under a constraint enriched by $\varphi$. In our notation: $p
\downarrow \tau \longrightarrow \exists \bar{\alpha} [\varphi] \Gamma$. The
fresh type variables $\bar{\alpha}$ are in fact introduced inside $\varphi$ in
$\tmop{DML} (X)$. For expression $e$, type $\tau$, formula $\varphi$, type
environment $\Delta$ assigning sorts to type variables, and environment
$\Gamma$ assigning types to expression variables, we write $\varphi ;
{\nobreak} \Delta ; {\nobreak} \Gamma \vdash e : {\nobreak} \tau$ to mean that
expression $e$ has type $\tau$ in the context of the environments $\Delta,
\Gamma$ and the constraint formula $\varphi$. In our notation: $\varphi ;
\Gamma \vdash e : \tau$, because we discriminate the sorts of type variables
by the names of the variables.

One of the most interesting rules of $\tmop{DML} (X)$ connects these two forms
of judgments:
\[ \frac{p \downarrow \tau_1 \vartriangleright (\varphi^{\prime} ;
   \Gamma^{\prime}) \text{ \ \ \ } \varphi \wedge \varphi^{\prime} ; \Delta ;
   \Gamma \Gamma^{\prime} \vdash e : \tau}{\varphi ; \Delta ; \Gamma \vdash p
   \Rightarrow e : \tau_1 \Rightarrow \tau_2} \]
The notation $\tau_1 \Rightarrow \tau_2$ is used instead of a function type
$\tau_1 \rightarrow \tau_2$ to indicate an ``internal'' use of a pattern
matching branch: it is always a part of a pattern matching syntactic
construct. The premise $\varphi \wedge \varphi^{\prime} ; \Delta ; \Gamma
\Gamma^{\prime} \vdash e : \tau$ means that for type-checking the body of a
pattern matching branch, we have available not only the pattern variables from
$\Gamma^{\prime}$, but also properties $\varphi^{\prime}$ of their types.

Another interesting aspect of $\tmop{DML} (X)$ is the presence of both
universal and existential types. Universal types are as in system F. Rules for
existential types:
\[ \frac{\varphi ; \Delta ; \Gamma \vdash e : \tau [\alpha \assign i] \text{ \
   \ \ } \varphi \vdash i : \gamma}{\varphi ; \Delta ; \Gamma \vdash \langle i
   \middle| e \rangle : \exists \alpha : \gamma . \tau} \text{ \ introduction}
\]
\[ \frac{\varphi ; \Delta ; \Gamma \vdash e_1 : \exists \alpha : \gamma .
   \tau_1 \text{ \ \ \ } \varphi \wedge \alpha : \gamma ; \Gamma \{ x : \tau_1
   \} \vdash e_2 : \tau_2}{\varphi ; \Delta ; \Gamma \vdash
   \text{{\tmstrong{let}}}  \langle \alpha \middle| x \rangle = e_1 
   \text{{\tmstrong{in}}} e_2 : \tau_2} \text{ \ \ \ elimination} \]
where $\gamma$ is a sort in the multisorted logic of $X$. The $\langle i
\middle| e \rangle$ construct pairs an expression $e$ with a witness $i$ for
the existential type $\exists \alpha : \gamma . \tau$. In fact, the concrete
language of DML does not have this construct. Rather, it is introduced by type
inference, as discussed below.

\subsection{Bidirectional Type Inference}

To facilitate type inference for the concrete language of DML, the system is
equipped with two mutually recursive kinds of judgements called
{\tmem{elaboration judgments}}: the {\tmem{synthesizing}} judgment $\varphi ;
\Gamma \vDash e \uparrow \tau \Rightarrow e^{\ast}$, and the {\tmem{checking}}
judgment $\varphi ; \Gamma \vDash e \downarrow \tau \Rightarrow e^{\ast}$. The
synthesizing judgments introduce fresh type variables when needed, which are
solved from the constraints $\varphi$. The types for lambda expressions and
recursive definition expressions are not synthesized. To see the interaction
between synthesizing and checking judgments, consider the rule for function
application:
\[ \frac{\varphi ; \Gamma \vdash e_1 \uparrow \tau_1 \rightarrow \tau_2
   \Rightarrow e_1^{\ast} \text{ \ \ \ } \varphi ; \Gamma \vdash e_2
   \downarrow \tau_1 \Rightarrow e_2^{\ast}}{\varphi ; \Gamma \vdash e_1 e_2
   \uparrow \tau_2 \Rightarrow e_1^{\ast} e_2^{\ast}} \]
Since $\tau_1$ is synthesized as part of the function type by $\varphi ;
\Gamma \vdash e_1 \uparrow \tau_1 \rightarrow \tau_2 \Rightarrow e_1$,
subsequently it only needs to be checked for the argument $e_2$. Synthesis for
variables is done by referring to the environment, and for data constructors
by referring to the respective datatype definition. Checking for variables and
data constructors refers back to synthesis, and we employ constraint solving
to decide whether the types agree:
\[ \frac{\varphi ; \Gamma \vdash x \uparrow \tau_1 \Rightarrow e^{\ast} \text{
   \ \ \ } \varphi \vDash \tau_1 \dot{=} \tau_2}{\varphi ; \Gamma \vdash x
   \downarrow \tau_2 \Rightarrow e^{\ast}} \]
With existential types, the situation is further complicated by the need to
introduce, during elaboration, ``witnesses'' $i$ into expressions of
existential type. DML relies on a judgment $\tmop{coerce}$, whose one of most
interesting rules is:
\[ \frac{\varphi \vdash \tmop{coerce} (\tau_1, \tau [\alpha \assign i])
   \Rightarrow E \text{ \ \ \ } \varphi \vdash i : \gamma}{\varphi \vdash
   \tmop{coerce} (\tau_1, \Sigma (\alpha : \gamma) . \tau) \Rightarrow \langle
   i \middle| E \rangle} \]
where we see existential witness introduced when an existential type is
encountered. The $\tmop{coerce}$ judgment is connected to synthesizing
judgments via checking rules. Let us look at the other rule for function
application:
\[ \frac{\varphi ; \Gamma \vdash e_1 e_2 \uparrow \tau_1 \Rightarrow e^{\ast}
   \text{ \ \ \ } \varphi \vdash \tmop{coerce} (\tau_1, \tau_2) \Rightarrow
   E}{\varphi ; \Gamma \vdash e_1 e_2 \downarrow \tau_2 \Rightarrow E
   [e^{\ast}]} \]
where $E$ is an expression context with a single hole.

\subsection{Relevance}

The work on $\tmop{DML} (X)$ is foundational to {\tmname{InvarGenT}}. While we
chose to build on the more elegant formalism of $\tmop{HMG} (X)$, we share
some of the goals of DML, diverging on one:
\begin{enumerate}
  \item The formalism behind {\tmname{InvarGenT}} is parametrized by the
  domain of constraints, just as $\tmop{DML} (X)$. The implementation handles
  linear arithmetic constraints, and can be extended to other domains,
  similarly to DML.
  
  \item Both DML and {\tmname{InvarGenT}} place emphasis on the seamless use
  of existential types, aided by type inference. However, existential types
  are not synthesized by DML.
  
  \item DML and {\tmname{InvarGenT}} share concern with covering a sufficient
  portion of ML-family language features, in particular pattern matching with
  deep patterns.
  
  \item DML requires type annotations on all functions that cannot be typed in
  the Hindley-Milner type system. In exchange, $\tmop{DML} (X)$ has the
  principal type property and type inference in DML is efficient.
  {\tmname{InvarGenT}} takes the opposite stance, not requiring any type
  annotations.
\end{enumerate}
Unlike DML, {\tmname{InvarGenT}} does not currently support first-class
universal types, i.e. function arguments which can be used polymorphically,
with different types within a function. Introducing inferred universal types
to {\tmname{InvarGenT}} is left as future work. The type inference for
universal types of function arguments would work similarly to how type
inference for recursive definitions (which admits polymorphic recursion) works
currently.

DML performs A-transformation, which is a source level transformation
introducing $\text{{\tmstrong{let}}}$ bindings for subexpressions. In this
way, all occurrences of expressions with existential types are unpacked, with
the existential type eliminated. In interests of type inference times,
{\tmname{InvarGenT}} does not perform A-transformation. However, passing
expressions of an existential type as arguments to functions is prohibited in
{\tmname{InvarGenT}}'s type system, therefore A-transformation would be a
conservative extension. This captures errors where the programmer forgets to
unpack the existential type, without limiting expressivity.

In {\tmname{InvarGenT}}, introductions of existential types need to be marked
in the source by adding a single character to the corresponding concrete
syntax keywords: \tmverbatim{function}, \tmverbatim{match} and
\tmverbatim{if}. In DML, existential type introductions do not figure in the
concrete syntax, but existential types need to be spelled out in type
annotations. In {\tmname{InvarGenT}}, inference of existential types will find
out how many existential parameters are needed, including the case when none
are needed, i.e. the type is not in fact existential. Therefore, a more
sophisticated approach might introduce existential types automatically, at
cost of increased type inference times.

\section{The $\tmop{HMG} (X)$ Formalization}\label{HMGSection}

In parallel to the development of DML, Martin Odersky, Martin Sulzmann and
Martin Wehr in {\cite{HMX}} developed a general framework $\tmop{HM} (X)$ for
expressing type systems extending the Hindley-Milner type system and
parameterized by a constraint domain. Hongwei Xi and collaborators further
developed the ideas behind DML into an extension of type systems for languages
in the ML family -- Standard ML, Haskell, OCaml -- they originally called
{\tmem{Guarded Recursive Data Types}}, see {\cite{XiGRDTs}}, later to become
known as {\tmem{Generalized Algebraic Data Types}} or GADTs. GADTs were
independently investigated at that time by James Cheney and Ralf Hinze as
{\tmem{first-class phantom types}} in {\cite{CheneyHinzePhantomTypes}}, and by
Tim Sheard as {\tmem{equality qualified types}} in {\cite{Sheard04}}. Vincent
Simonet and Fran{\c c}ois Pottier in {\cite{simonet-pottier-hmg-toplas}}
brought these two lines of research, $\tmop{HM} (X)$ and GADTs, together, by
designing an elegant type system $\tmop{HMG} (X)$.

\subsection{The Untyped Calculus}

While the {\tmem{indexed types}} of Christoph Zenger {\cite{IndexedTypes}} and
dependent types of DML refine ML but do not make more ML programs typeable,
GADTs extend ML, i.e. make some programs typeable -- provided appropriate
datatype definitions -- that would not type-check otherwise.

We present the call-by-value $\lambda$-calculus behind $\tmop{HMG} (X)$ since
we will defer the semantics of {\tmname{InvarGenT}} to $\tmop{HMG} (X)$. The
rest of this section follows closely {\cite{simonet-pottier-hmg-toplas}},
starting with pages 12-14.\begin{table}[t]
  $p \assign x \; | \; 0 \; | \; 1 \; | \; p \wedge p \; | \; p \vee p \; | \;
  K \bar{p}
  
  c \assign p.e
  
  e \assign x \; | \; K \bar{e} \; | \; \text{{\tmstrong{let}}} x = e
  \text{{\tmstrong{in}}} e \; | \; e e \; | \; \lambda (\bar{c}) \; | \; \mu
  x.v
  
  v \assign K \bar{v} \; | \; \lambda (\bar{c})$
  
  \label{CalcSynHMG}
  \caption{Syntax of the calculus for $\tmop{HMG} (X)$}
\end{table}

Let $x$ and $K$ range over disjoint denumerable sets of variables and data
constructors, respectively. For every data constructor $K$, we assume a fixed
nonnegative arity. The syntax of patterns, expressions, clauses, and values is
given in Figure \ref{CalcSynHMG}. Patterns include the empty pattern $0$, the
wildcard pattern $1$, variables, conjunction and disjunction patterns, and
data constructor applications. Defined program variables $\tmop{dpv} (p)$ are
the unique variables in $p$. The pattern $p$ is considered ill-formed for
nonlinear patterns (i.e. patterns with repeating variables). Expressions
include variables, functions, data constructor or function applications,
recursive definitions, and local variable definitions. Functions are defined
by cases: a $\lambda$-abstraction, written $\lambda (c_1, \ldots, c_n)$,
consists of a sequence of clauses. A clause $c$ is made up of a pattern $p$
and an expression $e$ and is written $p.e$; the variables in $\tmop{dpv} (p)$
are bound within $e$. We occasionally use $\tmop{ce}$ to stand for a clause or
an expression. Values include functions and applications of a data constructor
to values. Within patterns, expressions, and values, all applications of a
data constructor must respect its arity: data constructors cannot be partially
applied.\begin{table}[t]
  \begin{eqnarray*}
    {}[0 \assign v] & = & \text{undefined}\\
    {}[1 \assign v] & = & \varnothing\\
    {}[p_1 \wedge p_2 \assign v] & = & [p_1 \assign v] \otimes [p_2 \assign
    v]\\
    {}[p_1 \vee p_2 \assign v] & = & [p_1 \assign v] \oplus [p_2 \assign v]\\
    {}[K p_1 \cdots p_n \assign K v_1 \cdots v_n] & = & [p_1 \assign v_1]
    \otimes \cdots \otimes [p_n \assign v_n]
  \end{eqnarray*}
  \label{SubstHMG}
  \caption{Extended substitution for $\tmop{HMG} (X)$}
\end{table}

Whether a pattern $p$ matches a value $v$ is defined by an extended
substitution $[p \assign v]$ that is either undefined, which means that $p$
does not match $v$, or a mapping of $\tmop{dpv} (p)$ to values, which means
that $p$ does match $v$ and describes how its variables become bound. Of
course, when $p$ is a variable $x$, the extended substitution $[x \assign v]$
coincides with the ordinary substitution $[x \assign v]$, which justifies our
abuse of notation. Extended substitution for other pattern forms is defined in
Figure \ref{SubstHMG}. Let us briefly review the definition. The pattern $0$
matches no value, so $[0 \assign v]$ is always undefined. Conversely, the
pattern $1$ matches every value, but binds no variables, so $[1 \assign v]$ is
the empty substitution. In the case of conjunction patterns, $\otimes$ stands
for (disjoint) set-theoretic union, so that the bindings produced by $p_1
\wedge p_2$ are the union of those independently produced by $p_1$ and $p_2$.
The operator $\otimes$ is strict---that is, its result is undefined if either
of its operands is undefined---which means that a conjunction pattern matches
a value if and only if both of its members do. In the case of disjunction
patterns, $\oplus$ stands for a nonstrict, angelic choice operator with left
bias: when $o_1$ and $o_2$ are two possibly undefined mathematical objects
that belong to the same space when defined, $o_1 \oplus o_2$ stands for $o_1$
if it is defined and for $o_2$ otherwise. As a result, a disjunction pattern
matches a value if and only if either of its members does. The set of bindings
thus produced is that produced by $p_1$, if defined, otherwise that produced
by $p_2$. Last, the pattern $K p_1 \cdots p_n$ matches values of the form $K
v_1 \cdots v_n$ only; it matches such a value if and only if $p_i$ matches
$v_i$ for every $i \in \{ 1, \ldots, n \}$.\begin{table}[t]
  \begin{eqnarray*}
    \lambda (p_1 .e_1 \cdots p_n .e_n) v & \rightarrow & \bigoplus_{i = 1}^n
    e_i [p_i \assign v]\\
    \mu x.v & \rightarrow & v [x \assign \mu x.v]\\
    \text{{\tmstrong{let}}} x = v \text{{\tmstrong{in}}} e & \rightarrow & e
    [x \assign v]\\
    E [e] & \rightarrow & E [e^{\prime}] \text{ \ if \ } e \rightarrow
    e^{\prime}\\
    &  & \\
    E & \assign & K \bar{v}  []  \bar{e} \; | \; [] e \; | \; v [] \; | \;
    \text{{\tmstrong{let}}} x = []  \text{{\tmstrong{in}}} e
  \end{eqnarray*}
  \label{SemantHMG}
  \caption{Operational semantics for $\tmop{HMG} (X)$}
\end{table}

The call-by-value small-step semantics, written $\rightarrow$, is defined by
the rules of Figure \ref{SemantHMG}. It is standard. The first rule governs
function application and pattern-matching: $\lambda (p_1 .e_1 \cdots p_n
.e_n)$ reduces to $e_i [p_i \assign v_i]$, where $i$ is the least element of
$\{ 1, \ldots, n \}$ such that $p_i$ matches $v$. Note that this expression is
stuck (does not reduce) when no such $i$ exists. The last rule lifts reduction
to arbitrary evaluation contexts.

\subsection{The $\tmop{HMG} (X)$ Type System}

$\tmop{HMG} (X)$ supports subtyping. The language of constraints is as
follows:
\begin{eqnarray*}
  \tau & \assign & \alpha \; | \; \tau \rightarrow \tau \; | \; \varepsilon
  (\tau, \ldots, \tau)\\
  \pi & \assign & \preccurlyeq \; | \; \ldots\\
  C, D & \assign & \pi \bar{\tau} \; | \; C \wedge C \; | \; C \vee C \; | \;
  \exists \alpha .C \; | \; \forall \alpha .C \; | \; C \Rightarrow C
\end{eqnarray*}
where $\tau$ are types, $\pi$ are constraint relations and $\preccurlyeq$ is a
subtyping relation, $C, D$ are constraint formulas. Environments $\Gamma$
assign variables to type schemes, e.g. $\{ x \mapsto \sigma \}$, where type
schemes are triples:
\[ \sigma \assign \forall \bar{\alpha} [C] . \tau \]
indicating a value of type $\tau$ polymorphic wrt. variables $\bar{\alpha}$,
constrained by formula $C$. A novel concept is that of an {\tmem{environment
fragment}}, a triple:
\[ \Delta \assign \exists \bar{\beta} [D] \Gamma \]
where $\Gamma$ is a {\tmem{simple environment}} which assigns variables to
types (not type schemes). Environment fragments are used to describe the
static knowledge that is gained by successfully matching a value against a
pattern. We write $\exists \bar{\alpha} [C] \Delta$ for $\exists \bar{\alpha}
\bar{\beta} [C \wedge D] \Gamma$, and $\Delta_1 \times \Delta_2$ for:
\[ \exists \bar{\beta}_1 \bar{\beta}_2 [D_1 \wedge D_2] (\Gamma_1 \cup
   \Gamma_2) \]
Structures in which types can be interpreted have to meet three requirements,
see {\cite{simonet-pottier-hmg-toplas}} pages 18-19, of which we give two.
Every constraint of the form $\tau_1 \rightarrow \tau_2 \preccurlyeq
\varepsilon (\bar{\tau})$ or $\varepsilon (\bar{\tau}) \preccurlyeq \tau_1
\rightarrow \tau_2$ or $\varepsilon (\bar{\tau}) \preccurlyeq
\varepsilon^{\prime} (\bar{\tau}^{\prime})$ for $\varepsilon \neq
\varepsilon^{\prime}$, is unsatisfiable. $\tau_1 \rightarrow \tau_2
\preccurlyeq \tau_1^{\prime} \rightarrow \tau_2^{\prime}$ entails
$\tau_1^{\prime} \preccurlyeq \tau_1 \wedge \tau_2 \preccurlyeq
\tau_2^{\prime}$.

Judgments about expressions retain the same form as in $\tmop{HM} (X)$: they
are written $C, {\nobreak} \Gamma \vdash e : \sigma$, where $C$ represents an
assumption about the judgment's free type variables, $\Gamma$ assigns type
schemes to variables, and $\sigma$ is the type scheme assigned to $e$.

Judgments about patterns are written $C \vdash p : \tau \rightsquigarrow
\exists \bar{\beta} [D] \Gamma$, where the domain of $\Gamma$ is $\tmop{dpv}
(p)$. Such a judgment can be read: under assumption $C$, it is legal to match
a value of type $\tau$ against $p$; furthermore, if successful, this test
guarantees that there exist types $\bar{\beta}$ that satisfy $D$ such that
$\Gamma$ is a valid description of the values that the variables in
$\tmop{dpv} (p)$ receive. $D$ carries the information unpacked from guarded
data constructors $K$:
\[ \frac{\forall i \text{ \ \ } C \wedge D \vdash p_i : \tau_i
   \rightsquigarrow \Delta_i \text{ \ \ \ } K \colons \forall \bar{\alpha}
   \bar{\beta} [D] . \tau_1 \times \cdots \times \tau_n \rightarrow
   \varepsilon (\bar{\alpha}) \text{ \ \ \ } \bar{\beta} \# \tmop{FV} (C)}{C
   \vdash K p_1 \cdots p_n : \varepsilon (\bar{\alpha}) \rightsquigarrow
   \exists \bar{\beta} [D] (\Delta_1 \times \cdots \times \Delta_n)} \]
Subsequently, the unpacked information is available when type-checking the
expression in the corresponding branch:
\begin{equation}
  \frac{C \vdash p : \tau^{\prime} \rightsquigarrow \exists \bar{\beta} [D]
  \Gamma^{\prime} \text{ \ \ \ } C \wedge D, \Gamma \Gamma^{\prime} \vdash e :
  \tau \text{ \ \ \ } \bar{\beta} \# \tmop{FV} (C, \Gamma, \tau)}{C, \Gamma
  \vdash p.e : \tau^{\prime} \rightarrow \tau} \label{HMGClauseRule}
\end{equation}
The type judgments involving clauses are used to derive types of
$\lambda$-abstractions, i.e. anonymous functions defined by cases:
\[ \frac{\forall i \text{ \ \ } C, \Gamma \vdash c_i : \tau}{C, \Gamma \vdash
   \lambda (c_1 \cdots c_n) : \tau} \]
where $\tau$ is always of the form $\tau_1 \rightarrow \tau_2$.

Let us look at a selection of remaining rules. Patterns in a disjunction need
to agree on the information they bring about:
\[ \frac{\forall i \text{ \ \ } C \vdash p_i : \tau \rightsquigarrow \Delta}{C
   \vdash p_1 \vee p_2 : \tau \rightsquigarrow \Delta} \]
However, we can make them agree by discarding irrelevant information:
\[ \frac{C \vdash p : \tau \rightsquigarrow \Delta^{\prime} \text{ \ \ \ } C
   \Vdash \Delta^{\prime} \leqslant \Delta}{C \vdash p : \tau \rightsquigarrow
   \Delta} \]
where $C \Vdash \Delta^{\prime} \leqslant \Delta$ is defined in terms of
interpreting environment fragments as sets of ground environments.
Fortunately, we can equivalently define $C \Vdash \Delta^{\prime} \leqslant
\Delta$ as $C \wedge D^{\prime} \Vdash \exists \bar{\beta} .D \wedge
\Gamma^{\prime} \preccurlyeq \Gamma$, where $\Gamma^{\prime} \preccurlyeq
\Gamma$ is a shorthand for $\tmop{Dom} (\Gamma^{\prime}) = \tmop{Dom} (\Gamma)
\wedge_{x \in \tmop{Dom} (\Gamma^{\prime})} \Gamma^{\prime} (x) \preccurlyeq
\Gamma (x)$ and $\Delta = \exists \bar{\beta} [D] \Gamma$, $\Delta^{\prime} =
\exists \overline{\beta^{\prime}} [D^{\prime}] \Gamma^{\prime}$. Note that $C
\Vdash D$ is defined as $\vDash C \Rightarrow D$, i.e. holding in the model,
rather than $\vdash C \Rightarrow D$, i.e. provability.

To construct a value of a GADT, we need to check that the guard, or invariant,
holds:
\[ \frac{\forall i \text{ \ \ } C, \Gamma \vdash e_i : \tau_i \text{ \ \ \ } K
   \colons \forall \bar{\alpha} \bar{\beta} [D] . \tau_1 \cdots \tau_n
   \rightarrow \varepsilon (\bar{\alpha}) \text{ \ \ \ } C \Vdash D}{C, \Gamma
   \vdash K e_1 \cdots e_n : \varepsilon (\bar{\alpha})} \]
The use of a variable is type-checked in the same manner, the variable is
looked up in the environment $\Gamma$. The following rule can be derived in
$\tmop{HMG} (X)$ and is a bit more clear than the original:
\begin{equation}
  \frac{\Gamma (x) = \forall \bar{\alpha} [D] . \tau \text{ \ \ } C \Vdash
  \exists \bar{\alpha} .D}{C, \Gamma \vdash x : \tau} \label{HMGVarRule}
\end{equation}
The rule for recursive definition is initially given in the Milner-Mycroft
style to simplify proofs related to semantics rather than type inference:
\[ \frac{C, \Gamma \{ x \mapsto \sigma \} \vdash v : \sigma}{C, \Gamma \vdash
   \mu x.v : \sigma} \]
A $v$ instead of an $e$ is used to prevent diverging definitions. The final
form of $\tmop{HMG} (X)$, as employed in proofs of equivalence with the
constraint derivation-based presentation, enforces the type annotation: $e
\assign \ldots \; | \; \mu (x : \exists \bar{\beta} . \sigma) .v$, where
$\tmop{ftv} (\sigma) \subseteq \bar{\beta}$.

A closed (i.e. without unbound variables) expression $e$ is
{\tmem{well-typed}} if and only if $C, \varnothing \vdash e : \sigma$ holds
for some satisfiable constraint $C$. We have the following type soundness
results:

\begin{theorem}
  {\tmem{(Subject reduction).}} $C, \varnothing \vdash e : \sigma$ and $e
  \rightarrow e^{\prime}$ imply $C, \varnothing \vdash e^{\prime} : \sigma$.
\end{theorem}

\begin{theorem}
  {\tmem{(Progress).}} If $e$ is well-typed and contains exhaustive case
  analyses only, then it is either reducible or a value.
\end{theorem}

\begin{theorem}
  {\tmem{(Type soundness).}} If $e$ is well-typed and contains exhaustive case
  analyses only, then it does not reduce to a stuck
  expression.\label{HMGSoundness}
\end{theorem}

\subsection{Constraint Derivation}

The type inference for $\tmop{HMG} (X)$, as well as for {\tmname{InvarGenT}},
starts by generating a constraint for an expression, whose satisfiability
determines whether the expression is well-typed. The rules, or equations, for
deriving the constraints form an alternative specification of the type system.
Both specifications are readable, the specification by constraint derivation
is actually more concise. As with the type system rules, we have constraint
derivation equations for patterns, expressions, and pattern matching clauses.

The constraint $\llbracket p \downarrow \tau \rrbracket$ asserts that it is
legal to match a value of type $\tau$ against $p$, while the environment
fragment $\llbracket p \uparrow \tau \rrbracket$ represents knowledge about
the bindings that arise when such a test succeeds. (Note that our use of
$\downarrow$ and $\uparrow$ has nothing to do with bidirectional type
inference.) The rules for $\llbracket p \downarrow \tau \rrbracket$ and
$\llbracket p \uparrow \tau \rrbracket$ are the same in {\tmname{InvarGenT}}
as in $\tmop{HMG} (X)$. They are described in
{\cite{simonet-pottier-hmg-toplas}} on pages 32-33.\begin{table}[t]
  \begin{eqnarray*}
    \llbracket \Gamma \vdash x : \tau \rrbracket & = & \Gamma (x) \preccurlyeq
    \tau\\
    &  & \\
    \llbracket \Gamma \vdash \lambda \bar{c} : \tau \rrbracket & = & \exists
    \alpha_1 \alpha_2 . \wedge_c \llbracket \Gamma \vdash c : \alpha_1
    \rightarrow \alpha_2 \rrbracket \wedge \alpha_1 \rightarrow \alpha_2
    \preccurlyeq \tau\\
    &  & \\
    \llbracket \Gamma \vdash e_1 e_2 : \tau \rrbracket & = & \exists \alpha .
    \llbracket \Gamma \vdash e_1 : \alpha \rightarrow \tau \rrbracket \wedge
    \llbracket e_2 : \alpha \rrbracket\\
    &  & \\
    \llbracket K e_1 \cdots e_n : \tau \rrbracket & = & \exists \bar{\alpha}
    \bar{\beta} . \wedge_i \llbracket \Gamma \vdash e_i : \tau_i \rrbracket
    \wedge D \wedge \varepsilon (\bar{\alpha}) \preccurlyeq \tau\\
    &  & \text{where } K \colons \forall \bar{\alpha} \bar{\beta} [D] .
    \tau_1 \times \cdots \times \tau_n \rightarrow \varepsilon
    (\bar{\alpha})\\
    &  & \\
    \llbracket \Gamma \vdash \mu (x : \exists \bar{\beta} . \sigma) .e : \tau
    \rrbracket & = & \exists \bar{\beta} . \llbracket \Gamma \{ x \mapsto
    \sigma \} \vdash e : \sigma \rrbracket \wedge \sigma \preccurlyeq \tau\\
    &  & \\
    \llbracket \Gamma \vdash e : \forall \bar{\gamma} [C] . \tau \rrbracket &
    = & \forall \bar{\gamma} .C \Rightarrow \llbracket \Gamma \vdash e : \tau
    \rrbracket\\
    &  & \\
    \left\llbracket \text{{\tmstrong{let}}} x = e_1  \text{{\tmstrong{in}}}
    e_2 : \tau \right\rrbracket & = & \llbracket \Gamma \{ x \mapsto \forall
    \alpha [C] . \alpha \} \vdash e_2 : \tau \rrbracket \wedge \exists \alpha
    .C\\
    &  & \text{where } C \text{ is } \llbracket \Gamma \vdash e_1 : \alpha
    \rrbracket\\
    &  & \\
    \llbracket \Gamma \vdash p.e : \tau_1 \rightarrow \tau_2 \rrbracket & = &
    \llbracket p \downarrow \tau_1 \rrbracket \wedge \forall \bar{\beta} .D
    \Rightarrow \llbracket \Gamma \Gamma^{\prime} \vdash e : \tau_2
    \rrbracket\\
    &  & \text{where } \exists \bar{\beta} [D] \Gamma^{\prime} \text{ is }
    \llbracket p \uparrow \tau_1 \rrbracket
  \end{eqnarray*}
  \label{CnDerivHMG}
  \caption{Constraint derivation for $\tmop{HMG} (X)$: expressions and
  clauses}
\end{table}

We provide the $\tmop{HMG} (X)$ constraint derivation rules for expressions
and clauses in full in Table \ref{CnDerivHMG}, because they are analogous to
those in {\tmname{InvarGenT}}, but more concise. The novelty of $\tmop{HMG}
(X)$ compared to $\tmop{HM} (X)$ resides in the last rule, which deals with
clauses. The following description comes from
{\cite{simonet-pottier-hmg-toplas}} page 35. First, the function's domain type
is required to match the pattern's type, via the constraint $\llbracket p
\downarrow \tau_1 \rrbracket$. Then, the clause's right-hand side $e$ is
required to have type $\tau_2$ under a context extended with new abstract
types $\bar{\beta}$ and a new typing hypothesis $D$ and under an extended
environment $\Gamma$, all three of which are obtained by evaluating
$\llbracket p \uparrow \tau_1 \rrbracket$.

\subsection{Relevance}

The specification of the type system in terms of constraints can be seen as
even more declarative than the specification by natural deduction rules.
Instead of requiring that the reader builds understanding by analysing the
possible derivations, the constraint derivation rules reduce the meaning of
type judgments to that of formulas, whose semantics is already known.

There are three major differences between $\tmop{HMG} (X)$ and
{\tmname{InvarGenT}}. To simplify the already daunting task,
{\tmname{InvarGenT}} does not support subtyping. Recursive definitions do not
carry type annotations -- we perform type inference for polymorphic recursion.
The third difference is the requirement in {\tmname{InvarGenT}} that the
guards, or invariants, in type schemes and in value constructor definitions
are existentially quantified conjunctions of atoms rather than arbitrary
formulas. This is necessary to limit the space of candidate solutions to
predicate variables in the task of synthesizing types and invariants of
recursive definitions. But the restriction is also motivated by the need that
the invariants be readily understandable to the programmer. Logically complex
formulas, especially involving implications, are likely to not be sufficiently
self-explanatory.

As a consequence of restricting type scheme invariants to conjunctive
formulas, we cannot use the approach from Table \ref{CnDerivHMG} to
let-polymorphism. Instead, we perform ``division of labor'': we have
monomorphic bindings with a pattern on the left-hand-side
$\text{{\tmstrong{let}}} p = e \text{{\tmstrong{in}}} e$, and polymorphic
bindings, possibly recursive $\text{{\tmstrong{let rec}}} x = e
\text{{\tmstrong{in}}} e$.

\section{The $\tmop{OutsideIn} (X)$ System}\label{OutsideInSection}

The design of the Schrijvers, Peyton Jones, Sulzmann  and  Vytiniotis
$\tmop{OutsideIn} (X)$ {\cite{OutsideIn}} type system is a follow-up to the
{\tmname{OutsideIn}} algorithm from {\cite{ComplDecidableGADTs}} in a broader
context parameterized by a constraint logic. The type judgments in
$\tmop{OutsideIn} (X)$ are similar to $\tmop{HMG} (X)$. $C, \Gamma \vdash e :
\tau$ means that in a context where the constraint $C$ is available, and in
type environment $\Gamma$, the term $e$ has type $\tau$. The important
difference is that $C$ is a conjunction of atoms rather than an arbitrary
formula. Therefore, we cannot formulate an alternative, constraint derivation
based specification of the type system, with the elegant connection expressed
by results like Theorem \ref{InfCorrExpr}. Another difference between
$\tmop{OutsideIn} (X)$ and $\tmop{HMG} (X)$ is that $\tmop{OutsideIn} (X)$ is
defined against a background of a proof system for checking satisfiability of
formulas rather than against an abstract model. The relevance is that having a
model gives more flexibility for the implementer of type inference, while
having a proof system (i.e. a logic) gives more flexibility for the designer
of the type system. The third, related difference is that $\tmop{OutsideIn}
(X)$ uses quantifiers sparingly. Instead, it introduces a distinction into
{\tmem{skolem}} type variables and {\tmem{unification}} type variables.

$\tmop{OutsideIn} (X)$ has rules for type checking programs as well as single
expressions. Notably, the rule for un-annotated bindings makes use of
abduction:
\[ \frac{C_1, \Gamma \vdash e : \tau \text{ \ \ \ } \bar{\alpha} = \tmop{FV}
   (C, \tau) \text{ \ \ \ } \mathcal{C} \wedge C \Vdash C_1 \text{ \ \ \ }
   \mathcal{C}, \Gamma \{ f \mapsto \forall \bar{\alpha} [C] . \tau \} \vdash
   \tmop{prog}}{\text{ \ \ \ } \mathcal{C}, \Gamma \vdash f = e, \tmop{prog}}
\]
We determine the constraint $C_1$ which is required to make e typeable with
type $\tau$ in $\Gamma$. Next, we allow the invariant of $f$ be a simplified
version of $C_1$, namely $C$. Intuitively, $C$ is the ``extra information'',
not deducible from $\mathcal{C}$, that is needed to show the required
constraint $C_1$ (see {\cite{OutsideIn}} page 15). In our terms, $C$ is an
answer to an abduction problem $\mathcal{C} \Rightarrow C_1$.

Constraint generation in $\tmop{OutsideIn} (X)$ is similar to that in
$\tmop{HMG} (X)$, but the formulation is less concise, and the generated
constraint does not have alternating quantifiers. In our opinion, it makes the
semantics of the constraints less clear. The constraints are also more
restrictive than if quantifier scope was used to determine which solutions are
consistent, even for programs without type families or GADTs.

\subsection{Type Inference}

{\cite{OutsideIn}} page 20 defines a {\tmem{sound solution}}, which is
equivalent to our notion of an answer to a simple abduction problem, and a
{\tmem{guess-free solution}}, which is equivalent to our notion of a fully
maximal answer to a simple abduction problem. Type inference in
$\tmop{OutsideIn} (X)$ uses a solver of simple abduction problems, but the
abduction answers are not allowed to participate in the inferred types.
Instead, for type inference to succeed, the abduction answers to implications
have to be limited to substitutions of local variables, which are subsequently
ignored.

To make the exposition more concrete, we rewrite the solver infrastructure
specification from {\cite{OutsideIn}} page 41. Let $A, D_i, C_i$ be
conjunctions of atoms. Let $A \in \tmop{Abd}_{\mathcal{C}} (D_i, C_i)$ mean
that $A$ is an answer to the simple abduction problem $D_i \Rightarrow C_i$
under theory $\mathcal{C}$, i.e. $\mathcal{C} \wedge D_i \wedge A \Vdash C_i$.
Let $C \equiv \tmop{Simple} (C) \wedge \tmop{Implic} (C)$ be a decomposition
of a constraint $C$ into a conjunction of atoms $\tmop{Simple} (C)$ and a
conjunction of existentially quantified implications $\tmop{Implic} (C)$. For
a substitution $R = [\bar{\alpha} \assign \bar{\tau}] = [\alpha_1 \assign
\tau_1, \ldots, \alpha_n \assign \tau_n]$, let $\dot{R} = \bar{\alpha} \dot{=}
\bar{\tau} = \alpha_1 \dot{=} \tau_1 \wedge \ldots \wedge \alpha_n \dot{=}
\tau_n$. We define the solver recursively:
\[ \frac{\begin{array}{l}
     C_r \wedge \dot{R} \in \tmop{Abd}_{\mathcal{C}} (C_g, \tmop{Simple} (C))
     \text{ \ \ \ } \tmop{Dom} (R) \subseteq \bar{\alpha}\\
     \forall (\exists \bar{\alpha}_i .D_i \Rightarrow C_i \in \tmop{Implic}
     (C)) . \tmop{Solve} (\mathcal{C}; C_g \wedge C_r \wedge D_i ;
     \bar{\alpha}_i ; C_i) \rightsquigarrow (\varnothing, R_i)
   \end{array}}{\tmop{Solve} (\mathcal{C}; C_g ; \bar{\alpha}_i ; C_w)
   \rightsquigarrow (C_r, R)} \]
While it appears that the type inference solution is built out of abduction
answers $(C_r, R)$, the interesting simple abduction problems $D_i \Rightarrow
C_i$ are required to have abduction answer ($\dot{R_i}$) limited to variables
$\bar{\alpha}_i$, which can be subsequently ignored.

Motivated by considerations of efficiency and avoiding ambiguity,
{\cite{OutsideIn}} restricts the abduction algorithms to those only searching
for fully maximal answers. {\cite{OutsideIn}} conjectures that the algorithm
they provide is complete wrt. fully maximal answers to simple abduction
problems. $\tmop{OutsideIn} (X)$ might not be able to use the abduction
algorithm complete wrt. fully maximal answers to simple {\tmem{constraint}}
abduction problems from {\cite{AbductionSolvMaher}}, even when limited to
GADTs, because the $\tmop{OutsideIn} (X)$ type system is based on a logic
instead of on fixed models like Herbrand structures.

\subsection{Relevance}

The $\tmop{OutsideIn} (X)$ project exposition {\cite{OutsideIn}} lists
multiple challenges as being in its focus. Of these, {\tmname{InvarGenT}}
addresses GADTs, and is open to address {\tmem{units of measure}}, although
the corresponding sort has not been implemented. The other challenges fall
outside of the scope of {\tmname{InvarGenT}}: type-class constraints,
multi-parameter type classes with functional dependencies, and type families
with type family axioms. Some of these type system features, when implemented
directly, violate the requirement on the interpretation of types that we
preserve from $\tmop{HMG} (X)$: ``Every constraint of the form $\tau_1
\rightarrow \tau_2 \preccurlyeq \varepsilon (\bar{\tau})$ or $\varepsilon
(\bar{\tau}) \preccurlyeq \tau_1 \rightarrow \tau_2$ or $\varepsilon
(\bar{\tau}) \preccurlyeq \varepsilon^{\prime} (\bar{\tau}^{\prime})$ for
$\varepsilon \neq \varepsilon^{\prime}$, is unsatisfiable.'' (In
{\tmname{InvarGenT}}, we have equality $\dot{=}$ instead of subtyping
$\preccurlyeq$.)

{\cite{OutsideIn}} argues strongly against selecting an arbitrary correct type
when a definition does not have a most general type. {\tmname{InvarGenT}} does
select an arbitrary type without guarantees, although the implementation is
designed to pick useful types; e.g. if possible, with the return type of a
function sharing a parameter with an argument type. {\tmname{InvarGenT}} is
intended to provide type signatures for toplevel definitions to the
programmer. The programmer would then either accept the signature, modify the
program and re-generate the signature, or modify the signature directly.

$\tmop{OutsideIn} (X)$ shares with {\tmname{InvarGenT}} the restriction of
type scheme invariants, and guards of data constructors, to conjunctions of
atoms. On other points of difference between $\tmop{OutsideIn} (X)$ and
$\tmop{HMG} (X)$, {\tmname{InvarGenT}} follows $\tmop{HMG} (X)$. Both
$\tmop{OutsideIn} (X)$ and {\tmname{InvarGenT}} infer types for toplevel
definitions one at a time.

Finally, the type inference algorithm of $\tmop{OutsideIn} (X)$ is much weaker
than that of {\tmname{InvarGenT}}, as illustrated in our presentation of the
central step of the algorithm in terms of abduction. While our algorithm uses
abduction to search for the solution of the type inference problem,
$\tmop{OutsideIn} (X)$ only uses abduction to verify a partial solution found
by unification.

\section{Pointwise GADTs}\label{PointwiseSection}

Chuan-kai Lin in {\cite{PracticalGADTsInfer}}, see also Lin  and  Sheard
{\cite{PointwiseGADTs}}, presents a type inference algorithm for GADTs. Just
as {\tmname{InvarGenT}}, the type system Pointwise GADTs and the algorithm
$\mathcal{P}$ from {\cite{PracticalGADTsInfer}} does not require type
annotations. The Pointwise GADTs is not a constraint-based type system. Also,
it does not support deep patterns. Let us look at the two most important type
system rules. The recursive definition rule captures polymorphic recursion:
\[ \frac{\begin{array}{lll}
     \Gamma \{x \mapsto \forall \bar{\alpha} . \tau^{\prime} \} \vdash e_1 :
     \tau^{\prime} & \bar{\alpha} \# \tmop{FV} (\Gamma) & \Gamma \{x \mapsto
     \sigma\} \vdash e_2 : \tau
   \end{array}}{\Gamma \vdash \text{{\tmstrong{let rec}}} x = e_1 
   \text{{\tmstrong{in}}} e_2 : \tau} \]
The rule for pattern matching expressions is straightforward, nearly the same
as the corresponding derived rule in $\tmop{HMG} (X)$. We turn to the pattern
matching branch rule:
\[ \frac{\begin{array}{lll}
     K \colons \forall \bar{\alpha} . \bar{r} \longrightarrow \varepsilon
     \bar{s} &  & \bar{\alpha} \# \tmop{FV} (\Gamma, \bar{u}, \tau)\\
     S =\tmmathbf{P}\tmmathbf{U} (\varepsilon \bar{u} \dot{=} \varepsilon
     \bar{s}) &  & S (\Gamma \{ \bar{x} \assign \bar{r} \}) \vdash e : S
     (\tau)
   \end{array}}{\Gamma \vdash_p K \bar{x} .e : \varepsilon \bar{u} \rightarrow
   \tau} \]
$\tmmathbf{P}\tmmathbf{U}$ stands for {\tmem{pointwise unification}}. It is a
limited form of unification. If $U =\tmmathbf{P}\tmmathbf{U} (s \dot{=} t)$,
then $U (s) = U (t)$, but also: for $\alpha \in \tmop{Dom} (U)$, if $s
\downharpoonright_p = \alpha$, i.e. $s$ has variable $\alpha$ at position $p$,
then $t \downharpoonright_p = U (\alpha)$; similarly, if $t
\downharpoonright_p = \alpha$, then $s \downharpoonright_p = U (\alpha)$.

The type inference algorithm $\mathcal{P}$ is deliberately more restrictive
than the type system Pointwise GADTs. In particular, it uses
{\tmem{anti-unification}} to infer a tight type for the pattern-matched
expression, so that the pattern matching has a better chance of being
exhaustive given the type.

\subsection{Type Inference}

Algorithm $\mathcal{P}$ is based on Robin Milner's algorithm $\mathcal{W}$ as
in {\cite{Milner78}}, and its modification by Alan Mycroft to handle
polymorphic recursion as in {\cite{Mycroft84}}. See
{\cite{PracticalGADTsInfer}} pages 152-184 for detailed presentation. We
reproduce details of the algorithm $\mathcal{P}$, because \ algorithm
$\mathcal{P}$ provides a base-line wrt. which type inference for GADTs not
relying on type annotations can be compared. We defer the reader who finds the
current section cryptic to Chuan-kai Lin {\cite{PracticalGADTsInfer}}.

It might be worthwhile to collect the major pieces of algorithm $\mathcal{P}$
together:
\begin{eqnarray*}
  \tmop{infer} (\Gamma, e_1 e_2) & = & \\
  (S_1, \tau_1) & = & \tmop{infer} (\Gamma, e_1)\\
  (S_2, \tau_2) & = & \tmop{infer} (\Gamma, e_2)\\
  S_3 & = & \tmmathbf{U} (\tau_1 \dot{=} \tau_2 \rightarrow \beta) \text{, }
  \beta \text{ fresh}\\
  S & = & \tmmathbf{C} (S_1, S_2, S_3)\\
  &  & (S, S (\beta))\\
  &  & \\
  \tmop{infer} \left( \Gamma, \text{{\tmstrong{let rec}}} x = e \right) & = &
  \tmop{polyrec} (\{ \}, \forall \alpha . \alpha)\\
  \tmop{polyrec} (S_1, \sigma) & = & \\
  (S_2, \tau) & = & \tmop{infer} (\Gamma \{ u \mapsto \sigma \}, e)\\
  \bar{\beta} & = & \tmop{FV} (\tau) \backslash \tmop{FV} (\Gamma \{ u \mapsto
  \sigma \})\\
  \sigma^{\prime} & = & \forall \bar{\beta} . \tau\\
  \text{if} &  & S_2 (\sigma) = \sigma^{\prime}\\
  \text{then} &  & (S_2 S_1, \sigma^{\prime})\\
  \text{else} &  & \tmop{polyrec} (S_2 S_1, \sigma^{\prime})\\
  &  & \\
  \tmop{infer} (\Gamma, \tmmathbf{\tmop{case}} e \tmmathbf{\tmop{of}} 
  \overline{p_i .e_i}) & = & \\
  (S_0, u_0) & = & \tmop{infer} (\Gamma, e)\\
  (S_i, u_i, \tau_i) & = & \tmop{infer}_{\tmop{ALT}} (\Gamma, \beta, p_i .e_i)
  \text{, } \beta \text{ fresh}\\
  u & = & \tmop{LUB} (\tau_1, \ldots, \tau_n)\\
  S & = & \tmmathbf{U} ((u, \ldots, u) \dot{=} (u_0, u_1, \ldots, u_n))\\
  s & = & S (u)\\
  R_i & = & \tmmathbf{C} (S, S_0, S_i, \tmmathbf{U} (u_i \dot{=} \tau_i))\\
  \bar{\alpha} & = & \tmop{FV} (s) \cap (\cup_i \tmop{Dom} (R_i))\\
  \bar{\beta} & = & \tmop{FV} (s) \backslash \bar{\alpha}\\
  \tmop{trt} & = & \tmop{tabulate} (\bar{\alpha}, \overline{R_i})\\
  \tmop{btt} & = & \tmop{tabulate} (\beta \tmop{FV} (\Gamma),
  \overline{R_i})\\
  R & = & \tmop{reconcile} (\bar{\beta}, \tmop{trt}, \tmop{btt})\\
  &  & (R, R (\beta))\\
  &  & \\
  \tmop{infer}_{\tmop{ALT}} (\Gamma, \beta, K \bar{x} .e) & = & \\
  \forall \bar{\alpha} . \bar{r} \longrightarrow \varepsilon \bar{s} & = &
  \tmop{lookup} (K) \text{ \ where } \bar{\alpha} \text{ fresh}\\
  (S, \tau) & = & \tmop{infer} (\Gamma \{ \overline{x \mapsto r} \}, e)\\
  \bar{\alpha}^{\prime} & = & \bar{\alpha} \cap (\tmop{Dom} (S) \cup \tmop{FV}
  (\tmop{Rng} (S)) \cup \tmop{FV} (\tau))\\
  \bar{\gamma} & = & \{ \gamma | \gamma \in \bar{\alpha} \wedge \tmop{retain}
  (\bar{\alpha}, S, \gamma) \}\\
  \text{if} &  & \bar{\alpha}^{\prime} \nsubseteq \tmop{FV} (\bar{s})\\
  \text{then} &  & \bot\\
  \text{else if} &  & \bar{\gamma} = \varnothing\\
  \text{then} &  & ([\beta \assign \tau] S, \varepsilon \bar{\gamma}^{\prime},
  \varepsilon \bar{s}) \text{ where } \bar{\gamma}^{\prime} \text{ fresh}\\
  \text{else} &  & ([\beta \assign \tau] S, S (\tmop{transcb} (\bar{\gamma},
  \varepsilon \bar{s})), \varepsilon \bar{s})\\
  &  & \\
  \tmop{retain} (\bar{\alpha}, S, \gamma) & = & \\
  S (\gamma) = \varepsilon^{\prime}  \bar{r} & \rightarrow & \tmmathbf{T}\\
  S (\gamma) = \alpha & \rightarrow & \exists \beta \in \bar{\alpha} . \beta
  \neq \gamma \wedge \alpha \in \tmop{FV} (S (\beta))\\
  &  & \\
  \tmop{transcb} (\bar{\gamma}, \tau) & = & \\
  \text{if} &  & \bar{\gamma} \# \tmop{FV} (\tau)\\
  \text{then} &  & \beta \text{ where } \beta \text{ fresh}\\
  \text{else} &  & \\
  \tau = \varepsilon^{\prime}  \bar{r} & \rightarrow & \varepsilon^{\prime} 
  \overline{\tmop{transcb} (\bar{\gamma}, r)}\\
  \tau = \alpha & \rightarrow & \alpha
\end{eqnarray*}
The operation $\tmmathbf{C}$ is a combination of substitutions having the
effect of unification of the corresponding equations:
\[ \tmmathbf{C} (S_1, \ldots, S_n) =\tmmathbf{U} (\wedge_i \dot{S}_i)
   =\tmmathbf{U} ((\overline{a_1}, \ldots, \overline{\alpha_n}) \dot{=} (S_1
   (\overline{\alpha_1}), \ldots, S_n (\overline{\alpha_n}))) \text{ where }
   \overline{\alpha_i} = \tmop{Dom} (S_i) \]
The algorithm is designed after algorithm $\mathcal{W}$. However, it is made
more modular by the use of $\tmmathbf{C}$. This can be seen as a ``concealed''
form of constraint derivation based type inference, where the partial
constraints are solved during collection. The case $\tmop{infer} \left(
\Gamma, \text{{\tmstrong{let rec}}} x = e \right)$ is the Mycroft's
modification of algorithm $\mathcal{W}$ to handle polymorphic recursion. It
introduces iteration of type inference, for each recursive definition
separately, till the definition's type converges. The $\tmop{LUB} (\tau_1,
\ldots, \tau_n)$ subroutine used in $\tmop{infer} \left( \Gamma, {\nobreak}
\tmmathbf{\tmop{case}} e \tmmathbf{\tmop{of}}  \overline{p_i .e_i} \right)$ is
{\tmem{anti-unification}}, a special case of what we will call
{\tmem{constraint generalization}}. It computes the type $u$ for the
expression $e$ as specific as possible while agreeing with all types imposed
by the pattern matching branches. The complications in handling variables in
$\tmop{infer}_{\tmop{ALT}} (\Gamma, \beta, K \bar{x} .e)$, in particular the
failure condition $\bar{\alpha}^{\prime} \nsubseteq \tmop{FV} (\bar{s})$, are
analogous to the fact that the variables $\bar{\alpha}$ would be universally
quantified by $\tmop{HMG} (X)$, and thus cannot be part of, for example, the
domain of the resulting substitution.

The aspect of algorithm $\mathcal{P}$ that is closest to the intricacies of
term constraint abduction is the joint use of the $\tmop{tabulate}$ and
$\tmop{reconcile}$ subroutines (not presented above). They decompose
respectively the pattern-matched expression type, and the pattern-matching
branch body type, when the types start with the same type constructor
$\varepsilon$ across branches. When the types across branches do not agree,
but a $\tmop{btt}$ column corresponding to a subterm of the result type, i.e.
branch bodies type, is unifiable (branch-wise) with exactly one $\tmop{trt}$
column, then the subterm of the result type is replaced by the corresponding
subterm of the pattern-matched expression type.

\subsection{Relevance}

Algorithm $\mathcal{P}$ is designed to infer types for GADT programs without
relying on type annotations. This is the major goal of {\tmname{InvarGenT}},
alongside inference of existential types and ease of extension with novel
sorts. The type system behind algorithm $\mathcal{P}$ is not constraint based,
and in particular, it has no mechanism for extension with, for example,
numerical invariants. The Pointwise GADTs type system is much more restrictive
than the plain GADTs type system around which {\tmname{InvarGenT}} is
designed. Algorithm $\mathcal{P}$ is very efficient compared to
{\tmname{InvarGenT}}, but still quite capable. Some inspirations for
{\tmname{InvarGenT}} can be drawn from algorithm $\mathcal{P}$.

The examples from {\cite{PracticalGADTsInfer}} motivated the extension of
simple constraint abduction algorithm for terms in {\tmname{InvarGenT}} to
non-fully-maximal answers. The extension allows guessing equations between
invariant parameters. The use of anti-unification in algorithm $\mathcal{P}$
is interesting and the prospect of using constraint generalization in
{\tmname{InvarGenT}} to infer a tighter type when otherwise a pattern matching
expression would not be exhaustive, is intriguing, but we leave it as future
work. Finally, table-based approach to {\tmem{GADT type refinements}} in
algorithm $\mathcal{P}$, via the $\tmop{tabulate}$ and $\tmop{reconcile}$
subroutines, might inspire future work on augmenting the sequential approaches
to joint constraint abduction, by ``simultaneous'' stages, operating jointly
on all branches.

\section{Liquid Types}\label{LiquidSection}

Tim Freeman and Frank Pfenning's {\tmem{Refinement Types for ML}}
{\cite{RefinementFreemanPfenning}} predates even the Odersky, Sulzmann and
Wehr {\cite{HMX}}. Refinement types are so called because they refine the ML
type system, rather than extending it as GADTs do -- they do not make new
programs typeable. Refinement types became the universally and existentially
quantified types of DML. More recently, they were independently developed in
an inference-friendly way by Cormac Flanagan {\cite{Flanagan06}} and Knowles
and Flanagan {\cite{KnowlesF07}}. The work was continued by Ranjit Jhala,
Patrick Rondon and collaborators, and their systems {\tmname{DSolve}}
{\cite{DSolvePaper}} and {\tmname{HSolve}} {\cite{AbstractLiquid}} have become
more popular. They are fully focused on automatic invariant inference. The
work behind {\tmname{DSolve}} ({\cite{DSolvePaper}}, expanded in
{\cite{DSolveTechRep}} and {\cite{DSolveThesis}}) is the topic of this
section. {\cite{DSolvePaper}} cites the work on predicate abstraction
{\cite{AssertionGraphs}}, {\cite{AbstractStateGraphs}} as precursory for this
kind of invariant inference; and Jeffrey Scott Foster on Type Qualifiers
{\cite{TypeQualifiers}}, as inspiring some aspects of inference.

\subsection{The Type System}

Type refinements can be seen as preconditions (or postconditions) when they
refine the argument type (or the result type) of a function. The term
{\tmem{Liquid Types}} introduced by {\cite{DSolvePaper}} comes from
``logically qualified types'', a restriction of an otherwise more
depent-types-like type system, limiting type refinements to conjuctions of
atoms from a sort of linear inequalities and uninterpreted functions. The
refinements are over base (i.e. non-function) types, and are written: $\{ \nu
: B | e \}$, where $B$ is a base type like \tmverbatim{int}, $e$ is a
conjunction of atoms, and $\nu$ is a singled-out variable constrained by $e$.
(In the formalism used by {\tmname{InvarGenT}}, the letter $\delta$ will play
the role of $\nu$ here.) Recall the typing judgment form $C, \Gamma \vdash e :
\tau$ we encountered in systems described earlier: assuming constraint $C$
holds, in environment $\Gamma$, expression $e$ has type $\tau$. Although the
Liquid Types typing judgment $\Gamma \vdash e : T$ (where $T$ stands for a
liquid type or type scheme) seems simpler, here $C$ is a part of $\Gamma$. The
assumed constraint is:
\[ \llbracket \Gamma \rrbracket \equiv \wedge_{e \in \Gamma} e \wedge_{x : \{
   \nu : B | e \} \in \Gamma} e [\nu \assign x] \]
The Liquid Types type system is based on subtyping rather than type equality.
Here, subtyping implicitly captures the requirements on the assumed
constraint. The subtyping judgment has the form $\Gamma \vdash T_1 < : T_2$,
which means: $\llbracket \Gamma \rrbracket$ implies that $T_1$ is a subtype of
$T_2$.

The function type arguments are indexed by the variable corresponding to the
$\lambda$-abstraction introducing the type: $\Gamma \vdash \lambda x.e : (x :
T_x \rightarrow T)$. Together with the following rules:
\[ \frac{\Gamma (x) = \{ \nu : B | e \}}{\Gamma \vdash x : \{ \nu : B | \nu
   \dot{=} x \}} \]
\[ \frac{\Gamma \vdash e_1 : S \text{ \ \ } \Gamma ; x : S \vdash e_2 :
   T}{\Gamma \vdash \text{{\tmstrong{let}}} x = e_1  \text{{\tmstrong{in}}}
   e_2} \]
this avoids the use of existential quatification in inference constraints. The
crucial rules relevant for invariants are:
\[ \frac{\Gamma \vdash e_1 : \tmop{Bool} \text{ \ \ } \Gamma ; e_1 \vdash e_2
   : T \text{ \ \ } \Gamma ; \neg e_1 \vdash e_3 : T}{\Gamma \vdash
   \text{{\tmstrong{if}}} e_1  \text{{\tmstrong{then}}} e_2 
   \text{{\tmstrong{else}}} e_3 : T} \]
\[ \frac{\tmop{Valid} (\llbracket \Gamma \rrbracket \wedge e_1 \Rightarrow
   e_2)}{\Gamma \vdash \{ \nu : B | e_1 \} < : \{ \nu : B | e_2 \}} \]
In the first rule, the expressions $e_1$ are limited to well-formed formulas
of the constraint domain. Note how the former rule resembles the $\tmop{HMG}
(X)$ Rule \ref{HMGClauseRule} in that the local constraint is enhanced by the
information pertaining to the conditional branch. The latter rule resembles
the Rule \ref{HMGVarRule} in that the local constraint is used as a premise to
ensure the validity of the use of a value.

\subsection{Liquid Types Inference}

Similarly to {\tmname{{\tmname{InvarGenT}}}}, the {\tmname{DSolve}} system
first generates the inference constraint for the whole toplevel expression,
and then proceeds to solve the constraint. {\tmname{DSolve}} uses an oracle to
solve the Hindley-Milner polymorphic type inference subproblems, and the
resulting type shapes guide the construction of the constraint. In this way,
the constraint describes purely the subtyping issues, rather than both the
typing and subtyping issues. Each subtyping requirement contributes what would
be an implication in a normalized {\tmname{InvarGenT}} constraint. The Liquid
Type Inference Algorithm, including constraint generation, can be found in
{\cite{DSolvePaper}}, page 9, Figure 4. Liquid Types inference introduces
{\tmem{liquid type variables}} $\kappa$, standing for unknown invariants. (In
{\tmname{InvarGenT}}, we call them {\tmem{predicate variables}} $\chi$.) The
inference task is to find a substitution for the liquid type variables such
that all the implications corresponding to subtyping constraints are valid. We
present the liquid type variable solving part of the inference algorithm:
\begin{eqnarray}
  \tmop{Weaken} (\Gamma \vdash \{ \nu : B | \theta \cdot \kappa \}) A & = & A
  \circ \left[ \kappa \assign \\
  \text{ \ \ } \{ q \in A (\kappa) | \Gamma ; \nu : B \vdash \theta (q) :
  \tmop{Bool} \} \right]  \label{LiqAlgR1}\\
  \tmop{Weaken} (\Gamma \vdash \{ \nu : B | \rho \} < : \{ \nu : B | \theta
  \cdot \kappa \}) A & = & A \circ \left[ \kappa \assign \\
  \text{ \ \ } \{ q \in A (\kappa) | \llbracket A (\Gamma) \rrbracket \wedge
  A (\rho) \Rightarrow \theta (q) \} \right]  \label{LiqAlgR2}\\
  \tmop{Weaken} \_A & = & \bot  \label{LiqAlgR3}\\
  &  &  \nonumber\\
  \tmop{Solve} C A \text{ when } \exists c \in C \\
  \text{where } A (c) \text{ is not valid} & = & \tmop{Solve} C
  (\tmop{Weaken} c A) \nonumber\\
  \tmop{Solve} \_A & = & A \nonumber
\end{eqnarray}
The pair $\theta \cdot \kappa$ of a substitution $\theta$ and liquid type
variable $\kappa$ is a {\tmem{delayed substitution}}, see Knowles and Flanagan
{\cite{KnowlesF07}}. Entry \ref{LiqAlgR1} above ensures that the solution
formulas are well-formed. Entry \ref{LiqAlgR2} filters out the atoms which are
not implied by the premise of the implication considered. Entry \ref{LiqAlgR3}
is the failure case: the offending implication cannot be made valid. It cannot
be weakened, as it does not contain a liquid type variable in the conclusion,
and the premise is already as strong as possible. The initial call to
$\tmop{Solve}$ passes as the initial solution $A$, for each liquid type
variable $\kappa$, a conjunction of {\tmem{all}} atoms over the potential
parameters of invariants appearing in constraints.

\subsection{Relevance}

It will be illuminating for Chapter \ref{ChapSolv} to compare the solver of
{\tmname{DSolve}} with that of {\tmname{InvarGenT}}. In a single iteration of
the main algorithm of: {\tmname{DSolve}}, a single invariant (i.e. $\kappa$
substitution) is updated; {\tmname{InvarGenT}}, all invariants (i.e. $\chi$
substitutions) are updated. This makes: in {\tmname{DSolve}}, a single
implication valid, given before-update substitution of liquid type variables
in premises; in {\tmname{InvarGenT}}, all implications valid, given
before-update substitution of ``liquid type'', i.e. predicate variables in
conclusions. {\tmname{DSolve}} focuses on the role liquid type variables play
in conclusions, and ignores their role in premises, other than to verify
validity of implications. {\tmname{InvarGenT}} focuses on the role all
predicate variables play in premises, but also, using a different mechanism,
on postcondition predicate variables in conclusions. {\tmname{DSolve}} starts
with a full initial solution (a conjuction of all atoms, trivially
contradictory); {\tmname{InvarGenT}} starts with an empty initial solution (an
empty conjunction, trivially satisfiable). {\tmname{DSolve}} finds the most
specific (least general) solution for all invariants. {\tmname{InvarGenT}}
finds the least specific (most general) solution for preconditions, and the
most specific solution for postconditions given these preconditions.

Here are three arguments in favor of {\tmname{InvarGenT}} over
{\tmname{DSolve}}. The major limitation of {\tmname{DSolve}} is the need to
generate all the atoms as the initial solution; in {\tmname{InvarGenT}}, the
addition of atoms to the solution is driven by the conclusions that need to be
explained. Least general preconditions are less useful than most general
preconditions. Furthermore, in principle, there can be solutions that are
overlooked by solving one invariant at a time.

There are two shortcomings of {\tmname{InvarGenT}} relative to
{\tmname{DSolve}} that should inspire future work; we start with the major
one. {\tmname{InvarGenT}} prohibits passing values of an existentially
quantified type as arguments, and it does not allow type schemes as argument
types. Thus functions like \tmverbatim{ffor} from the \tmverbatim{fft}
example, see Appendix \ref{ExamplesCode} Subsection \ref{FFTExample}, require
manual packing and unpacking for the argument of the function passed to the
higher-order function. {\tmname{DSolve}} manages to infer types for such
higher-order functions, and this ability is further improved in Abstract
Refinement Types: Vazou, Rondon and Jhala {\cite{AbstractLiquid}}.

The other shortcoming is that, while {\tmname{InvarGenT}} is faster with
simple input programs, {\tmname{DSolve}} is faster with more complex programs
for which it is able to find a solution. There are two reasons. One is that
{\tmname{DSolve}} entirely eliminates variables which cannot be invariant
parameters before starting the solution process, i.e. intermediate variables
of the Hindley-Milner inference process. {\tmname{InvarGenT}} deploys the
general mechanism of abduction to solve for all variables. The other reason is
that {\tmname{DSolve}} only checks all constraints for contradiction once
every update of an invariant $\kappa$. {\tmname{InvarGenT}} checks all
constraints for contradiction for every candidate atom to be added during an
update of invariants.

\section{Herbrand Constraint Abduction}\label{HerbSection}

We have invoked ``abduction'' multiple times already, it is time to define the
term. {\tmem{Abduction}} is a form of inference where we find an explanation
of a target formula (often referred to as an observation) given background
knowledge. Given a signature $\Sigma$ and a set of variables $\tmop{Vars}$, a
{\tmem{constraint domain}} is a pair $(\mathcal{D}, \mathcal{L})$ where
$\mathcal{D}$ is a $\Sigma$-structure and $\mathcal{L}$ (the language of
constraints) is a set of $\Sigma$-formulas closed under conjunction and
renaming of free variables. {\tmem{Constraint abduction}} is a formalization
of abduction in the context of a constraint domain. {\tmem{Simple Constraint
Abduction}} is the task of solving an implication $B \Rightarrow C$, where $B,
C \in \mathcal{L}$ are conjunctions of atoms. The constraint abduction answer
$A \in \mathcal{L}$ is a solution to the implication $B \Rightarrow C$ if and
only if $\mathcal{D} \vDash (A \wedge B) \Rightarrow C$ and $\mathcal{D}
\vDash \exists \tmop{FV} (A, B) .A \wedge B$, where $\tmop{FV} (\cdummy)$
finds the free variables of an expression. $\mathcal{D}$ and $B$ play the role
of background knowledge -- $\mathcal{D}$ general, and $B$ context specific.
$C$ plays the role of observation, and $A$ of explanation. If $B \Rightarrow
C$ is a conjunct in a type inference constraint, then $B$ contains the
background information about a particular location in a program's source code,
coming from the pattern matching patterns that ``are the case'' -- the
location is in their scope. $C$ contains the requirements imposed by the
source at that location, for example the preconditions of the functions that
are called. The answer $A$ explains what needs to be the case for the
requirements to be met. {\tmem{Joint Constraint Abduction}} is the task of
solving implications $\wedge_i (B_i \Rightarrow C_i)$, where $B_i, C_i \in
\mathcal{L}$ are conjunctions of atoms and $\wedge_i \varphi_i$ stands for
$\varphi_1 \wedge \ldots \wedge \varphi_n$. The answer $A \in \mathcal{L}$ to
this Joint Constraint Abduction problem has to meet the conditions:
$\mathcal{D} \vDash (A \wedge B_i) \Rightarrow C_i$ and $\mathcal{D} \vDash
\exists \tmop{FV} (A, B_i) .A \wedge B_i$, for all $i$. The suitability of
joint constraint abduction for GADTs type inference was probably first
observed by Martin Sulzmann, Tom Schrijvers and Peter J. Stuckey, see
{\cite{SulzmannAbdGADTsRep}} (originally {\cite{SulzmannAbdGADTs}}) and
{\cite{AbdEADTs}}.

Michael Maher has studied constraint abduction for terms
({\cite{AbductionMaher}}, {\cite{AbductionSolvMaher}}) and linear arithmetic
({\cite{AbductionLinArithMaher}}). Michael Maher and Ge Huang
{\cite{AbductionSolvMaher}} provided the basis for our Simple Constraint
Abduction algorithm for terms, which we describe in this section. An abduction
answer $A$ is {\tmem{maximally general}}, when for every other abduction
answer $A^{\prime}$, if $\mathcal{D} \vDash A \Rightarrow A^{\prime}$, then
$\mathcal{D} \vDash A^{\prime} \Rightarrow A$. An abduction answer $A$ to $B
\Rightarrow C$ is {\tmem{fully maximal}}, when it is maximally general and
$\mathcal{D} \vDash (A \wedge B) \Leftrightarrow (B \wedge C)$. A fully
maximal answer does not ``guess'' any fact not entailed by the formulas
considered.\begin{table}[t]
  \begin{tabular}{l}
    $\tmop{pos} (t, A)$ \ \ returns the set of positions of the term $t$ in
    the right-hand-side of $A$\\
    $\tmop{repl} (S, A)$ \ replaces all terms in the right-hand-side of $A$
    occurring at a position\\
    \ \ \ \ \ \ \ \ \ \ \ \ \ \ in $S$ by a new variable (that is
    existentially quantified)\\
    $\tmop{next} (A)$ \ \ \ is the set $\{ \tmop{repl} (S, A) | \varnothing
    \subsetneq S \subset \tmop{pos} (t, A), t \nin \tmop{Vars} \}$\\
    \\
    {\tmstrong{algorithm}} $\tmop{FMA} (B, C)$\\
    \ \ \begin{tabular}{l}
      {\tmstrong{if}} $\vDash B \Rightarrow C$ {\tmstrong{then return}}
      $\top$\\
      {\tmstrong{let}} $A$ be the standard form of $B \wedge C$\\
      {\tmstrong{do}}\\
      \ \ \begin{tabular}{l}
        {\tmstrong{let}} $\mathcal{A}$ be $\tmop{next} (A)$\\
        {\tmstrong{if }}$(\forall A^{\prime} \in \mathcal{A}. \nvDash
        A^{\prime} \wedge B \Rightarrow C)$ {\tmstrong{then return}} $A$\\
        {\tmstrong{choose}} $A \in \mathcal{A}$ such that $\vDash A \wedge B
        \Rightarrow C$
      \end{tabular}
    \end{tabular}\\
    \\
    {\tmstrong{algorithm JCA-Solve}}\\
    \ \ \begin{tabular}{l}
      $\mathcal{M} \assign \{ \wedge_{i = 1}^n A_i | A_i \text{ is a maximally
      general answer of the } i \text{'th SCA problem} \}$\\
      {\tmstrong{while }}exist $A \in \mathcal{M}$ and $i$ s.t. $A \wedge B_i$
      is unsatisfiable in $\mathcal{D}$, $\mathcal{M} \assign
      \mathcal{M}\backslash \{ A \}$\\
      {\tmstrong{if }}$\mathcal{M}= \varnothing$ {\tmstrong{then return
      }}$\bot$\\
      {\tmstrong{while}} exist $A, A^{\prime} \in \mathcal{M}$ s.t.
      $\mathcal{D} \vDash A \Rightarrow A^{\prime}$ and $A \neq A^{\prime}$,
      $\mathcal{M} \assign \mathcal{M}\backslash \{ A \}$\\
      {\tmstrong{return }}$\mathcal{M}$
    \end{tabular}
  \end{tabular}\label{HerbAbdAlg}
  \caption{SCA algorithm for computing fully maximal answers, JCA algorithm}
\end{table}

Let $\mathcal{D}=\mathcal{T}= T (\Sigma, \tmop{Vars})$ be the free algebra of
terms over signature $\Sigma$ with variables $\tmop{Vars}$ and
$\mathcal{L}=\mathcal{F}\mathcal{T}_{\exists}$ be existentially quantified
conjunctions of equations. In Table \ref{HerbAbdAlg} we quote the fully
maximal abduction algorithm from {\cite{AbductionSolvMaher}}, Figure 4, page
13. The following theorem is Theorem 6 from {\cite{AbductionSolvMaher}}, page
14.

\begin{theorem}
  \label{MaherHerbSolve}Algorithm $\tmop{FMA}$ outputs all fully maximal
  answers to the SCA problem over $\mathcal{F}\mathcal{T}_{\exists}$ and
  terminates.
\end{theorem}

Now let us turn to the joint problem. The following proposition is Proposition
8 from {\cite{AbductionMaher}}, page 9. In Table \ref{HerbAbdAlg} we quote the
corresponding algorithm {\tmstrong{JCA-Solve}}.

\begin{proposition}
  \label{MaherJCASolve}Consider a joint constraint abduction problem composed
  of $n$ SCA component problems. Let $\mathcal{A}= \{ \wedge_{i = 1}^n A_i |
  A_i \text{ is a maximally general answer of the } i \text{'th SCA problem}
  \}$. If $A$ is a maximally general answer to the JCA problem then $A \in
  \mathcal{A}$. The JCA problem has no answers iff no constraint in
  $\mathcal{A}$ is an answer.
\end{proposition}

\subsection{Relevance}

Sulzmann, Schrijvers and Stuckey {\cite{SulzmannAbdGADTs}} and
{\cite{AbdEADTs}} inspired the approach taken by the {\tmname{InvarGenT}}
project. Maher and Huang {\cite{AbductionSolvMaher}} forms a basis of our
simple constraint abduction solver for the Herbrand domain. It turns out
however, that fully maximal answers are insufficient for type and invariant
inference. We go beyond the algorithms from Table \ref{HerbAbdAlg} as follows.

We extend the $\tmop{FMA} (B, C)$ algorithm by considering atoms $x \dot{=}
y$, where $x, y \in \tmop{Vars}$, $x \dot{=} t_x$ and $y \dot{=} t_y$ belong
to the solved form of $B \wedge C$, $t_x \nin \tmop{Vars}$ and $t_y \nin
\tmop{Vars}$, and $t_x \dot{=} t_y$ is satisfiable. These conditions limit the
number of pairs of variables to consider. There are many potential SCA
problems where we still will not be able to find an answer, but they appear
not to be relevant: we have not encountered a practical example where the
algorithm would need further extension.

We extend the {\tmstrong{JCA-Solve}} algorithm by using the partial solution,
i.e. starting from $B \wedge C \wedge_{i = 1}^{k - 1} A_i$ instead of $B
\wedge C$, when solving the $k$th component SCA problem.

We propose a novel algorithm (not based on {\cite{AbductionLinArithMaher}})
for constraint abduction for the linear arithmetic domain.{\newpage}

\chapter{The Type System}

The main idea of this work is that useful properties of functions can be
generated by solving ordinary-looking constraints under a quantifier prefix
(without resorting to Herbrandization), derived via a GADTs-based type system.
The content of Sections \ref{TypSys} and \ref{TypInfCns} is based on Simonet
and Pottier {\cite{simonet-pottier-hmg-toplas}}. To our knowledge, the idea
behind the {\tmname{NegClause}} rule is novel. The content of Sections
\ref{ExTypes} and \ref{ExTypesCns} is novel.

We start by introducing notation. By the bar $\bar{e}$ we denote a sequence
(or a set, depending on context) of elements $e$, by $\#$ we mean
disjointedness. With a free index $i$, $\overline{e_i}$ means $(e_1, \ldots,
e_n)$ for some $n$ associated with the index $i$; i.e. $\overline{e_i}$ or
$\bar{e}$ is a sequence $(e_1, \ldots, e_n)$ and $e_i$ is an $i$th element of
the sequence. Similarly, $\wedge_i \Phi_i$ denotes $\Phi_1 \wedge \ldots
\wedge \Phi_n$. For convenience, we treat a conjunction of atoms $\wedge_i
c_i$ as a set of atoms $\{ c_1, \ldots, c_n \}$.

In some contexts, for a quantifier prefix $\mathcal{Q}$ we write $\mathcal{Q}$
to denote the set of variables quantified by $\mathcal{Q}$. Let $\tmop{FV}$ be
a generic function returning the free variables of any expression. For a
quantifier prefix $\mathcal{Q}$ and variables $x, y$ in $\mathcal{Q}$, by $x
<_{\mathcal{Q}} y$ we denote that $x$ is to the left of $y$ in $\mathcal{Q}$
and they are separated by a quantifier alternation, by $x
\leqslant_{\mathcal{Q}} y$ that it is not the case that $y <_{\mathcal{Q}} x$.

By $\Phi [\bar{\alpha} \assign \bar{t}]$, $\Phi [\overline{\alpha \assign
t}]$, or $\Phi [\alpha_1 \assign t_1 ; \ldots ; \alpha_n \assign t_n]$, we
denote a substitution of terms $\bar{t}$ for corresponding variables
$\bar{\alpha}$ in the formula $\Phi$ (where $\bar{\alpha}$ and $\bar{t}$ are
finite sequences of the same length). Equality with a dot $\dot{=}$ is an
object-level equality, i.e. a relation in the language $\mathcal{L}$
introduced below. Equality $=$ is a meta-level equality, usually syntactic
equality on terms or formulas. By $\bar{s} \dot{=} \bar{t}$ we denote
$\wedge_i s_i \dot{=} t_i$, where $\bar{s} = (s_1, \ldots, s_n)$ and $\bar{t}
= (t_1, \ldots, t_n)$ for some $n$. We use letters $R, S, U$ to denote
substitutions. For a substitution $S = [\bar{\alpha} \assign \bar{t}]$, we
write substitution application as $S (\Phi) = \Phi [\bar{\alpha} \assign
\bar{t}]$; we write $\dot{S} = \bar{\alpha} \dot{=} \bar{t}$; and we denote
the substitution $S$ corresponding to a formula $A = \dot{S} = \bar{\alpha}
\dot{=} \bar{t}$ by $\tilde{A}$. We say that a substitution $[\bar{\alpha}
\assign \bar{t}]$ agrees with a quantifier prefix $\mathcal{Q}$, when $\vDash
\mathcal{Q}. \bar{\alpha} \dot{=} \bar{t}$ and in case of $\alpha_1 \dot{=}
\alpha_2 \in \bar{\alpha} \dot{=} \bar{t}$ for variables $\alpha_1, \alpha_2$,
we have $\alpha_2 \leqslant_{\mathcal{Q}} \alpha_1$. Syntactically, this means
that $\alpha_i$ is not to the right of variables in $t_i$: $\forall \beta \in
\tmop{FV} (t_i)$, $\alpha_i \nless_{\mathcal{Q}} \beta$. We use letters $\tau,
r, s, t, u$ to denote terms and letters $\alpha, \beta, v, x, y, z$ to denote
variables.

\section{The Language of Constraints}

We are interested in a multi-sorted first order language $\mathcal{L}_1$ with
equality, interpreted in a given, fixed model $\mathcal{M}$. Even when we
write $\vDash \Phi$, it is usually a shortcut notation for validity of a
formula $\Phi$ in the model $\mathcal{M}$: $\mathcal{M} \vDash \Phi$. The sort
of terms or ``types proper'', denoted $s_{\tmop{type}}$ and \tmverbatim{type}
in the concrete syntax of {\tmname{InvarGenT}}, plays a special role. In the
current presentation, we abstract from details of the language, posing the
necessary properties as assumptions.\begin{table}[t]
  \[ \begin{array}{lll}
       \tau_{\tmop{type}} & \assign & v_{\tmop{type}} \; | \;
       \tau_{\tmop{type}} \rightarrow \tau_{\tmop{type}} \; | \; \varepsilon
       (\bar{\tau}) \; | \; \tmop{Num} (\tau_{\tmop{num}})\\
       \tau_{\tmop{num}} & \assign & v_{\tmop{num}} \; | \; k
       \tau_{\tmop{num}} \; | \; \tau_{\tmop{num}} + \tau_{\tmop{num}} \; | \;
       k\\
       k & \assign & 0 \; | \; 1 \; | \; - 1 \; | \; 2 \; | \; - 2 \; | \;
       \ldots\\
       \tau & \assign & \tau_{\tmop{type}} \; | \; \tau_{\tmop{num}}\\
       a_{\tmop{type}} & \assign & \tau_{\tmop{type}} \dot{=}
       \tau_{\tmop{type}}\\
       a_{\tmop{num}} & \assign & \tau_{\tmop{num}} \dot{=} \tau_{\tmop{num}}
       \; | \; \tau_{\tmop{num}} \leq \tau_{\tmop{num}}\\
       & | & v_{\tmop{num}} \dot{=} \max (\tau_{\tmop{num}},
       \tau_{\tmop{num}}) \; | \; v_{\tmop{num}} \dot{=} \min
       (\tau_{\tmop{num}}, \tau_{\tmop{num}})\\
       & | & v_{\tmop{num}} \leq \max (\tau_{\tmop{num}}, \tau_{\tmop{num}})
       \; | \; \min (\tau_{\tmop{num}}, \tau_{\tmop{num}}) \leq
       v_{\tmop{num}}\\
       a & \assign & a_{\tmop{type}} \; | \; a_{\tmop{num}}\\
       v & \assign & v_{\tmop{type}} \; | \; v_{\tmop{num}} = \alpha_i \; | \;
       \beta_i \; | \; \ldots\\
       \sigma & \assign & \forall v_{\tmop{type}}^1 [\exists \bar{\alpha} .
       \bar{a}] .v_{\tmop{type}}^1 \; | \; \forall \bar{v} [\bar{a}] .
       \tau_{\tmop{type}}
     \end{array} \]
  \label{TypeSyn}
  \caption{Abstract syntax of types and type schemes}
\end{table}

In Appendix \ref{ChapProofs}, we introduce a Henkin semantics for existential
second order logic $\mathcal{L}$ extending $\mathcal{L}_1$ by predicate
symbols $\chi (\cdummy)$ that we call {\tmem{predicate variables}}.
$\mathcal{L}$ is tailored to our needs of invariant and postcondition
inference. Let $\tmop{PV} (\cdummy)$ be the set of predicate variables in any
expression. We define {\tmem{solved form formulas}} to be existentially
quantified conjunctions of atoms $\exists \bar{\alpha} .A$ without predicate
variables. An interpretation of predicate variables $\mathcal{I}$ substitutes
predicate variables by solved form formulas.

In Table \ref{TypeSyn}, we present a particular instance of $\mathcal{L}_1$,
introducing a numerical constraint domain. The terms of $\mathcal{L}_1$ are
$\tau$, and $a$ are the atomic formulas, as currently implemented in
{\tmname{InvarGenT}}. There are two sorts: $s_{\tmop{type}}$ with terms
$\tau_{\tmop{type}}$ and relations $a_{\tmop{type}}$, and $s_{\tmop{num}}$
with terms $\tau_{\tmop{num}}$ and relations $a_{\tmop{num}}$. The remaining
entry $\sigma$ is built on top of $\mathcal{L}_1$ rather than part of it. The
notation $v_{\tmop{type}}^1$ stands for the same variable of sort
$s_{\tmop{type}}$ in both occurrences. Besides $\tau$, we also use the letter
$t$ for terms of $\mathcal{L}_1$. Rarely, in interests of readability we use
the letters $x, y, z$ besides $\alpha, \beta$ to denote variables of
$\mathcal{L}_1$.

\section{The Type System with GADTs}\label{TypSys}

By {\tmem{types}} $\tau$ we mean terms of sort $s_{\tmop{type}}$. Define
{\tmem{type schemes}} $\sigma$ as $\forall \beta [D] . \beta$, where $D$ is
either a solved form formula $\exists \bar{\alpha} .E$ or a predicate variable
$\chi (\beta)$, and $\beta$ is a variable of sort $s_{\tmop{type}}$. A
{\tmem{simple environment}} (or {\tmem{monomorphic environment}}) maps
variables $x$ to types $\tau$. An {\tmem{environment}} (or {\tmem{polymorphic
environment}}) maps variables $x$ to type schemes $\sigma$. When a simple
environment is appended to an environment, we identify $\tau$ and $\forall
\beta [\beta \dot{=} \tau] . \beta$ for $\beta \nin \tmop{FV} (\tau)$. When
operations pertaining to formulas are applied to a type scheme $\forall \beta
[\exists \bar{\alpha} .E] . \beta$ or $\forall \beta [\chi (\beta)] . \beta$,
they are performed on the formula $\exists \bar{\alpha} .E$ or $\chi (\beta)$.
$\forall \beta \bar{\alpha} [b \dot{=} \tau \wedge E] . \tau$ is a notational
variant of $\forall \beta [\exists \bar{\alpha} .E] . \beta$. When operations
pertaining to type schemes (types) are applied to (simple) environments
$\Gamma$, they are performed on the image of $\Gamma$. Define
{\tmem{environment fragments}} $\Delta$ to be triples $\exists \bar{\alpha}
[D] . \Gamma$ of variables $\bar{\alpha}$, atomic conjunctions $D$ in
$\mathcal{L}$ and simple environment $\Gamma$.

Table \ref{ExprSyn} presents expressions currently in {\tmname{InvarGenT}},
more domain-specific $\tmmathbf{\tmop{assert}}$ and $\tmmathbf{\tmop{when}}$
clauses can be added. The table defines several expression languages,
underlying the corresponding type systems: $\tmop{MMG} (X)$, its supersets --
in terms of expressions $e$ and system rules -- $\tmop{MMG} (\tmop{num})$ and
$\tmop{MMG}_{\exists} (X)$, and $\tmop{MMG}_{\exists} (\tmop{num})$ which is a
union of $\tmop{MMG} (\tmop{num})$ and $\tmop{MMG}_{\exists} (X)$. The
domain-independent languages $\tmop{MMG} (X)$ and $\tmop{MMG}_{\exists} (X)$
are extended with some expressions specific to the numerical sort in
$\tmop{MMG} (\tmop{num})$ and $\tmop{MMG}_{\exists} (\tmop{num})$. The
language $\mathcal{L}_1$ and model $\mathcal{M}$ for $\tmop{MMG} (X)$ and
$\tmop{MMG}_{\exists} (X)$ are arbitrary. We fix the language $\mathcal{L}_1$
for $\tmop{MMG} (\tmop{num})$ and $\tmop{MMG}_{\exists} (\tmop{num})$ as
having besides $s_{\tmop{type}}$ only a linear arithmetic sort. We discuss the
type system $\tmop{MMG} (\tmop{num})$ shortly, and extend it to
$\tmop{MMG}_{\exists} (\tmop{num})$ in Section \ref{ExTypes}. Disjunctive
patterns are not yet implemented in {\tmname{InvarGenT}}, we omit them in the
presentation for brevity.\begin{table}[t]
  $p \assign x \; | \; 0 \; | \; 1 \; | \; p \wedge p \; | \; K \bar{p}
  
  c \assign
  
  \text{ \ } \tmop{MMG} (X) :
  
  \text{ \ \ \ \ \ \ } p.e
  
  \text{ \ } \tmop{MMG} (\tmop{num}) :
  
  \text{ \ \ \ \ \ } | \; p \text{{\tmstrong{when}}}  \overline{e \leqslant e}
  .e
  
  e \assign
  
  \text{ \ } \tmop{MMG} (X) :
  
  \text{ \ \ \ \ \ } x \; | \; K \bar{e} \; | \; \text{{\tmstrong{let}}} p = e
  \text{{\tmstrong{in}}} e \; | \; e e \; | \; \lambda (\bar{c}) \; | \;
  \text{{\tmstrong{let rec}}} x = e \text{{\tmstrong{in}}} e
  
  \text{ \ \ \ \ } | \; \text{{\tmstrong{assert false}}}
  
  \text{ \ \ \ \ } | \; \text{{\tmstrong{assert type}}} e = e ; e
  
  \text{ \ } \tmop{MMG} (\tmop{num}) :
  
  \text{ \ \ \ \ } | \; \text{{\tmstrong{assert num}}} e \leqslant e ; e
  
  \text{ \ } \tmop{MMG}_{\exists} (X) :
  
  \text{ \ \ \ \ } | \; \lambda [K] \bar{c}
  
  \text{ \ \ \ \ } | \; \lambda_K \bar{c}$
  
  \label{ExprSyn}
  \caption{Abstract syntax of expressions}
\end{table}

First, we present the type system in the standard, natural deduction style.
The {\tmem{type judgment}} $C, \Gamma \vdash e : \tau$ or $C, \Gamma \vdash e
: \sigma$ is composed of a formula $C$ without predicate variables, an
environment $\Gamma$, an expression $e$ and a type $\tau$ or type scheme
$\sigma$. Not mentioned explicitly is a set of data constructors $\Sigma$,
which is fixed when typing an expression.\begin{table}[t]
  Syntax-directed:
  
  \begin{tabular}{lll}
    {\tmname{p-Empty}} & {\tmname{p-Wild}} & {\tmname{p-Var}}\\
    $C \vdash 0 : \tau \longrightarrow \exists \varnothing [\tmmathbf{F}]
    \{\}$ & $C \vdash 1 : \tau \longrightarrow \exists \varnothing
    [\tmmathbf{T}] \{\}$ \ \ \ \ \  & $C \vdash x : \tau \longrightarrow
    \exists \varnothing [\tmmathbf{T}] \{x \mapsto \tau\}$\\
    &  & \\
    {\tmname{p-Cstr}} &  & {\tmname{p-And}}\\
    $\large{\frac{\begin{array}{llll}
      \forall i & C \wedge D \vdash p_i : \tau_i \longrightarrow \Delta_i & K
      \colons \forall \bar{\alpha} \bar{\beta} [D] . \tau_1 \times \ldots
      \times \tau_n \rightarrow \varepsilon (\bar{\alpha}) & \bar{\beta} \#
      \tmop{FV} (C)
    \end{array}}{C \vdash K p_1 \ldots p_n : \varepsilon (\bar{\alpha})
    \longrightarrow \exists \bar{\beta} [D] (\Delta_1 \times \ldots \times
    \Delta_n)}}$ &  & $\large{\frac{\forall i \text{ \ } \; C \vdash p_i :
    \tau \longrightarrow \Delta_i}{C \vdash p_1 \wedge p_2 : \tau
    \longrightarrow \Delta_1 \times \Delta_2}}$
  \end{tabular}
  
  \
  
  \
  
  Non-syntax-directed:
  
  \begin{tabular}{lllll}
    {\tmname{p-EqIn}} &  & {\tmname{p-SubOut}} &  & {\tmname{p-Hide}}\\
    $\large{\frac{\begin{array}{l}
      C \vdash p : \tau^{\prime} \longrightarrow \Delta\\
      C \vDash \tau \dot{=} \tau^{\prime}
    \end{array}}{C \vdash p : \tau \longrightarrow \Delta}}$ &  \ \  &
    $\large{\frac{\begin{array}{l}
      C \vdash p : \tau \longrightarrow \Delta^{\prime}\\
      C \vDash \Delta^{\prime} \leqslant \Delta
    \end{array}}{C \vdash p : \tau \longrightarrow \Delta}}$ &  \ \  &
    $\large{\frac{\begin{array}{l}
      C \vdash p : \tau \longrightarrow \Delta\\
      \bar{\alpha} \# \tmop{FV} (\tau, \Delta)
    \end{array}}{\exists \bar{\alpha} .C \vdash p : \tau \longrightarrow
    \Delta}}$
  \end{tabular}\label{TypRulPat}
  \caption{Typing rules for $\tmop{MMG} (X) : \tmop{patterns}$}
\end{table}\begin{table}[t]
  Syntax-directed:
  
  \begin{tabular}{lll}
    {\tmname{Var}} &  & {\tmname{Cstr}}\\
    $\large{\frac{\begin{array}{ll}
      \Gamma (x) = \forall \beta [\exists \bar{\alpha} .D] . \beta & C \vDash
      D
    \end{array}}{C, \Gamma \vdash x : \beta}}$ &  &
    {\large{$\frac{\begin{array}{l}
      \forall i C, \Gamma \vdash e_i : \tau_i \text{ \ \ \ \ \ } C \vDash D\\
      K \colons \forall \bar{\alpha} \bar{\beta} [D] . \tau_1 \ldots \tau_n
      \rightarrow \varepsilon (\bar{\alpha})
    \end{array}}{C, \Gamma \vdash K e_1 \ldots e_n : \varepsilon
    (\bar{\alpha})}$}}\\
    &  & \\
    {\tmname{App}} &  & {\tmname{Abs}}\\
    {\large{$\frac{\begin{array}{l}
      C, \Gamma \vdash e_1 : \tau^{\prime} \rightarrow \tau\\
      C, \Gamma \vdash e_2 : \tau^{\prime}
    \end{array}}{C, \Gamma \vdash e_1 e_2 : \tau}$}} &  &
    {\large{$\frac{\forall i C, \Gamma \vdash c_i : \tau_1 \rightarrow
    \tau_2}{C, \Gamma \vdash \lambda (c_1 \ldots c_n) : \tau_1 \rightarrow
    \tau_2}$}}\\
    &  & \\
    {\tmname{LetIn}} &  & {\tmname{LetRec}}\\
    {\large{$\frac{C, \Gamma \vdash \lambda (p.e_2) e_1 : \tau}{C, \Gamma
    \vdash \text{{\tmstrong{let}}} p = e_1  \text{{\tmstrong{in}}} e_2 :
    \tau}$}} &  & $\large{\frac{\begin{array}{ll}
      C, \Gamma^{\prime} \vdash e_1 : \sigma & C, \Gamma^{\prime} \vdash e_2 :
      \tau\\
      \sigma = \forall \beta [\exists \bar{\alpha} .D] . \beta &
      \Gamma^{\prime} = \Gamma \{x \mapsto \sigma\}
    \end{array}}{C, \Gamma \vdash \text{{\tmstrong{let rec}}} x = e_1 
    \text{{\tmstrong{in}}} e_2 : \tau}}$\\
    &  & \\
    {\tmname{AssertLeq}} &  & {\tmname{AssertEqty}}\\
    $\large{\frac{\begin{array}{lll}
      C, \Gamma \vdash e_1 : \tmop{Num} (\tau_1) &  & C \vDash \tau_1 \leq
      \tau_2\\
      C, \Gamma \vdash e_2 : \tmop{Num} (\tau_2) &  & C, \Gamma \vdash e_3 :
      \tau
    \end{array}}{C, \Gamma \vdash \text{{\tmstrong{assert num}}} e_1 \leqslant
    e_2 ; e_3 : \tau_{}}}$ &  & $\large{\frac{\begin{array}{lll}
      C, \Gamma \vdash e_1 : \tau_1 &  & C \vDash \tau_1 \dot{=} \tau_2\\
      C, \Gamma \vdash e_2 : \tau_2 &  & C, \Gamma \vdash e_3 : \tau
    \end{array}}{C, \Gamma \vdash \text{{\tmstrong{assert type}}} e_1 = e_2 ;
    e_3 : \tau_{}}}$\\
    &  & \\
    {\tmname{AssertFalse}} &  & {\tmname{RuntimeFailure}}\\
    {\large{$\frac{\begin{array}{l}
      C \vDash \tmmathbf{F}
    \end{array}}{C, \Gamma \vdash \text{{\tmstrong{assert false}}} :
    \tau_{}}$}} &  & {\large{$\frac{C, \Gamma \vdash s : \tmop{String}}{C,
    \Gamma \vdash \text{{\tmstrong{runtime failure}}} s : \tau_{}}$}}
  \end{tabular}
  
  \
  
  \
  
  Non-syntax-directed:
  
  \begin{tabular}{lllllll}
    {\tmname{Gen}} &  & {\tmname{Inst}} &  & {\tmname{Hide}} &  &
    {\tmname{Equ}}\\
    $\large{\frac{\begin{array}{l}
      C \wedge D, \Gamma \vdash e : \beta\\
      \beta \bar{\alpha} \# \tmop{FV} (\Gamma, C)
    \end{array}}{C \wedge \exists \beta \bar{\alpha} .D, \Gamma \vdash e :
    \forall \beta [\exists \bar{\alpha} .D] . \beta}}$ &  &
    $\large{\frac{\begin{array}{l}
      C, \Gamma \vdash e : \forall \bar{\alpha} [D] . \tau^{\prime}\\
      C \vDash D [\bar{\alpha} \assign \bar{\tau}]
    \end{array}}{C, \Gamma \vdash e : \tau^{\prime} [\bar{\alpha} \assign
    \bar{\tau}]}}$ &  & $\large{\frac{\begin{array}{l}
      C, \Gamma \vdash e : \tau\\
      \bar{\alpha} \# \tmop{FV} (\Gamma, \tau)
    \end{array}}{\exists \bar{\alpha} .C, \Gamma \vdash e : \tau}}$ &  &
    $\large{\frac{\begin{array}{l}
      C, \Gamma \vdash e : \tau\\
      C \vDash \tau \dot{=} \tau^{\prime}
    \end{array}}{C, \Gamma \vdash e : \tau^{\prime}}}$\\
    &  &  &  &  &  & \\
    {\tmname{DisjElim}} &  & {\tmname{FElim}} &  &  &  & \\
    $\large{\frac{\begin{array}{lll}
      C, \Gamma \vdash e : \tau &  & D, \Gamma \vdash e : \tau
    \end{array}}{C \vee D, \Gamma \vdash e : \tau}}$ &  &
    $\large{\frac{}{\tmmathbf{F}, \Gamma \vdash e : \tau}}$ &  &  &  & 
  \end{tabular}
  
  \label{TypRulExpr}
  \caption{Typing rules for $\tmop{MMG} (\tmop{num})$: expressions}
\end{table} If alternative sets of constructors are considered, we make them
explicit by writing $C, \Gamma, \Sigma \vdash e : \tau$. By $\mathcal{I}$ we
denote substitutions of predicate variables $\chi$. By interpretations we mean
substitutions leading to ground formulas, which have truth value in the model.
The intended meaning of the type judgment $C, \Gamma, \Sigma \vdash e : \tau$
is: for every interpretation $\mathcal{I}, R$, if $\mathcal{I}, R \vDash C$,
then the expression $e$ has a ground type $R (\tau)$ in a ground environment
$R (\mathcal{I}(\Gamma))$; and with constructors $\mathcal{I} (\Sigma)$ but
this only becomes relevant starting from Section \ref{ExTypes}. We define
derivability of type judgments in Tables \ref{TypRulExpr} and
\ref{TypRulClaus}, where we use pattern-related derivations from Table
\ref{TypRulPat}. For example, the rule {\tmname{Var}}:
$\frac{\begin{array}{ll}
  \Gamma (x) = \forall \beta [\exists \bar{\alpha} .D] . \beta & C \vDash D
\end{array}}{C, \Gamma \vdash x : \beta}$ means that we can derive a type
$\beta$ for $x$, with properties $D$, if the properties $D$ of $\beta$ are
implied by the judgement constraint $C$. And the standard rule {\tmname{App}}:
$\frac{C, \Gamma \vdash e_1 : \tau^{\prime} \rightarrow \tau \hspace{1.8em} C,
\Gamma \vdash e_2 : \tau^{\prime}}{C, \Gamma \vdash e_1 e_2 : \tau}$ means
that a type for an application $e_1 e_2$ can be derived if a function type can
be derived for $e_1$ and the corresponding argument type can be derived for
$e_2$.\begin{table}[t]
  Clauses
  
  \begin{tabular}{ll}
    {\tmname{Clause}} & \\
    $\large{\frac{\begin{array}{ll}
      C \wedge D, \Gamma \Gamma^{\prime} \vdash m_i : \tmop{Num} (\tau_{m_i})
      & e \neq \tmmathbf{\tmop{runtime} \tmop{failure}} \wedge e \neq
      \text{{\tmstrong{assert false}}} \wedge \ldots \wedge\\
      C \wedge D, \Gamma \Gamma^{\prime} \vdash n_i : \tmop{Num} (\tau_{n_i})
      & e \neq \lambda \left( p^{\prime} \ldots \lambda \left(
      p^{\prime\prime} . \text{{\tmstrong{assert false}}} \right) \right)\\
      C \vdash p : \tau_1 \longrightarrow \exists \bar{\beta} [D]
      \Gamma^{\prime} \text{ } \bar{\beta} \# \tmop{FV} (C, \Gamma, \tau_2) &
      C \wedge D \wedge_i \tau_{m_i} \leq \tau_{n_i}, \Gamma \Gamma^{\prime}
      \vdash e : \tau_2
    \end{array}}{C, \Gamma \vdash p \text{{\tmstrong{when}}} \wedge_i m_i
    \leqslant n_i .e : \tau_1 \rightarrow \tau_2}}$ & \\
    & \\
    {\tmname{NegClause}} & \\
    $\large{\frac{\begin{array}{ll}
      C \wedge D, \Gamma \Gamma^{\prime} \vdash m_i : \tmop{Num} (\tau_{m_i})
      & e = \text{{\tmstrong{assert false}}} \vee \ldots \vee\\
      C \wedge D, \Gamma \Gamma^{\prime} \vdash n_i : \tmop{Num} (\tau_{n_i})
      & e = \lambda \left( p^{\prime} \ldots \lambda \left( p^{\prime\prime} .
      \text{{\tmstrong{assert false}}} \right) \right)\\
      C \vdash p : \tau_3 \longrightarrow \exists \bar{\beta} [D]
      \Gamma^{\prime} \text{ } \bar{\beta} \# \tmop{FV} (C, \Gamma, \tau_2) &
      C \wedge D \wedge \tau_1 \dot{=} \tau_3 \wedge_i \tau_{m_i} \leq
      \tau_{n_i}, \Gamma \Gamma^{\prime} \vdash e : \tau_2
    \end{array}}{C, \Gamma \vdash p \text{{\tmstrong{when}}} \wedge_i m_i
    \leqslant n_i .e : \tau_1 \rightarrow \tau_2}}$ & \\
    & \\
    {\tmname{FailClause}} & \\
    $\large{\frac{\begin{array}{ll}
      C \wedge D, \Gamma \Gamma^{\prime} \vdash m_i : \tmop{Num} (\tau_{m_i})
      & C \vdash p : \tau_3 \longrightarrow \exists \bar{\beta} [D]
      \Gamma^{\prime} \text{ } \bar{\beta} \# \tmop{FV} (C, \Gamma)\\
      C \wedge D, \Gamma \Gamma^{\prime} \vdash n_i : \tmop{Num} (\tau_{n_i})
      & C \wedge D \wedge \tau_1 \dot{=} \tau_3 \wedge_i \tau_{m_i} \leq
      \tau_{n_i}, \Gamma \Gamma^{\prime} \vdash s : \tmop{String}
    \end{array}}{C, \Gamma \vdash p \text{{\tmstrong{when}}} \wedge_i m_i
    \leqslant n_i . \text{{\tmstrong{runtime failure}}} s : \tau_1 \rightarrow
    \tau_2}}$ & 
  \end{tabular}
  
  \label{TypRulClaus}
  \caption{Typing rules for $\tmop{MMG} (\tmop{num})$: clauses}
\end{table}

Top-level definitions introduce new entries into a global value constructor
environment $\Sigma$ or a global variable environment $\Gamma$. Here we
present the most interesting case of a recursive definition with a test
clause. The type information flow from the test $e_2$ to the definition $e_1$
is through the shared type scheme $\sigma$ in $\Gamma^{\prime}$. Other cases
are analogous. In particular, toplevel \tmverbatim{let} definitions, unlike
\tmverbatim{let}$\ldots$\tmverbatim{in} expressions, are polymorphic, and,
like \tmverbatim{let}$\ldots$\tmverbatim{in} expressions, allow binding to a
pattern.
\[ \frac{\begin{array}{lll}
     C, \Gamma^{\prime} \vdash e_1 : \sigma & C, \Gamma^{\prime} \vdash e_2 :
     \tmop{Bool} & \\
     \sigma = \forall \beta [\exists \bar{\alpha} .D] . \beta &
     \Gamma^{\prime} = \Gamma \{x \mapsto \sigma\} & \mathcal{M} \vDash C
   \end{array}}{C, \Gamma \vdash \text{{\tmstrong{let rec}}} x = e_1 
   \text{{\tmstrong{test}}} e_2 \longrightarrow \Gamma^{\prime}} \]
A {\tmem{data constructor}} $K$ for a {\tmem{datatype}} $\varepsilon$ (recall
that the sort $s_{\tmop{type}}$ holds two categories of elements: datatypes
and function types) has definition $K \colons \forall \bar{\alpha} \bar{\beta}
[D] . \tau_1 \times \ldots \times \tau_n \rightarrow \varepsilon
(\bar{\alpha})$ where $\tmop{FV} (D, \tau_1, \ldots, \tau_n) \subseteq
\bar{\alpha} \bar{\beta}$. $D$ is a conjunction of atoms.

At this point the construction {\tmname{LetIn}} is a syntactic sugar for
single branch patterns -- if polymorphic \tmverbatim{let} is needed, use
{\tmname{LetRec}}. We identify clauses $p.e$ of $\tmop{MMG} (X)$ with $p
\text{{\tmstrong{when}}} .e$. Note that {\tmname{Gen}} and {\tmname{DisjElim}}
are unrelated to Constraint Generalization we introduce in the next chapter.
Most of the rules in Tables \ref{TypRulPat} and \ref{TypRulExpr} are close to
$\tmop{HMG} (X)$ rules well described in Simonet and Pottier
{\cite{simonet-pottier-hmg-toplas}}. A salient difference is the rule
{\tmname{NegClause}} for clauses whose right-hand-side is an
\tmverbatim{assert false} expression. Rather than unifying the type of the
function argument with the type of the pattern, we want the discrepancy
between these types be a possible reason of unreachability of the pattern
matching clause. Without the distinct treatment of clauses with
\tmverbatim{assert false} as right-hand-side, the \tmverbatim{assert}-based
variant of the \tmverbatim{equal} example would fail. The similar rule
{\tmname{FailClause}} is an ``escape hatch'' for cases which cannot be ruled
out by the type system.

An expression $e$ is {\tmem{well typed}} given $\Gamma, \Sigma$ when
$\tmop{PV} (\Gamma, \Sigma) = \varnothing$ and $C, \Gamma, \Sigma \vdash e :
{\nobreak} \sigma$ holds for some satisfiable constraint $C$. For simplicity,
{\tmname{InvarGenT}} only admits type and invariant annotations from the user
on toplevel definitions.

\section{Type Inference Constraints}\label{TypInfCns}

In Tables \ref{IntermInfPat}, \ref{IntermInf} and \ref{IntermInfClaus}, we
present type judgments declaratively by reducing them to constraints. The
presentation is a little bit heavy due to explicit capture-avoidance
conditions. For example, $\llbracket \Gamma \vdash x : \tau \rrbracket =
\exists \beta^{\prime} \bar{\alpha}^{\prime} .D [\beta \bar{\alpha} \assign
\beta^{\prime} \bar{\alpha}^{\prime}] \wedge \beta^{\prime} \dot{=} \tau
\text{ where } \Gamma (x) = \forall \beta [\exists \bar{\alpha} .D] . \beta,
\beta^{\prime} \bar{\alpha}^{\prime} \# \tmop{FV} (\Gamma, {\nobreak} \tau)$
introduces a constraint over the type $\tau$ required by the properties of $x$
stored in the environment $\Gamma$. The constraint $\llbracket \Gamma \vdash
e_1 e_2 : \tau \rrbracket = \exists \alpha . \llbracket \Gamma \vdash e_1 :
\alpha \rightarrow \tau \rrbracket \wedge \llbracket \Gamma \vdash e_2 :
\alpha \rrbracket, \text{ \ } \alpha \# \tmop{FV} (\Gamma, \tau)$ for typing
the result of application of $e_1$ to $e_2$, imposes a function type on $e_1$,
and a type on $e_2$ -- represented by the fresh variable $\alpha$ -- that
matches the argument type of $e_1$. The constraint for a pattern matching
clause without the \tmverbatim{when} guard becomes more readable: $\llbracket
\Gamma \vdash p.e : \tau_1 \rightarrow \tau_2 \rrbracket = \llbracket \vdash p
\downarrow \tau_1 \rrbracket \wedge \forall \bar{\beta} .D \Rightarrow
\llbracket \Gamma \Gamma^{\prime} \vdash e : \tau_2 \rrbracket$ where $\exists
\bar{\beta} [D] \Gamma^{\prime} = \llbracket \vdash p \uparrow \alpha_3
\rrbracket$. The premise $D$ derived from the pattern $p$ provides
context-specific information for typing the expression $e$.

The two presentations are equivalent, in the sense of the
{\tmem{correctness}} and {\tmem{completeness}} theorems. The proofs are in
Section \ref{GADTProofs}.

\begin{theorem}
  \label{InfCorrExpr}{\tmem{Correctness}} (expressions). For all environments
  $\Gamma$, expressions $e$ and types $\tau$, $\llbracket \Gamma \vdash e :
  \tau \rrbracket, \Gamma \vdash e : \tau$ is derivable.
\end{theorem}

\begin{theorem}
  \label{InfComplExpr}{\tmem{Completeness}} (expressions). Let $\mathcal{L}_1$
  be a first order language with equality $\dot{=}$ and a model $\mathcal{M}$.
  Let $\mathcal{L}$ be an extension of $\mathcal{L}_1$ with predicate
  variables $\chi (\cdummy)$. Let $\Gamma$ be an environment, $C$ a formula in
  $\mathcal{L}$, $e$ an expression and $\tau$ a type. If $\tmop{PV} (C,
  \Gamma) = \varnothing$ and $C, \Gamma \vdash e : \tau$ is derivable, then
  there exists an interpretation of predicate variables $\mathcal{I}$ such
  that $\mathcal{M}, \mathcal{I} \vDash C \Rightarrow \left\llbracket \Gamma
  \vdash e : {\nobreak} \tau \right\rrbracket$.
\end{theorem}

In Theorem \ref{InfComplExpr}, we explicitly introduce all symbols used. We
often state propositions and theorems in a more compact manner, with symbols
introduced implicitly.

\begin{corollary}
  \label{InfComplCoro}If $C, \Gamma \vdash e : \forall \bar{\alpha} [D] .
  \tau$ and $\bar{\alpha} \# \tmop{FV} (\Gamma)$, then there is an
  interpretation $\mathcal{I}$ such that $\mathcal{M}, \mathcal{I} \vDash C
  \Rightarrow (\forall \bar{\alpha} .D \Rightarrow \llbracket \Gamma \vdash e
  : \tau \rrbracket)$.
\end{corollary}

\begin{table}[t]
  Constraint generation:
  \begin{eqnarray*}
    \llbracket \vdash 0 \downarrow \tau \rrbracket \nobracket \nobracket & = &
    \tmmathbf{T}\\
    &  & \\
    \llbracket \vdash 1 \downarrow \tau \rrbracket \nobracket \nobracket & = &
    \tmmathbf{T}\\
    &  & \\
    \llbracket \vdash x \downarrow \tau \rrbracket \nobracket \nobracket & = &
    \tmmathbf{T}\\
    &  & \\
    \llbracket \vdash p_1 \wedge p_2 \downarrow \tau \rrbracket \nobracket
    \nobracket & = & \llbracket \vdash p_1 \downarrow \tau \rrbracket \wedge
    \llbracket \vdash p_2 \downarrow \tau \rrbracket \nobracket \nobracket
    \nobracket \nobracket\\
    &  & \\
    \llbracket \vdash K p_1 \ldots p_n \downarrow \tau \rrbracket \nobracket
    \nobracket & = & \exists \bar{\alpha}^{\prime} . \varepsilon
    (\bar{\alpha}^{\prime}) \dot{=} \tau \wedge \forall \bar{\beta}^{\prime}
    .D [\bar{\alpha} \bar{\beta} \assign \bar{\alpha}^{\prime}
    \bar{\beta}^{\prime}] \Rightarrow \wedge_i \llbracket p_i \downarrow
    \tau_i [\bar{\alpha} \bar{\beta} \assign \bar{\alpha}^{\prime}
    \bar{\beta}^{\prime}] \rrbracket \nobracket \nobracket\\
    &  & \\
    &  & \text{where } K \colons \forall \bar{\alpha} \bar{\beta} [D] .
    \tau_1 \times \ldots \times \tau_n \rightarrow \varepsilon (\bar{\alpha}),
    \bar{\alpha}^{\prime} \bar{\beta}^{\prime} \# \tmop{FV} (\Sigma, \tau)
  \end{eqnarray*}
  Environment fragment generation:
  \begin{eqnarray*}
    \llbracket \vdash 0 \uparrow \tau \rrbracket \nobracket \nobracket & = &
    \exists \varnothing [\tmmathbf{F}] \{\}\\
    &  & \\
    \llbracket \vdash 1 \uparrow \tau \rrbracket \nobracket \nobracket & = &
    \exists \varnothing [\tmmathbf{T}] \{\}\\
    &  & \\
    \llbracket \vdash x \uparrow \tau \rrbracket \nobracket \nobracket & = &
    \exists \varnothing [\tmmathbf{T}] \{x \mapsto \tau\}\\
    &  & \\
    \llbracket \vdash p_1 \wedge p_2 \uparrow \tau \rrbracket \nobracket
    \nobracket & = & \llbracket \vdash p_1 \uparrow \tau \rrbracket \times
    \llbracket \vdash p_2 \uparrow \tau \rrbracket \nobracket \nobracket
    \nobracket \nobracket\\
    &  & \\
    \llbracket \vdash K p_1 \ldots p_n \uparrow \tau \rrbracket \nobracket
    \nobracket & = & \exists \bar{\alpha}^{\prime} \bar{\beta}^{\prime}
    [\varepsilon (\bar{\alpha}^{\prime}) \dot{=} \tau \wedge D [\bar{\alpha}
    \bar{\beta} \assign \bar{\alpha}^{\prime} \bar{\beta}^{\prime}]] (\times_i
    \llbracket p_i \uparrow \tau_i [\bar{\alpha} \bar{\beta} \assign
    \bar{\alpha}^{\prime} \bar{\beta}^{\prime}] \rrbracket) \nobracket\\
    &  & \\
    &  & \text{where } K \colons \forall \bar{\alpha} \bar{\beta} [D] .
    \tau_1 \times \ldots \times \tau_n \rightarrow \varepsilon (\bar{\alpha}),
    \bar{\alpha}^{\prime} \bar{\beta}^{\prime} \# \tmop{FV} (\Sigma, \tau)
  \end{eqnarray*}
  \label{IntermInfPat}
  \caption{Type inference for patterns}
\end{table}\begin{table}[t]
  \begin{eqnarray*}
    \llbracket \Gamma \vdash x : \tau \rrbracket & = & \tmmathbf{F} \text{ \
    when \ } x \nin \tmop{Dom} (\Gamma)\\
    &  & \\
    \llbracket \Gamma \vdash x : \tau \rrbracket & = & \exists \beta^{\prime}
    \bar{\alpha}^{\prime} .D [\beta \bar{\alpha} \assign \beta^{\prime}
    \bar{\alpha}^{\prime}] \wedge \beta^{\prime} \dot{=} \tau \text{ \ where }
    \Gamma (x) = \forall \beta [\exists \bar{\alpha} .D] . \beta,
    \beta^{\prime} \bar{\alpha}^{\prime} \# \tmop{FV} (\Gamma, \tau)\\
    &  & \\
    \left\llbracket \Gamma \vdash \text{{\tmstrong{assert false}}} : \tau
    \right\rrbracket & = & \tmmathbf{F}\\
    &  & \\
    \left\llbracket \Gamma \vdash \text{{\tmstrong{runtime failure}}} s : \tau
    \right\rrbracket & = & \llbracket \Gamma \vdash s : \tmop{String}
    \rrbracket\\
    &  & \\
    \left\llbracket \Gamma \vdash \text{{\tmstrong{assert num}}} e_1 \leqslant
    e_2 ; e_3 : \tau \right\rrbracket & = & \exists \alpha^1 \alpha^2 .
    \llbracket \Gamma \vdash e_1 : \tmop{Num} (\alpha^1) \rrbracket \wedge
    \llbracket \Gamma \vdash e_2 : \tmop{Num} (\alpha^2) \rrbracket \wedge
    \alpha^1 \leq \alpha^2 \wedge \left\llbracket \Gamma \vdash e_3 :
    {\nobreak} \tau \right\rrbracket, \alpha_1 \alpha_2 \# \tmop{FV} (\Gamma,
    {\nobreak} \tau)\\
    &  & \\
    \left\llbracket \Gamma \vdash \text{{\tmstrong{assert type}}} e_1 = e_2 ;
    e_3 : \tau \right\rrbracket & = & \exists \alpha^1 \alpha^2 . \llbracket
    \Gamma \vdash e_1 : \alpha^1 \rrbracket \wedge \llbracket \Gamma \vdash
    e_2 : \alpha^2 \rrbracket \wedge \alpha^1 \dot{=} \alpha^2 \wedge
    \llbracket \Gamma \vdash e_3 : \tau \rrbracket, \alpha_1 \alpha_2 \#
    \tmop{FV} (\Gamma, \tau)\\
    &  & \\
    \llbracket \Gamma \vdash \lambda \bar{c} : \tau \rrbracket & = & \exists
    \alpha_1 \alpha_2 . \llbracket \Gamma \vdash \bar{c} : \alpha_1
    \rightarrow \alpha_2 \rrbracket \wedge \alpha_1 \rightarrow \alpha_2
    \dot{=} \tau, \text{} \alpha_1 \alpha_2 \# \tmop{FV} (\Gamma, \tau)\\
    &  & \\
    \llbracket \Gamma \vdash e_1 e_2 : \tau \rrbracket & = & \exists \alpha .
    \llbracket \Gamma \vdash e_1 : \alpha \rightarrow \tau \rrbracket \wedge
    \llbracket \Gamma \vdash e_2 : \alpha \rrbracket, \text{ \ } \alpha \#
    \tmop{FV} (\Gamma, \tau)\\
    &  & \\
    \llbracket \Gamma \vdash K e_1 \ldots e_n : \tau \rrbracket & = & \exists
    \bar{\alpha}^{\prime} \bar{\beta}^{\prime} . (\wedge_i \llbracket \Gamma
    \vdash e_i : \tau_i [\bar{\alpha} \bar{\beta} \assign
    \bar{\alpha}^{\prime} \bar{\beta}^{\prime}] \rrbracket \wedge D
    [\bar{\alpha} \bar{\beta} \assign \bar{\alpha}^{\prime}
    \bar{\beta}^{\prime}] \wedge \varepsilon (\bar{\alpha}^{\prime}) \dot{=}
    \tau)\\
    &  & \text{where } \Sigma \ni K \colons \forall \bar{\alpha} \bar{\beta}
    [D] . \tau_1 \times \ldots \times \tau_n \rightarrow \varepsilon
    (\bar{\alpha}), \text{ \ } \bar{\alpha}^{\prime} \bar{\beta}^{\prime} \#
    \tmop{FV} (\Gamma, \tau)\\
    &  & \\
    \left\llbracket \Gamma \vdash \text{{\tmstrong{let rec}}} x = e_1
    \text{{\tmstrong{in}}} e_2 : \tau \right\rrbracket & = & (\forall \beta
    (\chi (\beta) \Rightarrow \llbracket \Gamma \{x \mapsto \forall \beta
    [\chi (\beta)] . \beta\} \vdash e_1 : \beta \rrbracket)) \wedge \\
    (\exists \alpha . \chi (\alpha)) \wedge \llbracket \Gamma \{x \mapsto
    \forall \beta [\chi (\beta)] . \beta\} \vdash e_2 : \tau \rrbracket\\
    &  & \text{where } \beta \# \tmop{FV} (\Gamma, \tau), \chi \# \tmop{PV}
    (\Gamma)\\
    &  & \\
    \llbracket \Gamma \vdash \overline{c_i} : \tau_1 \rightarrow \tau_2
    \rrbracket & = & \wedge_i \llbracket \Gamma \vdash c_i : \tau_1
    \rightarrow \tau_2 \rrbracket\\
    &  & \\
    \llbracket \Gamma \vdash e : \forall \bar{\alpha} [D] . \tau \rrbracket &
    = & \forall \bar{\alpha}^{\prime} .D [\bar{\alpha} \assign
    \bar{\alpha}^{\prime}] \Rightarrow \llbracket \Gamma \vdash e : \tau
    [\bar{\alpha} \assign \bar{\alpha}^{\prime}] \rrbracket, \text{ \ }
    \bar{\alpha}^{\prime} \# \tmop{FV} (\Gamma)
  \end{eqnarray*}
  \label{IntermInf}
  \caption{Type inference for expressions}
\end{table}\begin{table}[t]
  \begin{eqnarray*}
    \left\llbracket \Gamma \vdash p \text{{\tmstrong{when}}} \wedge_i m_i
    \leqslant n_i .e : \tau_1 \rightarrow \tau_2 \right\rrbracket & = &
    \llbracket \vdash p \downarrow \tau_1 \rrbracket \wedge \forall
    \bar{\beta} . \exists \overline{\alpha^1_i \alpha^2_i} .D \Rightarrow \\
    \text{ \ \ \ \ \ } \wedge_i \llbracket \Gamma \Gamma^{\prime} \vdash m_i
    : \tmop{Num} (\alpha^1_i) \rrbracket \wedge_i \llbracket \Gamma
    \Gamma^{\prime} \vdash n_i : \tmop{Num} (\alpha^2_i) \rrbracket\\
    \text{when } e \neq \tmmathbf{\tmop{runtime} \tmop{failure}} \wedge &  &
    \text{ \ \ \ \ \ } \wedge (\wedge_i \alpha^1_i \leq \alpha^2_i \Rightarrow
    \llbracket \Gamma \Gamma^{\prime} \vdash e : \tau_2 \rrbracket)\\
    e \neq \text{{\tmstrong{assert false}}} \wedge \ldots \wedge &  & \\
    e \neq \lambda \left( p^{\prime} \ldots \lambda \left( p^{\prime\prime} .
    \text{{\tmstrong{assert false}}} \right) \right) &  & \\
    &  & \text{where } \exists \bar{\beta} [D] \Gamma^{\prime} \text{ is }
    \llbracket \vdash p \uparrow \tau_1 \rrbracket, \bar{\beta}
    \overline{\alpha^1_i \alpha^2_i} \# \tmop{FV} (\Gamma, \tau_1, \tau_2)\\
    &  & \\
    \left\llbracket \Gamma \vdash p \text{{\tmstrong{when}}} \wedge_i m_i
    \leqslant n_i .e : \tau_1 \rightarrow \tau_2 \right\rrbracket & = &
    \exists \alpha_3 . \llbracket \vdash p \downarrow \alpha_3 \rrbracket
    \wedge \forall \bar{\beta} . \exists \overline{\alpha^1_i \alpha^2_i} .D
    \Rightarrow \\
    \text{ \ \ \ \ \ } \wedge_i \llbracket \Gamma \Gamma^{\prime} \vdash m_i
    : \tmop{Num} (\alpha^1_i) \rrbracket \wedge_i \llbracket \Gamma
    \Gamma^{\prime} \vdash n_i : \tmop{Num} (\alpha^2_i) \rrbracket\\
    \text{when } e = \text{{\tmstrong{assert false}}} \vee \ldots \vee &  &
    \text{ \ \ \ \ \ } \wedge (\alpha_3 \dot{=} \tau_1 \wedge_i \alpha^1_i
    \leq \alpha^2_i \Rightarrow \llbracket \Gamma \Gamma^{\prime} \vdash e :
    \tau_2 \rrbracket)\\
    e = \lambda \left( p^{\prime} \ldots \lambda \left( p^{\prime\prime} .
    \text{{\tmstrong{assert false}}} \right) \right) &  & \\
    &  & \text{where } \exists \bar{\beta} [D] \Gamma^{\prime} \text{ is }
    \llbracket \vdash p \uparrow \alpha_3 \rrbracket, \bar{\beta} \alpha_3
    \overline{\alpha^1_i \alpha^2_i} \# \tmop{FV} (\Gamma, \tau_1, \tau_2)\\
    &  & \\
    \left\llbracket \Gamma \vdash p \text{{\tmstrong{when}}} \wedge_i m_i
    \leqslant n_i .e : \tau_1 \rightarrow \tau_2 \right\rrbracket & = &
    \exists \alpha_3 . \llbracket \vdash p \downarrow \alpha_3 \rrbracket
    \wedge \forall \bar{\beta} . \exists \overline{\alpha^1_i \alpha^2_i} .D
    \Rightarrow \\
    \text{ \ \ \ \ \ } \wedge_i \llbracket \Gamma \Gamma^{\prime} \vdash m_i
    : \tmop{Num} (\alpha^1_i) \rrbracket \wedge_i \llbracket \Gamma
    \Gamma^{\prime} \vdash n_i : \tmop{Num} (\alpha^2_i) \rrbracket\\
    \text{when } e = \text{{\tmstrong{runtime failure}}} s &  & \text{ \ \ \ \
    \ } \wedge (\alpha_3 \dot{=} \tau_1 \wedge_i \alpha^1_i \leq \alpha^2_i
    \Rightarrow \llbracket \Gamma \Gamma^{\prime} \vdash s : \tmop{String}
    \rrbracket)\\
    &  & \text{where } \exists \bar{\beta} [D] \Gamma^{\prime} \text{ is }
    \llbracket \vdash p \uparrow \alpha_3 \rrbracket, \bar{\beta} \alpha_3
    \overline{\alpha^1_i \alpha^2_i} \# \tmop{FV} (\Gamma, \tau_1, \tau_2)
  \end{eqnarray*}
  \label{IntermInfClaus}
  \caption{Type inference for clauses}
\end{table}\begin{table}[t]
  $\begin{array}{lll}
    l (x) & = & \tmmathbf{F}\\
    l (\lambda \bar{c}) & = & \tmmathbf{F}\\
    l (e_1 e_2) & = & l (e_1)\\
    l (K e_1 \ldots e_n) & = & \tmmathbf{F}\\
    l \left( \text{{\tmstrong{let rec}}} x = e_1  \text{{\tmstrong{in}}} e_2
    \right) & = & l (e_2)\\
    l (\lambda [K] \overline{p_i .e_i}) & = & \tmmathbf{T}\\
    l \left( \text{{\tmstrong{let}}} p = e_1  \text{{\tmstrong{in}}} e_2
    \right) & = & \tmmathbf{T}\\
    l \left( \text{{\tmstrong{assert num}}} e_1 \leqslant e_2 ; e_3 \right) &
    = & l (e_3)\\
    l \left( \text{{\tmstrong{assert type}}} e_1 = e_2 ; e_3 \right) & = & l
    (e_3)\\
    l \left( \text{{\tmstrong{runtime failure}}} s \right) & = &
    \tmmathbf{T}\\
    l \left( \text{{\tmstrong{assert false}}} \right) & = & \tmmathbf{T}
  \end{array}$\label{IsExBrs}
  \caption{Expressions that directly handle an existential type}
\end{table}\begin{table}[t]
  $\begin{array}{lll}
    \mathcal{E} (x, v) & = & \varnothing\\
    \mathcal{E} (\lambda \bar{c}, v) & = & \cup \overline{\mathcal{E} (c,
    \tmmathbf{F})}\\
    \mathcal{E} (e_1 e_2, v) & = & \mathcal{E} (e_1, v) \cup \mathcal{E} (e_2,
    \tmmathbf{F})\\
    \mathcal{E} (K e_1 \ldots e_n, v) & = & \cup_i \mathcal{E} (e_i,
    \tmmathbf{F})\\
    \mathcal{E} \left( \text{{\tmstrong{let rec}}} x = e_1 
    \text{{\tmstrong{in}}} e_2, v \right) & = & \mathcal{E} (e_1,
    \tmmathbf{F}) \cup \mathcal{E} (e_2, v)\\
    \mathcal{E} \left( p \text{{\tmstrong{when}}}  \overline{m_i \leqslant
    n_i} .e, v \right) & = & \cup_i \mathcal{E} (m_i, \tmmathbf{F}) \cup_i
    \mathcal{E} (n_i, \tmmathbf{F})\\
    &  & \cup \mathcal{E} (e, v)\\
    \mathcal{E} (\lambda [K] \bar{c}, \tmmathbf{F}) & = & \{ K \} \cup
    \overline{\mathcal{E} (c, \tmmathbf{T})}\\
    \mathcal{E} (\lambda [K] \bar{c}, \tmmathbf{T}) & = & \cup
    \overline{\mathcal{E} (c, \tmmathbf{T})}\\
    \mathcal{E} \left( \text{{\tmstrong{let}}} p = e_1  \text{{\tmstrong{in}}}
    e_2, v \right) & = & \mathcal{E} (e_1, \tmmathbf{F}) \cup \mathcal{E}
    (e_2, v)\\
    \mathcal{E} \left( \text{{\tmstrong{assert num}}} e_1 \leqslant e_2 ; e_3,
    K^{\prime} \right) & = & (e_1, \tmmathbf{F}) \cup (e_2, \tmmathbf{F}) \cup
    \mathcal{E} (e_3, v)\\
    \mathcal{E} \left( \text{{\tmstrong{assert type}}} e_1 = e_2 ; e_3,
    K^{\prime} \right) & = & (e_1, \tmmathbf{F}) \cup (e_2, \tmmathbf{F}) \cup
    \mathcal{E} (e_3, v)
  \end{array}$\label{CollectK}
  \caption{Collect introduced value constructors}
\end{table}

In Definition \ref{SolTypInfer}, we describe a class of type judgments that is
a better target for type inference. First, we introduce a normal form of
constraints that simplifies formal considerations.

\begin{proposition}
  \label{NormalFormProp}For any $\Gamma, e, \tau$, the formula $\llbracket
  \Gamma \vdash e : \tau \rrbracket$ is equivalent to a prenex-normal,
  normalized form: there is a quantifier prefix $\mathcal{Q}$ and pairs of
  conjunctions of atoms $D_i, C_i$, such that
  \[ \mathcal{M} \vDash \llbracket \Gamma \vdash e : \tau \rrbracket
     \Leftrightarrow \mathcal{Q}. \wedge_i (D_i \Rightarrow C_i) \]
\end{proposition}

We put $\tmop{NF} (\llbracket \Gamma \vdash e : \tau \rrbracket) =\mathcal{Q}.
\wedge_i (D_i \Rightarrow C_i)$.

\begin{definition}
  \label{SolTypInfer}We call the formula $\llbracket \Gamma \vdash e : \tau
  \rrbracket$ the {\tmem{type inference problem}} for environment $\Gamma$,
  expression $e$ and type $\tau$. Let $\mathcal{M} \vDash \left\llbracket
  \Gamma \vdash e : {\nobreak} \tau \right\rrbracket \Leftrightarrow
  \mathcal{Q}. \Phi_N$ with $\Phi_N = \wedge_i (D_i \Rightarrow C_i)$ as in
  Proposition \ref{NormalFormProp}. Let $F_{\tmop{res}}$ be a conjunction of
  atoms without predicate variables. We call $\exists
  \bar{\alpha}_{\tmop{res}} .F_{\tmop{res}}, \mathcal{I}$ a {\tmem{solution to
  the type inference problem}} $\left\llbracket \Gamma \vdash e : {\nobreak}
  \tau \right\rrbracket$ when $\mathcal{I}, \mathcal{M} \vDash F_{\tmop{res}}
  \Rightarrow \Phi_N$ (i.e. $\mathcal{M} \vDash F_{\tmop{res}} \Rightarrow
  \mathcal{I} (\Phi_N)$), $\mathcal{M} \vDash \mathcal{Q}.F_{\tmop{res}}
  [\bar{\alpha}_{\tmop{res}} \assign \bar{t}]$ for some $\bar{t}$, and for
  every implication in $\Phi_N$, if $\mathcal{M}, \mathcal{I} \vDash \exists
  \tmop{FV} (D_i) .D_i$ then $\mathcal{M}, \mathcal{I} \vDash \exists
  \tmop{FV} (D_i, F_{\tmop{res}}) .D_i \wedge F_{\tmop{res}}$.
\end{definition}

We are more interested in finding an unambiguous, directly given solution to a
type inference problem, as introduced in Definition \ref{SolTypInfer}, than in
derivability of a judgement $C, \Gamma \vdash e : \tau$ for some satisfiable
$C$. By Theorems \ref{InfCorrExpr} and \ref{InfComplExpr}, these notions are
close. Therefore, unlike in $\tmop{HMG} (X)$, we say that an expression $e$ is
{\tmem{well-typed}} in $\tmop{MMG} (X)$ under environment $\Gamma$ with type
$\tau$, when the type inference problem $\llbracket \Gamma \vdash e : \tau
\rrbracket$ has a solution. The class of programs admitting a solution to the
type inference problem can sometimes be expanded by enriching the languages of
constraint domains.

Note that both the derivability of $C, \Gamma \vdash e : \tau$ and the
existence of a solution to the type inference problem $\left\llbracket \Gamma
\vdash e : {\nobreak} \tau \right\rrbracket$ allows for dead code in $e$, i.e.
unreachable cases of pattern matching. We offer an option to reject programs
with dead code, according to the following definition: We call $\exists
\bar{\alpha}_{\tmop{res}} .F_{\tmop{res}}, \mathcal{I}$ a {\tmem{solution
without dead code to the type inference problem}} $\left\llbracket \Gamma
\vdash e : {\nobreak} \tau \right\rrbracket$ when $\mathcal{I}, \mathcal{M}
\vDash F_{\tmop{res}} \Rightarrow \Phi_N$ (i.e. $\mathcal{M} \vDash
F_{\tmop{res}} \Rightarrow \mathcal{I} (\Phi_N)$), $\mathcal{M} \vDash
\mathcal{Q}.F_{\tmop{res}} [\bar{\alpha}_{\tmop{res}} \assign \bar{t}]$ for
some $\bar{t}$, and for every implication in $\Phi_N$, if $C_i \neq
\tmmathbf{F}$ then $\mathcal{M}, \mathcal{I} \vDash \exists \tmop{FV} \left(
D_i, {\nobreak} F_{\tmop{res}} \right) .D_i \wedge {\nobreak} F_{\tmop{res}}$.
However, some datatype-related dead code cases are rejected regardless of this
option, due to constraints introduced by $\llbracket \vdash p \downarrow \tau
\rrbracket \nobracket \nobracket$ for patterns $p$.

\section{Existential Types}\label{ExTypes}

In context of GADTs, existential types play a prominent role, beyond the
traditional role of abstraction in software engineering. Without existential
types, computations would need to express parameters of the output datatype
invariant as a function of parameters of the input datatype invariant. Since
GADTs are introduced to curtail the expressiveness of types compared to full
dependent type systems, opportunities for such functional dependency are rare
by design.

We need the capacity in the type system to express whatever relations it can
of the resulting datatype parameters to the input datatype parameters.
Traditionally in GADTs we package the result into a custom datatype. This is
tedious and contrary to the benefits of type inference. We automate this
process, in effect introducing inferred existential types to our type system.
Since the modification of the type system is minimal, formal guarantees carry
over to it and it will be familiar to users of GADTs.

Existential quantifiers in (contravariant) argument positions of function
types are redundant: they can be lifted to be traditional, polymorphic
variables constrained by the invariant of the function. However, they may
sometimes be needed in nested positions of higher-order function types. We
prohibit the use of inferred existential types in argument positions in the
current version of the type system.

We introduce a new expression construct $\lambda [K] \bar{c}$, where $K$ is a
value constructor, but is not available in concrete syntax, and $\bar{c}$ are
pattern matching clauses. In the implementation, the parser introduces a fresh
$K$ and forms $\lambda [K] \bar{c}$ for \tmverbatim{efunction $\bar{c}$}. We
also introduce a corresponding construct $\lambda_K \bar{c}$. $\lambda [K]
\bar{c}$ is either eliminated or replaced by $\lambda_K \bar{c}$ in a
normalization step. When $K \colons \forall \bar{\alpha} \bar{\beta} \gamma
[E] . \gamma \rightarrow \varepsilon_K (\bar{\alpha}) \in \Sigma$ is such a
data constructor absent from concrete syntax, the pretty-printer for types
prints $\varepsilon_K (\bar{\tau})$ as $\exists \bar{\beta} \gamma [E
[\bar{\alpha} \assign \bar{\tau}]] . \gamma$, or $\exists \bar{\beta} [E
[\bar{\alpha} \assign \bar{\tau}]] . \tau_e$ when $E$ implies $\gamma \dot{=}
\tau_e$. We also parse $\exists \bar{\beta} [E] . \tau_e$ generating a fresh
$K \colons \forall \bar{\alpha} \bar{\beta} [E] . \tau_e \rightarrow
\varepsilon_K (\bar{\alpha}) \in \Sigma$ in the toplevel $\Sigma$, where
$\bar{\alpha} = {\nobreak} \tmop{FV} (\tau_e) \backslash \bar{\beta}$.

Let $l (e)$ defined in Table \ref{IsExBrs} determine whether an expression
introduces or eliminates an existential type. It is used in the normalization
process described below.

Let all occurrences of $\lambda [K]$ in $e$ use distinct $K$. Let $n (e)
\assign n (e, \bot)$, defined in Table \ref{NormExpr}, flatten nested
introductions of existential types. Let $\mathcal{E} (e) \assign \mathcal{E}
(e, \tmmathbf{F})$, defined in Table \ref{CollectK}, collect value
constructors introduced for existential types. The normalization-related
functions $n (\cdummy, v)$ and $\mathcal{E} (\cdummy, v)$ are defined for both
expressions and clauses.

The $\tmop{MMG}_{\exists} (X)$ typing rules that differ from $\tmop{MMG} (X)$
are provided in Table \ref{ExTypingRul}. We put the normalization step into
the type system as a rule {\tmname{ExIntro}}. W.l.o.g. {\tmname{ExIntro}} can
be used once at the beginning of a derivation. {\tmname{ExIntro}} performs a
normalization of the expression and shares the job of introducing existential
types with the rule {\tmname{ExAbs}}. {\tmname{ExLetIn}} is the elimination
rule for existential types. Although {\tmname{LetIn}} and {\tmname{ExLetIn}}
resemble ``syntactic sugar'', their application is non-deterministic. We
include value constructor environment in judgments to facilitate the
completeness proof. We modify the rule {\tmname{App}} to exclude function
arguments that have existential types. We achieve that by introducing a new
atomic predicate $\not{E}$ to the sort $s_{\tmop{type}}$: $\not{E} (\tau)$ iff
$\wedge_K \neg {\nobreak} \exists \bar{\alpha} . \tau \dot{=} \varepsilon_K
(\bar{\alpha})$, i.e. $\tau$ is not an existential type.\begin{table}[t]
  \begin{eqnarray*}
    n (e, K^{\prime}) & = & \text{{\tmstrong{let}}} x = n (e, \bot) 
    \text{{\tmstrong{in}}} K^{\prime} x\\
    \text{when } K^{\prime} \neq \bot \wedge l (e) =\tmmathbf{F} &  & \\
    n (x, \bot) & = & x\\
    n (\lambda \bar{c}, \bot) & = & \lambda (\overline{n (c, \bot)})\\
    n (e_1 e_2, K^{\prime}) & = & n (e_1, K^{\prime}) n (e_2, \bot)\\
    n (K e_1 \ldots e_n, \bot) & = & K n (e_1, \bot) \ldots n (e_n, \bot)\\
    n \left( \text{{\tmstrong{let rec}}} x = e_1  \text{{\tmstrong{in}}} e_2,
    K^{\prime} \right) & = & \text{{\tmstrong{let rec}}} x = n (e_1, \bot) 
    \text{{\tmstrong{in}}} n (e_2, K^{\prime})\\
    n \left( p \text{{\tmstrong{when}}}  \overline{m_i \leqslant n_i} .e,
    K^{\prime} \right) & = & p \text{{\tmstrong{when}}}  \overline{n (m_i,
    \bot) \leqslant n (n_i, \bot)} . \text{{\tmstrong{runtime failure}}} n (s,
    \bot)\\
    \text{when } e = \text{{\tmstrong{runtime failure}}} s &  & \\
    n \left( p \text{{\tmstrong{when}}}  \overline{m_i \leqslant n_i} .e,
    K^{\prime} \right) & = & p \text{{\tmstrong{when}}}  \overline{n (m_i,
    \bot) \leqslant n (n_i, \bot)} .n (e, K^{\prime})\\
    \text{when } e \neq \text{{\tmstrong{runtime failure}}} s &  & \\
    n (\lambda [K] \bar{c}, \bot) & = & \lambda_K (\overline{n (c, K)})\\
    n (\lambda [K] \bar{c}, K^{\prime}) & = & \lambda (\overline{n (c,
    K^{\prime})})\\
    \text{when } K^{\prime} \neq \bot &  & \\
    n \left( \text{{\tmstrong{let}}} p = e_1  \text{{\tmstrong{in}}} e_2,
    K^{\prime} \right) & = & \text{{\tmstrong{let}}} p = n (e_1, \bot) 
    \text{{\tmstrong{in}}} n (e_2, K^{\prime})\\
    n \left( \text{{\tmstrong{assert num}}} e_1 \leqslant e_2 ; e_3,
    K^{\prime} \right) & = & \text{{\tmstrong{assert num}}} n (e_1, \bot)
    \leqslant n (e_2, \bot) ; n (e_3, K^{\prime})\\
    n \left( \text{{\tmstrong{assert type}}} e_1 = e_2 ; e_3, K^{\prime}
    \right) & = & \text{{\tmstrong{assert type}}} n (e_1, \bot) = n (e_2,
    \bot) ; n (e_3, K^{\prime})
  \end{eqnarray*}
  \label{NormExpr}
  \caption{Flatten nested introductions of existential types}
\end{table}\begin{table}[t]
  \begin{itemize}
    \begin{tabular}{llll}
      {\tmname{App}} &  & {\tmname{ExLetIn}} & \\
      $\large{\frac{\begin{array}{ll}
        C, \Gamma, \Sigma \vdash e_1 : \tau^{\prime} \rightarrow \tau & \\
        C, \Gamma, \Sigma \vdash e_2 : \tau^{\prime} & C \vDash \not{E}
        (\tau^{\prime})
      \end{array}}{C, \Gamma, \Sigma \vdash e_1 e_2 : \tau}}$ &  &
      {\large{$\frac{\begin{array}{ll}
        \varepsilon_K (\bar{\alpha}) \text{ in } \Sigma & \text{ } C, \Gamma,
        \Sigma \vdash e_1 : \tau^{\prime}\\
        C, \Gamma, \Sigma \vdash K p.e_2 : \tau^{\prime} \rightarrow \tau & 
      \end{array}}{C, \Gamma, \Sigma \vdash^{} \text{{\tmstrong{let}}} p = e_1
      \text{{\tmstrong{in}}} e_2 : \tau}$}} & \\
      &  &  & \\
      {\tmname{ExIntro}} &  & {\tmname{ExAbs}} & \\
      $\large{\frac{\begin{array}{l}
        \tmop{Dom} (\Sigma^{\prime}) \backslash \tmop{Dom} (\Sigma)
        =\mathcal{E} (e)\\
        C, \Gamma, \Sigma^{\prime} \vdash n (e) : \tau
      \end{array}}{C, \Gamma, \Sigma \vdash e : \tau}}$ &  &
      {\large{$\frac{\begin{array}{l}
        C \vDash \tmop{RetType} (\tau, \varepsilon_K (\bar{\tau}^{\prime}))\\
        \forall i C, \Gamma \vdash \lambda \bar{c} : \tau_{}
      \end{array}}{C, \Gamma \vdash \lambda_K \bar{c} : \tau}$}} & 
    \end{tabular}
    
    \ 
  \end{itemize}
  \label{ExTypingRul}
  \caption{Typing rules added by $\tmop{MMG}_{\exists} (X)$}
\end{table}\begin{table}[t]
  \begin{eqnarray*}
    \llbracket \Gamma \vdash e_1 e_2 : \tau \rrbracket & = & \exists \alpha .
    \llbracket \Gamma \vdash e_1 : \alpha \rightarrow \tau \rrbracket \wedge
    \llbracket \Gamma \vdash e_2 : \alpha \rrbracket \wedge \not{E} (\alpha),
    \alpha \# \tmop{FV} (\Gamma, \tau)\\
    &  & \\
    \llbracket \Gamma, \Sigma_0 \vdash e : \tau \rrbracket & = & \llbracket
    \Gamma, \Sigma \vdash n (e) : \tau \rrbracket\\
    \text{when } \mathcal{E} (e) \neq \varnothing &  & \text{where } \Sigma =
    \Sigma_0 \overline{K \colons \forall \alpha_K \gamma_K [\chi_K (\gamma_K,
    \alpha_K)] . \gamma_K \rightarrow \varepsilon_K (\alpha_K)}_{K \in
    \mathcal{E} (e)}\\
    &  & \\
    \left\llbracket \Gamma \vdash \text{{\tmstrong{let}}} p = e_1 
    \text{{\tmstrong{in}}} e_2 : \tau \right\rrbracket & = & \exists \alpha_0
    . \llbracket \Gamma \vdash e_1 : \alpha_0 \rrbracket \wedge \left(
    \llbracket \Gamma \vdash p.e_2 : \alpha_0 \rightarrow \tau \rrbracket
    \wedge \not{E} (\alpha_0) \\
    \text{ \ \ \ \ \ } \vee_{\mathcal{E}} \llbracket \Gamma \vdash K p.e_2 :
    \alpha_0 \rightarrow \tau \rrbracket \right)\\
    &  & \\
    &  & \text{where } \mathcal{E}= \{ K | K \colons \forall \bar{\alpha}
    \bar{\beta} [E] . \tau \rightarrow \varepsilon_K (\bar{\alpha}) \in \Sigma
    \}\\
    &  & \\
    \llbracket \Gamma \vdash \lambda_K \bar{c} : \tau \rrbracket & = & \exists
    \alpha_0 . \tmop{RetType} (\tau, \alpha_0) \wedge (\exists \alpha_1 .
    \alpha_0 \dot{=} \varepsilon_K (\alpha_1) \wedge \chi_K (\alpha_1)) \wedge
    \llbracket \Gamma \vdash \lambda \bar{c} : \tau \rrbracket, \text{}
    \alpha_0 \alpha_1 \# \tmop{FV} (\Gamma, \tau)
  \end{eqnarray*}
  \label{ExIntermInf}
  \caption{Type inference for the added expressions}
\end{table}

\begin{definition}
  \label{WellTypedEx}Let $\Sigma = \Sigma_0 \cup \Sigma_e$ and
  $\Sigma^{\prime} = \Sigma_0 \cup \Sigma_e^{\prime}$ be sets of value
  constructors related to each other as follows:
  \begin{itemize}
    \item $\tmop{PV}^2 (\Sigma_0) = \varnothing$,
    
    \item $\Sigma_e = \overline{K \colons \forall \alpha_K \gamma_K [\chi_K
    \left( \gamma_K, {\nobreak} \alpha_K \right)] . \gamma_K \rightarrow
    \varepsilon_K (\alpha_K)}$,
    
    \item and $\Sigma_e^{\prime} = \overline{K \colons \forall
    \bar{\alpha}_K^{\prime} \bar{\beta}_K^{\prime} \gamma_K [E_K] . \gamma_K
    \rightarrow \varepsilon_K (\bar{\alpha}_K^{\prime})}$
  \end{itemize}
  where $\exists \bar{\alpha}_K^{\prime} \bar{\beta}_K^{\prime} \gamma_K .E_K$
  are solved form formulas. Define $\Sigma^{\prime} / \Sigma =\mathcal{I}_e =
  [\overline{\chi_K} \assign \overline{\exists \bar{\alpha}_K .F_K}]$, where
  $F_K = E_K \wedge \alpha_K \dot{=} \overrightarrow{\alpha_K}^{\prime}$ and
  $\bar{\alpha}_K = \bar{\alpha}_K^{\prime} \bar{\beta}_K^{\prime}$.
\end{definition}

Note that we do not lose generality by using single-argument datatypes
$\varepsilon_K (\alpha)$ rather than the general form $\varepsilon_K
(\bar{\alpha})$. Multiple parameters can be captured by providing a constraint
$\exists \alpha_1 \ldots \alpha_n . \alpha \dot{=} \alpha_1 \rightarrow \ldots
\rightarrow \alpha_n$, with existential variables $\alpha_1 \ldots \alpha_n$,
as part of the constraints of constructors of $\varepsilon_K (\alpha)$. In the
implementation of {\tmname{InvarGenT}}, we recover the multiple parameter
representation $\varepsilon_K (\bar{\alpha})$ of existential types, at the end
of inference.

Normalization defined in Table \ref{NormExpr} is responsible for introduction
of existential types, but it also ensures that existential types never
directly wrap around other existential types. This flattening enables the use
of all information available to derive the postcondition, i.e. the existential
type. To flatten nested existential types, we rename constructors $K$ to
$K^{\prime}$ in $n (\lambda [K] \bar{c}, K^{\prime})$, and eliminate potential
existential type before introducing one in $n (e, K^{\prime})$ when
$K^{\prime} \neq \bot \wedge l (e) =\tmmathbf{F}$. Note the frequent need to
use \tmverbatim{efunction} (and its syntax sugar forms \tmverbatim{ematch} and
\tmverbatim{eif}) for all branches of a definition, as for example in function
\tmverbatim{bsearch2} in Section \ref{ExBsearch2}. The branches bound by
\tmverbatim{match} or \tmverbatim{if} would be required to have the exact same
type, e.g. the same number \tmverbatim{Num}, height of a tree or length of a
list.

\section{Type Inference Constraints for Existential Types}\label{ExTypesCns}

The type inference uses predicate variables to determine the information
bundled inside existential types. The constraint derivation rules related to
existential types are provided in Table \ref{ExIntermInf}. For the initial
call to $\llbracket \cdummy \rrbracket$ (i.e. case $\mathcal{E} (e) \neq
\varnothing$), we normalize the expression. We shorten $\llbracket \Gamma,
\Sigma \vdash \cdummy : \tau \rrbracket$ to $\llbracket \Gamma \vdash \cdummy
: \tau \rrbracket$.

Our tools for solving second order constraints only handle conjunctions of
implications. We solve disjunctions early, which is problematic as selecting a
disjunct may require information hidden in other disjunctions or in predicate
variables. For example, in the normalization of constraints we need to
associate each unary predicate variable with at most one inferred existential
type that can occur as return type in its solution. The pragmatics we adopt in
{\tmname{InvarGenT}} is that whenever the $\llbracket \Gamma \vdash p.e_2 :
\alpha_0 \rightarrow \tau \rrbracket$ disjunct coming from the
{\tmname{LetIn}} rule is satisfiable with the rest of the constraint, we
select it for the solution. One can turn the pragmatics into semantics by
adding the premise $C, \Gamma, \Sigma \nvdash \lambda (p.e_2) e_1 : \tau$ to
the {\tmname{ExLetIn}} rule, but it would make the formalism more complex. The
proof of the following theorem is in Section \ref{ExGADTsProofs}.

\begin{theorem}
  \label{InfCorrComplEx}Theorems \ref{InfCorrExpr} (Correctness) and
  \ref{InfComplExpr} (Completeness) hold for the type system extended with
  {\tmname{ExIntro}} and {\tmname{ExLetIn}} in the following sense. Recall the
  notation from Theorem \ref{InfComplExpr}.
  
  Correctness: For all $\Gamma, \Sigma, e$ and $\tau$, $\llbracket \Gamma,
  \Sigma \vdash e : \tau \rrbracket, \Gamma, \Sigma \vdash e : \tau$ is
  derivable.
  
  Completeness: For all $\Gamma, \Sigma, e$ and $\tau$, if $\tmop{PV} (C,
  \Gamma, \Sigma) = \varnothing$ and $C, \Gamma, \Sigma \vdash e : \tau$, then
  there exist interpretations of predicate variables $\mathcal{I}_u,
  \mathcal{I}_e$ such that $\tmop{Dom} (\mathcal{I}_u)$ are unary, $\tmop{Dom}
  (\mathcal{I}_e) = \{ \chi_K | K \in \mathcal{E} (e) \}$, and $\mathcal{M},
  \mathcal{I}_u \vDash C \Rightarrow \mathcal{I}_e (\llbracket \Gamma, \Sigma
  \vdash e : \tau \rrbracket) [\overline{\varepsilon_K (\vec{\tau})} \assign
  \overline{\varepsilon_K (\bar{\tau})}]$.
\end{theorem}

The set of value constructors is updated in {\tmname{InvarGenT}} after a
toplevel definition with a well typed body: from $\Sigma_0$ to
$\Sigma^{\prime}$, using the notation from Definition \ref{WellTypedEx}.

\begin{example}
  Consider the function \tmverbatim{filter} for lists with length. The
  \tmverbatim{eif}$\ldots$\tmverbatim{then}$\ldots$\tmverbatim{else} syntax is
  a syntactic sugar for the \tmverbatim{ematch}$\ldots$\tmverbatim{with True
  ->}$\ldots$\tmverbatim{ \textbar  False ->}$\ldots$ syntax, which in turn is
  a syntactic sugar for \tmverbatim{(efunction True ->}$\ldots$\tmverbatim{
  \textbar  False ->}$\ldots$\tmverbatim{)}$\ldots$ expressions.
  \begin{tmcode}
  datatype List : type * num
datacons LNil : a. List(a, 0)
datacons LCons :
  n, a [0n]. a * List(a, n)  List(a, n+1)

let rec filter = fun f ->
  efunction LNil -> LNil
  | LCons (x, xs) ->
    eif f x
    then let ys = filter f xs in LCons (x,ys)
    else filter f xs
  \end{tmcode}
  As an aside, an example interaction with {\tmname{InvarGenT}} might look as
  follows:
  \begin{tmcode}
  $ ./invargent -inform examples/filter.gadt
val filter :
  n, a.
  (a  Bool)  List (a, n)  k[0  k  k  n].List (a, k)
InvarGenT: Generated file examples/filter.gadti
InvarGenT: Generated file examples/filter.ml
InvarGenT: Command "ocamlc -w -25 -c examples/filter.ml" exited with code 0
  \end{tmcode}
  Below we show the corresponding type inference constraint, slightly
  simplified by hand for readability. Constructors such as $\varepsilon_{K_2}
  (\alpha_6)$ identify an occurrence of an existential type, here $\exists
  \delta . \chi_1 (\delta, \alpha_6)$ in a constructor environment with $K_2
  \colons \forall \delta \alpha [\chi_1 (\delta, \alpha)] . \delta \rightarrow
  \varepsilon_{K_2} (\alpha)$.
  \begin{eqnarray}
    \exists \alpha_0 . \forall \beta_0 . \exists n_1 \alpha_1 \alpha_2
    \alpha_3 . \forall k_1 \beta_1 \beta_2 . \exists \alpha_5 \alpha_6
    \alpha_7 \alpha_8 \alpha_9 . \forall \beta_3 \beta_4 \beta_5 \beta_6 .
    \exists n_2 \alpha_{10} \alpha_{11} \alpha_{12} \alpha_{13} . &  & 
    \nonumber\\
    &  & \chi_2 (\alpha_0)  \label{ExCnBr1}\\
    & \wedge &  \nonumber\\
    \chi_2 (\beta) & \Longrightarrow & \alpha_2 \dot{=} \tmop{List} (\alpha_4,
    n_1) \wedge \beta \dot{=} \alpha_1 \rightarrow \tmop{List} (\alpha_4, n_1)
    \rightarrow \alpha_3  \label{ExCnBr2}\\
    & \wedge &  \nonumber\\
    \tmop{List} (\beta_1, 0) \dot{=} \tmop{List} (\alpha_4, n_1) \wedge \chi_2
    (\beta) & \Longrightarrow & \alpha_3 \dot{=} \varepsilon_{K_2} (\alpha_6)
    \wedge \chi_1 (\tmop{List} (\alpha_5, 0), \alpha_6)  \label{ExCnBr3}\\
    & \wedge &  \nonumber\\
    \tmop{List} (\beta_2, k_1 + 1) \dot{=} \tmop{List} (\alpha_4, n_1) \wedge
    & \Longrightarrow & \alpha_1 \dot{=} \beta_2 \rightarrow \tmop{Bool}
    \wedge  \label{ExCnBr4}\\
    0 \leq k_1 \wedge \chi_2 (\beta) &  & \chi_2 (\alpha_1 \rightarrow
    \tmop{List} (\beta_2, k_1) \rightarrow \varepsilon_{K_2} (\alpha_8))
    \wedge \nonumber\\
    &  & \chi_2 (\alpha_1 \rightarrow \tmop{List} (\beta_2, k_1) \rightarrow
    \varepsilon_{K_2} (\alpha_9)) \nonumber\\
    & \wedge &  \nonumber\\
    \varepsilon_{K_2} (\beta_4) \dot{=} \varepsilon_{K_2} (\alpha_8) \wedge
    \chi_1 (\beta_3, \beta_4) \wedge & \Longrightarrow & \alpha_3 \dot{=}
    \varepsilon_{K_2} (\alpha_8) \wedge \chi_1 (\alpha_{10}, \alpha_{11})
    \wedge \beta_3 \dot{=} \tmop{List} (\beta_2, n_2) \wedge \nonumber\\
    \tmop{List} (\beta_2, k_1 + 1) \dot{=} \tmop{List} (\alpha_4, n_1) \wedge
    &  & \alpha_{10} \dot{=} \tmop{List} (\beta_2, n_2 + 1) \wedge 0 \leq n_2 
    \label{ExCnBr5}\\
    0 \leq k_1 \wedge \chi_2 (\beta) &  &  \nonumber\\
    & \wedge &  \nonumber\\
    \varepsilon_{K_2} (\beta_6) \dot{=} \varepsilon_{K_2} (\alpha_9) \wedge
    \chi_1 (\beta_5, \beta_6) \wedge & \Longrightarrow & \alpha_{12} \dot{=}
    \beta_5 \wedge \alpha_3 \dot{=} \varepsilon_{K_2} (\alpha_{13}) \wedge
    \chi_1 (\alpha_{12}, \alpha_{13})  \label{ExCnBr6}\\
    \tmop{List} (\beta_2, k_1 + 1) \dot{=} \tmop{List} (\alpha_4, n_1) \wedge
    &  &  \nonumber\\
    0 \leq k_1 \wedge \chi_2 (\beta) &  &  \nonumber
  \end{eqnarray}
  Branch \ref{ExCnBr1} ensures that the invariant is satisfiable. Branch
  \ref{ExCnBr2} decomposes the type of the recursive definition and ensures
  that the second argument is a list. Branch \ref{ExCnBr3} is the case of
  empty list. Branch \ref{ExCnBr4} handles the two recursive calls. Branch
  \ref{ExCnBr5} is the case of kept element: the recursive call result
  $\beta_3$ has length one-shorter than the result $\alpha_{10}$. Branch
  \ref{ExCnBr6} is the case of a dropped element: the recursive call result
  $\beta_5$ and the function result $\alpha_{12}$ are the same.
\end{example}

\section{Semantics by Reduction to $\tmop{HMG} (X)$}

The typing rules that importantly differ between $\tmop{HMG} (X)$ and
$\tmop{MMG}_{\exists} (X)$ are: {\tmname{NegClause}}, {\tmname{FailClause}},
{\tmname{App}}, {\tmname{LetIn}}, {\tmname{ExLetIn}}, {\tmname{ExIntro}},
{\tmname{ExAbs}}. The remaining, crucial difference lies in having predicate
variables among constraints. However:
\begin{itemize}
  \item {\tmname{ExIntro}}, {\tmname{LetIn}} and {\tmname{ExLetIn}} can be
  seen as ``desugaring'' a program to a core language,
  
  \item pattern matching branches that invoke {\tmname{NegClause}} are
  unreachable and so can be dropped,
  
  \item pattern matching branches that invoke {\tmname{FailClause}} can be
  dropped, because they lead to runtime failure, computationally equivalent to
  a match failure,
  
  \item {\tmname{App}} in $\tmop{MMG}_{\exists} (X)$ is a restricted form of
  its variant from $\tmop{HMG} (X)$.
\end{itemize}
Therefore, if an inference problem constraint $\llbracket \Gamma, \Sigma
\vdash e : \tau \rrbracket$ holds, then we can build environment
$\Gamma^{\prime}$, constructor environment $\Sigma^{\prime}$ and an expression
$e^{\prime}$ such that $C^{\prime}, \Gamma^{\prime} \vdash e^{\prime} : \tau$
is derivable in $\tmop{HMG} (X)$ for a satisfiable constraint $C^{\prime}$.
For the reduction to be a convincing argument for soundness of the type
system, we also need to show that the ``desugared'' expression $e^{\prime}$
preserves the intended meaning of $e$. For the details of the $\tmop{HMG} (X)$
type system, we refer the reader to Simonet and Pottier
{\cite{simonet-pottier-hmg-toplas}}.

\begin{definition}
  \label{SemNormForm}Consider a constructor environment $\Sigma$ and an
  $\tmop{HMG} (X)$ expression $e$. A {\tmem{tag erasure}} of $e$ w.r.t.
  $\Sigma$ is an expression achieved by iteratively replacing each subpattern
  $K p_1$ of $e$ with $K \nin \tmop{Dom} (\Sigma)$ by $p_1$ and each
  subexpression $K e_1$ of $e$ with $K \nin \tmop{Dom} (\Sigma)$ by $e_1$.
  
  Consider an $\tmop{MMG}_{\exists} (X)$ expression $e$. An {\tmem{HMG-form}}
  of $e$ is an $\tmop{HMG} (X)$ expression achieved by iteratively replacing
  each subexpression $\lambda [K] \bar{c}$ and $\lambda_K \bar{c}$ of $e$ by
  $\lambda \bar{c}$, each subexpression $\text{{\tmstrong{assert type}}} e_1 =
  e_2 ; e_3$ of $e$ by $e_3$, each subexpression $\text{{\tmstrong{runtime
  failure}}} s$ by a function call $\tmop{failwith} s$, and each subexpression
  $\text{{\tmstrong{let rec}}} x = e_1  \text{{\tmstrong{in}}} e_2$ in $e$ by
  $\text{{\tmstrong{let}}} x = \mu x.e_1  \text{{\tmstrong{in}}} e_2$.
\end{definition}

The HMG-form from Definition \ref{SemNormForm} captures the intended meaning
of programs. We provide semantics for the $\lambda$-calculus underlying
$\tmop{MMG}_{\exists} (X)$, by reducing each expression to its HMG-form and
referring to the semantics for $\tmop{HMG} (X)$ provided in
{\cite{simonet-pottier-hmg-toplas}}. In terms of concrete syntax, the HMG-form
of a program is the program with \tmverbatim{assert type} clauses dropped,
keywords \tmverbatim{efunction} replaced by \tmverbatim{function}, and
\tmverbatim{ematch} replaced by \tmverbatim{match}.

\begin{definition}
  \label{CEquiv}For our purposes, let {\tmem{computational equivalence}} be
  the smallest congruence relation $\sim_C$ such that $\text{{\tmstrong{let}}}
  x = e_1  \text{{\tmstrong{in}}} e_2 \sim_C (\lambda x.e_2) e_1$ for any
  expressions $e_1, e_2$, and $\lambda (\overline{p_i .e_i}) \sim_C \lambda
  (\overline{p_i .e_i}_{i : \neg \tmop{unreach} (e_i) \wedge \neg
  \tmop{failure} (e_i)})$ for any sequences of clauses $\overline{p_i .e_i}$.
  The condition is defined by: $\tmop{unreach} (e) \assign e =
  \text{{\tmstrong{assert false}}} \vee (e = \lambda (p_1 .e_1) \wedge
  \tmop{unreach} (e_1))$ and $\tmop{failure} (e) \assign e =
  \text{{\tmstrong{runtime failure}}} s \vee (e = \lambda (p_1 .e_1) \wedge
  \tmop{failure} (e_1))$.
\end{definition}

Expressions $\text{{\tmstrong{let}}} x = e_1  \text{{\tmstrong{in}}} e_2$ and
$(\lambda x.e_2) e_1$ have the same result of reduction at the outermost
context: $e_2 [x \assign e_1]$. Subexpressions $\text{{\tmstrong{assert
false}}}$ in a well-typed program are guaranteed to be unreachable in both
$\tmop{MMG}_{\exists} (X)$ and $\tmop{HMG} (X)$ type systems. Subexpressions
$\text{{\tmstrong{runtime failure}}} s$ are intended to halt a program, so we
assume they are computationally equivalent to match failures (matching against
a value not covered by any pattern matching case). The proof of the follwing
theorem is in Section \ref{SemanticProofs}.

\begin{theorem}
  \label{SemRedHMG}Let $\Gamma, \Sigma \vdash e : \tau$ be a typing judgment
  in $\tmop{MMG}_{\exists} (X)$. If $\llbracket \Gamma, \Sigma \vdash e : \tau
  \rrbracket$ is satisfiable, we can construct a satisfiable constraint
  $C^{\prime}$, an $\tmop{HMG} (X)$ environment $\Gamma^{\prime}$, constructor
  environment $\Sigma^{\prime}$ and an expression $e^{\prime}$ such that
  $C^{\prime}, \Gamma^{\prime} \vdash e^{\prime} : \tau$ is derivable in
  $\tmop{HMG} (X)$, and the tag erasure of $e^{\prime}$ w.r.t. $\Sigma$ is
  computationally equivalent to the HMG-form of $e$ (see Definitions
  \ref{SemNormForm}, \ref{CEquiv}).
\end{theorem}

An expression is {\tmem{closed}} when it does not contain variables not bound
by a $\text{{\tmstrong{let}}}$- or $\text{{\tmstrong{let rec}}}$-expression,
or a pattern. A closed expression is {\tmem{stuck}} if it is neither reducible
nor a value, as defined in {\cite{simonet-pottier-hmg-toplas}}.

\begin{corollary}
  \label{TypeSoundness}{\tmem{Type soundness}}. If a closed expression $e$ is
  well-typed under some constructor environment, then its HMG-form does not
  reduce to a stuck expression.
\end{corollary}

\chapter{Solving Second Order Constraints}\label{ChapSolv}

Least Upper Bounds and Greatest Lower Bounds computations are the standard
tools for finding unknowns involved in an order structure, see e.g. Knowles
and Flanagan {\cite{KnowlesF07}} Section 6.3. In case of implicational
constraints, constraint abduction (see Maher and Huang
{\cite{AbductionSolvMaher}}, Sulzmann, Schrijvers and Stuckey
{\cite{AbdEADTs}}) and constraint generalization belong to this tool-set.
{\tmem{Simple Constraint Abduction}} under Quantifier Prefix is the task of
finding for an implication $\mathcal{Q}.D \Rightarrow C$, where $\mathcal{Q}$
is a quantifier prefix and $D, C$ are conjunctions of atoms, a weakest solved
form formula $\exists \bar{\alpha} .A$ such that $\mathcal{M} \vDash (\exists
\bar{\alpha} .A) \Rightarrow (D \Rightarrow C)$, equivalently $\mathcal{M}
\vDash (\exists \bar{\alpha} .A) \wedge D \Rightarrow C$, $\mathcal{M} \vDash
\exists \tmop{FV} (A, D, C) .A \wedge D$ and \ $\mathcal{M} \vDash
\mathcal{Q}.A [\bar{\alpha} \assign \bar{t}]$ for some $\bar{t}$. {\tmem{Joint
Constraint Abduction}} under Quantifier Prefix handles several implications,
i.e. $\mathcal{Q}. \wedge_i (D_i \Rightarrow C_i)$, simultaneously: each
condition holds for each $D_i, C_i$ pair. It is used to solve for a predicate
variable appearing in multiple premises -- where we are interested in the GLB
of the corresponding conclusions, ``modulo'' the corresponding premises.
{\tmem{Constraint Generalization}} answer to a disjunction $\vee_i D_i$ of
conjunctions of atoms is a solved form formula $\exists \bar{\alpha} .A$ such
that for each $i$, $\mathcal{M} \vDash D_i \Rightarrow \exists \bar{\alpha}
.A$. It is used to solve for a predicate variable appearing in multiple
conclusions -- where we are interested in the LUB of the corresponding
premises.

\section{Overview of Solving for Predicate Variables}\label{PredVarOverview}

Equipped with these tools, consider first solving an un-normalized constraint
$\Phi$ for invariants -- unary predicates $\chi (\cdummy)$. We want the
invariants to be as weak as possible, to make the use of the corresponding
$\tmop{MMG} (X)$ definitions (toplevel expressions) as easy as possible: the
weaker the invariant, the more general the type of the definition. We solve
for the predicate variables in multiple steps, maintaining partial solutions
$\overline{\exists \bar{\beta}^k_{\chi} .F_{\chi}^k}$ and iterating till the
solutions converge (see e.g. Cousot and Cousot {\cite{Cousot77}}). We start
with $k = 0$, $\bar{\beta}^0_{\chi} = \varnothing$, $F^0_{\chi} =\tmmathbf{T}$
for $\chi \in \tmop{PV} (\Phi)$. In each step, we solve the first order
constraint $\mathcal{Q}^k . \wedge_i (D_i^k \Rightarrow C_i^k)$ achieved by
reducing $\Phi \left[ \bar{\chi} \assign \overline{\exists
\bar{\beta}^k_{\chi} . \overline{F^k_{\chi}}} \right]$ to prenex-normal form
and duplicating shared premises. Let $\bar{\beta}^k =
\overline{\bar{\beta}^k_{\chi}}$. We perform joint constraint abduction under
prefix $\mathcal{Q}^k$ (with parameters $\bar{\beta}^k$), let $\exists
\bar{\alpha}^k .A^k$ be one of the abduction answers. We divide the atoms of
the answer $\exists \bar{\alpha}^k .A^k$ into solutions to predicate variables
$A_{\chi}^k$ and a remainder, or residuum, $A_{\tmop{res}}^k = A^k \backslash
\cup_{\chi} A_{\chi}^k$, depending on the variables in the atoms and so that
the residuum holds under the quantifiers: $\vDash \mathcal{Q}^k
.A_{\tmop{res}}^k$. Let $\exists \bar{\beta}^{k + 1}_{\chi} .F^{k + 1}_{\chi}
= \exists \bar{\beta}^k_{\chi} \bar{\alpha}^k .F^k_{\chi} \wedge A^k_{\chi}
[\beta_{\chi} \assign \delta]$.

If in $k$th iteration, $A_{\tmop{res}}^k = A^k$, we are done. From $\vDash
\mathcal{Q}^k .A^k$ and $\vDash A^k \Rightarrow \wedge_i (D_i^k \Rightarrow
C_i^k)$, it follows that the constraint holds: $\mathcal{I} \vDash \Phi$ for
$\mathcal{I}= \left[ \bar{\chi} \assign \overline{\exists \bar{\beta}^k .
\overline{F^k_{\chi}}} \right]$.

When $A^k_{\chi} \neq \tmmathbf{T}$ for some $\chi$, we need to perform
another iteration. It might be that the constraints added by a positive
occurrence of $\chi$ -- a recursive call -- cannot all fit in next iteration's
$\vDash \mathcal{Q}.A_{\tmop{res}}^{k + 1}$ and have to be a part of next
iteration's $A_{\chi}^{k + 1}$.

For postconditions, we want the strongest possible solutions, because a
stronger postcondition provides more information at use sites of a definition.
Therefore, in each iteration, we use constraint generalization to solve for
binary predicate variables $\chi_K (\cdummy, \cdummy)$ without ``hurting'' the
constraint. The generalization results $\exists \bar{\alpha}_{\chi_K}^k
.G_{\chi_K}^k$ are used to get the first order constraints for abduction
$\left[ \bar{\chi} \assign \overline{\exists \bar{\beta}^k_{\chi} .F^k_{\chi}}
; \overline{\chi_K} \assign {\nobreak} \overline{\exists
\bar{\alpha}_{\chi_K}^k .G_{\chi_K}^k} \right]$. However, we have more work
when due to recursive calls of functions returning values of existential
types, $A_{\chi_K}^k \neq \tmmathbf{T}$. Note that postconditions of each
iteration are constructed from scratch, they do not accumulate. The way we
utilize answers $A_{\chi_K}^k$ is by performing second-stage abduction. Let
$\vee_i D_i^{k, \chi_K}$ be the constraint generalization problem resulting in
the postcondition $G_{\chi_K}^k$. We form the joint constraint abduction
problem $\mathcal{Q}^k . \wedge_i (D_i^{k, \chi_K} \Rightarrow A_{\chi_K}^k)$,
and split its answer $A^{k, \chi_K}$ in the same way as for the first-stage
abduction described above. We include the result in the partial solutions:
$\exists \bar{\beta}^{k + 1}_{\chi} .F^{k + 1}_{\chi} = \exists
\bar{\beta}^k_{\chi} \bar{\alpha}^k .F^k_{\chi} \wedge (A^k_{\chi} \wedge_K
A_{\chi}^{k, \chi_K}) [\beta_{\chi} \assign \delta]$. This way, $D_i^{k + 1,
\chi_K}$ may be strong enough to imply $A_{\chi_K}^k$, and hopefully
$A_{\chi_K}^{k + 1} =\tmmathbf{T}$.

But we are still in trouble. Without the postcondition already in place,
branches with recursive calls have too little information to imply what the
postcondition should be. Therefore, during initial iterations ($k = 0$ and $k
= 1$), we only include non-recursive branches in the constraint generalization
process. This leads to ``overshooting'': the initial postconditions may
reflect properties specific to base cases. These properties get relaxed during
successive iterations, and we only end with success after reaching a
fix-point.

The termination condition of the iteration might not be met. Actually,
{\tmname{InvarGenT}} reports type inference failure if the solutions
$F_{\chi}^k$ not converge after a fixed number of iterations. We have only
observed diverging solutions due to numerical constraints.

\section{Constraint Abduction}\label{SecAbduction}\label{AbdAlg}

Abduction is a reasoning technique concerned with the search for explanations.
An abduction problem is usually given by a background theory $\Theta$ and a
formula $C$, and the answer is a formula $A$ such that $\Theta \cup \{A\}
\vDash C$ (relevance), $\Theta \nvDash \neg A$ (consistency), and $A$ has some
restricted syntactical form, see Marta Mayer and Fiora Pirri
{\cite{AbductionSequents}}. We are however interested in {\tmem{constraint
abduction}} problems, with constraints expressed over a fixed model
$\mathcal{M}$. A constraint abduction problem is then given by a pair of
formulas $D, C$, and $A$ is its answer when (at least) $\mathcal{M} \vDash (D
\wedge A) \Rightarrow C$ (relevance) and $\mathcal{M} \vDash \exists \tmop{FV}
(D, A) .D \wedge A$ (consistency), see Michael Maher{\newblock}
{\tmem{Herbrand constraint abduction}} {\cite{AbductionMaher}}. We extend the
formulation of joint constraint abduction (see {\cite{AbductionMaher}}) to
constraints with quantifiers. In general abduction, where a model is built as
a solution, the quantifiers could be eliminated by Herbrandization. However,
Herbrandization (and Skolemization) are operations on the level of a logic,
not within a model -- they do not return equivalent formulas and do not work
in a fixed model. The introduced functions are uninterpreted. We opt to keep
the model $\mathcal{M}$ and handle quantification explicitly. We develop a
combination procedure for abduction algorithms when their domains of
constraints are combined in a very simple way (yet sufficing for
type-inference-driven invariant generation).

\subsection{Formulating the Joint Constraint Abduction Problem}

In joint, also called simultaneous, problems, we expect a single answer to
solve several problems, here: to solve a conjunction of implications.

A {\tmem{solved form}} is a syntactically specified class of formulas,
associated with a class of problems, for which satisfiability is trivial to
check. We restrict solved forms to existentially quantified conjunctions of
atoms, $\exists \bar{\alpha} .A$. The variables $\bar{\alpha}$ of a solved
form $\exists \bar{\alpha} .A$ that is an abduction problem answer, are ``free
parameters'' of the answer, they are required to be ``unconstrained''.

Due to the iterative nature of the main algorithm, we need to account for
prior parameters $\bar{\beta}$ similar to the new parameters $\bar{\alpha}$.

The constraints that we need to solve form a {\tmem{Joint Constraint Abduction
under a Quantifier Prefix problem}} ({\tmem{JCAQP}} problem for short) of the
form $\mathcal{Q}. \wedge_i (D_i \Rightarrow C_i)$, where $D_i$ and $C_i$ are
conjunctions of atomic formulas, and $\mathcal{Q}$ is an arbitrary quantifier
prefix. We assume that $\tmop{FV} (\wedge_i (D_i \Rightarrow C_i)) \subseteq
\mathcal{Q}$. We also have a set $\bar{\beta}$ of {\tmem{parameters}} of the
invariants. The conditions on abduction answers $\exists \bar{\alpha} .A$ are
as follows:
\begin{enumerate}
  \item {\tmem{relevance condition}}: $\mathcal{M} \vDash \wedge_i (D_i \wedge
  A \Rightarrow C_i)$,
  
  \item {\tmem{validity condition}}: $\mathcal{M} \vDash \mathcal{Q}.A
  [\bar{\alpha} \bar{\beta} \assign \bar{t}]$ for some $\bar{t}$,
  
  \item {\tmem{consistency condition}}: $\mathcal{M} \vDash \wedge_i \exists
  \tmop{FV} (D_i \wedge A) .D_i \wedge A$,
  
  \item {\tmem{no escaping variables}} condition: for all atoms $c \in A$ such
  that $\mathcal{M} \nvDash \mathcal{Q}.c$ and $\tmop{FV} (c) \cap \bar{\beta}
  \neq \varnothing$, and for all $\beta_1 \in \tmop{FV} (c)$ such that
  $(\forall \beta_1) \in \mathcal{Q}$, there exists $\beta_2 \in \tmop{FV} (c)
  \cap \bar{\beta}$ such that $\beta_1 \leqslant_{\mathcal{Q}} \beta_2$.
  \begin{enumerate}
    \item {\tmem{Strong no escaping variables condition}} is a stronger
    variant we use for sorts whose solved forms are substitutions. For all
    atoms $\beta_2 \dot{=} t \in A$ such that $\beta_2 \in \bar{\beta}$, if
    $\beta_1 \in \tmop{FV} (c)$ such that $(\forall \beta_1) \in \mathcal{Q}$,
    then $\beta_1 \leqslant_{\mathcal{Q}} \beta_2$.
  \end{enumerate}
\end{enumerate}
\begin{definition}
  \label{AbdAnsDef}$\exists \bar{\alpha} .A$ is an answer to a JCAQP problem
  $\mathcal{Q}. \wedge_i (D_i \Rightarrow C_i)$ with parameters $\bar{\beta}$
  for model $\mathcal{M}$, written $\exists \bar{\alpha} .A \in \tmop{Abd}
  (\mathcal{Q}, \bar{\beta}, \overline{D_i, C_i})$, when $A$ is a conjunction
  of atoms, $\bar{\alpha} \# \tmop{FV} (\wedge_i (D_i \Rightarrow C_i),
  \bar{\beta})$, meeting the relevance condition, validity condition,
  consistency condition and {\tmem{no escaping variables}} condition.
  
  We can also consider Joint Abduction under a Quantifier Prefix problems for
  a logic, checking relevance condition: $\vDash \wedge_i (D_i \wedge A
  \Rightarrow C_i)$ and validity condition: $\vDash \mathcal{Q}.A
  [\bar{\alpha} \bar{\beta} \assign \bar{t}]$. The consistency conditions: for
  all $i$, $D_i \nvDash \neg A$, are always met.
\end{definition}

The natural setting for constraint abduction problems is with a fixed model.
Above we extend the definition to the case where instead of a model just a
logic is given, just to shed light on relations between constraint abduction,
general abduction, and decision (i.e. validity or satisfiability) problems.

We call a JCAQP problem $\mathcal{Q}.D \Rightarrow C$ (i.e. a non-simultaneous
problem) a {\tmem{simple constraint abduction under a quantifier prefix}}
problem SCAQP. We write JCAQP$_{\mathcal{M}}$ to indicate the model.

\begin{proposition}
  If the JCAQP\tmrsub{$\mathcal{M}$} problem $\mathcal{Q}. \wedge_i (D_i
  \Rightarrow C_i)$ without parameters has an answer, then $\mathcal{M} \vDash
  \mathcal{Q}. \wedge_i (D_i \Rightarrow C_i)$.
\end{proposition}

We say that JCAQP\tmrsub{$\mathcal{M}$} answer $\exists \bar{\alpha} .A$ is
{\tmem{more general}} than $\exists \bar{\beta} .B$, when there exist terms
$\bar{t}$ such that $\mathcal{M} \vDash B \Rightarrow A [\bar{\alpha} \assign
\bar{t}]$.

\begin{example}
  For a free term algebra $T (F)$ over signature $F$ containing binary
  functors $f, g$ and constants $a, b$, and JCAQP\tmrsub{$T (F)$} problem $_{}
  \exists x, y, z. (y \dot{=} f (a, x) \Rightarrow z \dot{=} a) \wedge (y
  \dot{=} f (b, x) \Rightarrow z \dot{=} b)$, the most general answer is
  $\exists \alpha .y \dot{=} f (z, \alpha)$. Indeed, $\forall \alpha \exists
  x, y, z.y \dot{=} f (z, \alpha) \wedge y \dot{=} f (a, x)$ holds with $x =
  \alpha, y = f (a, \alpha), z = a$ and $\forall \alpha \exists x, y, z.y
  \dot{=} f (z, \alpha) \wedge y \dot{=} f (b, x)$ holds with $x = \alpha, y =
  f (b, \alpha), z = b$. For the SCAQP\tmrsub{$T (F)$} problem $_{} \exists x,
  y, z. (y \dot{=} f (a, x) \Rightarrow z \dot{=} a)$, the set of maximally
  general answers is $\{\exists \alpha .y \dot{=} f (z, \alpha), z \dot{=}
  a\}$.
\end{example}

Consider a conjunction of atoms $C$ interpreted in any free term algebra $T
(F)$ over a signature $F$. If $C$ is satisfiable, let $\tmmathbf{U} (C)$ be a
conjunction of equations whose left-hand-sides are variables not occurring in
any of the right-hand-sides, such that $T (F) \vDash C \Leftrightarrow
\tmmathbf{U} (C)$, otherwise let $\tmmathbf{U} (C) = \bot$. Let $\tmmathbf{U}
(\mathcal{Q}.C)$ in case $T (F) \not{\vDash} \mathcal{Q}.C$ be $\bot$, and
otherwise be as before but with equations directed so that variables later in
the prefix are on the left. $\tmmathbf{U} (\mathcal{Q}.C)$ can be computed by
unification with linear constant restrictions, see Baader and Schulz
{\cite{UnificationBaader}}. In effect: if for some $x \dot{=} t \in
\tmmathbf{U} (C)$ such that $(\exists t) \nin \mathcal{Q}$ (e.g. $t$ is not a
variable) and $(\forall x) \in \mathcal{Q}$, then $\tmmathbf{U}
(\mathcal{Q}.C) = \bot$; if for some $x \dot{=} t \in \tmmathbf{U} (C)$ such
that $(\exists x) \in \mathcal{Q}$, there is a $y \in \tmop{FV} (t)$,
$(\forall y) \in \mathcal{Q}$ and $x \leqslant_{\mathcal{Q}} y$, then
$\tmmathbf{U} (\mathcal{Q}.C) = \bot$; otherwise $\tmmathbf{U} (\mathcal{Q}.C)
=\tmmathbf{U} (C)$. Let $\tmmathbf{U}_{\bar{\alpha}} (\mathcal{Q}.C)$ be
$\tmmathbf{U}_{\bar{\alpha}} ((\mathcal{Q}\backslash \{ \forall \bar{\alpha}
\}) \exists \bar{\alpha} .C)$, i.e. $\tmmathbf{U}_{\bar{\alpha}}
(\mathcal{Q}.C) =\tmmathbf{U} (\mathcal{Q}^{\prime} .C)$ where variables
$\bar{\alpha}$ are existentially quantified in $\mathcal{Q}^{\prime}$ and
otherwise $\mathcal{Q}^{\prime}$ is like $\mathcal{Q}$. Computing
$\tmmathbf{U}_{\bar{\alpha} \bar{\beta}} (\mathcal{Q}.A)$ decides the validity
condition for $\mathcal{M}= T (F)$.

\begin{definition}
  \label{CorComplAbd}An abduction algorithm for JCAQP\tmrsub{$\mathcal{M}$}
  with parameters $\bar{\beta}$, $\tmop{Abd} (\mathcal{Q}, \bar{\beta},
  \overline{D_i, C_i})$, generates a sequence of quantified conjunctions of
  atoms $\overline{\exists \bar{\alpha}_j .A_j}$, possibly infinite. The
  algorithm is {\tmem{correct}} if for every $j$, $\exists \bar{\alpha}_j
  .A_j$ is an answer to the JCAQP\tmrsub{$\mathcal{M}$} problem $\mathcal{Q}.
  \wedge_i (D_i \Rightarrow C_i)$ with parameters $\bar{\beta}$. It is
  {\tmem{complete}} if for every JCAQP\tmrsub{$\mathcal{M}$} answer $\exists
  \bar{\alpha} .A$, there is a $j$ and some $\bar{t}$ such that $\mathcal{M}
  \vDash A \Rightarrow A_j [\bar{\alpha}_j \assign \bar{t}]$ (with variables
  renamed so that $\bar{\alpha} \# \tmop{FV} (A_j)$). If the sequence is
  empty, we write $\tmop{Abd} (\mathcal{Q}, \bar{\beta}, \overline{D_i, C_i})
  = \bot$.
\end{definition}

Note that we do not require that only maximally general answers are returned.
Although this would be preferrable, it is costly to guarantee in many
constraint domains even with incomplete algorithms.

\subsection{Abduction Algorithm for The Combination of
Domains}\label{JCAAlienMain}

We provide a plug-in architecture where to add a new sort to the logic it is
enough to give an algorithm solving the JCAQP problem. Define an {\tmem{alien
subterm}} (cf. {\cite{UnificationBaader}}) of a term $\tau$ of sort
$s_{\tmop{type}}$ to be a maximally large subterm $t$ of $\tau$ of sort $s
\neq s_{\tmop{type}}$.

Let $\mathcal{L}_{\tmop{ty}} = T (F, \cup_{s \in \tmop{sorts}} X_s)$ be a
language interpreted in a multi-sorted free term algebra $T = T (F \cup_{s \in
\tmop{usorts}} D_s)$ which, besides the term variables $X_{s_{\tmop{type}}}$,
has {\tmem{alien subterm variables}} $\cup_{s \in \tmop{usorts}} X_s$.

Let $\mathcal{Q}$ be a quantifier prefix and $D_i, C_i$ be atomic conjunctions
in $\mathcal{L}$ that form a joint abduction problem $\mathcal{Q}. \wedge_i
(D_i \Rightarrow C_i)$. For the purpose of Theorem \ref{MultisortAbdThm}, let
$\tmop{Abd}_s$ be complete JCAQP algorithms for $s \in \tmop{usorts}$ and
$\tmop{Abd}_T$ be a complete JCAQP algorithm for the free term algebra $T$. We
start the multi-sorted abduction procedure $\tmop{Abd} (\mathcal{Q},
\bar{\beta}, \overline{D_i, C_i})$ by performing $\tmop{Abd}_T$ on the
$s_{\tmop{type}}$ part of constraints with alien subterms replaced by
variables. For each JCAQP solution $\exists \bar{\alpha}_j .A_j$, we replace
the $s_{\tmop{type}}$ part of the $i$th premise by the ``residual'' formula
$A_p^i_j$ and of $i$th conclusion by ``residual'' formula $A_c^i_j$, defined
in Table \ref{MultisortAbd}. We split the resulting JCAQP problem into
single-sort problems for each sort $s \in \tmop{usorts}$, and we solve them
using $\tmop{Abd}_s$ algorithms. Finally, we build answers as conjunctions of
answers for each sort.

Let $\mathcal{Q}$ be a quantifier prefix and $D_i, C_i$ be atomic conjunctions
in $\mathcal{L}$ that form a joint abduction problem $\mathcal{Q}. \wedge_i
(D_i \Rightarrow C_i)$ with parameters $\bar{\beta}$, and $\overline{\forall
\beta} \subset \mathcal{Q}$. Let $\tmop{Abd}_s$ be correct and complete JCAQP
algorithms for $s \in \tmop{usorts}$ and $\tmop{Abd}_T$ be a correct and
complete JCAQP algorithm for language $\mathcal{L}_{\tmop{ty}} = T (F, \cup_{s
\in \tmop{sorts}} X_s)$ and model $T (F \cup_{s \in \tmop{usorts}} D_s)$. We
define the algorithm $\tmop{Abd} (\mathcal{Q}, \bar{\beta}, \overline{D_i,
C_i})$ in Table \ref{MultisortAbd}.\begin{table}[t]
  \begin{tabular}{l}
    {\tmstrong{let}} $D_i \equiv \wedge_s D_i^s$ and $C_i \equiv \wedge_s
    C_i^s$, where, for $s \in \tmop{sorts}$, $D_i^s, C_i^s$ are atomic
    conjunctions in $\mathcal{L}_s$.\\
    {\tmstrong{let}} $D_i^t, C_i^t$ be formulas $D_i^{s_{\tmop{ty}}},
    C_i^{s_{\tmop{ty}}}$ with subterms $\bar{r}_i$ replaced with fresh
    variables $\overline{\alpha_{r_i}}$\\
    \ \begin{tabular}{l}
      (such that $\alpha_r \in X_s$ for $r$ of sort $s$)\\
      where $\bar{r}_i$ are all alien subterms in $D_i^{s_{\tmop{type}}},
      C_i^{s_{\tmop{type}}}$
    \end{tabular}\\
    {\tmstrong{let}} $\overline{\bar{r}} = \overline{r_1} \ldots
    \overline{r_n}$, $\overline{\overline{\alpha_r}} = \overline{\alpha_{r_1}}
    \ldots \overline{\alpha_{r_n}}$\\
    {\tmstrong{if}} $D_i^t$ is not satisfiable for some $i$, {\tmstrong{then}}
    $\tmop{Abd} (\mathcal{Q}, \bar{\beta}, \overline{D_i, C_i}) \assign
    \tmop{Abd} (\mathcal{Q}, \bar{\beta}, \overline{D_j, C_j}_{j \neq i})$\\
    {\tmstrong{else if}} $C_i^t$ is not satisfiable for some $i$,
    {\tmstrong{then return}} $\tmop{Abd} (\mathcal{Q}, \bar{\beta},
    \overline{D_i, C_i}) \assign \bot$\\
    {\tmstrong{else let}} $\wedge_s D^t_{i, s} =\tmmathbf{U} (D^t_i)$,
    $\wedge_s C^t_{i, s} =\tmmathbf{U} (C^t_i)$ be solved forms of $D^t_i$,
    $C^t_i$\\
    Similarly, discard branches $i$ for which $D^t_{i, s} \wedge D_i^s$ is not
    satisfiable for some $s \in \tmop{usorts}$.\\
    {\tmstrong{if}} $\tmop{Abd}_T (\mathcal{Q}, \overline{\overline{\alpha_r}}
    \bar{\beta}, \overline{D_{i, s_{\tmop{type}}}^t, C_{i,
    s_{\tmop{type}}}^t}) = \bot$, {\tmstrong{then}} $\tmop{Abd} (\mathcal{Q},
    \bar{\beta}, \overline{D_i, C_i}) \assign \bot$\\
    {\tmstrong{else let}} $\overline{\exists \overline{\alpha_j^T} .A_T^j} =
    \tmop{Abd}_T (\mathcal{Q}, \overline{\overline{\alpha_r}} \bar{\beta},
    \overline{D_{i, s_{\tmop{type}}}^t, C_{i, s_{\tmop{type}}}^t})$\\
    {\tmstrong{let}} $A_T^{j\prime}$ be $A_T^j$ with alien subterms replaced
    by distinct fresh variables $\bar{\alpha}_r^j$,
    $\overline{\alpha_j^T}^{\prime} = \overline{\alpha_j^T}
    \bar{\alpha}_r^j$\\
    {\tmstrong{let}} $A^i_{p j} = \{ x \dot{=} t \in \tmmathbf{U} (D_{i,
    s_{\tmop{type}}}^t \wedge A_T^{j\prime}) |x \in X_s, s \neq
    s_{\tmop{type}} \}$\\
    {\tmstrong{let}} $A^i_{c j} = \{ x \dot{=} t \in \tmmathbf{U} (D_{i,
    s_{\tmop{type}}}^t \wedge C_{i, s_{\tmop{type}}}^t \wedge A_T^{j\prime})
    |x \in X_s, s \neq s_{\tmop{type}} \}$\\
    {\tmstrong{let}} $A^i_{p / c, j} = \wedge_s A_{p / c, j, s}^i$, $A_{p / c,
    j, s}^i \in \mathcal{L}_s$\\
    {\tmstrong{let}} $J_s = \left\{ j \middle| \tmop{Abd}_s (\mathcal{Q},
    \bar{\alpha}_r^j \bar{\beta}, \overline{D_i^s \wedge (D^t_{i, s} \wedge
    A^i_{p, j, s}) [\overline{\overline{\alpha_r}} \assign
    \overline{\bar{r}}], C_i^s \wedge (C_{i, s}^t \wedge A^i_{c, j, s})
    [\overline{\overline{\alpha_r}} \assign \overline{\bar{r}}]}) \neq \bot
    \right\}$\\
    {\tmstrong{let}} $\overline{\exists \overline{\alpha_s^{k_j^s}}
    .A_s^{k_j^s}} = \tmop{Abd}_s (\mathcal{Q}, \bar{\alpha}_r^j \bar{\beta},
    \overline{D_i^s \wedge (D^t_{i, s} \wedge A^i_{p, j, s})
    [\overline{\overline{\alpha_r}} \assign \overline{\bar{r}}], C_i^s \wedge
    (C_{i, s}^t \wedge A^i_{c, j, s}) [\overline{\overline{\alpha_r}} \assign
    \overline{\bar{r}}]})$,\\
    \ \ \ \ \ \ $j \in \cap_{s \in \tmop{usorts}} J_s$\\
    {\tmstrong{if}} for some $s \in \tmop{usorts}$, $J_s = \varnothing$,
    {\tmstrong{then}} $\tmop{Abd} (\mathcal{Q}, \bar{\beta}, \overline{D_i,
    C_i}) \assign \bot$\\
    {\tmstrong{return}} $\tmop{Abd} (\mathcal{Q}, \bar{\beta}, \overline{D_i,
    C_i}) \assign \overline{\exists \overline{\alpha_{\overline{k_j^s}}}
    .A_T^{j\prime} \wedge_s A_s^{k_j^s}}_{\overline{k_j^s} : j \in \cap_s
    J_s}$,\\
    \ \begin{tabular}{l}
      where $\overline{\alpha_{\overline{k_j^s}}}$ are
      $\overline{\alpha_j^T}^{\prime} \overline{\overline{\alpha_s^{k_j^s}}}
      \cap \tmop{FV} (A_T^{j\prime} \wedge_s A_s^{k_j^s})$,\\
      $\overline{\overline{\alpha_s^{k_j^s}}}$ is a concatenation of
      $\overline{\alpha_s^{k_j^s}}$ for $s \in \tmop{usorts}$,\\
      for $j \in \cap_s J_s$, $\overline{k_j^s}$ span the cartesian product
      $\times_{s \in \tmop{usorts}} \tmop{Abd}_s$ of solutions for fixed $j$
    \end{tabular}
  \end{tabular}\label{MultisortAbd}
  \caption{Complete multi-sorted abduction $\tmop{Abd} (\mathcal{Q},
  \bar{\beta}, \overline{D_i, C_i})$}
\end{table} In this algorithm, we separate formulas into single-sorted
formulas, we perform abduction for terms, and we perform abduction for other
sorts with additional equations due to abduction for terms. The proof of the
following theorem is in Section \ref{SortCombProof}.

\begin{theorem}
  \label{MultisortAbdThm}Assume the conditions listed above. Then $\tmop{Abd}$
  meets Definition \ref{CorComplAbd} {\tmem{correct}} and {\tmem{complete}}
  algorithm conditions.
\end{theorem}

Introducing new variables $\bar{\alpha}_r^j$ puts additional burden on
abduction in other sorts. In the current implementation of
{\tmname{InvarGenT}}, we only replace $A_T^j$ with $A_T^{j\prime}$ in the
first iteration, when abduction in other sorts is not performed.

For the details of the implementation, see Section \ref{JCAAlienAlg}.

\subsection{Joint Constraint Abduction}\label{JCAMain}

In Table \ref{MultisortAbd}, we assumed abduction algorithms are returning
sequences of answers, from which the answers of interest are extracted. We
implement abduction algorithms differently. We maintain a {\tmem{discard
list}} -- a list of answers to avoid, and the algorithms return the first
answer they find which does not imply any answer in the discard list.

Besides the discard list, the simple abduction algorithms take another
argument: the partial answer. By starting the search for an answer to an
implication from the joint solution to already solved implications, we ensure
by construction that the final answer to the joint constraint abduction
problem meets validity and consistency conditions. The relevance and no
escaping variables conditions would be met regardless of starting from the
partial answer. However, the search is greatly facilitated, because simple
constraint abduction does not need to rediscover relevant parts of the answers
to already solved implications.

Sometimes such rediscovery is not possible. Abduction answer $\exists
\bar{\alpha} .A$ to $D \Rightarrow C$ is {\tmem{fully maximal}} when
$\mathcal{M} \vDash (\exists \bar{\alpha} .D \wedge A) \Leftrightarrow D
\wedge C$. The notion was introduced by Michael Maher, see e.g.
{\cite{AbductionMaher}}, to curb the difficulty of finding all abduction
answers. Even equipped only with fully maximal abduction algorithms, by
accumulating answers across implications we can still solve problems where
some implications do not have fully maximal answers. To this effect, we vary
the order in which implications are solved, so that when an implication $D
\Rightarrow C$ without fully maximal answers is encountered, the answer
$\exists \bar{\alpha}_p .A_p$ to previous implications is such that $D \wedge
A_p \Rightarrow C$ has a fully maximal answer. In fact, our simple constraint
abduction algorithms go beyond fully maximal answers.

Rather than testing permutations blindly, we use a search scheme which might
not capture all opportunities to solve a JCA problem, but detects unsolvable
JCA problems earlier. We set aside branches that do not have any answer
extending the partial answer so far. After all branches have been tried and
the partial answer is not an empty conjunction (i.e. not $\tmmathbf{T}$), we
retry the set-aside branches. If during the retry, any of the set-aside
branches fails, we add the partial answer to discarded answers -- which are
avoided during simple abduction -- and restart. Restart puts the set-aside
branches to be tried first. If, when left with set-aside branches only, the
partial answer is an empty conjunction, i.e. all the answer-contributing
branches have been set aside, we fail -- return $\bot$ from the joint
abduction.

Before starting joint constraint abduction, we separate out negative
constraints, as explained in Section \ref{AbdNegConstrn}. To ensure the
overall consistency and validity conditions without needless backtracking, we
pass a validation suite to SCA algorithms. The validation ensures that the
partial answers are consistent with all implications of the constraint. That
is, we reject the partial answers $A$ such that $\mathcal{M} \nvDash \wedge_i
\exists \tmop{FV} (D_i \wedge C_i \wedge A) .D_i \wedge C_i \wedge A$.

For the details of the algorithm, see Section \ref{JCAAlg}.

\subsection{Simple Constraint Abduction}

JCA problems for most interesting domains are unknown to be decidable, because
the corresponding SCA problems are not known to have effective
characterizations of the sets of answers. Maher {\cite{AbductionMaher}}
introduces subsets of answers which are more amenable to search: abduction
answer $\exists \bar{\alpha} .A$ to SCA problem $D \Rightarrow C$ is fully
maximal when $\mathcal{M} \vDash (\exists \bar{\alpha} .D \wedge A)
\Leftrightarrow D \wedge C$. We can search for fully maximal answers using
various forms of a brute-force approach: starting from an {\tmem{initial
candidate}} $D \wedge C$ and generalizing until we find a formula, implied by
$D \wedge C$, which meets all conditions for a correct SCA answer and all its
further generalizations do not imply $C$. It turns out that fully maximal
answers are insufficient. We have come up with two additional ways to
introduce initial candidates. One is to add -- guess -- atoms constraining
variables which are already significantly constrained by the premise $D$. The
``significant constraint'' condition limits the number of guesses to try. The
other way is to decompose the SCA problem $D \Rightarrow C$ into subproblems
$d \Rightarrow c$ for atoms $d \in D, c \in C$ and add their answers to
initial candidates. In many domains, all maximally general answers to $d
\Rightarrow c$ can be easily enumerated. Even considering all initial
candidates described in this paragraph does not ensure finding all maximally
general answers.

We preprocess the initial SCA answer candidate $C_a$ by trying to eliminate
universally quantified variables, using the premise of the SCA problem. We
define this preprocessing, separately for the different sorts, as
$\tmop{Rev}_{\forall} (\mathcal{Q}, \bar{\beta}, D, C_a)$. It makes the
valididty condition of the resulting answer, $\mathcal{M} \vDash
\mathcal{Q}.A$, easier to meet. See Dillig, Dillig, Li and McMillan
{\cite{InductLoopInv}} for a similar approach.

\subsubsection{Abduction for Terms}\label{SCATermsMain}

The JAQP problem for first order logic with function symbols and equality is
undecidable, because it is equivalent, by Herbrandization, to simultaneous
rigid E-unification (see Degtyarev and Voronkov
{\cite{simulRigidEUnifUndec}}): the substitution that is a solution to
simultaneous rigid E-unification when expressed as a conjunction of equations
has the same properties as a JCAQP answer (therefore the existence of JCAQP
answers coincides with intuitionistic satisfiability). Remember that outside
of Definition \ref{AbdAnsDef}, we use $\vDash \Phi$ as a shorthand for
$\mathcal{M} \vDash \Phi$.

The decision problem $T (F) \vDash \mathcal{Q}. \wedge_i (D_i \Rightarrow
C_i)$ is decidable, see Comon {\cite{DisunificationComon}}. Actually,
{\cite{DisunificationComon}} provides a disjunction as a solution, each
disjunct meeting the relevance and validity conditions of the JCAQP\tmrsub{$T
(F)$} problem. It is often the case though that each disjunct does not meet
the consistency condition, despite the JCAQP\tmrsub{$T (F)$} problem
considered having answers.

The JCAQP problem for free term algebra $T (F)$ is unknown to be decidable. A
limited form of JCA is to find the fully maximal answers, introduced in
{\cite{AbductionMaher}}, using non-simultaneous abduction algorithm from Maher
and Huang {\cite{AbductionSolvMaher}}. {\cite{AbductionMaher}} refers to
{\cite{MaherParamSubstitutions}} as establishing that there are finitely many
fully maximal answers. {\cite{AbductionSolvMaher}} gives an algorithm finding
fully maximal answers for simple (i.e. non-simultaneous) constraint abduction
problems.

Let us recall why Herbrandization is insufficient and therefore we cannot
apply the algorithm from {\cite{AbductionSolvMaher}} without some adjustments.
Take a formula $\exists x. (a \dot{=} b (x) \Rightarrow x \dot{=} b (x))
\wedge \varphi (x)$. In the original model $T (F)$, the formula is equivalent
to $\exists x. \varphi (x)$, because $a \dot{=} b (x)$ is equivalent to
falsehood. However, it is a Herbrandization of $\exists x. (\forall b.a
\dot{=} b \Rightarrow x \dot{=} b) \wedge \varphi (x)$, which is equivalent to
$\varphi (a)$.

Fully maximal answers are not sufficient for the practical application of
joint constraint abduction. Consider a constructor $\text{\tmverbatim{Pair}} :
\forall \alpha \beta . \tmop{Term} (\alpha) \times \tmop{Term} (\beta)
\longrightarrow \tmop{Term} ((\alpha, \beta))$ and a pattern matching branch
that we could add to our \tmverbatim{eval} function: \tmverbatim{\textbar 
Pair (y1, y2) -> (eval y1, eval y2)}. It leads to an implication:
\[ \alpha \dot{=} \tmop{Term} ((\alpha^{\prime}, \beta^{\prime})) \Rightarrow
   \beta \dot{=} (\alpha^{\prime\prime}, \beta^{\prime\prime}) \]
where $\alpha^{\prime}, \beta^{\prime}$ are universally quantified and
$\alpha^{\prime\prime}, \beta^{\prime\prime}$ are existentially quantified.
The expected abduction answer $\alpha \dot{=} \tmop{Term} (\beta) \wedge
\alpha^{\prime\prime} \dot{=} \alpha^{\prime} \wedge \beta^{\prime\prime}
\dot{=} \beta^{\prime}$ is not fully maximal because:
\[ \alpha \dot{=} \tmop{Term} ((\alpha^{\prime}, \beta^{\prime})) \wedge \beta
   \dot{=} (\alpha^{\prime\prime}, \beta^{\prime\prime}) \nRightarrow
   \alpha^{\prime} \dot{=} \alpha^{\prime\prime} \wedge \beta^{\prime} \dot{=}
   \beta^{\prime\prime} \]
In the case of \tmverbatim{eval}, the inference problem is solved by our JCA
scheme even based on fully maximal abduction. But examples from Chuan-kai Lin
{\cite{PracticalGADTsInfer}} (for example the \tmverbatim{zip2} and
\tmverbatim{zip1} functions) motivated a development that goes beyond fully
maximal answers: guessing equations between variables.

To show how we eliminate universally quantified variables, we slightly abuse
notation:
\begin{eqnarray*}
  S & = & [\overline{t_u} \assign \overline{t^{\prime}_u}] \text{ for }
  \tmop{FV} (t_u) \cap \overline{\beta_u} \neq \varnothing, \overline{\forall
  \beta_u} \subset \mathcal{Q} \text{ such that } \mathcal{M} \vDash D
  \Rightarrow \dot{S},\\
  S^{\prime} & = & [\bar{u} \assign \overline{t^{\prime}_u}] \text{ for }
  \bar{u} \subset \overline{\beta_u}, \overline{\forall \beta_u} \subset
  \mathcal{Q} \text{ such that } \mathcal{M} \vDash D \wedge C \Rightarrow
  \dot{S}^{\prime},\\
  \tmop{Rev}_{\forall} (\mathcal{Q}, \bar{\beta}, D, C) & = & \{ c^{\prime} |
  c = x \dot{=} t \in C, \text{if } x = S^{\prime} (t) \text{ then }
  c^{\prime} = S (c) \text{ else } c^{\prime} = SS^{\prime} (c) \}
\end{eqnarray*}
$S$ is a substitution of subterms rather than a regular substitution of
variables.

To solve $D \Rightarrow C$ the algorithm from {\cite{AbductionSolvMaher}}
page 13, reproduced in Table \ref{HerbAbdAlg}, starts with $\tmmathbf{U} (D
\wedge C)$ and iteratively replaces subterms by fresh variables $\alpha \in
\bar{\alpha}$ for a final solution $\exists \bar{\alpha} .A$. If the same
subterm occurs at multiple positions, we try replacing by the same variable at
subsets of these positions. We start from $\tmop{Rev}_{\forall} (\mathcal{Q},
\bar{\beta}, \tmmathbf{U} (D \wedge A_p), \tmmathbf{U} (A_p \wedge D \wedge
C))$, where $\exists \bar{\alpha}_p .A_p$ is the solution to previous problems
solved by the joint abduction algorithm. Optionally, we also substitute-out in
the initial candidates, constants $\tau \assign \alpha$ when $\alpha \dot{=}
\tau \in \tmmathbf{U} (D \wedge A_p)$. The modification going beyond fully
maximal answers, is to consider candidate atoms not implied by $D \wedge C$,
even in the case without an initial partial answer: $A_p =\tmmathbf{T}$. A
natural choice is to consider equations $\beta_1 \dot{=} \beta_2$ for
parameters $\beta_1 \beta_2 \subseteq \bar{\beta}$. To curtail the search
space, we limit the choices of pairs $\beta_1, \beta_2$ to cases $A_p \wedge D
\wedge C \Rightarrow \beta_1 \dot{=} \tau_1 \wedge \beta_2 \dot{=} \tau_2$
where $\tau_1$ and $\tau_2$ are not variables but are unifiable.

For the details of the algorithm, see Section \ref{SCATermsAlg}.

\subsubsection{Abduction for Linear Arithmetic}\label{SCANumMain}

We start with some insights into abduction for linear arithmetic, and then
discuss our algorithm. Unlike in term abduction, there is no need to introduce
variables.

\begin{proposition}
  The domain of linear arithmetic has the following {\tmem{quantifier
  elimination}} property: for every constraint (i.e. conjunction of atoms) $A$
  and variables $\bar{\alpha}$ there exists a constraint $A^{\prime}$ such
  that $\mathcal{M} \vDash (\exists \bar{\alpha} .A) \Leftrightarrow
  A^{\prime}$.
\end{proposition}

With inequalities, it can happen that there are no maximally general answers,
there can also be infinitely many maximally general answers outside of fully
maximal answers.

The {\tmem{implicit equalities}} $E$ of a conjunction $C$ is a conjunction of
equations of biggest rank such that $\mathcal{M} \vDash C \Rightarrow E$ and
$E$ are linearly independent of equations in $C$.

\begin{proposition}[{\cite{AbductionLinArithMaher}} p. 14 Lemma 10]
  Suppose $D \wedge C$ has implicit equalities $E$. Then $\vDash A \wedge
  {\nobreak} D \Rightarrow C$ iff $\vDash A \wedge D \Rightarrow E$ and
  $\vDash \tilde{E} (A \wedge D) \Rightarrow \tilde{E} (C)$.
\end{proposition}

Consider the SCA problem $x \dot{=} 2 r \wedge y \dot{=} 2 s \Rightarrow z
\dot{=} 2 t$ and a maximally general answer $A = x + {\nobreak} y \dot{=} z
\wedge r + s \dot{=} t$. It is not fully maximal, because $C \wedge D$ does
not imply any relation between $x, y$ and $z$.

To simplify the search in presence of a quantifier prefix, we preprocess the
initial candidates by trying to eliminate universally quantified variables:
\begin{eqnarray*}
  S & = & [\overline{\beta_u} \assign \overline{t_u}] \text{ for }
  \overline{\forall \beta_u} \subset \mathcal{Q} \text{ such that }
  \mathcal{M} \vDash D \Rightarrow \dot{S},\\
  \tmop{Rev}_{\forall} (\mathcal{Q}, \bar{\beta}, D, C) & = & \{ c^{\prime} |
  c \in C, \text{if } \mathcal{M} \vDash \mathcal{Q}.c [\bar{\beta} \assign
  \bar{t}] \text{ for some } \bar{t} \text{ then } c^{\prime} = c \text{ else
  } c^{\prime} = S (c) \}
\end{eqnarray*}
We approach solving for equations $c = t_1 \dot{=} t_2 \in C$ and inequalities
$c = t_1 \leqslant t_2 \in C$ differently. For equations, we construct initial
candidates as linear combinations of $c$ and selected equations from $D$. For
inequalities, we construct the initial candidates by finding the abudction
answers to $d \Rightarrow c$ for each inequality $d \in D$.

To find the abduction answers to $d \Rightarrow c$, pick a common variable
$\alpha \in \tmop{FV} (d) \cap \tmop{FV} (c)$ or the constant $\alpha = 1$. We
have four possibilities:
\begin{enumerate}
  \item $d \Leftrightarrow \alpha \leqslant d_{\alpha}$ and $c \Leftrightarrow
  \alpha \leqslant c_{\alpha}$: the abduction answers are $c$ and $d_{\alpha}
  \leqslant c_{\alpha}$,
  
  \item $d \Leftrightarrow \alpha \leqslant d_{\alpha}$ and $c \Leftrightarrow
  c_{\alpha} \leqslant \alpha$: the abduction answer is only $c$,
  
  \item $d \Leftrightarrow d_{\alpha} \leqslant \alpha$ and $c \Leftrightarrow
  \alpha \leqslant c_{\alpha}$: the abduction answer is only $c$,
  
  \item $d \Leftrightarrow d_{\alpha} \leqslant \alpha$ and $c \Leftrightarrow
  c_{\alpha} \leqslant \alpha$: the abduction answers are $c$ and $c_{\alpha}
  \leqslant d_{\alpha}$.
\end{enumerate}
Thanks to cases (1) and (4) above, the abduction algorithm can find some
answers which are not fully maximal.

To check whether $A \Rightarrow B$, we check for each $b \in B$:
\begin{itemize}
  \item if $b = x \dot{=} y$, that $A (x) = A (y)$, where $A (\cdummy)$ is the
  substitution corresponding to equations and implicit equalities in $A$;
  
  \item if $b = x \dot{\leqslant} y$, that $A \wedge y \dot{<} x$ is not
  satisfiable.
\end{itemize}
For more details of the algorithm, see Section \ref{SCANumAlg}.

\section{Constraint Generalization}\label{ConstrnGenMain}

We define {\tmem{constraint generalization}} as the task of finding common
consequences, as specific as possible, of conjunctions of atomic formulas.
Notice that there is at most one maximally specific constraint generalization
answer. Therefore the completeness condition for constraint generalization
reduces to requiring most specific answers.

\begin{definition}
  \label{CorComplLUB}A quantified conjunction of atoms $\exists \bar{\alpha}
  .A \in \mathcal{F}$ is a {\tmem{constraint generalization answer}} to
  $\vee_i D_i$, where $D_i$ are conjunctions of atoms in $\mathcal{L}$, when
  $\mathcal{M} \vDash \wedge_i (D_i \Rightarrow \exists \bar{\alpha} .A)$. A
  constraint generalization answer is {\tmem{most specific}} when: for every
  solution $\exists \bar{\alpha}_s .A_s$ with $\mathcal{M} \vDash \wedge_i
  (D_i \Rightarrow \exists \bar{\alpha}_s .A_s)$, $\mathcal{M} \vDash A
  \Rightarrow \exists \bar{\alpha}_s .A_s$ (with variables renamed so that
  $\bar{\alpha}_s \# \tmop{FV} (A)$). By $\tmop{LUB} (\overline{D_i})$ we
  denote the most general constraint generalization answer to $\vee_i D_i$, by
  $\tmop{LUB} (\cdot)$ an algorithm that computes it.
\end{definition}

Consider an \tmverbatim{eval}-like example, where the disjuncts include the
conjunctions of premises and conclusions of a type inference constraint. We
are interested in finding
\begin{eqnarray*}
  \tmop{LUB} (\overline{D_i \wedge C_i}) & = \text{ } \tmop{LUB} & \left(
  \beta \dot{=} \alpha \rightarrow \beta_1 \wedge \alpha \dot{=} \tmop{Term}
  (\tmop{Int}) \wedge \beta_1 \dot{=} \tmop{Int}, \\
  \beta \dot{=} \alpha \rightarrow \beta_1 \wedge \alpha \dot{=} \tmop{Term}
  (\tmop{Bool}) \wedge \beta_1 \dot{=} \tmop{Bool}, \\
  \beta \dot{=} \alpha \rightarrow \beta_1 \wedge \alpha \dot{=} \tmop{Term}
  (\gamma) \wedge \beta_1 \dot{=} \gamma \right)
\end{eqnarray*}
for which the solution is $\exists \alpha_1 . \beta \dot{=} \alpha \rightarrow
\beta_1 \wedge \alpha \dot{=} \tmop{Term} (\alpha_1) \wedge \beta_1 \dot{=}
\alpha_1$.

We discuss issues specific to postcondition inference in Section
\ref{ConstrnGenAlg}.

\subsection{Algorithm for Combining Domains}

First order most specific anti-unifiers are closely related to constraint
generalization.

\begin{definition}
  A sequence of substitutions $\overline{S_i}$ is an {\tmem{anti-unifier}} of
  a same-length sequence of terms $\overline{t_i}$ when there is a term $t_G$
  such that $\forall i, S_i (t_G) = t_i$. The term $t_G$ is called a
  {\tmem{common generalization}} of $\overline{t_i}$.
  
  An anti-unifier $\overline{S_i}$ is a {\tmem{most general anti-unifier}} of
  $\overline{t_i}$ when given any other anti-unifier $\overline{\eta_i}$ there
  exists a substitution $\rho$ such that $\eta_i = S_i \rho$. We call $t_G$
  such that $S_i (t_G) = t_i$ the {\tmem{most specific generalization}} (MSG)
  of $\overline{t_i}$. We require, without loss of generality, that variables
  used by anti-unification are fresh: $\{ x \in X | S_i (x) \neq x \} \#
  \tmop{FV} (\overline{t_i})$.
\end{definition}

We define an {\tmem{alien subterm}} (cf. Baader and Schulz
{\cite{UnificationBaader}}) of a term $\tau$ of sort $s_{\tmop{type}}$ to be a
maximally large subterm $t$ of $\tau$ of sort $s \neq s_{\tmop{type}}$. Let
$\tmop{LUB}_s (\overline{D_i^s})$ give most specific constraint generalization
answers to $\vee_i D_i^s$ for conjunctions of atoms of sort $s$ for $s \neq
s_{\tmop{type}}$. We define $\tmop{LUB} (\overline{D_i}) = \exists
\bar{\alpha} .A$ in Table \ref{LUBAlg}.\begin{table}[t]
  \begin{tabular}{l}
    {\tmstrong{let}} $\wedge_s D_i^s \equiv D_i$ where $D_i^s$ is of sort
    $s$\\
    {\tmstrong{let}} $D_i^t$ be $D_i^{s_{\tmop{type}}}$ with all alien
    subterms $\overline{a^j_i}$ replaced by fresh variables
    $\overline{\alpha^j_i}$ of corresp. sorts\\
    {\tmstrong{let}} $D^a_i = \overline{\alpha^j_i} \dot{=} \overline{a^j_i}$
    and $\wedge_s D_i^{t, s} \equiv D_i^a$ where $D_i^{t, s}$ is in sort $s$\\
    {\tmstrong{let}} $\wedge_s D^t_{i, s} \equiv \tmmathbf{U} (D^t_i)$ where
    $D^t_{i, s}$ is of sort $s$\\
    For the sort $s_{\tmop{type}}$:\\
    \ \begin{tabular}{l}
      {\tmstrong{let}} $V = \{ x_j, \overline{t_{i, j}} | \forall i \exists
      t_{i, j} .x_j \dot{=} t_{i, j} \in D_{i, s_{\tmop{type}}}^t \}$\\
      {\tmstrong{let}} $G = \{ \bar{\alpha}_j, g_j, \overline{S_{i, j}} |
      S_{i, j} = [\bar{\alpha}_j \assign \bar{g}_j^i], S_{i, j} (g_j) = t_{i,
      j} \}$\\
      \ \ \ \ \ be the most specific anti-unifiers of $\{ \overline{_{} t_{i,
      j}} | (x_j, \overline{t_{i, j}}) \in V \}$ for each $j$\\
      {\tmstrong{let}} $D_i^u = \wedge_j \bar{\alpha}_j \dot{=} \bar{g}_j^i$
      and $D_i^g = D_{i, s_{\tmop{type}}}^t \wedge D_i^u$\\
      {\tmstrong{let}} $D^v_i = \{ x \dot{=} y | x \dot{=} t_1 \in D_i^g, y
      \dot{=} t_2 \in D_i^g, \mathcal{M} \vDash D_i^g \Rightarrow t_1 \dot{=}
      t_2 \}$\\
      {\tmstrong{let}} $A_{s_{\tmop{type}}} = \wedge_j x_j \dot{=} g_j \wedge
      \bigcap_i (D_i^g \wedge D_i^v)$\\
      \ where equations are ordered so that only one of $a \dot{=} b, b
      \dot{=} a$ appears anywhere\\
      {\tmstrong{let}} $\bar{\alpha}_{s_{\tmop{type}}} =
      \overline{\bar{\alpha}_j}$
    \end{tabular}\\
    {\tmstrong{let}} $\wedge_s D^u_{i, s} \equiv D^u_i$ for $D^u_{i, s}$ of
    sort $s$\\
    For sorts $s \neq s_{\tmop{type}}$:\\
    \ \ {\tmstrong{let}} $\exists \bar{\alpha}_s .A_s = \tmop{LUB}_s \left(
    \overline{D_i^s \wedge \widetilde{D^a_i} (D^t_{i, s} \wedge D^u_{i, s})}
    \right)$\\
    {\tmstrong{return}} $\exists \overline{\overline{\alpha^j_i}}
    \overline{\bar{\alpha}_s} . \wedge_s A_s$
  \end{tabular}\label{LUBAlg}
  \caption{LUB algorithm $\tmop{LUB} (\overline{D_i}) = \exists \bar{\alpha}
  .A$}
\end{table} The algorithm separates formulas into single-sorted formulas,
computes anti-unification of terms equal to the same variable in each
disjunct, and calls constraint generalization for other sorts. It adds to the
disjuncts passed to generalization for other sorts, equations binding the
generalization variables (introduced by anti-unification) of the particular
sort. To find the anti-unifiers, we adapt the anti-unification algorithm
provided in {\O}stvold {\cite{AntiUnifAlg}}, fig. 2. The proof of the
following theorem is in Section \ref{LUBProofs}.

\begin{theorem}
  \label{CorComplLUBThm}$\tmop{LUB} (\overline{D_i})$ from Table \ref{LUBAlg}
  meets the Definition \ref{CorComplLUB} conditions for most specific
  constraint generalization answer to $\vee_i D_i$.
\end{theorem}

In our actual implementation, we do not separate alien subterms, i.e. we do
not compute $D_i^t$ and $D_i^a$. Rather, we rely on our implementation of
unification $\tmmathbf{U}$ to separate out equations in sorts other than
$s_{\tmop{type}}$.

\subsection{Linear Arithmetic}

Equality under premises within linear arithmetic: $D \vDash \alpha \dot{=}
\beta$ can be decided by checking whether the polytopes $D \wedge \alpha <
\beta$ and $D \wedge \beta < \alpha$ are empty, which can be done by
Fourier-Motzkin elimination.

In contrast to the free term algebra, for the logic of linear inequalities
over real or rational numbers we have the following quantifier elimination
property. For all $\exists \bar{\alpha} .A \in \mathcal{F}$, there is a
conjunction of atoms $A^{\prime}$ such that $\mathcal{M} \vDash A^{\prime}
\Leftrightarrow \exists \bar{\alpha} .A$. Therefore, we need not introduce new
variables.

A system of linear inequalities is {\tmem{full-dimensional}} when the set of
solutions is not contained in an affine subspace of dimension smaller than the
number of variables. In other words, a full-dimensional system of inequalities
does not have implicit equalities. The task of constraint generalization
$\tmop{LUB} (\overline{D_i})$ for full-dimensional inequalities is equivalent
to finding the convex hull of the half-space represented, possibly unbounded
polytopes $D_i$. Fukuda, Liebling and L{\"u}tolf {\cite{ConvexHull}} provide a
polynomial-time algorithm to find the half-space represented convex hull of
closed polytopes. The algorithm can be generalized to unbounded polytopes.

When all variables of an equation $a \dot{=} b$ appear in all branches $D_i$,
we can turn the equation $a \dot{=} b$ into pair of inequalities $a \leqslant
b \wedge b \leqslant a$. We eliminate all equations and implicit equalities
which contain a variable not shared by all $D_i$, by substituting out such
variables. We conjecture that the resulting algorithm meets the correctness
and completeness conditions.

More details of the algorithm implemented in {\tmname{InvarGenT}} can be found
in Section \ref{NumConstrnGenAlg}.

\subsection{Abductive Constraint Generalization}\label{AbdConstrnGenMain}

It turns out that constraint generalization as we defined it is insufficient.
Consider replacing \tmverbatim{function} by \tmverbatim{efunction} in a
function of type $\tau_1 \rightarrow \tau_2$ for which type inference works
fine. We expect the resulting type to have the form $\tau_1 \rightarrow
\exists . \tau_2$, i.e. without existential type variables. But while
unification reconstructs $\tau_2$ from the result types of the branches,
constraint generalization will fail to generate $\tau_2$ when a branch
under-constrains the result, does not ``care'' about the specific type. To
mitigate this problem we introduce {\tmem{abductive constraint
generalization}} and modify our anti-unification algorithm. Constraint
generalization in the numerical domain does not need such modification.

We extend the notion of constraint generalization:

\begin{definition}
  \label{CorAbdLUB}Substitution $U$ and solved form formula $\exists
  \bar{\alpha} .A$ are an answer to {\tmem{abductive constraint
  generalization}} problem $\overline{D_i}$ given a quantifier prefix
  $\mathcal{Q}$ when:
  \begin{enumerate}
    \item $(\forall i) \mathcal{M} \vDash U (D_i) \Rightarrow \exists
    \bar{\alpha} \backslash \tmop{FV} (U) .A$;
    
    \item If $\alpha \in \tmop{Dom} (U)$, then $(\exists \alpha) \in
    \mathcal{Q}$ -- variables substituted by $U$ are existentially quantified;
    
    \item $(\forall i) \mathcal{M} \vDash \forall (\tmop{Dom} (U)) \exists
    (\tmop{FV} (D_i) \backslash \tmop{Dom} (U)) .D_i$.
  \end{enumerate}
\end{definition}

The sort-integrating algorithm changes only minimally. We adapt the
anti-unification algorithm to compute the substitution $U$ from Definition
\ref{CorAbdLUB}. For details of the algorithms, see Section
\ref{AbdConstrnGenAlg}.

\section{Negative Constraints}\label{NegElimMain}

We collect implication branches with conclusion $C_i =\tmmathbf{F}$ and do not
pass them to abduction or constraint generalization.

For the numerical sort, optionally but by default, we try to turn the negative
constraint $D_i^k \Rightarrow \tmmathbf{F}$ into a positive contribution to
constraints under assumption that the numerical domain is integer numbers
rather than rationals. We convert $\neg D_i^k$ into disjunctive normal form,
and replace strict inequalities $w > 0$ with weak inequalities $- w \leq - 1$,
assuming w.l.o.g. integer coefficients in $w$. We eliminate each disjunct
inconsistent with some branch $D_j^k \wedge C_j^k$. If exactly one disjunct
remains, it is the solution to the negative constraint $D_i^k \Rightarrow
\tmmathbf{F}$.

After a joint constraint abduction answer has been found, we check whether all
negated constraints, i.e. $D_i^k$ for $C_i^k =\tmmathbf{F}$, are inconsistent.
If not, we backtrack: redo abduction with $A^k$ put into the discard list.
Moreover, we include in simple constraint abduction a heuristic to prefer
partial answers which make more negated constraints inconsistent.

For details of the algorithm, see Section \ref{NegElimAlg}.

\section{Details of Solving for Predicate Variables}\label{PredVarMain}

Faced with an inference problem $\Phi$, we solve it using iterated abduction.
Let
\[ \exists \delta \Phi \equiv \exists \delta \mathcal{Q}. \wedge_i (D_i
   \Rightarrow C_i) = \exists \delta \tmop{NF} (\Phi) \]
where $D_i, C_i$ are conjunctions of atoms, be quantifier
alternation-minimizing normalization of $\Phi$ (the variable $\delta$ is used
just to bring existentially quantified variables to the front, if possible).
Using the prenex-normal form simplifies formal presentation.

\begin{definition}
  \label{AtomizedForm}We will say that a quantifier prefix and a solved form
  formula $\exists \bar{\alpha} .A$ are in {\tmem{atomized form with respect
  to parameters $\bar{\beta}$}}, when: variables $\bar{\beta}$ are in the same
  quantifier alternation in $\mathcal{Q}$, $\mathcal{M} \vDash \mathcal{Q}.A
  [\bar{\alpha} \bar{\beta} \assign \bar{t}]$ for some $\bar{t}$, and the
  following two conditions hold:
  \begin{enumerate}
    \item there are no properties expressed in $A$ in a way constraining
    parameters $\bar{\alpha} \bar{\beta}$ that can also be expressed without
    constraining them: if there are conjunctions of atoms $C, D$ such that
    $\mathcal{M} \vDash D \wedge C \Rightarrow A$ and $\mathcal{M} \vDash
    \mathcal{Q}.D$, then there exists a conjunction of atoms $B$ such that $B
    \subset A$, $\mathcal{M} \vDash \mathcal{Q}.B$ and $\mathcal{M} \vDash C
    \Rightarrow (A\backslash B)$;
    
    \item all properties expressed in $A$ in a way constraining parameters
    $\bar{\alpha} \bar{\beta}$ are also expressed using only variables
    $\mathcal{Q}_{< \beta}$ (for any $\beta \in \bar{\beta}$) and
    $\bar{\alpha} \bar{\beta}$: if there is a conjunction of atoms $C \subset
    A$ such that $\mathcal{M} \nvDash \mathcal{Q} [(\forall \bar{\beta})
    \assign (\forall \bar{\alpha} \bar{\beta})] .C$, then there is a
    conjunction of atoms $D \subset A$ such that $\tmop{FV} (D) \subset
    \mathcal{Q}_{< \beta} \bar{\alpha} \bar{\beta}$ and $\mathcal{M} \vDash
    \mathcal{Q} [(\forall \bar{\beta}) \assign (\forall \bar{\alpha}
    \bar{\beta})] .D \Rightarrow C$.
  \end{enumerate}
  The atomized form is {\tmem{with respect to parameters
  $\overline{\bar{\beta}^{\chi}}$}} when it is an atomized form with respect
  to parameters $\bar{\beta}^{\chi}$ for all $\bar{\beta}^{\chi} \in
  \overline{\bar{\beta}^{\chi}}$.
\end{definition}

\begin{proposition}
  Let $\mathcal{Q}. \bar{\alpha}_1 \dot{=} \bar{t}_1 \wedge \bar{\alpha}_2
  \dot{=} \bar{t}_2$ be a formula in $\mathcal{L}$ with sorts
  $s_{\tmop{type}}$ and the linear arithmetic sort. If $[\bar{\alpha}_1
  \assign \bar{t}_1 ; \bar{\alpha}_2 \assign \bar{t}_2]$ is a substitution,
  $[\bar{\alpha}_1 \assign \bar{t}_1]$ agrees with quantifier prefix
  $\mathcal{Q}$ and $\bar{\alpha}_2 \subseteq \bar{\alpha} \bar{\beta}$, then
  $\exists \bar{\alpha} . \bar{\alpha}_1 \dot{=} \bar{t}_1 \wedge
  \bar{\alpha}_2 \dot{=} \bar{t}_2$ is in atomized form w.r.t. parameters
  $\bar{\beta}$.
\end{proposition}

Take a conjunction of atoms $A \in \mathcal{L}$, a quantifier prefix
$\mathcal{Q}$, and variables $\bar{\alpha}, \overline{\bar{\beta}^{\chi}}$,
where $\chi$ varies over $\tmop{PV} (\mathcal{Q})$. Let us discuss the routine
$\tmop{Split} (\mathcal{Q}, \bar{\alpha}, A, \overline{\bar{\beta}^{\chi}},
\overline{A_{\chi}^0})$ defined in Table \ref{SplitAlg}\begin{table}[t]
  \label{SplitAlg}{\tmstrong{let}}
  \begin{eqnarray*}
    \alpha \prec \beta & = & \alpha <_{\mathcal{Q}} \beta \vee (\alpha
    \leqslant_{\mathcal{Q}} \beta \wedge \beta \nless_{\mathcal{Q}} \alpha
    \wedge \alpha \in \overline{\bar{\beta}^{\chi}} \wedge \beta \nin
    \overline{\bar{\beta}^{\chi}})\\
    A_{\alpha \beta} & = & \{ \beta \dot{=} \alpha \in A | \beta \in
    \overline{\bar{\beta}^{\chi}} \wedge (\exists \alpha) \in \mathcal{Q}
    \wedge \beta \prec \alpha \}\\
    A_0 & = & A\backslash A_{\alpha \beta}\\
    A^1_{\chi} & = & \left\{ c \in A_0 \middle| \forall \alpha \in \tmop{FV}
    (c) . (\exists \alpha) \in \mathcal{Q} \vee \\
    \text{ \ \ \ } \alpha <_{\mathcal{Q}} \beta_{\chi} \wedge \alpha \nin
    \tmop{PrimCV} (c) \vee \alpha \in \overline{\bar{\beta}^{\chi}} \wedge
    \alpha \in \tmop{PrimCV} (c) \right\}\\
    A^2_{\chi} & = & \tmop{Atomized} (\overline{\bar{\beta}^{\chi}},
    A_{\chi}^1)\\
    A^3_{\chi} & = & A^2_{\chi} \backslash \cup_{\chi^{\prime}}
    A^1_{\chi^{\prime}}
  \end{eqnarray*}
  {\tmstrong{if}} $\mathcal{M} \nvDash \mathcal{Q}. (A \setminus \cup_{\chi}
  A^2_{\chi}) [\bar{\alpha} \assign \bar{t}] \text{ \ for all } \bar{t}$
  
  {\tmstrong{then return}} $\bot$
  
  {\tmstrong{for all}} $\overline{A_{\chi}^+} \text{ min. w.r.t. } \subset
  \text{ s.t. } \wedge_{\chi} (A_{\chi}^+ \subset A^2_{\chi}) \wedge
  \mathcal{M} \vDash \mathcal{Q}. (\cup_{\chi} A_{\chi}^+ \Rightarrow A)
  [\bar{\alpha} \assign \bar{t}] \text{ \ for some } \bar{t}$:
  
  \ \ {\tmstrong{if}} $\tmop{Strat} (A^+_{\chi}, \bar{\beta}^{\chi}) \text{ \
  returns } \bot \text{ for some } \chi$
  
  \ \ {\tmstrong{then return}} $\bot$
  
  \ \ {\tmstrong{else let}}
  \begin{eqnarray*}
    \bar{\alpha}_+^{\chi}, A_{\chi}^L, A^R_{\chi} & = & \tmop{Strat}
    (A^+_{\chi}, \bar{\beta}^{\chi}) \text{ \ \ \ \ \ \ \ \ \ \ \ \ \ \ \ \ \
    \ \ \ \ \ \ \ \ \ \ \ \ \ \ \ \ \ \ \ \ \ \ \ \ \ \ \ \ \ \ \ \ \ \ \ }\\
    A_{\chi} & = & A_{\chi}^0 \cup A_{\chi}^L\\
    \bar{\alpha}^{\chi}_0 & = & \bar{\alpha} \cap \tmop{FV} (A_{\chi})\\
    \bar{\alpha}^{\chi} & = & \left( \bar{\alpha}^{\chi}_0 \setminus
    \bigcup_{\chi^{\prime} <_{\mathcal{Q}} \chi}
    \bar{\alpha}^{\chi^{\prime}}_0 \right) \bar{\alpha}_+^{\chi}\\
    A_+ & = & \cup_{\chi} A_{\chi}^R\\
    A_{\tmop{res}} & = & A_+ \cup \widetilde{A_+} (A \setminus \cup_{\chi}
    A_{\chi}^+)
  \end{eqnarray*}
  \ \ \ \ \ {\tmstrong{if}} $\cup_{\chi} \bar{\alpha}^{\chi} \neq \varnothing
  \vee \cup_{\chi} A^3_{\chi} \neq \varnothing$
  
  \ \ \ \ \ {\tmstrong{then}}
  
  \ \ \ \ \ \ \ \ $\mathcal{Q}^{\prime}, A_{\tmop{res}}^{\prime},
  \overline{\exists \bar{\alpha}^{\prime \chi} .A_{\chi}^{\prime}} \in$
  
  \ \ \ \ \ \ \ \ \ \ \ $\tmop{Split} (\mathcal{Q} [\overline{\forall
  \bar{\beta}^{\chi}} \assign \overline{\forall (\bar{\beta}^{\chi} \cup
  \bar{\alpha}^{\chi})}], \bar{\alpha} \setminus \cup_{\chi}
  \bar{\alpha}^{\chi}, A_{\tmop{res}} \wedge_{\chi} A^3_{\chi},
  \overline{\bar{\beta}^{\chi} \cup \bar{\alpha}^{\chi}},
  \overline{A_{\chi}})$
  
  \ \ \ \ \ \ \ \ {\tmstrong{return}} $\mathcal{Q}^{\prime}, A_{\alpha \beta}
  \wedge A_{\tmop{res}}^{\prime}, \overline{\exists \bar{\alpha}^{\chi}
  \bar{\alpha}^{\prime \chi} .A_{\chi}^{\prime}}$
  
  \ \ \ \ \ {\tmstrong{else return}} $\mathcal{Q} \exists (\bar{\alpha}
  \setminus \cup_{\chi} \bar{\alpha}^{\chi}), A_{\alpha \beta} \wedge
  A_{\tmop{res}}, \overline{\exists \bar{\alpha}^{\chi} .A_{\chi}}$
  
  {\tmstrong{where}} $\tmop{Strat} (A, \bar{\beta}^{\chi})$ is computed as
  follows:
  \begin{enumerate}
    \item for every $c \in A$, and for every $\beta_2 \in \tmop{FV} (c)$ such
    that $\beta_1 <_{\mathcal{Q}} \beta_2$ for $\beta_1 \in
    \bar{\beta}^{\chi}$,
    
    \item if $\beta_2$ is universally quantified in $\mathcal{Q}$, then return
    $\bot$;
    
    \item otherwise, introduce a fresh variable $\alpha_f$, replace $c \assign
    c [\beta_2 \assign \alpha_f]$,
    
    \item add $\beta_2 \dot{=} \alpha_f$ to $A^R_{\chi}$ and $\alpha_f$ to
    $\bar{\alpha}_+^{\chi}$,
    
    \item after replacing all such $\beta_2$ add the resulting $c$ to
    $A_{\chi}^L$.
  \end{enumerate}
  
  \caption{Split routine $\tmop{Split} (\mathcal{Q}, \bar{\alpha}, A,
  \overline{\bar{\beta}^{\chi}}, \overline{A_{\chi}})$}
\end{table}. $\tmop{PrimCV} (c)$ stands for {\tmem{primary constrained
variables}} of an atom $c$. These are the variables that can be separated to
one side of the atom $c$. The routine $\tmop{Split}$ separates an answer
formula into components that need to become (parts of) the predicate variable
solutions, and the residual answer $A_{\tmop{res}}$. The residual answer is a
satisfiable formula, it contains solutions to the existentially quantified
variables. In particular, the types of expressions of interest can be
extracted from the final residual answer. Note that due to existential types
predicates, we actually compute $\tmop{Split} \left( \mathcal{Q},
\bar{\alpha}, A, \overline{\bar{\beta}^{\beta_{\chi}}},
\overline{A_{\beta_{\chi}}^0} \right)$, i.e. we index by $\beta_{\chi}$ (which
can be multiple for a single $\chi$) rather than $\chi$. We retain the
notation indexing by $\chi$ as it better conveys the intent. Let $\tmop{Split}
(\mathcal{Q}, \bar{\alpha}, A, \overline{\bar{\beta}^{\chi}})$ be
$\tmop{Split} (\mathcal{Q}, \bar{\alpha}, A, \overline{\bar{\beta}^{\chi}},
\bar{\tmmathbf{T}})$.

In Table \ref{MainAlgIter}, we describe a single step of solving for predicate
variables. The iteration of the algorithm starts by substituting the partial
solutions to predicate variables from the previous iteration (Table
\ref{MainAlgIter}, Eq. \ref{MainAlgEnt1}). Next, it performs constraint
generalization, to find the current solutions for postcondition-related
predicate variables (Eq. \ref{MainAlgEnt2}). The found postconditions are then
substituted for predicate variables in premises, i.e. at places where a
definition returning an existential type is used recursively (Eq.
\ref{MainAlgEnt3}). Having stronger premises simplifies, and sometimes
enables, the subsequent abduction (Eq. \ref{MainAlgEnt4}). The abduction
answer is split into parts (Eq. \ref{MainAlgEnt5}), that are added to the
partial solutions to predicate variables (Eqs. \ref{MainAlgEnt6},
\ref{MainAlgEnt8}). When abduction shows we need richer postconditions, an
auxiliary round of abduction is performed (Eq. \ref{MainAlgEnt7}). It infers
additional preconditions that enable the required postconditions. We add these
preconditions to the partial solutions (Eq. \ref{MainAlgEnt8}). If the
resulting partial solutions differ from the initial partial solutions, we
continue with another iteration of the main algorithm.\begin{table}[t]
  \label{MainAlgIter}
  \begin{eqnarray}
    \overline{\exists \bar{\beta}^{\chi, k} .F_{\chi}} & = & S_k \nonumber\\
    &  & \text{In iteration 2, remove non-term-sort atoms} \nonumber\\
    &  & \text{containing outer-scope parameters from } S_k . \nonumber\\
    D_K^{\alpha} \Rightarrow C_K^{\alpha} \in R_k^- S_k (\Phi) & = & \text{all
    such that } \chi_K (\alpha, \alpha_{\alpha}^K) \in C_K^{\alpha}, 
    \label{MainAlgEnt1}\\
    &  & \overline{C^{\alpha}_j} = \{ C | D \Rightarrow C \in S_k (\Phi)
    \wedge D \subseteq D^{\alpha}_K \} \nonumber\\
    \alpha_K & : & \chi_K (\alpha_K) \text{ is a positive atom in } \Phi
    \nonumber\\
    U_{\chi_K}, \exists \bar{\alpha}^{\chi_K}_g .G_{\chi_K}, \overline{B_d^K}
    & = & \tmop{DisjElim} \left( \alpha_K, \overline{\delta \dot{=} \alpha
    \wedge D^{\alpha}_K \wedge_j C^{\alpha}_j}_{\alpha \in
    \overline{\alpha_3^{i, K}}} \right)  \label{MainAlgEnt2}\\
    \vec{\tau}_{\varepsilon_K} & = & \overrightarrow{\{ \alpha \in \tmop{FV}
    (G_{\chi_K}) \backslash \delta \delta^{\prime} \bar{\alpha}^{\chi_K}_g |
    (\exists \alpha) \in \mathcal{Q} \vee \alpha \in
    \overline{\bar{\beta}^{\chi}} \backslash \bar{\beta}^{\chi_K} \}}
    \nonumber\\
    \Xi (\exists \bar{\alpha}^{\chi_K}_g .G_{\chi_K}) & = & \exists \tmop{FV}
    (\vec{\tau}_{\varepsilon_K}, G_{\chi_K}) \backslash \delta \delta^{\prime}
    . \delta^{\prime} \dot{=} \vec{\tau}_{\varepsilon_K} \wedge G_{\chi_K}
    \nonumber\\
    (\exists \bar{\alpha}^{\chi_K} .F_{\chi_K}), \rho & = & H (R_k (\chi_K),
    \Xi (\exists \bar{\alpha}_g^{\chi_K} .G_{\chi_K})) \nonumber\\
    R_g (\chi_K) & = & \exists \bar{\alpha}^{\chi_K} .F_{\chi_K} \nonumber\\
    P_g (\chi_K (\delta, \delta^{\prime})) = P_g (\chi_K (\delta^{\prime})) &
    = & \delta^{\prime} \dot{=} \vec{\tau}_{\varepsilon_K} \nonumber\\
    F_{\chi}^{\prime} & = & F_{\chi} [\varepsilon_K (\vec{\tau}_{\tmop{old}})
    \assign \varepsilon_K (\vec{\tau}_{\varepsilon_K}) [\tmop{recover}
    (\vec{\tau}_{\varepsilon_K}, \vec{\tau}_{\tmop{old}})]] \nonumber\\
    S_k^{\prime} & = & \overline{\exists \bar{\beta}^{\chi, k} \{ \alpha \in
    \tmop{FV} (F^{\prime}_{\chi}) | \beta_{\chi} <_{\mathcal{Q}} \alpha \}
    .F^{\prime}_{\chi}^{}} \nonumber\\
    \mathcal{Q}^{\prime} . \wedge_i (D_i \Rightarrow C_i) \wedge_j (D_j^-
    \Rightarrow \tmmathbf{F}) & = & R_g^- P_g^+ S_k^{\prime} (\Phi
    \wedge_{\chi_K} U_{\chi_K})  \label{MainAlgEnt3}\\
    \exists \bar{\alpha} .A_0 & = & \tmop{Abd} \left( \mathcal{Q}^{\prime},
    \bar{\beta} = \overline{\beta_{\chi} \overline{\beta^{}}^{\chi}},
    \overline{D_i, C_i} \right)  \label{MainAlgEnt4}\\
    A & = & A_0 \wedge_j \tmop{NegElim} (\neg \tmop{Simpl} (\tmop{FV} (D_j^-)
    \backslash \overline{\bar{\beta}^{\chi}} .D_j^-), A_0, \overline{D_i,
    C_i}) \nonumber\\
    &  & \text{In later iterations, check negative constraints.} \nonumber\\
    \left( \mathcal{Q}^{k + 1}, A_{\tmop{res}}, \overline{\exists
    \bar{\alpha}^{\beta_{\chi}} .A_{\beta_{\chi}}} \right) & = & \tmop{Split}
    \left( \mathcal{Q}^{\prime}, \bar{\alpha}, A, \overline{\beta_{\chi}
    \overline{\beta^{}}^{\chi}} \right)  \label{MainAlgEnt5}\\
    \vec{\tau}_{\varepsilon_K}^{\prime} & = & \overrightarrow{\tmop{FV}
    (\widetilde{A_{\tmop{res}}} (\vec{\tau}_{\varepsilon_K}))} \nonumber\\
    R_{k + 1} (\chi_K) & = & \exists \bar{\beta}^{\chi_K, k}
    \overline{\bar{\alpha}^{\beta_{\chi_K}}} \bar{\alpha}^{\chi_K} \backslash
    \tmop{FV} (\vec{\tau}_{\varepsilon_K}^{\prime}) . \delta^{\prime} \dot{=}
    \vec{\tau}_{\varepsilon_K}^{\prime} \wedge \widetilde{A_{\tmop{res}}}
    (F_{\chi_K} \backslash \delta^{\prime} \dot{=} \ldots) \\
    \text{ \ \ \ } \wedge_{\chi_K} A_{\beta_{\chi_K}} \left[
    \overline{\beta_{\chi_K} \overline{\beta^{}}^{\beta_{\chi_K}}} \assign
    \overline{\delta \overline{\beta^{}}^{\chi_K, k}} \right] 
    \label{MainAlgEnt6}\\
    A_K^d & = & \left\{ c \in A_{\beta_{\chi_K}} \left[
    \overline{\beta_{\chi_K} \overline{\beta^{}}^{\beta_{\chi_K}}} \assign
    \overline{\delta \overline{\beta^{}}^{\chi_K, k}} \right] \middle| \\
    \text{ \ \ \ \ \ } \tmop{FV} (c) \subseteq \tmop{FV} (\rho (B_d^K))
    \right\} \nonumber\\
    \exists \bar{\alpha}^{\prime}_K .A^{\prime}_K & = & \tmop{Abd} \left(
    \mathcal{Q}^{\prime}, \overline{\beta_{\chi} \overline{\beta^{}}^{\chi}}
    \backslash \overline{\beta_{\chi_K} \overline{\beta^{}}^{\chi_K}},
    \overline{\rho (B_d^K), A_K^d} \right)  \label{MainAlgEnt7}\\
    \left( \mathcal{Q}^{k + 1}, \varnothing, \overline{\exists
    \bar{\alpha}^{\beta_{\chi}\prime} .A_{\beta_{\chi}}^{\prime}} \right) & =
    & \tmop{Split} \left( \mathcal{Q}^{k + 1},
    \overline{\bar{\alpha}^{\prime}_K}, \wedge_K A^{\prime}_K,
    \overline{\beta_{\chi} \overline{\beta^{}}^{\chi}} \backslash
    \overline{\beta_{\chi_K} \overline{\beta^{}}^{\chi_K}} \right) \nonumber\\
    S_{k + 1} (\chi) & = & \exists \bar{\beta}^{\chi, k} . \tmop{Simpl} \left(
    \exists \overline{\bar{\alpha}^{\beta_{\chi}}} .F_{\chi}^{\prime} \\
    \text{ \ \ \ \ \ \ \ \ } \wedge A_{\beta_{\chi}} \left[
    \overline{\beta_{\chi} \overline{\beta^{}}^{\chi}} \assign
    \overline{\delta \overline{\beta^{}}^{\chi, k}} \right] \wedge
    A_{\beta_{\chi}}^{\prime} \right)  \label{MainAlgEnt8}
  \end{eqnarray}
  
  \caption{Iteration $k$ of solving for predicate variables}
\end{table}

Recall Definition \ref{SolTypInfer} of solutions to a type inference problem
$\Phi$.

\begin{proposition}
  If a solution to the inference problem $\Phi$ exists, then there is an
  interpretation of predicate variables $\mathcal{I}$ such that $\mathcal{M},
  \mathcal{I} \vDash \Phi$.
\end{proposition}

Define
\[ \Psi_0 = \{ (\tmmathbf{T}, \bar{\tmmathbf{T}}, \bar{\tmmathbf{T}}) \}
   \text{, \ } \Psi_{k + 1} = \bigcup_{\left( F_{\tmop{res}}^k,
   \overline{\exists \bar{\alpha}^{\chi, k} .F_{\chi}^k}, \overline{\exists
   \bar{\alpha}^{\chi_K, k} .F_{\chi_K}^k} \right) \in \Psi_k} \Psi \left( k,
   \Phi, \overline{\exists \bar{\alpha}^{\chi, k} .F_{\chi}^k},
   \overline{\exists \bar{\alpha}^{\chi_K, k} .F_{\chi_K}^k} \right) \]
where the operator $\Psi$ such that $(\exists \bar{\alpha}_{\tmop{res}}
.A_{\tmop{res}}, S_{k + 1}, R_{k + 1}) \in \Psi (k, \Phi, S_k, R_k)$ for $S_k
= \overline{\exists \bar{\alpha}^{\chi, k} .F_{\chi}^k}$, $R_k =
\overline{\exists \bar{\alpha}^{\chi_K, k} .F_{\chi_K}^k}$, is defined in
Table \ref{MainAlgIter}. The convergence condition is:
\begin{eqnarray*}
  &  & (\forall \chi) S_{k + 1} (\chi) \subseteq S_k (\chi),\\
  & \wedge & (\forall \chi_K) R_{k + 1} (\chi_K) = R_k (\chi_K),\\
  & \wedge & (\forall \beta_{\chi_K}) A_{\beta_{\chi_K}} =\tmmathbf{T},\\
  & \wedge & k > 1
\end{eqnarray*}
We argue for the truth of the following theorem in Section
\ref{PredVarProofs}.

\begin{theorem}
  \label{ThCorrect}{\tmem{Correctness}}. Let $\left( \exists
  \bar{\alpha}_{\tmop{res}} .F_{\tmop{res}}, \overline{\exists
  \bar{\alpha}^{\chi} .F_{\chi}}, \overline{\exists \bar{\alpha}^{\chi_K}
  .F_{\chi_K}} \right)$ be a potential solution to an inference problem
  $\Phi$. If $\left( \exists \bar{\alpha}_{\tmop{res}}^{\prime}
  .F_{\tmop{res}}^{\prime}, \overline{\exists \bar{\alpha}^{\chi} .F_{\chi}},
  \overline{\exists \bar{\alpha}^{\chi_K} .F_{\chi_K}} \right) \in \Psi (\Phi,
  \overline{\exists \bar{\alpha}^{\chi} .F_{\chi}}, \overline{\exists
  \bar{\alpha}^{\chi_K} .F_{\chi_K}})$, then $\mathcal{M}, \mathcal{I} \vDash
  \Phi$ for $\mathcal{I}= \left[ \bar{\chi} \assign \overline{\exists
  \bar{\alpha}^{\chi} .F_{\chi}} ; \overline{\chi_K} \assign \overline{\exists
  \bar{\alpha}^{\chi_K} .F_{\chi_K}} \right]$.
\end{theorem}

To prove Theorem \ref{InferCompleteness}, we have to assume that there are no
existential type predicate variables, i.e. for an inference problem $\Phi$,
$\tmop{PV} (\Phi) = \tmop{PV}^1 (\Phi)$; and also that $\tmop{Abd}$ in the
definition of $\Psi$ is a complete abduction algorithm returning an atomized
form formula. Since the completeness of abduction assumption is not met,
Theorem \ref{InferCompleteness} does not describe an actual implementation. We
argue for the truth of Theorem \ref{InferCompleteness} in Section
\ref{PredVarProofs}.

\begin{theorem}
  \label{InferCompleteness}Let $(\exists \bar{\alpha}^{\tmop{res}}_s
  .F_{\tmop{res}}^s, \overline{\exists \bar{\alpha}^{\chi}_s .F_{\chi}^s})$ be
  any solution to an inference problem $\Phi$ with $\bar{\chi} = \tmop{PV}
  (\Phi) = \tmop{PV}^1 (\Phi)$. Assume that $\tmop{Abd}$ in the definition of
  $\Psi$ is a complete abduction algorithm returning an atomized form formula.
  Then there is a chain $(\exists \bar{\alpha}^{\tmop{res}}_k
  .F_{\tmop{res}}^k, \overline{\exists \bar{\alpha}^{\chi, k} .F_{\chi}^k})
  \in \Psi_k$, with $(\overline{\exists \bar{\alpha}^{\chi, k + 1}
  .F_{\chi}^{k + 1}}) \in \Psi (k, \Phi, \overline{\exists \bar{\alpha}^{\chi,
  k} .F_{\chi}^k})$, such that for all $k \geqslant 1$, $\mathcal{M} \vDash
  \wedge_{\chi} F^s_{\chi} \Rightarrow (\wedge_{\chi} F^k_{\chi})
  [\overline{\bar{\alpha}^{\chi, k}} \assign \bar{t}]$ for some $\bar{t}$.
\end{theorem}

For more discussion on the implementation, see Section
\ref{PredVarAlg}.{\newpage}

\chapter{{\tmname{InvarGenT}}: Tests and Limitations}\label{Limitations}

In this chapter we present types and inference times of tested examples,
measured on a laptop with {\tmem{Intel Core i3}} processor at 2.30GHz, 3MB
cache. The implementation does not use parallelism.

\section{Benchmarks}

Table \ref{Experiments} contains examples using the plain GADTs type system,
without numerical constraints and existential types. Consider the examples
from Chuan-kai Lin's {\cite{PracticalGADTsInfer}}. We tested 7 longer examples
where algorithm $\mathcal{P}$ finds the expected type. {\tmname{InvarGenT}}
does not have problems with these cases either, except the \tmverbatim{head}
function. We also reproduce all 8 examples from {\cite{PracticalGADTsInfer}}
where algorithm $\mathcal{P}$ fails, adapting them only to accomodate to the
type system $\tmop{MMG} (X)$ behind {\tmname{InvarGenT}}. {\tmname{InvarGenT}}
fails on two examples. One, the function \tmverbatim{vary}, was contrived to
not have a clear intended type. The other, the function \tmverbatim{fd\_comp},
can be slightly modified to work with {\tmname{InvarGenT}}'s type inference.
There is also one example, function \tmverbatim{leq}, which requires
non-default parameter setting (passing a parameter \tmverbatim{-prefer\_guess}
to {\tmname{InvarGenT}}). In a future version, {\tmname{InvarGenT}} will
attempt multiple parameter settings before giving up. Overall, we consider the
behavior of {\tmname{InvarGenT}} on this test suite to be a
success.\begin{table}[t]
  \
  
  \begin{tabular}{|l|l|l|}
    \hline
    {\tmem{Class of examples}} & {\tmem{Function name: inferred type}} &
    {\tmem{Inf. time}}\\
    \hline
    incompl. ex. {\cite{OutsideIn}} & \tmverbatim{rx: $\forall$a.R
    a$\rightarrow$Int} & <0.01s\\
    \hline
    examples from {\cite{PointwiseGADTs}} within the scope of algorithm
    $\mathcal{P}$ & \tmverbatim{rotate:
    {\small{$\forall$a.Dir$\rightarrow$Int$\rightarrow$RoB (Black,
    a)$\rightarrow$Dir$\rightarrow$Int$\rightarrow \\
    $ RoB (Black, a)$\rightarrow$RoB (Red, a)$\rightarrow$RoB (Black, S a)}}}
    & 0.02s\\
    \hline
    & \tmverbatim{zip2: {\small{$\forall$a, b.Zip2 (B,
    a)$\rightarrow$b$\rightarrow$a}}} & 0.16s\\
    \hline
    & \tmverbatim{rotl: {\small{$\forall$a.AVL
    a$\rightarrow$Int$\rightarrow$AVL (S (S a))$\rightarrow$\\
    Choice (AVL (S (S a)), AVL (S (S (S a))))}}} & 0.03s\\
    \hline
    & \tmverbatim{ins: {\small{$\forall$a.Int$\rightarrow$AVL
    a$\rightarrow$Choice (AVL a, AVL (S a))}}} & 0.41s\\
    \hline
    & \tmverbatim{extract: {\small{$\forall$a, b.Path b$\rightarrow$Tree (b,
    a)$\rightarrow$a}}} & 0.06s\\
    \hline
    & \tmverbatim{run\_state: {\small{$\forall$a, b.b$\rightarrow$State (b,
    a)$\rightarrow$(b, a)}}} & 0.01s\\
    \hline
    & \tmverbatim{head: {\small{$\forall$a, b.List (a, S b)$\rightarrow$a}}}
    & $\infty$\\
    \hline
    examples from {\cite{PointwiseGADTs}} outside the scope of algorithm
    $\mathcal{P}$ & \tmverbatim{joint: {\small{$\forall$a.Split (a,
    a)$\rightarrow$a}}} & <0.01s\\
    \hline
    & \tmverbatim{rotr: {\small{$\forall$a.Int$\rightarrow$AVL
    a$\rightarrow$AVL (S (S a))$\rightarrow$\\
    Choice (AVL (S (S a)), AVL (S (S (S a))))}}} & 0.09s\\
    \hline
    & \tmverbatim{delmin: {\small{$\forall$a.AVL (S a)$\rightarrow$\\
    (Int, Choice (AVL a, AVL (S a)))}}} & 0.31s\\
    \hline
    & \tmverbatim{fd\_comp: {\small{$\forall$a, b, c.FunDesc (c,
    b)$\rightarrow$FunDesc (b, a)$\rightarrow$ FunDesc (c, a)}}} & 0.2s*,
    0.1s*\\
    \hline
    & \tmverbatim{zip1: {\small{$\forall$a, b.Zip1 (List b,
    a)$\rightarrow$b$\rightarrow$a}}} & 0.08s\\
    \hline
    & \tmverbatim{leq: {\small{$\forall$a.Nat a$\rightarrow$NatLeq (a, a)}}}
    & <0.01s**\\
    \hline
    & \tmverbatim{run\_state: {\small{$\forall$a, b.b$\rightarrow$State (b,
    a)$\rightarrow$(b, a)}}} & 0.03s\\
    \hline
    \begin{tabular}{l}
      run-time type\\
      representations
    \end{tabular} & \tmverbatim{eval: {\small{$\forall$a.Term
    a$\rightarrow$a}}} & 0.06s\\
    \hline
    & \tmverbatim{equal: {\small{$\forall$a,b.(Ty a, Ty
    b)$\rightarrow$a$\rightarrow$b$\rightarrow$Bool}}} & 0.3s, 2.1s\\
    \hline
  \end{tabular}
  
  * Slight meaning-preserving modification of test program
  
  ** Needs a non-default option \tmverbatim{-prefer\_guess}\label{Experiments}
  \caption{Examples of types and inference times, plain GADTs}
\end{table}

Table \ref{ExperimentsNum} contains examples using numerical constraints and
existential types. The examples are organized around several data structures:
lists and arrays with length, binary numbers, AVL balanced search trees. We
provide the inferred types in the hope that they are
self-explanatory.\begin{table}[t]
  \
  
  \begin{tabular}{|l|l|l|}
    \hline
    {\tmem{Class of examples}} & {\tmem{Function name: inferred type}} &
    {\tmem{Inf. time}}\\
    \hline
    lists with length and arrays with static bound checks & \tmverbatim{head:
    {\small{$\forall$n,a[1$\leqslant$n].List (a, n)$\rightarrow$a}}} &
    <0.01s\\
    \hline
    & \tmverbatim{append: {\small{$\forall$a,n,k.List (a, n)$\rightarrow$List
    (a, k)$\rightarrow$List(a, n+k)}}} & 0.02s\\
    \hline
    & \tmverbatim{flatten\_pairs: {\small{$\forall$n, a. List ((a, a), n)
    $\rightarrow$\\
    List (a, 2 n)}}} & 0.01s\\
    \hline
    & \tmverbatim{flatten\_quadrs: {\small{$\forall$n, a. List ((a, a, a, a),
    n) $\rightarrow$ List (a, 4 n)}}} & 0.06s\\
    \hline
    & \tmverbatim{filter: {\small{$\forall$n, a. (a $\rightarrow$ Bool)
    $\rightarrow$ List (a, n) $\rightarrow$\\
    $\exists$k[0 $\leqslant$ k $\wedge$ k $\leqslant$ n].List (a, k)}}} &
    0.18s\\
    \hline
    & \tmverbatim{zip: {\small{$\forall$a,b,n,k.(List (a,n), List
    (b,k))$\rightarrow$ $\exists$i[i=min(n, k)].List ((a,b),i)}}} & 0.38s\\
    \hline
    & \tmverbatim{bsearch2: {\small{$\forall$n, a[0 $\leqslant$ n]. a
    $\rightarrow$ Array (a, n) $\rightarrow$\\
    $\exists$k[0 $\leqslant$ k + 1 $\wedge$ k + 1 $\leqslant$ n].Num k}}} \
    {\small{(uses assertion)}} & 1.39s\\
    \hline
    & \tmverbatim{swap\_interval: {\small{$\forall$i, j, k, n, a[1
    $\leqslant$ j $\wedge$ 0 $\leqslant$ n $\wedge$\\
    0 $\leqslant$ k $\wedge$ i + k $\leqslant$ j $\wedge$ i + n $\leqslant$
    j]. Array (a, j) $\rightarrow$\\
    Num n $\rightarrow$ Num k $\rightarrow$ Num i $\rightarrow$ ()}}} &
    0.27s\\
    \hline
    & \tmverbatim{matmul: {\small{$\forall$i, j, k, n[0 $\leqslant$ n
    $\wedge$ 0 $\leqslant$ k $\wedge$ 0 $\leqslant$ j $\wedge$\\
    j $\leqslant$ i]. Matrix (n,j) $\rightarrow$ Matrix (i,k) $\rightarrow$
    Matrix (n,k)}}} & 0.34s\\
    \hline
    binary numbers & \tmverbatim{plus: {\small{$\forall$n,k,i.Carry
    i$\rightarrow$Binary k$\rightarrow$Binary n$\rightarrow$Binary (n+k+i)}}}
    & 0.66s\\
    \hline
    & \tmverbatim{increment: {\small{$\forall$n.Binary n$\rightarrow$Binary
    (n + 1)}}} & 0.01s\\
    \hline
    & \tmverbatim{bitwise\_or: {\small{$\forall$k, n. Binary k $\rightarrow$
    Binary n $\rightarrow$\\
    $\exists$i[k $\leqslant$ i $\wedge$ n $\leqslant$ i $\wedge$ i
    $\leqslant$ n + k].Binary i}}} & 1.21s\\
    \hline
    AVL trees, imbalance of 2 & \tmverbatim{create: {\small{$\forall$k, n, a[0
    $\leqslant$ n $\wedge$ 0 $\leqslant$ k $\wedge$ n $\leqslant$ k + 2
    $\wedge$\\
    k $\leqslant$ n + 2]. Avl (a, k) $\rightarrow$ a $\rightarrow$ Avl (a, n)
    $\rightarrow$\\
    $\exists$i[i=max (k + 1, n + 1)].Avl (a, i)}}} & 0.09s\\
    \hline
    & \tmverbatim{rotr: {\small{$\forall$k, n, a[0 $\leqslant$ n $\wedge$ n +
    2 $\leqslant$ k $\wedge$ k $\leqslant$ n + 3]. Avl (a, k) $\rightarrow$ a
    $\rightarrow$ Avl (a, n) $\rightarrow$\\
    $\exists$n[k $\leqslant$ n $\wedge$ n $\leqslant$ k + 1].Avl (a, n)}}} &
    1.07s\\
    \hline
    & \tmverbatim{add: {\small{$\forall$n, a. a $\rightarrow$ Avl (a, n)
    $\rightarrow$ $\exists$k[1 $\leqslant$ k $\wedge$ n $\leqslant$ k $\wedge$
    k $\leqslant$ n + 1].Avl (a, k)}}} & 0.69s\\
    \hline
    & \tmverbatim{remove\_min\_binding: {\small{$\forall$n, a[1 $\leqslant$
    n]. Avl (a, n) $\rightarrow$ $\exists$k[n $\leqslant$ k + 1 $\wedge$ k
    $\leqslant$ n $\wedge$ k + 2 $\leqslant$ 2 n].Avl (a, k)}}} & 0.59s\\
    \hline
    & \tmverbatim{merge: {\small{$\forall$k, n, a[n $\leqslant$ k + 2
    $\wedge$ k $\leqslant$ n + 2]. (Avl (a, n), Avl (a, k)) $\rightarrow$
    $\exists$i[n $\leqslant$ i $\wedge$ k $\leqslant$ i $\wedge$ i $\leqslant$
    n + k $\wedge$ i$\leqslant$max (k + 1, n + 1)].Avl (a, i)}}} & 0.93s\\
    \hline
    & \tmverbatim{remove: {\small{$\forall$n, a. a $\rightarrow$ Avl (a, n)
    $\rightarrow$ $\exists$k[n $\leqslant$ k + 1 $\wedge$\\
    0 $\leqslant$ k $\wedge$ k $\leqslant$ n].Avl (a, k)}}} & 0.38s\\
    \hline
    & Total time for \tmverbatim{add}, \tmverbatim{remove} and helper
    functions: & 4.92s\\
    \hline
  \end{tabular}\label{ExperimentsNum}
  \caption{Examples of types and inference times, with numerical constraints
  and existentials}
\end{table}

Table \ref{ExperimentsBounds} contains examples of static array (and matrix)
bounds checking. These tests originally date back to the Hongwei Xi's DML
system. They were selected, by {\cite{DSolvePaper}}, to demonstrate the
abilities of the Liquid Types system inference implementation
{\tmname{DSolve}}. The reported {\tmname{DSolve}} times come from
{\cite{DSolvePaper}} and {\cite{DSolveThesis}}. The example \tmverbatim{isort}
lists two times: the shorter time is inferred from a program without an unused
nested definition of \tmverbatim{vecswap}. We extracted the corresponding
computation as the example \tmverbatim{swap\_interval} in Table
\ref{ExperimentsNum}. The shorter time for the example \tmverbatim{simplex}
corresponds to a variant where some helper functions are separated out as
toplevel definitions, while the longer time is for a variant with a single
toplevel definition. For the example program \tmverbatim{gauss} without
assertions we again have two variants, but they both simplify the original
example. They remove a redundant level of nesting in a helper, nested
definition \tmverbatim{rowMax}. One of the variants requires passing a
non-default option \tmverbatim{-prefer\_bound\_to\_local} to guide inference,
the other slightly generalizes the algorithm by passing a different matrix
dimension in two places. Since our type system does not allow existential
types as function arguments, for the FFT examples we introduce an explicit
datatype \tmverbatim{Bounded}. The FFT example with assertion requires
non-default option \tmverbatim{-same\_with\_assertions}. On the whole, we
believe that the approach taken by {\tmname{InvarGenT}} shows promise in this
comparison. It is clear that {\tmname{InvarGenT}} faces increasing
difficulties when analysing more complex programs. We propose remedies to the
issues encountered as future work.\begin{table}[t]
  \
  
  \begin{tabular}{|l|l|l|}
    \hline
    {\tmem{Program}} & {\tmem{Inf. time}} & {\tmem{Time from
    {\cite{DSolvePaper}}}}\\
    \hline
    \tmverbatim{dotprod} & 0.05s & 0.31s\\
    \hline
    \tmverbatim{bcopy} & 0.03s & 0.15s\\
    \hline
    \tmverbatim{bsearch} & 0.07s & 0.46s\\
    \hline
    \tmverbatim{queen} & 0.42s & 0.7s\\
    \hline
    \tmverbatim{isort} & 0.3s, 0.37s & 0.88s\\
    \hline
    \tmverbatim{tower} no assertions & 0.84s & $\infty$\\
    \hline
    \tmverbatim{tower} with assertion & 3.93s & 3.33s\\
    \hline
    \tmverbatim{matmult} & 0.34s & 1.79s\\
    \hline
    \tmverbatim{heapsort} & 2.34s & 0.53s\\
    \hline
    \tmverbatim{fft} no assertions & 36.4s* & ?\\
    \hline
    \tmverbatim{fft} with assertion & 37.5s*,** & 9.13s\\
    \hline
    \tmverbatim{simplex} & 8.1s*, 31.4s & 7.73s\\
    \hline
    \tmverbatim{gauss} no assertions & 2.66s*, 1.02s*,*** & ?\\
    \hline
    \tmverbatim{gauss} with assertion & 2.72s & 3.17s\\
    \hline
  \end{tabular}
  
  * Slight meaning-preserving modification of test program
  
  ** Needs a non-default option \tmverbatim{-same\_with\_assertions}
  
  *** Needs a non-default option
  \tmverbatim{-prefer\_bound\_to\_local}\label{ExperimentsBounds}
  \caption{Examples of inference times, bound checking}
\end{table}

The code of examples from Tables \ref{Experiments}, \ref{ExperimentsNum} and
\ref{ExperimentsBounds} can be found in Appendix \ref{ExamplesCode}.

\section{{\tmname{InvarGenT}} Failure Cases}\label{FailureCases}

Type inference for the type system underlying {\tmname{InvarGenT}} is
undecidable. Constraint abduction itself is not known to be decidable. Rather
than trying to enumerate all possible answers, we devised constraint abduction
algorithms that only consider answers of restricted forms. The
\tmverbatim{vary} example from {\cite{PracticalGADTsInfer}} fails due to such
limitation. However, the practical problems we encountered are not of such
nature. Each subsection below describes a class of problematic cases,
exhausting the types of inference failures we encountered. Near discuss the
prospects of improving the situation.

\subsection{The Need to Expand Pattern Variables}

The one function from {\cite{PracticalGADTsInfer}} that needed modification to
be type-inferred is \tmverbatim{fd\_comp}. Take a look at the fragment of
original definition that needed modification:
\begin{tmcode}
let fd_comp = fun fd1 fd2 ->
  let o = fun f g x -> f (g x) in
  match fd1 with
    | FDI -> fd2
    | FDC b -> ...
\end{tmcode}
The modified form:
\begin{tmcode}
let fd_comp = fun fd1 fd2 ->
  let o = fun f g x -> f (g x) in
  match fd1 with
    | FDI ->
      (match fd2 with | FDI -> fd2 | FDC _ -> fd2 | FDG _ -> fd2)
    | FDC b -> ...
\end{tmcode}
Type inference needs to know the structure of both arguments to be able to
infer the resulting type. If this problem proves cumbersome, we can devise
heuristics for automatic expansion of pattern variables.

\subsection{Constraints Shared by Constructors of a Datatype}

The above problem is more pronounced in the numerical domain. Fortunately,
here it can have a principled solution. Consider the following list appending
example that does not type-check:
\begin{tmcode}
let rec append =
  function LNil -> (fun l -> l)
    | LCons (x, xs) -> (fun l -> LCons (x, append xs l))
\end{tmcode}
We can mitigate the problem by expanding the pattern, as in Example
\ref{BadFdComp}:
\begin{tmcode}
let rec append =
  function LNil -> (function LNil -> LNil | LCons (_,_) as l -> l)
    | LCons (x, xs) ->
      (function LNil -> LCons (x, append xs LNil)
        | LCons (_,_) as l -> LCons (x, append xs l))
\end{tmcode}
But we can also provide the ``hidden'' information explicitly, by either a
positive assertion, or a negative assertion as below:
\begin{tmcode}
let rec append =
  function
    | LNil ->
      (function l when (length l + 1)  assert false | l -> l)
    | LCons (x, xs) ->
      (function l when (length l + 1)  assert false
      | l -> LCons (x, append xs l))
\end{tmcode}
One can investigate a modification to the type system so that each datatype is
associated with an invariant common to all constructors of the datatype. In
case of lists, the shared invariant is that the length of a list is
non-negative. In the modified type system, when it becomes determined that a
pattern variable's type is a datatype, the corresponding invariant is made
available for inference.

\subsection{Negative Constraints in the Sort of Terms}

\begin{example}
  Consider the following variant of the \tmverbatim{head} example from
  {\cite{PracticalGADTsInfer}}:
  \begin{tmcode}
  datatype Z
datatype S : type
datatype List : type * num
datacons LNil : a. List(a, Z)
datacons LCons : a, b. a * List(a, b)  List(a, S b)

let head = function
  | LCons (x, _) -> x
  | LNil -> assert false
  \end{tmcode}
  We introduced the assertion to disallow the overly general type
  \tmverbatim{$\forall$a, b. List (a, b) $\rightarrow$ a}, which is the
  correct most general type for \tmverbatim{function LCons (x, \_) -> x},
  because the type system does not enforce exhaustiveness of pattern matching.
  The above function has type \tmverbatim{$\forall$a, b. List (a, S b)
  $\rightarrow$ a}, but type inference fails to find it. The problem stems
  from the inability to express disequations $t_1 \neq t_2$ in the sort of
  terms as conjunctions of atoms. We already implemented a workaround in
  {\tmname{InvarGenT}} for the case of datatypes with only constant
  constructors.
\end{example}

\subsection{Insufficient Context to Infer Postconditions}

\begin{example}
  In the following variant of the \tmverbatim{bsearch2} example it turns out
  to be too hard to infer the full postcondition.
  \begin{tmcode}
  datatype Array : type * num
external let array_make :
  n, a [0n]. Num n  a  Array (a, n) = "fun a b -> Array.make a b"
external let array_get :
n, k, a [0k  k+1n]. Array (a, n)  Num k  a =
  "fun a b -> Array.get a b"
external let array_length :
  n, a [0n]. Array (a, n)  Num n = "fun a -> Array.length a"
datatype LinOrder
datacons LE : LinOrder
datacons GT : LinOrder
datacons EQ : LinOrder
external let compare : a. a  a  LinOrder =
  "fun a b -> let c = Pervasives.compare a b in
              if c  0 then GT else EQ"
external let equal : a. a  a  Bool = "fun a b -> a = b"
external let div2 : n. Num (2 n)  Num n = "fun x -> x / 2"

let bsearch key vec =
  let rec look key vec lo hi =
    eif lo 
        let m = div2 (hi + lo) in
        let x = array_get vec m in
        ematch compare key x with
          | LE -> look key vec lo (m + (-1))
          | GT -> look key vec (m + 1) hi
          | EQ -> eif equal key x then m else -1
    else -1 in
  look key vec 0 (array_length vec + (-1))
  \end{tmcode}
  We get the result type \tmverbatim{$\exists$n[0 $\leqslant$ n + 1].Num n}
  instead of \tmverbatim{$\exists$k[k $\leqslant$ n $\wedge$ 0 $\leqslant$ k +
  1].Num k}. The inference of the intended type succeeds after we introduce an
  appropriate assertion, e.g. \tmverbatim{assert num -1 }. Alternatively, we
  can include a use case for \tmverbatim{bsearch} where the full postcondition
  is required. It would be challenging to guess the assertion or use-case
  automatically.
\end{example}

\subsection{Insufficient Search}

The need to pass options, in examples like the function \tmverbatim{leq} and
the function \tmverbatim{fft} with assertions, stems from the failure of our
search strategy for partial solutions across iterations of the main algorithm
(rather than within a single call to abduction). Our search can recover from
choices leading to contradictions in the following iteration, but not from
choices leading down blind alleys. The problem might be more serious than it
appears, because the heuristics (i.e. the default options) have been tuned in
favor of our test suite. The solution to the problem is {\tmem{beam search}}:
processing a fixed number of partial solutions, those with highest scores. A
score measures heuristically the quality of a partial solution: penalizing
complexity, rewarding locality of scope of parameters in an atom, etc.

\subsection{Nested Definitions with Tied Postconditions}

The variant of the \tmverbatim{gauss} program without assertions and with more
nesting illustrates a problem with nested definitions introducing existential
types. To arrive at the postcondition of the inner definition, we need to
propagate information from the use-site of the outer definition. The
postconditions are tied in the sense that the postcondition of the outer
definition depends on the postcondition of the inner definition. We need to
further investigate the issue to see what can be done.

\chapter{Conclusions}\label{ChapRelWork}

In this chapter, we summarize our work in its broader context.

\section{Related Work}\label{SecRelWork}

Perhaps Xi and Pfenning {\cite{DML99}} can be considered the earliest work on
GADTs with invariants over a constraint domain; see Section \ref{DMLSection}.
Its numerical constraint language contains {\tmem{min}} and {\tmem{max}}
operations enabling e.g. a concise definition of AVL trees datatype as in
Chapter \ref{Intro}. Its presentation of existential types is similar to ours.
{\cite{DML99}} opted to automate existential type elimination by
{\tmem{A-translation}}: $A_{\tmop{tr}} (e_1 e_2) = \left(
\text{{\tmstrong{let}}} x = A_{\tmop{tr}} (e_1)  \text{{\tmstrong{in}}} 
\text{{\tmstrong{let}}} y = A_{\tmop{tr}} (e_2)  \text{{\tmstrong{in}}} x y
\right)$. Thanks to the modified {\tmname{App}} rule, performing A-translation
during the normalization step {\tmname{ExIntro}} would make strictly more
programs typeable in $\tmop{MMG}_{\exists} (X)$. The DML language from
{\cite{DML99}} and its successor the ATS language do not perform type
inference for recursive functions.

We base our type system on the $\tmop{HMG} (X)$ type system from Simonet and
Pottier {\cite{simonet-pottier-hmg-toplas}}; see Section \ref{HMGSection}. In
the tradition of the Milner-Mycroft type system (see Henglein
{\cite{Henglein93}}), we drop the type specifications on recursive definitions
from program terms. We also naturally restrict $\tmop{HMG} (X)$ by limiting
the user-specified and inferred invariant constraints to use conjunction as
the only logical connective. We call the resulting type system $\tmop{MMG}
(X)$.

The traditional framework for loop invariant generation of Cousot and Cousot
{\cite{Cousot77}} inspired the iterative aspect of our solver. Mycroft's
modification in {\cite{Mycroft84}} of algorithm $\mathcal{W}$ to polymorphic
recursion is also iterative, but it solves each recursive definition
separately. We solve all nested definitions jointly in a single iteration.

Initially we were only aware of the work in Knowles and Flanagan
{\cite{KnowlesF07}} in the framework of {\tmem{refinement types}}, which
applies Dijkstra's weakest precondition calculus to general refinement types.
An interesting question is whether work similar to ours could be done by
application of the weakest precondition calculus to the Hoare logic of
Regis-Gianas and Pottier {\cite{regis-gianas-pottier-08}}, with the conditions
inserted by type inference. Work on {\tmem{Liquid Types}} by Rondon, Kawaguchi
and Jhala {\cite{DSolvePaper}} is a continuation of this approach closer to
{\tmname{InvarGenT}} in that it does not use weakest precondition calculus
explicitly; see Section \ref{LiquidSection}. {\tmem{Abstract Refinement
Types}} by Vazou, Rondon and Jhala {\cite{AbstractLiquid}} is a continuation
of Liquid Types in an interesting direction beyond the scope of current
{\tmname{InvarGenT}}, opening an area for future work.

Early work on applying abduction to type inference was an inspiration for us.
The work by Sulzmann, Schrijvers and Stuckey {\cite{AbdEADTs}} and
{\cite{SulzmannAbdGADTs}} closely relates to our work in their application of
fully-maximal constraint abduction to GADTs type inference. But without
annotations, this would ``throw the baby out with the bathwater'' by
rejecting, as not having an intuitive type, for example any variant of
\tmverbatim{eval} that computes for pairs: \tmverbatim{datacons Pair :
$\forall$a, b. Term a * Term b $\longrightarrow$ Term (a, b)},
\tmverbatim{datacons Fst : $\forall$a, b. Term (a, b) $\longrightarrow$ Term
a} and \tmverbatim{datacons Snd : $\forall$a, b. Term (a, b) $\longrightarrow$
Term b}. The corresponding implications do not have fully maximal answers.
{\cite{AbdEADTs}} and {\cite{SulzmannAbdGADTs}} use Herbrandization but solve
the resulting constraints in the free algebra of terms, which we find
problematic (see Section \ref{AbdUniVars}).

Maher {\cite{AbductionMaher}} and Maher and Huang {\cite{AbductionSolvMaher}}
introduce fully maximal constraint abduction; see Section \ref{HerbSection}.
Our algorithm for injective terms with multi-sorted {\tmem{alien subterms}} is
extensible to any constraint domain, given algorithms for the new domain, as
long as the domains do not interact. Allowing the domains to interact would
require considerable work, Baader and Schulz {\cite{UnificationBaader}} might
be relevant. Abduction algorithm for the term algebra is provided in
{\cite{AbductionSolvMaher}}, although further work driven by practical issues
was needed. The linear arithmetic constraint abduction in Maher
{\cite{AbductionLinArithMaher}} turned out not to be useful as a basis for an
algorithm, our algorithm is novel. Maher {\cite{HeytingAbdMaher}} studies
constraint domains where a (single) most general Simple Constraint Abduction
answer exists, and gives a closed formula for SCA answer for the case of
Boolean lattices.

The work in Unno and Kobayashi {\cite{InterpolationInfer}} can be seen as
extending Knowles and Flanagan {\cite{KnowlesF07}} with reasoning by Boolean
cases. Their programming language and type system is in several ways less
expressive than the ML language with polymorphic recursion and the full GADTs
type system: no inductive types (and therefore no pattern matching),
refinement predicates over integers only instead of over arbitrary domains
including types.

The recent {\tmem{counterexample-guided abstraction refinement}} work of Zhu
and Jagannathan {\cite{LightDML}} (building on Kobayashi, Sato and Unno
{\cite{CEGAR}}) pushes the envelope on both expressiveness and efficiency of
inference of a refinement types based approach. It is further developed in
Zhu, Nori and Jagannathan {\cite{SpecLearnArrays}}, the implementation is
called {\tmname{SpecLearn}}. The work includes features beyond the scope of
{\tmname{InvarGenT}}, for example effect tracking. Wrt. invariant inference,
{\tmname{SpecLearn}} is similar to {\tmname{InvarGenT}} in that it starts with
coarse invariants and refines them. However, rather than deriving the
invariants mostly from the definitions they describe, {\tmname{SpecLearn}}
generates test cases to refine the invariants. It still cannot solve, for
example, the AVL trees inference task, without assertions in the source code.
{\tmname{SpecLearn}} is slower than {\tmname{InvarGenT}}, at least for those
of the tests reported in {\cite{SpecLearnArrays}} which belong to our test
suite (Table \ref{ExperimentsBounds}).

Schrijvers and Bruynooghe {\cite{ADTReconstruct}} shares the objective with
our work: softening the learning curve and facilitating rapid prototyping. It
reconstructs algebraic data types. An interesting direction of future work
would be to reconstruct GADTs using the mechanisms already present in
{\tmname{InvarGenT}}. In fact, {\tmname{InvarGenT}} reconstructs GADT
constructors that serve as existential types.

There is a surge of work on type inference for GADTs over the last decade, not
contributing to our approach. Works such as Pottier and R{\'e}gis-Gianas
{\cite{StratifiedGADTs}} (older), Schrijvers, Peyton Jones, Sulzmann and
Vytiniotis {\cite{ComplDecidableGADTs}}, Lin and Sheard
{\cite{PointwiseGADTs}} (see also {\cite{PracticalGADTsInfer}}) modify the
GADTs type system to make it more amenable to type inference (rejecting some
reasonable programs as untypeable), and develop less declarative inference
algorithms. These works also do not allow other domains (than the free term
algebra) to express invariants, but Schrijvers, Peyton Jones, Sulzmann  and 
Vytiniotis {\cite{OutsideIn}} extends the {\tmname{OutsideIn}} system to
arbitrary domains, see Section \ref{OutsideInSection}.
{\cite{PracticalGADTsInfer}} and {\cite{ADTReconstruct}} stand out from our
point of view as they handle type inference for polymorphic recursion:
{\cite{PracticalGADTsInfer}} by iteration, see Section \ref{PointwiseSection},
and {\cite{ADTReconstruct}} using an algorithm by Henglein.

{\tmem{Inductive loop invariants}} from the recent work Dillig, Dillig, Li and
McMillan {\cite{InductLoopInv}} are closely related to fully maximal joint
constraint abduction answers. However, {\cite{InductLoopInv}} employs a richer
language of answers, including implications. We decided to limit solved forms,
thus preconditions and postconditions, to conjunctions of atoms, to not tax
with complexity both algorithms and users. Exploring {\cite{InductLoopInv}} in
context of {\tmname{InvarGenT}} may be fruitful, by improving on
{\tmname{InvarGenT}}'s abduction algorithm for linear arithmetic. One
improvement is to perform quantifier elimination of universal variables (as in
{\cite{InductLoopInv}}) when they cannot be substituted out. An improvement
would be constraint abduction that generates {\tmem{min}} and {\tmem{max}}
relations, which in current {\tmname{InvarGenT}} are only introduced by
constraint generalization, but here {\cite{InductLoopInv}} does not help
directly. {\cite{InductLoopInv}} would gain by incorporating ideas from our
abduction algorithm.

A related work by Deepak Kapur {\cite{LoopQuantElim}} is part of an
interesting research program of Deepak Kapur and Enric Rodr{\'i}guez
Carbonell, for example {\cite{LoopQuantElim}}, {\cite{RodriKapurICTAC04}},
{\cite{RodriKapurJSC07}}. {\cite{LoopQuantElim}} generates loop invariants for
imperative programs, by quantifier elimination. The variables that are not
eliminated are the invariant parameters. In fact, the generated invariant is
an abduction answer to the corresponding invariant inference constraint;
compare how {\tmname{InvarGenT}} uses abduction to solve for preconditions. On
the other hand, {\cite{RodriKapurICTAC04}} uses constraint generalization,
iterating it until convergence. This is in effect how {\tmname{InvarGenT}}
solves for postconditions. {\cite{RodriKapurICTAC04}} and
{\cite{RodriKapurJSC07}} use polynomial equations as the language of
invariants. This suggests, as future work, that including in
{\tmname{InvarGenT}} a domain of polynomial equations.

In case of the free algebra of terms, constraint generalization reduces to
anti-unification. Anti-unification was first introduced by Plotkin
{\cite{AntiUnifPlotkin}} and Reynolds {\cite{AntiUnifReynolds}}. Bulychev,
Kostylev and Zakharov {\cite{AntiUnifInv}} is a recent work on
anti-unification, with an example application to invariant inference. Our
constraint generalization algorithm for linear inequalities has inspirations
from Fukuda, Liebling  and  L{\"u}tolf {\cite{ConvexHull}}.

\section{Future Work}\label{SecFutWork}

Let us collect together proposed directions for future work.
\begin{itemize}
  \item Implement beam search, where several answers are maintained and
  inferior answers (less general precondition, less specific postcondition),
  or answers with worse heuristic score (more complex, etc.), are dropped.
  
  \item Address other issues discussed in Section \ref{FailureCases}: handle
  constraints shared by constructors of a datatype, handle use-site constraint
  propagation for nested definitions with tied postconditions.
  
  \item Explore extending the $\tmop{MMG}_{\exists} (X)$ type system to handle
  the use of invariants in arguments of higher-order functions, as in Rondon,
  Kawaguchi and Jhala {\cite{DSolvePaper}} and especially in abstract
  refinement types of Vazou, Rondon and Jhala {\cite{AbstractLiquid}}.
  
  \item Work on error reporting for {\tmname{InvarGenT}}. Trace use-site
  location and program path location associated with implications, separately;
  see Zhu and Jagannathan {\cite{LightDML}}.
  
  \item Integrate {\tmname{InvarGenT}} with an IDE.
  
  \item Add more constraint domains to {\tmname{InvarGenT}}.
\end{itemize}

\section{Summary}

We presented the type inference problem for $\tmop{MMG} (X)$, a Milner-Mycroft
style variant of the $\tmop{HMG} (X)$ type system without subtyping, as
satisfaction of second order constraints over a multi-sorted domain. We
provided a minimal extension $\tmop{MMG}_{\exists} (X)$ of this type system
that enables inference and easy use of existential types. Although insertions
of introduction and elimination of existential types are not automated by the
inference process, they are seamlessly integrated into expressions. We
demonstrated several use cases using the {\tmname{{\tmname{InvarGenT}}}}
system.

Our Joint Constraint Abduction under Quantifier Prefix algorithm builds on
the fully maximal Simple Constraint Abduction answers algorithm from Maher and
Huang {\cite{AbductionSolvMaher}}, extended to guess equalities of parameters.
Thanks to aggregating answers during JCA, it can find answers to some cases of
joint problems where it would fail to solve some of the implications
independently. Our SCA algorithm for linear arithmetic is novel.

Consider programs not using the numerical sort. Undecidability of type
inference for polymorphic recursion does not enforce an unbounded number of
iteration steps of the main algorithm, because there can be infinitely many
(maximally general but not fully maximal) term abduction answers. We
encountered examples that need two iteration steps: one to generate a
solution, and the following step to discard a wrong solution and accept the
correct one. However, with numerical sort constraints, a series of programs
can be written needing unbounded number of iteration steps.

We define the Constraint Generalization problem. In case of free terms it is
equivalent to anti-unification and in case of linear equations and
inequalities it is equivalent to finding extended convex hull. In the former
case, we extend the notion to Abductive Constraint Generalization, providing
more specific answers. As we do for abduction, we provide a
combination-of-domains algorithm for constraint generalization.

We implemented an algorithm solving for predicate variables of the existential
second order constraints generated for our type system. It uses abduction to
find requirements on invariants, augmented by constraint generalization to
find postconditions.

The inference times in {\tmname{InvarGenT}} are sufficient for interactive use
with toplevel definitions of small-to-moderate size.

\begin{thebibliography}{10}
  \bibitem[1]{AssertionGraphs}Tilak~Agerwala  and  Jayadev~Misra.{\newblock}
  Assertion graphs for verifying and synthesizing programs.{\newblock}
  Technical Report , University of Texas Technical Report 83, 1978.{\newblock}
  
  \bibitem[2]{Cayenne}Lennart~Augustsson.{\newblock} Cayenne -- a language
  with dependent types.{\newblock} In \tmtextit{Proceedings of the Third ACM
  SIGPLAN International Conference on Functional Programming}, ICFP '98, 
  pages  239--250. New York, NY, USA, 1998. ACM.{\newblock}
  
  \bibitem[3]{UnificationBaader}Franz~Baader  and  Klaus U.~Schulz.{\newblock}
  Unification in the union of disjoint equational theories: combining decision
  procedures.{\newblock} In \tmtextit{Proceedings of the 11th International
  Conference on Automated Deduction: Automated Deduction}, CADE-11,  pages 
  50--65. London, UK, 1992. Springer-Verlag.{\newblock}
  
  \bibitem[4]{ArithQuantElim}Sergey~Berezin, Vijay~Ganesh  and  David
  L.~Dill.{\newblock} An online proof-producing decision procedure for
  mixed-integer linear arithmetic.{\newblock} In \tmtextit{Proceedings of the
  9th international conference on Tools and algorithms for the construction
  and analysis of systems}, TACAS'03,  pages  521--536. Berlin, Heidelberg,
  2003. Springer-Verlag.{\newblock}
  
  \bibitem[5]{IdrisCurrent}Edwin~Brady.{\newblock} Idris, a general-purpose
  dependently typed programming language: design and
  implementation.{\newblock} \tmtextit{Journal of Functional Programming},
  23:552--593, 9 2013.{\newblock}
  
  \bibitem[6]{IdrisEarly}Edwin~Brady, Christoph A.~Herrmann  and 
  Kevin~Hammond.{\newblock} Lightweight invariants with full dependent
  types.{\newblock} In \tmtextit{Proceedings of the Nineth Symposium on Trends
  in Functional Programming, TFP 2008, Nijmegen, The Netherlands, May 26-28,
  2008.},  pages  161--177. 2008.{\newblock}
  
  \bibitem[7]{AntiUnifInv}Peter~Bulychev, Egor~Kostylev  and 
  Vladimir~Zakharov.{\newblock} Anti-unification algorithms and their
  applications in program analysis.{\newblock} In  Amir~Pnueli,
  Irina~Virbitskaite  and  Andrei~Voronkov , editors , \tmtextit{Perspectives
  of Systems Informatics},  volume  5947  of \tmtextit{Lecture Notes in
  Computer Science},  pages  413--423. Springer Berlin / Heidelberg,
  2010.{\newblock}
  
  \bibitem[8]{CheneyHinzePhantomTypes}James~Cheney  and 
  Ralf~Hinze.{\newblock} A lightweight implementation of generics and
  dynamics.{\newblock} In \tmtextit{Proceedings of the 2002 ACM SIGPLAN
  Workshop on Haskell}, Haskell '02,  pages  90--104. New York, NY, USA, 2002.
  ACM.{\newblock}
  
  \bibitem[9]{DisunificationComon}Hubert~Comon.{\newblock} Disunification: a
  survey.{\newblock} In \tmtextit{Computational Logic: Essays in Honor of Alan
  Robinson},  pages  322--359. MIT Press, 1991.{\newblock}
  
  \bibitem[10]{Cousot77}Patrick~Cousot  and  Radhia~Cousot.{\newblock}
  Automatic synthesis of optimal invariant assertions: mathematical
  foundations.{\newblock} \tmtextit{SIGPLAN Notices}, 12(8):1--12, aug
  1977.{\newblock}
  
  \bibitem[11]{simulRigidEUnifUndec}Anatoli~Degtyarev  and 
  Andrei~Voronkov.{\newblock} Simultaneous rigid e-unification is
  undecidable.{\newblock} In  Hans Kleine~B{\"u}ning , editor ,
  \tmtextit{Computer Science Logic},  volume  1092  of \tmtextit{Lecture Notes
  in Computer Science},  pages  178--190. Springer Berlin Heidelberg,
  1996.{\newblock}
  
  \bibitem[12]{InductLoopInv}Isil~Dillig, Thomas~Dillig, Boyang~Li  and 
  Ken~McMillan.{\newblock} Inductive invariant generation via abductive
  inference.{\newblock} In \tmtextit{Proceedings of the 2013 ACM SIGPLAN
  International Conference on Object Oriented Programming Systems Languages
  and Applications}, OOPSLA '13,  pages  443--456. New York, NY, USA, 2013.
  ACM.{\newblock}
  
  \bibitem[13]{Flanagan06}Cormac~Flanagan.{\newblock} Hybrid type
  checking.{\newblock} \tmtextit{SIGPLAN Not.}, 41(1):245--256, jan
  2006.{\newblock}
  
  \bibitem[14]{TypeQualifiers}Jeffrey Scott~Foster.{\newblock} \tmtextit{Type
  Qualifiers: Lightweight Specifications to Improve Software
  Quality}.{\newblock} PhD dissertation, University of California, Berkeley,
  Department of Computer Science, December 2002.{\newblock}
  
  \bibitem[15]{RefinementFreemanPfenning}Tim~Freeman  and 
  Frank~Pfenning.{\newblock} Refinement types for ml (extended
  abstract).{\newblock} In \tmtextit{Proc. ACM SIGPLAN '91 Conf. on
  Programming Language Design and Implementation, Toronto, Ontario},  pages 
  268--277. ACM Press, June 1991.{\newblock}
  
  \bibitem[16]{ConvexHull}Komei~Fukuda, Thomas M.~Liebling  and 
  Christine~L{\"u}tolf.{\newblock} Extended convex hull.{\newblock} In
  \tmtextit{Proceedings of the 12th Canadian Conference on Computational
  Geometry, Fredericton, New Brunswick, Canada, August 16-19, 2000}, 
  volume~20,  pages  13--23. 2001.{\newblock}
  
  \bibitem[17]{AbstractStateGraphs}Susanne~Graf  and 
  Hassen~Sa{\"i}di.{\newblock} Construction of abstract state graphs with
  pvs.{\newblock} In \tmtextit{Proceedings of the 9th International Conference
  on Computer Aided Verification}, CAV '97,  pages  72--83. London, UK, UK,
  1997. Springer-Verlag.{\newblock}
  
  \bibitem[18]{Henglein93}Fritz~Henglein.{\newblock} Type inference with
  polymorphic recursion.{\newblock} \tmtextit{ACM Trans. Program. Lang.
  Syst.}, 15(2):253--289, 1993.{\newblock}
  
  \bibitem[19]{LoopQuantElim}Deepak~Kapur.{\newblock} Automatically generating
  loop invariants using quantifier elimination.{\newblock} In  Franz~Baader,
  Peter~Baumgartner, Robert~Nieuwenhuis  and  Andrei~Voronkov , editors ,
  \tmtextit{Deduction and Applications},  number  05431  in  Dagstuhl Seminar
  Proceedings,  pages  10--10. Dagstuhl, Germany, 2006. Internationales
  Begegnungs- und Forschungszentrum f{\"u}r Informatik (IBFI), Schloss
  Dagstuhl, Germany.{\newblock}
  
  \bibitem[20]{KnowlesF07}Kenneth W.~Knowles  and  Cormac~Flanagan.{\newblock}
  Type reconstruction for general refinement types.{\newblock} In
  \tmtextit{ESOP},  volume  4421  of \tmtextit{Lecture Notes in Computer
  Science},  pages  505--519. Springer, 2007.{\newblock}
  
  \bibitem[21]{CEGAR}Naoki~Kobayashi, Ryosuke~Sato  and 
  Hiroshi~Unno.{\newblock} Predicate abstraction and cegar for higher-order
  model checking.{\newblock} In \tmtextit{ACM SIGPLAN Notices},  volume~46, 
  pages  222--233. ACM, 2011.{\newblock}
  
  \bibitem[22]{PracticalGADTsInfer}Chuan-kai~Lin.{\newblock}
  \tmtextit{Practical type inference for the GADT type system}.{\newblock} PhD
  dissertation, Portland State University, Department of Computer Science,
  2010.{\newblock}
  
  \bibitem[23]{PointwiseGADTs}Chuan-kai~Lin  and  Tim~Sheard.{\newblock}
  Pointwise generalized algebraic data types.{\newblock} In
  \tmtextit{Proceedings of the 5th ACM SIGPLAN workshop on Types in language
  design and implementation}, TLDI '10,  pages  51--62. New York, NY, USA,
  2010. ACM.{\newblock}
  
  \bibitem[24]{ExTyADTs}Konstantin~L{\"a}ufer  and  Martin~Odersky.{\newblock}
  Polymorphic type inference and abstract data types.{\newblock} \tmtextit{ACM
  Trans. Program. Lang. Syst.}, 16(5):1411--1430, sep 1994.{\newblock}
  
  \bibitem[25]{MaherParamSubstitutions}Michael~Maher.{\newblock} On
  parameterized substitutions.{\newblock} Technical Report , IBM Research
  Report RC 16042, 1990.{\newblock}
  
  \bibitem[26]{AbductionLinArithMaher}Michael~Maher.{\newblock} Abduction of
  linear arithmetic constraints.{\newblock} In  Maurizio~Gabbrielli  and 
  Gopal~Gupta , editors , \tmtextit{Logic Programming},  volume  3668  of
  \tmtextit{Lecture Notes in Computer Science},  pages  174--188. Springer
  Berlin Heidelberg, 2005.{\newblock}
  
  \bibitem[27]{AbductionMaher}Michael~Maher.{\newblock} Herbrand constraint
  abduction.{\newblock} In \tmtextit{LICS '05: Proceedings of the 20th Annual
  IEEE Symposium on Logic in Computer Science},  pages  397--406. Washington,
  DC, USA, 2005. IEEE Computer Society.{\newblock}
  
  \bibitem[28]{HeytingAbdMaher}Michael~Maher.{\newblock} Heyting domains for
  constraint abduction.{\newblock} In \tmtextit{Proceedings of the 19th
  Australian Joint Conference on Artificial Intelligence: Advances in
  Artificial Intelligence}, AI'06,  pages  9--18. Berlin, Heidelberg, 2006.
  Springer-Verlag.{\newblock}
  
  \bibitem[29]{AbductionSolvMaher}Michael~Maher  and  Ge~Huang.{\newblock} On
  computing constraint abduction answers.{\newblock} In  Iliano~Cervesato,
  Helmut~Veith  and  Andrei~Voronkov , editors , \tmtextit{Logic for
  Programming, Artificial Intelligence, and Reasoning},  volume  5330  of
  \tmtextit{Lecture Notes in Computer Science},  pages  421--435. Springer
  Berlin / Heidelberg, 2008.{\newblock}
  
  \bibitem[30]{MartinLof}Per~Martin-L{\"o}f.{\newblock} Constructive
  mathematics and computer programming.{\newblock} In \tmtextit{Proc. Of a
  Discussion Meeting of the Royal Society of London on Mathematical Logic and
  Programming Languages},  pages  167--184. Upper Saddle River, NJ, USA, 1985.
  Prentice-Hall, Inc.{\newblock}
  
  \bibitem[31]{AbductionSequents}Marta Cialdea~Mayer  and 
  Fiora~Pirri.{\newblock} First order abduction via tableau and sequent
  calculi.{\newblock} \tmtextit{Logic Journal of IGPL}, 1(1):99--117,
  1993.{\newblock}
  
  \bibitem[32]{Milner78}Robin~Milner.{\newblock} A theory of type polymorphism
  in programming.{\newblock} \tmtextit{Journal of Computer and System
  Sciences}, 17:348--375, 1978.{\newblock}
  
  \bibitem[33]{Mycroft84}Alan~Mycroft.{\newblock} Polymorphic type schemes and
  recursive definitions.{\newblock} In \tmtextit{Proceedings of the 6th
  Colloquium on International Symposium on Programming},  pages  217--228.
  London, UK, UK, 1984. Springer-Verlag.{\newblock}
  
  \bibitem[34]{HMX}Martin~Odersky, Martin~Sulzmann  and 
  Martin~Wehr.{\newblock} Type inference with constrained types.{\newblock} In
  \tmtextit{Fourth International Workshop on Foundations of Object-Oriented
  Programming (FOOL)}. 1997.{\newblock}
  
  \bibitem[35]{AntiUnifPlotkin}Gordon D.~Plotkin.{\newblock} A note on
  inductive generalization.{\newblock} \tmtextit{Machine Intelligence},
  5:153--163, 1970.{\newblock}
  
  \bibitem[36]{StratifiedGADTs}Fran{\c c}ois~Pottier  and 
  Yann~R{\'e}gis-Gianas.{\newblock} Stratified type inference for generalized
  algebraic data types.{\newblock} In \tmtextit{Conference record of the 33rd
  ACM SIGPLAN-SIGACT symposium on Principles of programming languages}, POPL
  '06,  pages  232--244. New York, NY, USA, 2006. ACM.{\newblock}
  
  \bibitem[37]{AntiUnifReynolds}John C.~Reynolds.{\newblock} Transformational
  systems and the algebraic structure of atomic formulas.{\newblock} In 
  Bernard~Meltzer  and  Donald~Michie , editors , \tmtextit{Machine
  Intelligence},  volume~5,  pages  135--151. Edinburgh, Scotland, 1969.
  Edinburgh University Press.{\newblock}
  
  \bibitem[38]{RodriKapurICTAC04}Enric~Rodrguez-Carbonell  and 
  Deepak~Kapur.{\newblock} Program verification using automatic generation of
  invariants.{\newblock} In \tmtextit{1st International Colloquium on
  Theoretical Aspects of Computing (ICTAC'04)},  volume  3407  of
  \tmtextit{Lecture Notes in Computer Science},  pages  325--340.
  Springer-Verlag, 2005.{\newblock}
  
  \bibitem[39]{RodriKapurJSC07}Enric~Rodriguez-Carbonell  and 
  Deepak~Kapur.{\newblock} Generating all polynomial invariants in simple
  loops.{\newblock} \tmtextit{Journal of Symbolic Computation}, 42(4):0,
  2007.{\newblock}
  
  \bibitem[40]{DSolveThesis}Patrick~Rondon.{\newblock} \tmtextit{Liquid
  Types}.{\newblock} PhD dissertation, University of California, San Diego,
  Department of Computer Science, 2012.{\newblock}
  
  \bibitem[41]{DSolvePaper}Patrick Maxim~Rondon, Ming~Kawaguchi  and 
  Ranjit~Jhala.{\newblock} Liquid types.{\newblock} In \tmtextit{Proceedings
  of the ACM SIGPLAN 2008 Conference on Programming Language Design and
  Implementation, Tucson, AZ, USA, June 7-13, 2008},  pages  159--169.
  2008.{\newblock}
  
  \bibitem[42]{DSolveTechRep}Patrick Maxim~Rondon, Ming~Kawaguchi  and 
  Ranjit~Jhala.{\newblock} Liquid types.{\newblock} Technical Report ,
  2008.{\newblock}
  
  \bibitem[43]{regis-gianas-pottier-08}Yann~R{\'e}gis-Gianas  and  Fran{\c
  c}ois~Pottier.{\newblock} A Hoare logic for call-by-value functional
  programs.{\newblock} In \tmtextit{Proceedings of the Ninth International
  Conference on Mathematics of Program Construction (MPC'08)},  volume  5133 
  of \tmtextit{Lecture Notes in Computer Science},  pages  305--335. Springer,
  JUL 2008.{\newblock}
  
  \bibitem[44]{ADTReconstruct}Tom~Schrijvers  and 
  Maurice~Bruynooghe.{\newblock} Polymorphic algebraic data type
  reconstruction.{\newblock} In \tmtextit{Proceedings of the 8th International
  ACM SIGPLAN Conference on Principles and Practice of Declarative
  Programming, July 10-12, 2006, Venice, Italy},  pages  85--96.
  2006.{\newblock}
  
  \bibitem[45]{ComplDecidableGADTs}Tom~Schrijvers, Simon Peyton~Jones,
  Martin~Sulzmann  and  Dimitrios~Vytiniotis.{\newblock} Complete and
  decidable type inference for gadts.{\newblock} In \tmtextit{Proceedings of
  the 14th ACM SIGPLAN international conference on Functional programming},
  ICFP '09,  pages  341--352. New York, NY, USA, 2009. ACM.{\newblock}
  
  \bibitem[46]{Sheard04}Tim~Sheard.{\newblock} Languages of the
  future.{\newblock} In \tmtextit{In OOPSLA '04: Companion to the 19th annual
  ACM SIGPLAN conference on Object-oriented programming systems, languages,
  and applications},  pages  116--119. ACM Press, 2004.{\newblock}
  
  \bibitem[47]{simonet-pottier-hmg-toplas}Vincent~Simonet  and  Fran{\c
  c}ois~Pottier.{\newblock} A constraint-based approach to guarded algebraic
  data types.{\newblock} \tmtextit{ACM Transactions on Programming Languages
  and Systems}, 29(1), JAN 2007.{\newblock}
  
  \bibitem[48]{SulzmannAbdGADTs}Martin~Sulzmann, Tom~Schrijvers  and  Peter
  J.~Stuckey.{\newblock} Type inference for GADTs via Herbrand constraint
  abduction.{\newblock} Manuscript, July 2006.{\newblock}
  
  \bibitem[49]{AbdEADTs}Martin~Sulzmann, Tom~Schrijvers  and  Peter
  J.~Stuckey.{\newblock} Type inference via constraint abduction for
  EADTs.{\newblock} Manuscript, April 2006.{\newblock}
  
  \bibitem[50]{SulzmannAbdGADTsRep}Martin~Sulzmann, Tom~Schrijvers  and  Peter
  J.~Stuckey.{\newblock} Type inference for GADTs via Herbrand constraint
  abduction.{\newblock} Technical Report , K. U. Leuven Dept. of Computer
  Science, 2008.{\newblock} Report CW 507.{\newblock}
  
  \bibitem[51]{InterpolationInfer}Hiroshi~Unno  and 
  Naoki~Kobayashi.{\newblock} Dependent type inference with
  interpolants.{\newblock} In \tmtextit{Proceedings of the 11th ACM SIGPLAN
  conference on Principles and practice of declarative programming}, PPDP '09,
  pages  277--288. New York, NY, USA, 2009. ACM.{\newblock}
  
  \bibitem[52]{AbstractLiquid}Niki~Vazou, Patrick Maxim~Rondon  and 
  Ranjit~Jhala.{\newblock} Abstract refinement types.{\newblock} In
  \tmtextit{Programming Languages and Systems - 22nd European Symposium on
  Programming, ESOP 2013, Held as Part of the European Joint Conferences on
  Theory and Practice of Software, ETAPS 2013, Rome, Italy, March 16-24, 2013.
  Proceedings},  pages  209--228. 2013.{\newblock}
  
  \bibitem[53]{OutsideIn}Dimitrios~Vytiniotis, Simon L. Peyton~Jones,
  Tom~Schrijvers  and  Martin~Sulzmann.{\newblock} Outsidein(x) modular type
  inference with local assumptions.{\newblock} \tmtextit{J. Funct. Program.},
  21(4-5):333--412, 2011.{\newblock}
  
  \bibitem[54]{DMLThesis}Hongwei~Xi.{\newblock} \tmtextit{Dependent Types in
  Practical Programming}.{\newblock} PhD thesis , Department of Mathematical
  Sciences, Carnegie Mellon University, December 1998.{\newblock}
  
  \bibitem[55]{XiGRDTs}Hongwei~Xi, Chiyan~Chen  and  Gang~Chen.{\newblock}
  Guarded recursive datatype constructors.{\newblock} In \tmtextit{Proceedings
  of the 30th ACM SIGPLAN-SIGACT Symposium on Principles of Programming
  Languages}, POPL '03,  pages  224--235. New York, NY, USA, 2003.
  ACM.{\newblock}
  
  \bibitem[56]{DMLArrays}Hongwei~Xi  and  Frank~Pfenning.{\newblock}
  Eliminating array bound checking through dependent types.{\newblock} In
  \tmtextit{Proceedings of the ACM SIGPLAN 1998 Conference on Programming
  Language Design and Implementation}, PLDI '98,  pages  249--257. New York,
  NY, USA, 1998. ACM.{\newblock}
  
  \bibitem[57]{DML99}Hongwei~Xi  and  Frank~Pfenning.{\newblock} Dependent
  types in practical programming.{\newblock} In \tmtextit{Proceedings of the
  26th ACM SIGPLAN Symposium on Principles of Programming Languages},  pages 
  214--227. San Antonio, January 1999.{\newblock}
  
  \bibitem[58]{IndexedTypes}Christoph~Zenger.{\newblock} Indexed
  types.{\newblock} \tmtextit{Theor. Comput. Sci.}, 187(1-2):147--165, nov
  1997.{\newblock}
  
  \bibitem[59]{LightDML}He~Zhu  and  Suresh~Jagannathan.{\newblock}
  Compositional and lightweight dependent type inference for ml.{\newblock} In
  Roberto~Giacobazzi, Josh~Berdine  and  Isabella~Mastroeni , editors ,
  \tmtextit{Verification, Model Checking, and Abstract Interpretation}, 
  volume  7737  of \tmtextit{Lecture Notes in Computer Science},  pages 
  295--314. Springer, 2013.{\newblock}
  
  \bibitem[60]{SpecLearnArrays}He~Zhu, Aditya V.~Nori  and 
  Suresh~Jagannathan.{\newblock} Dependent array type inference from
  tests.{\newblock} In  Deepak~D'Souza, Akash~Lal  and  Kim Guldstrand~Larsen
  , editors , \tmtextit{VMCAI '15: Verification, Model Checking and Abstract
  Interpretation},  volume  8931  of \tmtextit{Lecture Notes in Computer
  Science},  pages  412--430. Springer, January 2015.{\newblock}
  
  \bibitem[61]{AntiUnifAlg}Bjarte M.~{\O}stvold.{\newblock} A functional
  reconstruction of anti-unification.{\newblock} Technical Report , Norwegian
  Computing Center, Oslo, Norway, 2004.{\newblock}
\end{thebibliography}

\appendix\chapter{Proofs}\label{ChapProofs}

\section{The Type System}

\subsection{The Logic of Constraints}

We work with a first order language $\mathcal{L}_1$ interpreted in a model
$\mathcal{M}$, the language of constraints for our type inference problem. We
assume that the logic is multi-sorted, but the sorts are combined only via a
sort of finite trees, whose leaves can belong to other sorts. I.e., there is a
single sort $s_{\tmop{type}}$ whose terms can contain subterms from other
sorts and for any its terms \ $C \Rightarrow f (\overline{s_i}) \dot{=} g
(\overline{t_i})$ implies $f = g$ and $C \Rightarrow \wedge_i s_i \dot{=}
t_i$. The sorts $\mathcal{L}_s, \mathcal{M}_s$ for $s \neq s_{\tmop{type}}$
are single-sorted logics and $\mathcal{L}_1 = \dot{\cup}_s \mathcal{L}_s$.

Let $\rho$ be an interpretation of types, that is an assignment of elements of
$\mathcal{M}$ to variables in the corresponding sort, extended homomorphically
to terms in the standard way. In the main text, we used $R$ rather than $\rho$
to avoid unnecessary formality. For $\Phi \in \mathcal{L}_1$, let
$\mathcal{M}, \rho \vDash \Phi$ denote the interpretation of a formula $\Phi$
in the model $\mathcal{M}$ under the interpretation $\rho$, in the standard
way, for example $\mathcal{M}, \rho \vDash \pi (t)$ if and only if $\pi (\rho
(t))$ holds in $\mathcal{M}$, where predicate symbol $\pi$ in $\mathcal{L}_1$
corresponds to predicate $\pi$ in $\mathcal{M}$, etc.

Of the model $\mathcal{M}$ of $\mathcal{L}_1$ we require the following. For
the sort of types $s_{\tmop{type}}$:
\begin{enumerate}
  \item {\tmem{Type conservation.}} For any function symbols $f, g$ such that
  $f \neq g$, and arbitrary $\bar{s}, \bar{t}$, there is no interpretation
  $\rho$ such that $\mathcal{M}, \rho \vDash f (\bar{s}) \dot{=} g (\bar{t})$.
  
  \item {\tmem{Free generation.}} For any $f \in \{\rightarrow\} \cup
  \bar{\varepsilon}$ and arbitrary $\overline{s_i}_{1 \leqslant i \leqslant
  \tmop{ar} (f)}, \overline{t_i}_{1 \leqslant i \leqslant \tmop{ar} (f)}$, for
  any interpretation $\rho$, if $\mathcal{M}, \rho \vDash f (\bar{s}) \dot{=}
  f (\bar{t})$, then $\mathcal{M}, \rho \vDash \wedge_{1 \leqslant i \leqslant
  \tmop{ar} (f)} s_i \dot{=} t_i$.
\end{enumerate}
and every sort is nonempty.

Let $\mathcal{L}$ be $\mathcal{L}_1$ with: a set of unary predicates $\chi
(\cdummy)$, which stand for invariants of recursive definitions in the
constraints we will derive for type inference problems. And a set of binary
predicates $\chi_K (\cdummy, {\nobreak} \cdot)$, which will be put as
constraints of data constructors $K$ when we introduce inferred existential
types. We call $\chi$ and $\chi_K$ {\tmem{predicate variables}}. Let
$\tmop{PV}^1 (\cdummy)$, resp. $\tmop{PV}^2 (\cdummy)$ be the set of unary,
resp. binary predicate variables in any expression, and $\tmop{PV} (\Phi) =
\tmop{PV}^1 (\Phi) \cup \tmop{PV}^2 (\Phi)$. We define {\tmem{solved form
formulas}} to be existentially quantified conjunctions of atoms $\exists
\bar{\alpha} .A$ without predicate variables. Let $\delta, \delta^{\prime}$ be
fixed variables of sort $s_{\tmop{type}}$.

\begin{definition}
  \label{InterprPredVars}For a formula $\Phi$, let $\bar{\chi} = \tmop{PV}^1
  (\Phi)$, resp. $\overline{\chi_K} = \tmop{PV}^2 (\Phi)$, and let
  $\overline{\chi (\tau_{\chi, k})}$, resp. $\overline{\chi_K (\tau_{K, k},
  \tau_{K, k}^{\prime})}$ be all occurrences of $\chi$, resp. $\chi_K$ in
  $\Phi$. We call an assignment $\mathcal{I}= [\bar{\chi} \assign
  \overline{\exists \bar{\alpha}_{\chi} .F_{\chi}} ; \overline{\chi_K} \assign
  \overline{\exists \bar{\alpha}_K .F_K}]$ an {\tmem{interpretation of
  predicate variables for $\Phi$}} when
  \begin{enumerate}
    \item $\overline{\exists \bar{\alpha}_i .F_i} \overline{\exists
    \bar{\alpha}_j .F_j}$ are solved form formulas,
    
    \item $\delta \bar{\alpha}_{\chi} \# \tmop{FV} (\wedge_k \tau_{\chi, k})$
    and $\delta \delta^{\prime} \bar{\alpha}_K \# \tmop{FV} (\wedge_k \tau_{K,
    k} \wedge_k \tau_{K, k}^{\prime})$,
    
    \item for every variable $\beta \in \tmop{FV} (F_{\chi}) \backslash \delta
    \bar{\alpha}_{\chi}$, there is a quantifier that binds $\beta$ at every
    position of $\chi (\tau_{\chi, k})$ in $\Phi$,
    
    \item $\tmop{FV} (F_K) \subseteq \delta \delta^{\prime} \bar{\alpha}_K$.
  \end{enumerate}
\end{definition}

For an interpretation of predicate variables $\mathcal{I}= [\bar{\chi} \assign
\overline{\exists \bar{\alpha}_{\chi} .F_{\chi}} ; \overline{\chi_K} \assign
\overline{\exists \bar{\alpha}_K .F_K}]$, define the corresponding
substitution of predicate variables in a formula $\Phi$ by:
\begin{eqnarray*}
  \mathcal{I} (\chi (\tau_{\chi})) & = & \exists \bar{\alpha}_{\chi} .F_{\chi}
  [\delta \assign \tau_{\chi}]\\
  \mathcal{I} (\chi_K (\tau_{\chi}, \tau_{\chi}^{\prime})) & = & \exists
  \bar{\alpha}_{\chi} .F_{\chi} [\delta \assign \tau_{\chi} ; \delta^{\prime}
  \assign \tau_{\chi}^{\prime}]\\
  \mathcal{I} (a) & = & a \text{ \ for atom } a \in \mathcal{L}_1\\
  \mathcal{I} (\mathcal{Q}. \Phi) & = & \mathcal{Q}.\mathcal{I} (\Phi) \text{
  \ for any quantifier prefix } \mathcal{Q}\\
  \mathcal{I} (\Phi_1 \odot \Phi_2) & = & \mathcal{I} (\Phi_1) \odot
  \mathcal{I} (\Phi_2) \text{ \ for any logical connective } \odot
\end{eqnarray*}
Define a statement $\mathcal{M}, \mathcal{I}, \rho \vDash \Phi$ by:
$\mathcal{I}$ is an interpretation of predicate variables for $\Phi$, $\rho$
is an interpretation of types, and $\mathcal{M}, \rho \vDash \mathcal{I}
(\Phi)$. Define $\mathcal{M}, \mathcal{I} \vDash \Phi$ as: for all
interpretations of types $\rho$, $\mathcal{M}, \mathcal{I}, \rho \vDash \Phi$.
Define $\mathcal{M} \vDash \Phi$ as: for all interpretations of predicate
variables $\mathcal{I}$ for $\Phi$, $\mathcal{M}, \mathcal{I} \vDash \Phi$.
Sometimes we write $\mathcal{I} \vDash \Phi$, resp. $\vDash \Phi$, instead of
$\mathcal{M}, \mathcal{I} \vDash \Phi$, resp. $\mathcal{M} \vDash \Phi$, since
the model is fixed. We write $\mathcal{I}, C \vDash \Phi$, resp. $C \vDash
\Phi$, for $\mathcal{I} \vDash C \Rightarrow \Phi$, resp. $\vDash C
\Rightarrow \Phi$.

We say that a formula $\Phi$ is {\tmem{satisfiable}}, if and only if there
exists an interpretation of predicate variables $\mathcal{I}$ for $\Phi$, such
that $\mathcal{I} \vDash \exists \tmop{FV} (\Phi) . \Phi$. As seen above, we
extend the notion of substitution to handle predicate variable atoms, where
the replacement of each occurrence of a variable depends on the argument of
that variable. For interpretations of predicate variables $\mathcal{I}_1,
\mathcal{I}_2$ with disjoint domains, we write their composition
$\mathcal{I}_1 \mathcal{I}_2 (\cdummy) =\mathcal{I}_1 (\mathcal{I}_2
(\cdummy))$.

Above we in effect introduce a Henkin semantics for existential second order
logic, tailored to our needs of invariant and postcondition inference.

\subsection{The GADT Type System}\label{GADTProofs}

Set $\Delta \assign \exists \bar{\beta} [D] . \Gamma$ and $\Delta^{\prime}
\assign \exists \bar{\beta}^{\prime} [D^{\prime}] . \Gamma^{\prime}$ such that
$\bar{\beta} \# \tmop{FV} (\Gamma^{\prime})$, $\bar{\beta}^{\prime} \#
\tmop{FV} (\Delta)$ and $\bar{\beta}^{\prime} \#C$. Let $C \vDash
\Delta^{\prime} \leqslant \Delta$ denote $C \wedge D^{\prime} \vDash \exists
\bar{\beta} . (D \wedge_{x \in \tmop{Dom} (\Gamma)} \Gamma (x) \dot{=}
\Gamma^{\prime} (x))$ when $\tmop{Dom} (\Gamma) = \tmop{Dom}
(\Gamma^{\prime})$, and otherwise a falsehood (compare lemma 3.5 of
{\cite{simonet-pottier-hmg-toplas}}). Let $\Delta \times \Delta^{\prime}$
denote $\exists \bar{\beta} \overline{\beta^{\prime}} [D \wedge D^{\prime}] .
\Gamma \dot{\cup} \Gamma^{\prime}$, and $\exists \bar{\beta}^{\prime}
[D^{\prime}] \Delta$ denote $\exists \bar{\beta} \overline{\beta^{\prime}} [D
\wedge D^{\prime}] . \Gamma$.

\begin{proposition}
  \label{EnvFragProp}Properties of environment fragments (see
  {\cite{simonet-pottier-hmg-toplas}} lemma 3.15).
  \begin{description}
    \item[f-Hide] $\vDash \Delta \leqslant \exists \bar{\alpha} . \Delta$.
    
    \item[f-Imply] $C_1 \Rightarrow C_2 \vDash [C_1] \Delta \leqslant [C_2]
    \Delta$.
    
    \item[f-Enrich] $C \Rightarrow \Delta_1 \leqslant \Delta_2 \vDash [C]
    \Delta_1 \leqslant [C] \Delta_2$.
    
    \item[f-Ex] $\forall \bar{\alpha} . \Delta_1 \leqslant \Delta_2 \vDash
    (\exists \bar{\alpha} . \Delta_1) \leqslant (\exists \bar{\alpha} .
    \Delta_2)$.
    
    \item[f-And] $\Delta_1 \leqslant \Delta_2 \vDash \Delta \times \Delta_1
    \leqslant \Delta \times \Delta_2$.
  \end{description}
\end{proposition}

\begin{proposition}
  Constructor $K \colons \forall \bar{\alpha} \bar{\beta} [D] . \tau_1 \times
  \ldots \times \tau_n \rightarrow \varepsilon (\bar{\alpha})$ where $D =
  \exists \bar{\beta}^{\prime} .A$, is equivalent to $K \colons \forall
  \bar{\alpha} \overline{\gamma_i} [\exists \bar{\beta} \bar{\beta}^{\prime} .
  \overline{\gamma_i} \dot{=} \overline{\tau_i} \wedge A] . \gamma_1 \times
  \ldots \times \gamma_n \rightarrow \varepsilon (\bar{\alpha})$.
\end{proposition}

\begin{proposition}
  \label{collapse-ty-args}Constructors of the form $K \colons \forall
  \overline{\alpha_i} \bar{\beta} [D] . \tau_1 \times \ldots \times \tau_n
  \rightarrow \varepsilon (\overline{\alpha_i})$ where $D = \exists
  \bar{\beta}^{\prime} .A$, are equivalent to constructors of the form $K
  \colons \forall \alpha \bar{\beta} [\exists \overline{\alpha_i}
  \bar{\beta}^{\prime} . \alpha \dot{=} \alpha_1 \rightarrow \ldots
  \rightarrow \alpha_m \wedge A] . \gamma_1 \times \ldots \times \gamma_n
  \rightarrow \varepsilon (\alpha)$ when all uses of $\varepsilon (\tau_1,
  \ldots, \tau_m)$ are translated to $\varepsilon (\tau_1 \rightarrow \ldots
  \rightarrow \tau_m)$.
\end{proposition}

\begin{lemma}
  \label{Weakening}{\tmem{Weakening}} (patterns and expressions). Assume $C_1
  \vDash C_2$. If $C_2 \vdash p : \tau \longrightarrow \Delta$ (resp. $C_2,
  \Gamma \vdash e : \tau$, $C_2, \Gamma \vdash e : \sigma$) is derivable, then
  there exists a derivation of $C_1 \vdash p : \tau \longrightarrow \Delta$
  (resp. $C_1, \Gamma \vdash e : \tau$, $C_1, \Gamma \vdash e : \sigma$) of
  the same structure.
\end{lemma}

The lemma follows from transitivity of $\vDash$ ($A \vDash B$ and $B \vDash C$
imply $A \vDash C$) by induction on the structure of the derivation.

\begin{lemma}
  \label{MoreCstrs}If $\Sigma \subset \Sigma^{\prime}$ and $C \vdash p : \tau
  \longrightarrow \Delta$ (resp. $C, \Gamma \vdash e : \tau$, $C, \Gamma
  \vdash e : \sigma$) is derivable with constructors $\Sigma$, then the same
  derivation works with constructors $\Sigma^{\prime}$.
\end{lemma}

\begin{lemma}
  \label{CorrPat}{\tmem{Correctness}} (patterns). $\llbracket \vdash p
  \downarrow \tau \rrbracket \vdash p : \tau \longrightarrow \llbracket \vdash
  p \uparrow \tau \rrbracket$.
\end{lemma}

\begin{proof}
  By induction on the structure of $p$.
  \begin{itemize}
    \item Cases $0$, $1$ and $x$: follow directly from {\tmname{p-Empty}},
    {\tmname{p-Wild}} and {\tmname{p-Var}} respectively.
    
    \item Case $p_1 \wedge p_2$.
    \begin{enumerate}
      \item By the induction hypothesis, $\llbracket \vdash p_i \downarrow
      \tau \rrbracket \vdash p_i : \tau \longrightarrow \llbracket \vdash p_i
      \uparrow \tau \rrbracket$ for $i = 1, 2$.
      
      \item By weakening and {\tmname{p-And}} we have the goal.
    \end{enumerate}
    \item Case $K p_1 \ldots p_n$.
    \begin{enumerate}
      \item Let $\Sigma \ni K \colons \forall \bar{\alpha} \bar{\beta} [D] .
      \tau_1 \times \ldots \times \tau_n \rightarrow \varepsilon
      (\bar{\alpha})$.
      
      \item By the induction hypothesis, $\llbracket \vdash p_i \downarrow
      \tau_i \rrbracket \vdash p_i : \tau_i \longrightarrow \llbracket \vdash
      p_i \uparrow \tau_i \rrbracket$ for $i = 1, \ldots, n$.
      
      \item The {\tmname{p-Cstr}} rule says $\forall i (C \wedge D \vdash p_i
      : \tau_i \longrightarrow \Delta_i) /_{\text{p-Cstr}} C \vdash p :
      \varepsilon (\bar{\alpha}) \longrightarrow \exists \bar{\beta} [D]
      (\Delta_1 \times \ldots \times \Delta_n)$, where $\Delta_i \assign
      \llbracket \vdash p_i \uparrow \tau_i \rrbracket$. Applying it to (2) we
      get $C \vdash p : \varepsilon (\bar{\alpha}) \longrightarrow \exists
      \bar{\beta} [D] (\Delta_1 \times \ldots \times \Delta_n)$ as long as $C
      \wedge D \vDash \llbracket \vdash p_i \downarrow \tau_i \rrbracket$.
      
      \item Let $\bar{\alpha}^{\prime} \bar{\beta}^{\prime} \# \tmop{FV}
      (\Sigma, \tau)$ and $\tau_i^{\prime} \assign \tau_i [\bar{\alpha}
      \bar{\beta} \assign \bar{\alpha}^{\prime} \bar{\beta}^{\prime}],
      D^{\prime} \assign D [\bar{\alpha} \bar{\beta} \assign
      \bar{\alpha}^{\prime} \bar{\beta}^{\prime}]$. Let $\Delta_i^{\prime}$ be
      $\Delta_i$ with unbound occurrences of $\bar{\alpha} \bar{\beta}$
      renamed to $\bar{\alpha}^{\prime} \bar{\beta}^{\prime}$.
      
      \item By weakening and {\tmname{p-EqIn}}, (3) gives $\bar{\alpha}
      \bar{\beta} \dot{=} \bar{\alpha}^{\prime} \bar{\beta}^{\prime} \wedge
      \varepsilon (\bar{\alpha}^{\prime}) \dot{=} \tau \wedge C \vdash p :
      \tau \longrightarrow \exists \bar{\beta} [D] (\Delta_1 \times \ldots
      \times \Delta_n)$.
      
      \item By proposition \ref{EnvFragProp}, transitivity of $\leqslant$, and
      {\tmname{p-SubOut}}, we get $\bar{\alpha} \bar{\beta} \dot{=}
      \bar{\alpha}^{\prime} \bar{\beta}^{\prime} \wedge \varepsilon
      (\bar{\alpha}^{\prime}) \dot{=} \tau \wedge C \vdash p : \tau
      \longrightarrow \exists \bar{\alpha}^{\prime} \bar{\beta}^{\prime}
      [D^{\prime}] (\Delta_1^{\prime} \times \ldots \times
      \Delta_n^{\prime})$.
      
      \item By applying {\tmname{p-Hide}} to (6) with $C = \bar{\alpha}
      \dot{=} \bar{\alpha}^{\prime} \wedge \forall \bar{\beta}^{\prime}
      .D^{\prime} \Rightarrow \wedge_i \llbracket \vdash p_i \downarrow
      \tau_i^{\prime} \rrbracket$ and weakening, since w.l.o.g. $\bar{\alpha}
      \bar{\beta}$ do not appear unbound in the goal, and $C \wedge D \vDash
      \llbracket \vdash p_i \downarrow \tau_i \rrbracket$, we get the goal
      $\exists \bar{\alpha}^{\prime} . \varepsilon (\bar{\alpha}^{\prime})
      \dot{=} \tau \wedge \forall \bar{\beta}^{\prime} .D^{\prime} \Rightarrow
      \wedge_i \llbracket \vdash p_i \downarrow \tau_i^{\prime} \rrbracket
      \vdash p : \tau \longrightarrow \exists \bar{\alpha}^{\prime}
      \bar{\beta}^{\prime} [D^{\prime}] (\Delta_1^{\prime} \times \ldots
      \times \Delta_n^{\prime})$.
    \end{enumerate}
  \end{itemize}
\end{proof}

Proof of theorem \ref{InfCorrExpr}.

\begin{proof}
  By induction on the structure of $e$.
  \begin{itemize}
    \item Case $e$ is $x$.
    \begin{enumerate}
      \item If $x \nin \tmop{Dom} (\Gamma)$, then the goal follows by applying
      {\tmname{FElim}}. Otherwise, let $\Gamma (x)$ be $\forall \beta [\exists
      \bar{\alpha} .D] . \beta$. By {\tmname{Var}}, $D, \Gamma \vdash x :
      \beta$.
      
      \item Let $\beta^{\prime} \bar{\alpha}^{\prime} \# \tmop{FV} (\Gamma,
      \tau)$. By (1), weakening and {\tmname{Equ}}, $\beta \bar{\alpha}
      \dot{=} \beta^{\prime} \bar{\alpha}^{\prime} \wedge D^{\prime} \wedge
      \beta^{\prime} \dot{=} \tau, \Gamma \vdash x : \tau$, where $D^{\prime}
      \assign D [\beta \bar{\alpha} \assign \beta^{\prime}
      \bar{\alpha}^{\prime}]$.
      
      \item By {\tmname{Hide}} and weakening, since w.l.o.g. $\beta
      \bar{\alpha}$ do not appear unbounded in the goal, this implies the goal
      $\exists \beta^{\prime} \bar{\alpha}^{\prime} . (D^{\prime} \wedge
      \beta^{\prime} \dot{=} \tau), \Gamma \vdash x : \tau$.
    \end{enumerate}
    
    
    \item Case $e$ is $\text{{\tmstrong{assert false}}}$. The goal follows by
    {\tmname{FElim}}.
    
    \item Case $e$ is $\text{{\tmstrong{assert num}}} m \leq n ; e$.
    \begin{enumerate}
      \item Let $\alpha^1 \alpha^2 \# \tmop{FV} (\Gamma, \tau)$.
      
      \item By the induction hypothesis, we have $\llbracket \Gamma \vdash e_1
      : \tmop{Num} (\alpha^1) \rrbracket, \Gamma \vdash e_1 : \tmop{Num}
      (\alpha^1)$, $\llbracket \Gamma \vdash e_2 : \tmop{Num} (\alpha^2)
      \rrbracket, \Gamma \vdash e_2 : \tmop{Num} (\alpha^2)$ and $\llbracket
      \Gamma \vdash e_3 : \tau \rrbracket, \Gamma \vdash e_3 : \tau$.
      
      \item By weakening and {\tmname{AssertLeq}}, this yields $\llbracket
      \Gamma \vdash e_1 : \tmop{Num} (\alpha^1) \rrbracket \wedge \llbracket
      \Gamma \vdash e_2 : \tmop{Num} (\alpha^2) \rrbracket \wedge \alpha^1
      \leq \alpha^2 \wedge \left\llbracket \Gamma \vdash e_3 : {\nobreak} \tau
      \right\rrbracket, \Gamma \vdash \text{{\tmstrong{assert num}}} e_1 \leq
      e_2 ; e_3 : \tau$.
      
      \item By {\tmname{Hide}} using (1), this gives the goal.
    \end{enumerate}
    \item Case $e$ is $\text{{\tmstrong{assert type}}} e_1 = e_2 ; e_3$.
    \begin{enumerate}
      \item Let $\alpha^1 \alpha^2 \# \tmop{FV} (\Gamma, \tau)$.
      
      \item By the induction hypothesis, we have $\llbracket \Gamma \vdash e_1
      : \alpha^1 \rrbracket, \Gamma \vdash e_1 : \alpha^1$, $\llbracket \Gamma
      \vdash e_2 : \alpha^2 \rrbracket, \Gamma \vdash e_2 : \alpha^2$ and
      $\llbracket \Gamma \vdash e_3 : \tau \rrbracket, \Gamma \vdash e_3 :
      \tau$.
      
      \item By weakening and {\tmname{AssertEqty}}, this yields $\llbracket
      \Gamma \vdash e_1 : \alpha^1 \rrbracket \wedge \llbracket \Gamma \vdash
      e_2 : \alpha^2 \rrbracket \wedge \alpha^1 \dot{=} \alpha^2 \wedge
      \left\llbracket \Gamma \vdash e_3 : {\nobreak} \tau \right\rrbracket,
      \Gamma \vdash \text{{\tmstrong{assert type}}} e_1 = e_2 ; e_3 : \tau$.
      
      \item By {\tmname{Hide}} using (1), this gives the goal.
    \end{enumerate}
    \item Case $e$ is $\text{{\tmstrong{runtime failure}}} s$.
    \begin{enumerate}
      \item By the induction hypothesis, we have $\llbracket \Gamma \vdash s :
      \tmop{String} \rrbracket, \Gamma \vdash s : \tmop{String}$.
      
      \item The goal follows from {\tmname{RuntimeFailure}}, since
      $\left\llbracket \Gamma \vdash \text{{\tmstrong{runtime failure}}} s :
      {\nobreak} \tau \right\rrbracket = \llbracket \Gamma \vdash s :
      \tmop{String} \rrbracket$.
    \end{enumerate}
    \item Case $e$ is $\lambda \bar{c}$ where $\bar{c} = (c_1, \ldots, c_n)$.
    \begin{enumerate}
      \item Let $\alpha_1 \alpha_2 \# \tmop{FV} (\Gamma, \tau)$.
      
      \item Induction hypothesis yelds $\llbracket \Gamma \vdash c_i :
      \alpha_1 \rightarrow \alpha_2 \rrbracket, \Gamma \vdash c_i : \alpha_1
      \rightarrow \alpha_2$.
      
      \item By (2), weakening and {\tmname{Abs}}, $\llbracket \Gamma \vdash
      \bar{c} : \alpha_1 \rightarrow \alpha_2 \rrbracket, \Gamma \vdash
      \lambda \bar{c} : \alpha_1 \rightarrow \alpha_2$.
      
      \item By weakening and {\tmname{Equ}}, (3) implies $\llbracket \Gamma
      \vdash \bar{c} : \alpha_1 \rightarrow \alpha_2 \rrbracket \wedge
      \alpha_1 \rightarrow \alpha_2 \dot{=} \tau, \Gamma \vdash \lambda
      \bar{c} : \tau$.
      
      \item By (1) and {\tmname{Hide}}, this implies $\llbracket \Gamma \vdash
      \lambda \bar{c} : \tau \rrbracket, \Gamma \vdash \lambda \bar{c} :
      \tau$.
    \end{enumerate}
    \item Case $e$ is $e_1 e_2$.
    \begin{enumerate}
      \item Let $\alpha \# \tmop{FV} (\Gamma, \tau)$.
      
      \item By the induction hypothesis, we have $\llbracket \Gamma \vdash e_1
      : \alpha \rightarrow \tau \rrbracket, \Gamma \vdash e_1 : \alpha
      \rightarrow \tau$ and $\llbracket \Gamma \vdash e_2 : \alpha \rrbracket,
      \Gamma \vdash e_2 : \alpha$.
      
      \item By weakening and {\tmname{App}}, this yields $\llbracket \Gamma
      \vdash e_1 : \alpha \rightarrow \tau \rrbracket \wedge \llbracket \Gamma
      \vdash e_2 : \alpha \rrbracket, \Gamma \vdash e_1 e_2 : \tau$.
      
      \item By {\tmname{Hide}} using (1), $\llbracket \Gamma \vdash e_1 e_2 :
      \tau \rrbracket, \Gamma \vdash e_1 e_2 : \tau$.
    \end{enumerate}
    \item Case $e$ is $Ke_1 \ldots e_n$.
    \begin{enumerate}
      \item Let $\Sigma \ni K \colons \forall \bar{\alpha} \bar{\beta} [D] .
      \tau_1 \times \ldots \times \tau_n \rightarrow \varepsilon
      (\bar{\alpha})$.
      
      \item By induction hypothesis and weakening for each $i = 1, \ldots, n$
      \[ \wedge_j \llbracket \Gamma \vdash e_j : \tau_j \rrbracket \wedge D
         \wedge \varepsilon (\bar{\alpha}) \dot{=} \tau, \Gamma \vdash e_i :
         \tau_i \]
      \item Applying {\tmname{Cstr}} to (1) and (3) we obtain
      \[ \wedge_i \llbracket \Gamma \vdash e_i : \tau_i \rrbracket \wedge D
         \wedge \varepsilon (\bar{\alpha}) \dot{=} \tau, \Gamma \vdash Ke_1
         \ldots e_n : \varepsilon (\bar{\alpha}) \]
      \item Let $\bar{\alpha}^{\prime} \bar{\beta}^{\prime} \# \tmop{FV}
      (\Gamma, \tau)$ and $\tau_i^{\prime} \assign \tau_i [\bar{\alpha}
      \bar{\beta} \assign \bar{\alpha}^{\prime} \bar{\beta}^{\prime}]$,
      $D^{\prime} \assign D [\beta \bar{\alpha} \assign \beta^{\prime}
      \bar{\alpha}^{\prime}]$.
      \[ \bar{\alpha} \bar{\beta} \dot{=} \bar{\alpha}^{\prime}
         \bar{\beta}^{\prime} \wedge_i \llbracket \Gamma \vdash e_i :
         \tau_i^{\prime} \rrbracket \wedge D \wedge \varepsilon
         (\bar{\alpha}^{\prime}) \dot{=} \tau, \Gamma \vdash Ke_1 \ldots e_n :
         \varepsilon (\bar{\alpha}^{\prime}) \]
      \item By {\tmname{Equ}}, (1) {\tmname{Hide}} and weakening, since
      w.l.o.g. $\bar{\alpha} \bar{\beta}$ do not appear unbounded in the goal,
      $\llbracket \Gamma \vdash Ke_1 \ldots e_n : \tau \rrbracket, \Gamma
      \vdash Ke_1 \ldots e_n : \tau$.
    \end{enumerate}
    \item Case $e$ is $\text{{\tmstrong{let rec}}} x = e_1 
    \text{{\tmstrong{in}}} e_2$.
    \begin{enumerate}
      \item Let $\alpha \beta \# \tmop{FV} (\Gamma, \tau)$ and $\chi \#
      \tmop{PV} (\Gamma)$.
      
      \item Let $\sigma = \forall \beta [\chi (\beta)] . \beta$,
      $\Gamma^{\prime} = \Gamma \{x \mapsto \sigma\}$. By the induction
      hypothesis, $\llbracket \Gamma^{\prime} \vdash e_1 : \beta \rrbracket,
      \Gamma^{\prime} \vdash e_1 : \beta$ and $\llbracket \Gamma^{\prime}
      \vdash e_2 : \tau \rrbracket, \Gamma^{\prime} \vdash e_2 : \tau$.
      
      \item Let $D = \forall \beta . (\chi (\beta) \Rightarrow \llbracket
      \Gamma^{\prime} \vdash e_1 : \beta \rrbracket)$. Since $D \wedge \chi
      (\beta)$ implies $\left\llbracket \Gamma^{\prime} \vdash e_1 :
      {\nobreak} \beta \right\rrbracket$, by weakening of (2), we have $D
      \wedge \chi (\beta), \Gamma^{\prime} \vdash e_1 : {\nobreak} \beta$.
      From (1) we have $\alpha \# \tmop{FV} (D, \Gamma^{\prime}, \tau)$, by
      {\tmname{Gen}} we have $D \wedge \exists \beta . \chi (\beta),
      \Gamma^{\prime} \vdash e_1 : \forall \beta [\chi (\beta)] . \beta$, by
      (1) and renaming we have
      \[ D \wedge \exists \alpha . \chi (\alpha), \Gamma^{\prime} \vdash e_1 :
         \sigma . \]
      \item By weakening of both (2) and (3), and by {\tmname{LetRec}}, we
      have $\left\llbracket \Gamma \vdash \text{{\tmstrong{let rec}}} x = e_1 
      \text{{\tmstrong{in}}} e_2 : \tau \right\rrbracket, \Gamma \vdash
      \text{{\tmstrong{let rec}}} x = e_1  \text{{\tmstrong{in}}} e_2 : \tau$.
    \end{enumerate}
    \item Case $e$ is $p.e$.
    \begin{enumerate}
      \item $\tau$ is of the form $\tau_1 \rightarrow \tau_2$. Write
      $\llbracket \vdash p \uparrow \tau_1 \rrbracket$ as $\exists \bar{\beta}
      [D] \Gamma^{\prime}$, where $\bar{\beta} \# \tmop{FV} (\Gamma, \tau_1,
      \tau_2)$.
      
      \item By induction hypothesis, $\llbracket \Gamma \Gamma^{\prime} \vdash
      e : \tau_2 \rrbracket, \Gamma \Gamma^{\prime} \vdash e : \tau_2$.
      
      \item By lemma \ref{CorrPat} and (1), we have $\llbracket \vdash p
      \downarrow \tau_1 \rrbracket \vdash p : \tau_1 \longrightarrow \exists
      \bar{\beta} [D] \Gamma^{\prime}$.
      
      \item By instantiation of $\bar{\beta}$ and weakening, (2) implies
      \[ \llbracket \Gamma \vdash p.e : \tau \rrbracket \wedge D, \Gamma
         \Gamma^{\prime} \vdash e : \tau_2 \]
      \item By weakening, (3) implies $\llbracket \Gamma \vdash p.e : \tau
      \rrbracket \vdash p : \tau_1 \longrightarrow \exists \bar{\beta} [D]
      \Gamma^{\prime}$.
      
      \item By (4), (5), (1), and {\tmname{Clause}}, we obtain $\llbracket
      \Gamma \vdash p.e : \tau \rrbracket, \Gamma \vdash p.e : \tau$.
    \end{enumerate}
    \item Case $e$ is $p \text{{\tmstrong{when}}} \wedge_i m_i \leqslant n_i
    .e$ and $e \neq \text{{\tmstrong{assert false}}} \wedge \ldots \wedge e
    \neq \lambda \left( p^{\prime} \ldots \lambda \left( p^{\prime\prime} .
    \text{{\tmstrong{assert false}}} \right) \right)$.
    \begin{enumerate}
      \item Let $\bar{\beta} \overline{\alpha^1_i \alpha^2_i} \# \tmop{FV}
      (\Gamma, \tau_1, \tau_2)$. Recall that $\left\llbracket \Gamma \vdash p
      \text{{\tmstrong{when}}} \wedge_i m_i \leqslant n_i .e : \tau_1
      \rightarrow \tau_2 \right\rrbracket = \exists \overline{\alpha^1_i
      \alpha^2_i} . \Phi$, for $\Phi = \\
      \llbracket \vdash p \downarrow \tau_1 \rrbracket \wedge \forall
      \bar{\beta} .D \Rightarrow \wedge_i \llbracket \Gamma \Gamma^{\prime}
      \vdash m_i : \tmop{Num} (\alpha^1_i) \rrbracket \wedge_i \\
      \llbracket \Gamma \Gamma^{\prime} \vdash n_i : \tmop{Num} (\alpha^2_i)
      \rrbracket \wedge (\wedge_i \alpha^1_i \leq \alpha^2_i \Rightarrow
      \llbracket \Gamma \Gamma^{\prime} \vdash e : \tau_2 \rrbracket)$.
      
      \item $\tau$ is of the form $\tau_1 \rightarrow \tau_2$. Write
      $\llbracket \vdash p \uparrow \tau_1 \rrbracket$ as $\exists \bar{\beta}
      [D] \Gamma^{\prime}$.
      
      \item By induction hypothesis, $\llbracket \Gamma \Gamma^{\prime} \vdash
      e : \tau_2 \rrbracket, \Gamma \Gamma^{\prime} \vdash e : \tau_2$, and
      also $\llbracket \Gamma \Gamma^{\prime} \vdash m_i : \tmop{Num}
      (\alpha_i^1) \rrbracket, \Gamma \Gamma^{\prime} \vdash m_i : \tmop{Num}
      (\alpha_i^1)$ and $\llbracket \Gamma \Gamma^{\prime} \vdash n_i :
      \tmop{Num} (\alpha_i^2) \rrbracket, \Gamma \Gamma^{\prime} \vdash n_i :
      \tmop{Num} (\alpha_i^2)$.
      
      \item By lemma \ref{CorrPat} and (1), we have $\llbracket \vdash p
      \downarrow \tau_1 \rrbracket \vdash p : \tau_1 \longrightarrow \exists
      \bar{\beta} [D] \Gamma^{\prime}$.
      
      \item By instantiation of $\bar{\beta}$ and weakening, (3) implies
      \begin{eqnarray*}
        \Phi \wedge D \wedge_i \alpha^1_i \leq \alpha^2_i, \Gamma
        \Gamma^{\prime} & \vdash & e : \tau_2\\
        \Phi \wedge D, \Gamma \Gamma^{\prime} & \vdash & m_i : \tmop{Num}
        (\alpha^1_i)\\
        \Phi \wedge D, \Gamma \Gamma^{\prime} & \vdash & n_i : \tmop{Num}
        (\alpha^2_i)
      \end{eqnarray*}
      \item By weakening, (4) implies $\Phi \vdash p : \tau_1 \longrightarrow
      \exists \bar{\beta} [D] \Gamma^{\prime}$.
      
      \item By (5), (6), (1), and {\tmname{Clause}}, we obtain $\Phi, \Gamma
      \vdash p \text{{\tmstrong{when}}} \wedge_i m_i \leqslant n_i .e : \tau$.
      
      \item By (7) and {\tmname{Hide}}, we get the goal.
    \end{enumerate}
    \item Case $e$ is $p \text{{\tmstrong{when}}} \wedge_i m_i \leqslant n_i
    .e$ and $e = \text{{\tmstrong{assert false}}} \vee \ldots \vee e = \lambda
    \left( p^{\prime} \ldots \lambda \left( p^{\prime\prime} .
    \text{{\tmstrong{assert false}}} \right) \right)$. The proof is nearly
    identical to above.
    \begin{enumerate}
      \item Let $\alpha_3 \bar{\beta} \overline{\alpha^1_i \alpha^2_i} \#
      \tmop{FV} (\Gamma, \tau_1, \tau_2)$. Recall that $\left\llbracket \Gamma
      \vdash p \text{{\tmstrong{when}}} \wedge_i m_i \leqslant n_i .e : \tau_1
      \rightarrow \tau_2 \right\rrbracket = \exists \alpha_3
      \overline{\alpha^1_i \alpha^2_i} . \Phi$, for $\Phi = \\
      \llbracket \vdash p \downarrow \alpha_3 \rrbracket \wedge \forall
      \bar{\beta} .D \Rightarrow \wedge_i \llbracket \Gamma \Gamma^{\prime}
      \vdash m_i : \tmop{Num} (\alpha^1_i) \rrbracket \wedge_i \llbracket
      \Gamma \Gamma^{\prime} \vdash n_i : \tmop{Num} (\alpha^2_i) \rrbracket
      \\
      \wedge (\alpha_3 \dot{=} \tau_1 \wedge_i \alpha^1_i \leq \alpha^2_i
      \Rightarrow \llbracket \Gamma \Gamma^{\prime} \vdash e : \tau_2
      \rrbracket)$.
      
      \item $\tau$ is of the form $\tau_1 \rightarrow \tau_2$. Write
      $\llbracket \vdash p \uparrow \alpha_3 \rrbracket$ as $\exists
      \bar{\beta} [D] \Gamma^{\prime}$.
      
      \item By induction hypothesis, $\llbracket \Gamma \Gamma^{\prime} \vdash
      e : \tau_2 \rrbracket, \Gamma \Gamma^{\prime} \vdash e : \tau_2$, and
      also $\llbracket \Gamma \Gamma^{\prime} \vdash m_i : \tmop{Num}
      (\alpha_i^1) \rrbracket, \Gamma \Gamma^{\prime} \vdash m_i : \tmop{Num}
      (\alpha_i^1)$ and $\llbracket \Gamma \Gamma^{\prime} \vdash n_i :
      \tmop{Num} (\alpha_i^2) \rrbracket, \Gamma \Gamma^{\prime} \vdash n_i :
      \tmop{Num} (\alpha_i^2)$.
      
      \item By lemma \ref{CorrPat} and (1), we have $\llbracket \vdash p
      \downarrow \alpha_3 \rrbracket \vdash p : \alpha_3 \longrightarrow
      \exists \bar{\beta} [D] \Gamma^{\prime}$.
      
      \item By instantiation of $\bar{\beta}$ and weakening, (3) implies
      \begin{eqnarray*}
        \Phi \wedge D \wedge \alpha_3 \dot{=} \tau_1 \wedge_i \alpha^1_i \leq
        \alpha^2_i, \Gamma \Gamma^{\prime} & \vdash & e : \tau_2\\
        \Phi \wedge D, \Gamma \Gamma^{\prime} & \vdash & m_i : \tmop{Num}
        (\alpha^1_i)\\
        \Phi \wedge D, \Gamma \Gamma^{\prime} & \vdash & n_i : \tmop{Num}
        (\alpha^2_i)
      \end{eqnarray*}
      \item By weakening, (4) implies $\Phi \vdash p : \alpha_3
      \longrightarrow \exists \bar{\beta} [D] \Gamma^{\prime}$.
      
      \item By (5), (6), (1), and {\tmname{NegClause}}, we obtain $\Phi,
      \Gamma \vdash p \text{{\tmstrong{when}}} \wedge_i m_i \leqslant n_i .e :
      \tau$.
      
      \item By (7) and {\tmname{Hide}}, we get the goal.
    \end{enumerate}
  \end{itemize}
\end{proof}

$\Gamma^{\prime} \dot{=} \Gamma^{\prime\prime}$ stands for $\forall x \in
\tmop{Dom} (\Gamma^{\prime}) \cup \tmop{Dom} (\Gamma^{\prime\prime}) .
\Gamma^{\prime} (x) \dot{=} \Gamma^{\prime\prime} (x)$ and is false when
$\tmop{Dom} (\Gamma^{\prime}) \neq \tmop{Dom} (\Gamma^{\prime\prime})$. Recall
that for $\Delta \assign \exists \bar{\beta} [D] . \Gamma$ and
$\Delta^{\prime} \assign \exists \bar{\beta}^{\prime} [D^{\prime}] .
\Gamma^{\prime}$ such that $\bar{\beta} \# \tmop{FV} (\Gamma^{\prime})$,
$\bar{\beta}^{\prime} \# \tmop{FV} (\Delta)$ and $\bar{\beta}^{\prime} \#C$,
$C \vDash \Delta^{\prime} \leqslant \Delta$ denotes $C \wedge D^{\prime}
\vDash \exists \bar{\beta} .D \wedge \Gamma \dot{=} \Gamma^{\prime}$. Observe,
that $C \vDash \Delta^{\prime} \leqslant \Delta$ iff $C \vDash \forall
\bar{\beta}^{\prime} .D^{\prime} \Rightarrow \exists \bar{\beta} .D \wedge
\Gamma \dot{=} \Gamma^{\prime}$.

\begin{lemma}
  \label{CompletePat}{\tmem{Completeness}} (patterns). Let $\Delta = \exists
  \bar{\beta}^{\prime} [D^{\prime}] \Gamma^{\prime}$ and $\llbracket \vdash p
  \uparrow \tau \rrbracket = \exists \bar{\beta}^{\prime\prime}
  [D^{\prime\prime}] \Gamma^{\prime\prime} = \Delta^{\prime}$. $C \vdash p :
  \tau \longrightarrow \Delta$ implies $C \vDash \llbracket \vdash p
  \downarrow \tau \rrbracket$ and $C \vDash \forall \bar{\beta}^{\prime\prime}
  .D^{\prime\prime} \Rightarrow \exists \bar{\beta}^{\prime} . (D^{\prime}
  \wedge \Gamma^{\prime\prime} \dot{=} \Gamma^{\prime})$, i.e. $C \vDash
  \Delta^{\prime} \leqslant \Delta$.
\end{lemma}

\begin{proof}
  By induction on the derivation of $C \vdash p : \tau \longrightarrow
  \Delta$. To slightly simplify the proof, the induction is actually on the
  lexicographic ordering: (\# of applications of {\tmname{p-Cstr}}, \# of
  other rules applications).
  \begin{itemize}
    \item Cases {\tmname{p-Empty}}, {\tmname{p-Wild}}, {\tmname{p-Var}}.
    $\llbracket \vdash p \downarrow \tau \rrbracket =\tmmathbf{T}$.
    $\llbracket \vdash p \uparrow \tau \rrbracket$ and $\Delta$ coincide:
    $\Gamma^{\prime\prime} = \Gamma^{\prime}$, $D^{\prime} = D^{\prime\prime}
    =\tmmathbf{T}$ and $\vDash \exists \bar{\beta} . \Gamma^{\prime} \dot{=}
    \Gamma^{\prime}$ holds because sorts are nonempty.
    
    \item Case {\tmname{p-And}}. In this case $\Delta = \Delta_1 \times
    \Delta_2, \bar{\beta}^{\prime} = \bar{\beta}_1^{\prime}
    \bar{\beta}_2^{\prime}, D^{\prime} = D_1^{\prime} \wedge D_2^{\prime},
    \Gamma^{\prime} = \Gamma_1^{\prime} \dot{\cup} \Gamma_2^{\prime}$.
    \begin{enumerate}
      \item {\tmname{p-And}}'s premises are $C \vdash p_i : \tau
      \longrightarrow \Delta_i$, which by induction hypothesis gives $C \vDash
      \llbracket \vdash p_i \downarrow \tau \rrbracket$ and $C \vDash \forall
      \bar{\beta}^{\prime\prime}_i .D_i^{\prime\prime} \Rightarrow \exists
      \bar{\beta}_i^{\prime} . (D_i^{\prime} \wedge \Gamma_i^{\prime\prime}
      \dot{=} \Gamma_i^{\prime})$ for $i = 1, 2$.
      
      \item (1) gives $C \vDash \llbracket \vdash p_1 \wedge p_2 \downarrow
      \tau \rrbracket \nobracket \nobracket$ as $\llbracket \vdash p_1 \wedge
      p_2 \downarrow \tau \rrbracket = \llbracket \vdash p_1 \downarrow \tau
      \rrbracket \wedge \llbracket \vdash p_2 \downarrow \tau \rrbracket
      \nobracket \nobracket$.
      
      \item $\llbracket \vdash p_1 \wedge p_2 \uparrow \tau \rrbracket =
      \llbracket \vdash p_1 \uparrow \tau \rrbracket \times \llbracket \vdash
      p_2 \uparrow \tau \rrbracket = \exists \bar{\beta}^{\prime\prime}_1
      \bar{\beta}_2^{\prime\prime} [D_1^{\prime\prime} \wedge
      D_2^{\prime\prime}] \Gamma_1^{\prime\prime} \dot{\cup}
      \Gamma_2^{\prime\prime} \nobracket \nobracket \nobracket \nobracket
      \nobracket \nobracket$. We will show $C \vDash \forall
      \bar{\beta}_1^{\prime\prime} \bar{\beta}_2^{\prime\prime}
      .D_1^{\prime\prime} \wedge D_2^{\prime\prime} \Rightarrow \exists
      \bar{\beta}_1^{\prime} \bar{\beta}_2^{\prime} . (D_1^{\prime} \wedge
      D_2^{\prime} \wedge \Gamma_1^{\prime\prime} \dot{\cup}
      \Gamma_2^{\prime\prime} \dot{=} \Gamma_1^{\prime} \dot{\cup}
      \Gamma_2^{\prime})$.
      
      \item Assume w.l.o.g. $\bar{\beta}_1^{\prime} \#
      \bar{\beta}_2^{\prime}$, $\bar{\beta}_1^{\prime\prime} \#
      \bar{\beta}_2^{\prime\prime}$. Applying (1) for $i = 1, 2$ gives $C
      \vDash \forall \bar{\beta}_1^{\prime\prime} \bar{\beta}_2^{\prime\prime}
      .D_1^{\prime\prime} \wedge D_2^{\prime\prime} \Rightarrow \exists
      \bar{\beta}_1^{\prime} \bar{\beta}_2^{\prime} . (D_1^{\prime} \wedge
      D_2^{\prime} \wedge \Gamma_1^{\prime\prime} \dot{=} \Gamma_1^{\prime}
      \wedge \Gamma_2^{\prime\prime} \dot{=} \Gamma_2^{\prime})$, which
      completes the goal.
    \end{enumerate}
    \item Case {\tmname{p-Cstr}}. In this case $\Delta = \exists \bar{\beta}_0
    [D_0] (\Delta_1 \times \ldots \times \Delta_n)$, and $\tau = \varepsilon
    (\bar{\alpha}_0)$, where $D_0 \assign D_K [\bar{\alpha} \bar{\beta}
    \assign \bar{\alpha}_0 \bar{\beta}_0]$ for $\Sigma \ni K \colons \forall
    \bar{\alpha} \bar{\beta} [D_K] . \tau_1 \times \ldots \times \tau_n
    \rightarrow \varepsilon (\bar{\alpha})$ and $\bar{\beta}_0 \# \tmop{FV}
    (C)$.
    \begin{enumerate}
      \item {\tmname{p-Cstr}}'s premises are $C \wedge D_0 \vdash p_i : \tau_i
      [\bar{\alpha} \bar{\beta} \assign \bar{\alpha}_0 \bar{\beta}_0]
      \longrightarrow \Delta_i$.
      
      \item Let $\bar{\alpha}_0^{\prime} \bar{\beta}_0^{\prime} \# \tmop{FV}
      (\tau, \bar{\alpha} \bar{\beta}, \bar{\alpha}_0 \bar{\beta}_0, C)$.
      
      \item Let $\tau_i^{\prime} \assign \tau_i [\bar{\alpha} \bar{\beta}
      \assign \bar{\alpha}_0^{\prime} \bar{\beta}_0^{\prime}]$. By weakening
      and {\tmname{p-EqIn}}, (1) gives $C \wedge D_0 \wedge \bar{\alpha}_0
      \bar{\beta}_0 \dot{=} \bar{\alpha}_0^{\prime} \bar{\beta}_0^{\prime}
      \vdash p_i : \tau_i^{\prime} \longrightarrow \Delta_i$.
      
      \item By induction hypothesis we have $C \wedge D_0 \wedge
      \bar{\alpha}_0 \bar{\beta}_0 \dot{=} \bar{\alpha}_0^{\prime}
      \bar{\beta}_0^{\prime} \vDash \llbracket \vdash p_i \downarrow
      \tau_i^{\prime} \rrbracket$ and $C \wedge D_0 \wedge \bar{\alpha}_0
      \bar{\beta}_0 \dot{=} \bar{\alpha}_0^{\prime} \bar{\beta}_0^{\prime}
      \vDash \forall \bar{\beta}_i^{\prime\prime} .D_i^{\prime\prime}
      \Rightarrow \exists \bar{\beta}_i^{\prime} . (D_i^{\prime} \wedge
      \Gamma_i^{\prime\prime} \dot{=} \Gamma_i^{\prime})$ for $i = 1, \ldots,
      n$.
      
      \item Let $D_0^{\prime} \assign D_K [\bar{\alpha} \bar{\beta} \assign
      \bar{\alpha}_0^{\prime} \bar{\beta}_0^{\prime}]$. From (4) follows $C
      \wedge \bar{\alpha}_0 \bar{\beta}_0 \dot{=} \bar{\alpha}_0^{\prime}
      \bar{\beta}_0^{\prime} \vDash D_0^{\prime} \Rightarrow \wedge_i
      \llbracket \vdash p_i \downarrow \tau_i^{\prime} \rrbracket$.
      
      \item W.l.o.g. $\bar{\alpha}_0 \bar{\beta}_0 \# \tmop{FV} (D_0^{\prime}
      \Rightarrow \wedge_i \llbracket \vdash p_i \downarrow \tau_i^{\prime}
      \rrbracket)$. (5) gives $C \wedge \bar{\alpha}_0 \dot{=}
      \bar{\alpha}_0^{\prime} \vDash \forall \bar{\beta}_0^{\prime}
      .D_0^{\prime} \Rightarrow \wedge_i \llbracket \vdash p_i \downarrow
      \tau_i^{\prime} \rrbracket$ because we can drop $\bar{\beta}_0 \dot{=}
      \bar{\beta}_0^{\prime}$ from premises.
      
      \item (6) is equivalent to $C \wedge \bar{\alpha}_0 \dot{=}
      \bar{\alpha}_0^{\prime} \vDash \varepsilon (\bar{\alpha}_0) \dot{=}
      \varepsilon (\bar{\alpha}_0^{\prime}) \wedge \forall
      \bar{\beta}_0^{\prime} .D_0^{\prime} \Rightarrow \wedge_i \llbracket p_i
      \downarrow \tau_i^{\prime} \rrbracket$ which by the nonempty domain
      property implies $C \wedge \bar{\alpha}_0 \dot{=}
      \bar{\alpha}_0^{\prime} \vDash \exists \bar{\alpha}_0^{\prime} .
      \varepsilon (\bar{\alpha}_0) \dot{=} \varepsilon
      (\bar{\alpha}_0^{\prime}) \wedge \forall \bar{\beta}_0^{\prime}
      .D_0^{\prime} \Rightarrow \wedge_i \llbracket \vdash p_i \downarrow
      \tau_i^{\prime} \rrbracket$.
      
      \item Because by (6) we can drop $\bar{\alpha}_0 \dot{=}
      \bar{\alpha}_0^{\prime}$ from premises, (7) is equivalent to $C \vDash
      \exists \bar{\alpha}_0^{\prime} . \varepsilon (\bar{\alpha}_0) \dot{=}
      \varepsilon (\bar{\alpha}_0^{\prime}) \wedge \forall
      \bar{\beta}_0^{\prime} .D_0^{\prime} \Rightarrow \wedge_i \llbracket
      \vdash p_i \downarrow \tau_i^{\prime} \rrbracket$, which is the first
      part of the goal.
      
      \item From (4), $C \wedge \bar{\alpha}_0 \bar{\beta}_0 \dot{=}
      \bar{\alpha}_0^{\prime} \bar{\beta}_0^{\prime} \vDash D_0^{\prime}
      \Rightarrow \forall \bar{\beta}_1^{\prime\prime} \ldots
      \bar{\beta}_n^{\prime\prime} . \wedge_i D_i^{\prime\prime} \Rightarrow
      \exists \bar{\beta}_1^{\prime} \ldots \bar{\beta}_n^{\prime} . \wedge_i
      (D_i^{\prime} \wedge \Gamma_i^{\prime\prime} \dot{=}
      \Gamma_i^{\prime})$.
      
      \item From (9) by (2) and (6), $C \vDash \forall \bar{\alpha}_0^{\prime}
      \bar{\beta}_0^{\prime} . \bar{\alpha}_0 \bar{\beta}_0 \dot{=}
      \bar{\alpha}_0^{\prime} \bar{\beta}_0^{\prime} \wedge D_0^{\prime}
      \Rightarrow \forall \bar{\beta}_1^{\prime\prime} \ldots
      \bar{\beta}_n^{\prime\prime} . \wedge_i D_i^{\prime\prime} \Rightarrow
      \exists \bar{\beta}_1^{\prime} \ldots \bar{\beta}_n^{\prime} . \wedge_i
      (D_i^{\prime} \wedge \Gamma_i^{\prime\prime} \dot{=}
      \Gamma_i^{\prime})$, which is equivalent to
      \[ C \vDash \forall \bar{\alpha}_0^{\prime} \bar{\beta}_0^{\prime}
         \bar{\beta}_1^{\prime\prime} \ldots \bar{\beta}_n^{\prime\prime} .
         \bar{\alpha}_0 \bar{\beta}_0 \dot{=} \bar{\alpha}_0^{\prime}
         \bar{\beta}_0^{\prime} \wedge D_0^{\prime} \wedge_i
         D_i^{\prime\prime} \Rightarrow \text{ \ \ \ \ \ \ \ \ \ \ \ \ \ \ \ }
         \text{\\
         \ \ \ \ \ \ \ \ \ \ \ \ \ \ \ \ \ \ \ \ \ } \exists
         \bar{\beta}_1^{\prime} \ldots \bar{\beta}_n^{\prime} . \wedge_i
         \left( D_i^{\prime} \wedge {\nobreak} \Gamma_i^{\prime\prime} \dot{=}
         \Gamma_i^{\prime} \right) \]
      \item Observe, that w.l.o.g. $\bar{\beta}^{\prime\prime} \assign
      \bar{\alpha}_0^{\prime} \bar{\beta}_0^{\prime}
      \bar{\beta}_1^{\prime\prime} \ldots \bar{\beta}_n^{\prime\prime}$. Note
      by definition of $\llbracket \vdash p \uparrow \tau \rrbracket$, that
      $D^{\prime\prime} = \varepsilon (\bar{\alpha}_0) \dot{=} \varepsilon
      (\bar{\alpha}_0^{\prime}) \wedge D_0^{\prime} \wedge_i
      D_i^{\prime\prime}$. By the free generation property, $\vDash
      D^{\prime\prime} \Rightarrow \bar{\alpha}_0 \dot{=}
      \bar{\alpha}_0^{\prime}$.
      
      \item Observe, that $\Gamma^{\prime\prime} \dot{=} \Gamma^{\prime}
      \equiv \wedge_i (\Gamma_i^{\prime\prime} \dot{=} \Gamma_i^{\prime})$ and
      $D^{\prime} = D_0 \wedge_i D_i^{\prime}$. (10) and (11) imply
      \[ C \vDash \forall \bar{\beta}^{\prime\prime} . \bar{\beta}_0 \dot{=}
         \bar{\beta}_0^{\prime} \wedge D^{\prime\prime} \Rightarrow \exists
         \bar{\beta}_1^{\prime} \ldots \bar{\beta}_n^{\prime} .D^{\prime}
         \wedge \Gamma^{\prime\prime} \dot{=} \Gamma^{\prime} \]
      \item Also, $\bar{\beta}^{\prime} = \bar{\beta}_0 \bar{\beta}_1^{\prime}
      \ldots \bar{\beta}_n^{\prime}$. Because $\bar{\beta}_0 \# \tmop{FV}
      (D^{\prime\prime})$, because sorts are nonempty (12) gives $C \vDash
      \forall \bar{\beta}^{\prime\prime} . \bar{\alpha}_0 \dot{=}
      \bar{\alpha}_0^{\prime} \wedge D^{\prime\prime} \Rightarrow \exists
      \bar{\beta}^{\prime} .D^{\prime} \wedge \Gamma^{\prime\prime} \dot{=}
      \Gamma^{\prime}$, the other part of the goal.
    \end{enumerate}
    \item Case {\tmname{p-EqIn}}.
    \begin{enumerate}
      \item {\tmname{p-EqIn}}'s premises are: $C \vdash p : \tau^{\prime}
      \longrightarrow \Delta$, which by induction hypothesis gives $C \vDash
      \llbracket \vdash p \downarrow \tau^{\prime} \rrbracket$ and $C \vDash
      \Delta^{\prime}_1 \leqslant \Delta$, for $\Delta_1^{\prime} = \exists
      \bar{\beta}_1^{\prime\prime} [D_1^{\prime\prime}]
      \Gamma_1^{\prime\prime}$
      
      \item and $C \vDash \tau \dot{=} \tau^{\prime}$.
      
      \item Observe by induction on $p$, that $C \wedge \tau \dot{=}
      \tau^{\prime} \vDash \llbracket \vdash p \downarrow \tau^{\prime}
      \rrbracket$ iff $C \wedge \tau \dot{=} \tau^{\prime} \vDash \llbracket
      \vdash p \downarrow \tau \rrbracket$, which by (1) and (2) gives the
      first part of the goal.
      
      \item Observe by induction on $p$, that $C \wedge \tau \dot{=}
      \tau^{\prime} \vDash \llbracket \vdash p \uparrow \tau \rrbracket
      \leqslant \llbracket \vdash p \uparrow \tau^{\prime} \rrbracket$, i.e.
      $C \wedge \tau \dot{=} \tau^{\prime} \vDash \Delta^{\prime} \leqslant
      \Delta^{\prime}_1$, which by (1), (2) and transitivity of $\leqslant$,
      proves the second part of the goal.
    \end{enumerate}
    \item Case {\tmname{p-SubOut}} follows by transitivity of $\leqslant$.
    
    \item Case {\tmname{p-Hide}}.
    \begin{enumerate}
      \item {\tmname{p-Hide}}'s premises are $C^{\prime} \vdash p : \tau
      \longrightarrow \Delta$ and $\bar{\alpha}_0 \# \tmop{FV} (\tau, \Delta)$
      for $C = \exists \bar{\alpha}_0 .C^{\prime}$.
      
      \item By inductive hypothesis, $C^{\prime} \vDash \llbracket \vdash p
      \downarrow \tau \rrbracket$ and $C^{\prime} \vDash \Delta^{\prime}
      \leqslant \Delta$.
      
      \item By induction on $p$, $\tmop{FV} (\llbracket \vdash p \downarrow
      \tau \rrbracket) = \tmop{FV} (\tau)$.
      
      \item By (1), (2) and (3) we have $C \vDash \llbracket \vdash p
      \downarrow \tau \rrbracket$.
      
      \item By induction on $p$, $\tmop{FV} (D^{\prime\prime},
      \Gamma^{\prime\prime}) \subseteq \tmop{FV} (\tau) \cup
      \bar{\beta}^{\prime\prime}$.
      
      \item By (1), (2) and (3) we have $C \vDash \Delta^{\prime} \leqslant
      \Delta$.
    \end{enumerate}
  \end{itemize}
\end{proof}

\begin{lemma}
  \label{MonoEnvEq}Let $\Gamma$ be an environment and $\Gamma^{\prime},
  \Gamma^{\prime\prime}$ be simple (i.e. monomorphic) environments. For any
  $e, \tau$, $C \wedge \Gamma^{\prime} \dot{=} \Gamma^{\prime\prime}, \Gamma
  \Gamma^{\prime} \vdash e : \tau$ iff $C \wedge \Gamma^{\prime} \dot{=}
  \Gamma^{\prime\prime}, \Gamma \Gamma^{\prime\prime} \vdash e : \tau$.
\end{lemma}

\begin{proof}
  Consider a derivation of $C \wedge \Gamma^{\prime} \dot{=}
  \Gamma^{\prime\prime}, \Gamma \Gamma^{\prime} \vdash e : \tau$. The only
  case where $\Gamma^{\prime}$ is referred to, is in the {\tmname{Var}} rule,
  which for a monomorphic environment simplifies to: $\Gamma^{\prime} (x) =
  \tau^{\prime} / C, \Gamma \Gamma^{\prime} \vdash x : \tau^{\prime}$. Replace
  $\Gamma^{\prime}$ with $\Gamma^{\prime\prime}$ in judgments throughout the
  derivation. $\Gamma^{\prime} (x) = \tau^{\prime} /_{\tmop{Var}} C \wedge
  \Gamma^{\prime} \dot{=} \Gamma^{\prime\prime}, \Gamma \Gamma^{\prime\prime}
  \vdash x : \tau^{\prime}$ is not valid, correct it as $\Gamma^{\prime\prime}
  (x) = \tau^{\prime\prime} /_{\tmop{Var}} C \wedge \Gamma^{\prime} \dot{=}
  \Gamma^{\prime\prime}, \Gamma \Gamma^{\prime\prime} \vdash x :
  \tau^{\prime\prime} /_{\tmop{Equ}} C \wedge \Gamma^{\prime} \dot{=}
  \Gamma^{\prime\prime}, \Gamma \Gamma^{\prime\prime} \vdash x :
  \tau^{\prime}$. Analogically follows the other direction of the equivalence
  of $C \wedge \Gamma^{\prime} \dot{=} \Gamma^{\prime\prime}, \Gamma
  \Gamma^{\prime} \vdash e : \tau$ and $C \wedge \Gamma^{\prime} \dot{=}
  \Gamma^{\prime\prime}, \Gamma \Gamma^{\prime\prime} \vdash e : \tau$.
\end{proof}

Proof of theorem \ref{InfComplExpr}.

\begin{proof}
  We proceed by induction on the derivation of $C, \Gamma \vdash e : \tau$. To
  slightly simplify the proof, the induction is actually on the lexicographic
  ordering: (\# of structural rule applications {\tmname{Var}},
  {\tmname{Cstr}}, {\tmname{Abs}}, {\tmname{App}}, {\tmname{LetRec}},
  {\tmname{Clause}}; \# of nonstructural rule applications {\tmname{Equ}},
  {\tmname{Hide}}, {\tmname{FElim}}, {\tmname{DisjElim}}). (The rules
  {\tmname{FElim}} and {\tmname{DisjElim}} are not needed when deriving the
  syntax-directed rules.)
  \begin{itemize}
    \item Case {\tmname{Var}}.
    \begin{enumerate}
      \item {\tmname{Var}}'s first premise is $\Gamma (x) = \forall \beta
      [\exists \bar{\alpha} .D] . \beta$.
      
      \item {\tmname{Var}}'s second premise is $C \vDash D$.
      
      \item The goal is: $\mathcal{I}, C \vDash \exists \beta^{\prime}
      \bar{\alpha}^{\prime} . (D [\beta \bar{\alpha} \assign \beta^{\prime}
      \bar{\alpha}^{\prime}] \wedge \beta^{\prime} \dot{=} \tau)$, where
      w.l.o.g. $\beta^{\prime} \bar{\alpha}^{\prime} \# \tmop{FV} (C, \Gamma,
      \tau, \beta, \bar{\alpha})$.
      
      \item (3) follows from (2) by instantiating $\beta$ to $\tau$, because
      we assume that all sorts in $\mathcal{M}$ are non-empty. We can take an
      empty interpretation $\mathcal{I}= \epsilon$.
    \end{enumerate}
    \item Case {\tmname{AssertFalse}}. {\tmname{AssertFalse}}'s premise is $C
    \vDash \tmmathbf{F}$, we have the goal from $\vDash \tmmathbf{F}
    \Rightarrow \Phi$ holding for any $\Phi$, and transitivity of $\vDash$.
    
    \item Case {\tmname{AssertLeq}}.
    \begin{enumerate}
      \item {\tmname{AssertLeq}}'s premises are $C, \Gamma \vdash e_1 :
      \tmop{Num} (\tau_1)$, $C, \Gamma \vdash e_2 : \tmop{Num} (\tau_2)$, $C
      \vDash \tau_1 \leq \tau_2$ and $C, \Gamma \vdash e_3 : \tau$.
      
      \item Pick w.l.o.g. $\alpha^1 \alpha^2 \nin \tmop{FV} (C, \tau_1,
      \tau_2, \Gamma, \tau)$. By rule {\tmname{Equ}}, (1) implies $C \wedge
      \alpha^1 \dot{=} \tau_1 \wedge \alpha^2 \dot{=} \tau_2, \Gamma \vdash
      e_1 : \tmop{Num} (\alpha^1)$, $C \wedge \alpha^1 \dot{=} \tau_1 \wedge
      \alpha^2 \dot{=} \tau_2, \Gamma \vdash e_2 : \tmop{Num} (\alpha^2)$ and
      $C, \Gamma \vdash e_3 : \tau$.
      
      \item By induction hypothesis and weakening, (2) implies $\mathcal{I}_i,
      C \wedge \alpha^1 \dot{=} \tau_1 \wedge \alpha^2 \dot{=} \tau_2 \wedge
      \alpha^1 \leq \alpha^2 \vDash \Phi_i$ for $\Phi_i = \left\llbracket
      \Gamma \vdash e_i : {\nobreak} \tmop{Num} (\alpha^i) \right\rrbracket$,
      $i = 1, 2$, $\Phi_3 = \llbracket \Gamma \vdash e_3 : \tau \rrbracket$.
      
      \item By (2) and nonemptiness of sorts, we have $C \vDash \exists
      \alpha^1 \alpha^2 . (C \wedge \alpha^1 \dot{=} \tau_1 \wedge \alpha^2
      \dot{=} \tau_2 \wedge \alpha^1 \leq \alpha^2)$.
      
      \item By (3), the premise and because $C \vDash D$ implies $\exists
      \alpha .C \vDash \exists \alpha .D$, we have $\mathcal{I}_1
      \mathcal{I}_2 \mathcal{I}_3, \exists \alpha^1 \alpha^2 . (C \wedge
      \alpha^1 \dot{=} \tau_1 \wedge \alpha^2 \dot{=} \tau_2 \wedge \alpha^1
      \leq \alpha^2) \vDash \exists \alpha . (\Phi_1 \wedge \Phi_2 \wedge
      \Phi_3)$.
      
      \item By (4) and (5), we have the goal $\mathcal{I}, C \vDash
      \left\llbracket \Gamma \vdash \text{{\tmstrong{assert num}}} e_1 \leq
      e_2 ; e_3 : \tau \right\rrbracket$ with $\mathcal{I}=\mathcal{I}_1
      \mathcal{I}_2 \mathcal{I}_3$.
    \end{enumerate}
    \item Case {\tmname{AssertEqty}}.
    \begin{enumerate}
      \item {\tmname{AssertEqty}}'s premises are $C, \Gamma \vdash e_i :
      \tau_i$ for $i = 1, 2, 3, \tau_3 = \tau$ and $C \vDash \tau_1 \dot{=}
      \tau_2$.
      
      \item Pick w.l.o.g. $\alpha^1 \alpha^2 \nin \tmop{FV} (C, \tau_1,
      \tau_2, \Gamma, \tau)$. By rule {\tmname{Equ}}, (1) implies $C \wedge
      \alpha^1 \dot{=} \tau_1 \wedge \alpha^2 \dot{=} \tau_2, \Gamma \vdash
      e_i : \alpha^i$ for $i = 1, 2$.
      
      \item By induction hypothesis and weakening, (1) and (2) imply
      $\mathcal{I}_i, C \wedge \alpha^1 \dot{=} \tau_1 \wedge \alpha^2 \dot{=}
      \tau_2 \wedge \alpha^1 \dot{=} \alpha^2 \vDash \Phi_i$ for $\Phi_i =
      \left\llbracket \Gamma \vdash e_i : {\nobreak} \alpha^i
      \right\rrbracket$, $i = 1, 2$, $\Phi_3 = \llbracket \Gamma \vdash e_3 :
      \tau \rrbracket$.
      
      \item By (2) and nonemptiness of sorts, we have $C \vDash \exists
      \alpha^1 \alpha^2 . (C \wedge \alpha^1 \dot{=} \tau_1 \wedge \alpha^2
      \dot{=} \tau_2 \wedge \alpha^1 \dot{=} \alpha^2)$.
      
      \item By (3), the premise and because $C \vDash D$ implies $\exists
      \alpha .C \vDash \exists \alpha .D$, we have $\mathcal{I}_1
      \mathcal{I}_2 \mathcal{I}_3, \exists \alpha^1 \alpha^2 . (C \wedge
      \alpha^1 \dot{=} \tau_1 \wedge \alpha^2 \dot{=} \tau_2 \wedge \alpha^1
      \dot{=} \alpha^2) \vDash \exists \alpha . (\Phi_1 \wedge \Phi_2 \wedge
      \Phi_3)$.
      
      \item By (4) and (5), we have the goal $\mathcal{I}, C \vDash
      \left\llbracket \Gamma \vdash \text{{\tmstrong{assert type}}} e_1 = e_2
      ; e_3 : \tau \right\rrbracket$ with $\mathcal{I}=\mathcal{I}_1
      \mathcal{I}_2 \mathcal{I}_3$.
    \end{enumerate}
    \item Case {\tmname{RuntimeFailure}}.
    \begin{enumerate}
      \item {\tmname{RuntimeFailure}}'s premise is $C, \Gamma \vdash s :
      \tmop{String}$.
      
      \item By induction hypothesis, (1) implies $\mathcal{I}, C \vDash
      \llbracket \Gamma \vdash s : \tmop{String} \rrbracket$ and thus the
      goal.
    \end{enumerate}
    \item Case {\tmname{Cstr}}.
    \begin{enumerate}
      \item {\tmname{Cst`r}}'s premises are $C, \Gamma \vdash e_i : \tau_i$,
      $i = 1, \ldots, n$, $C \vDash D$ and $K \colons \forall \bar{\alpha}
      \bar{\beta} [D] . \tau_1 \ldots \tau_n \rightarrow \varepsilon
      (\bar{\alpha})$. $\tau = \varepsilon (\bar{\alpha})$.
      
      \item Let w.l.o.g. $\bar{\alpha}^{\prime} \bar{\beta}^{\prime} \#
      \tmop{FV} (C, \Gamma, \tau)$. By weakening and {\tmname{Equ}}, (1) gives
      $C \wedge \bar{\alpha}^{\prime} \bar{\beta}^{\prime} \dot{=}
      \bar{\alpha} \bar{\beta}, \Gamma \vdash e_i : \tau_i [\bar{\alpha}
      \bar{\beta} \assign \bar{\alpha}^{\prime} \bar{\beta}^{\prime}]$.
      
      \item Let $\Phi_i = \llbracket \Gamma \vdash e_i : \tau_i [\bar{\alpha}
      \bar{\beta} \assign \bar{\alpha}^{\prime} \bar{\beta}^{\prime}]
      \rrbracket$. By induction hypothesis, $\mathcal{I}_i, C \wedge
      \bar{\alpha}^{\prime} \bar{\beta}^{\prime} \dot{=} \bar{\alpha}
      \bar{\beta} \vDash \Phi_i$, $i = 1, \ldots, n$.
      
      \item Observe, that (1) and (3) imply $\mathcal{I}_i, C \wedge
      \bar{\alpha}^{\prime} \bar{\beta}^{\prime} \dot{=} \bar{\alpha}
      \bar{\beta} \vDash \wedge_i \Phi_i \wedge D [\bar{\alpha} \bar{\beta}
      \assign {\nobreak} \bar{\alpha}^{\prime} \bar{\beta}^{\prime}] \wedge
      \varepsilon (\bar{\alpha}^{\prime}) \dot{=} \varepsilon (\bar{\alpha})$.
      
      \item By non-emptiness of sorts and because the premise $\tmop{PV} (C,
      \Gamma) = {\nobreak} \varnothing$ gives disjoint domains for the
      $\mathcal{I}_i$, (4) and (2) imply $\mathcal{I}_1 \ldots \mathcal{I}_n,
      C \vDash \exists \bar{\alpha}^{\prime} \bar{\beta}^{\prime} . \wedge_i
      \Phi_i \wedge D [\bar{\alpha} \bar{\beta} \assign \bar{\alpha}^{\prime}
      \bar{\beta}^{\prime}] \wedge \varepsilon (\bar{\alpha}^{\prime}) \dot{=}
      \varepsilon (\bar{\alpha})$.
      
      \item By (1) and (5), $\mathcal{I}, C \vDash \llbracket \Gamma \vdash K
      e_1 \ldots e_n : \tau \rrbracket$ for $\mathcal{I}=\mathcal{I}_1 \ldots
      \mathcal{I}_n$.
    \end{enumerate}
    \item Case {\tmname{Abs}}. In this case, $\tau \assign \tau_1 \rightarrow
    \tau_2$.
    \begin{enumerate}
      \item {\tmname{Abs}}' premise is $C, \Gamma \vdash \bar{c} : \tau_1
      \rightarrow \tau_2$, which by induction hypothesis implies
      $\mathcal{I}_i, C \vDash \Phi_i$ for $\Phi_i = \llbracket \Gamma \vdash
      p_i .e_i : \tau_1 \rightarrow \tau_2 \rrbracket$, $i = 1, \ldots, n$.
      
      \item Let $\alpha_1 \alpha_2 \# \tmop{FV} (C, \tau_1, \tau_2)$. Then,
      because sorts are nonempty, $C \vDash \exists \alpha_1 \alpha_2 . (C
      \wedge \alpha_1 \dot{=} \tau_1 \wedge \alpha_2 \dot{=} \tau_2)$.
      
      \item (1) and the premise implies $\mathcal{I}_1 \mathcal{I}_2, C \wedge
      \alpha_1 \dot{=} \tau_1 \wedge \alpha_2 \dot{=} \tau_2 \vDash \wedge_i
      \Phi_i \wedge \alpha_1 \rightarrow \alpha_2 \dot{=} \tau_1 \rightarrow
      \tau_2$.
      
      \item Combining (2) and (3), $\mathcal{I}_1 \mathcal{I}_2, C \vDash
      \exists \alpha_1 \alpha_2 . (\wedge_i \Phi_i \wedge \alpha_1 \rightarrow
      \alpha_2 \dot{=} \tau_1 \rightarrow \tau_2)$.
      
      \item By (1) and (4), $\mathcal{I}_1 \mathcal{I}_2, C \vDash \llbracket
      \Gamma \vdash \lambda \bar{c} : \tau \rrbracket$.
    \end{enumerate}
    \item Case {\tmname{App}}.
    \begin{enumerate}
      \item {\tmname{App}}'s premises are $C, \Gamma \vdash e_1 :
      \tau^{\prime} \rightarrow \tau$ and $C, \Gamma \vdash e_2 :
      \tau^{\prime}$.
      
      \item Pick w.l.o.g. $\alpha \nin \tmop{FV} (C, \tau^{\prime}, \Gamma,
      \tau)$. By rule {\tmname{Equ}}, (1) implies $C \wedge \alpha \dot{=}
      \tau^{\prime}, \Gamma \vdash e_1 : \alpha \rightarrow \tau$ and $C
      \wedge \alpha \dot{=} \tau^{\prime}, \Gamma \vdash e_2 : \alpha$.
      
      \item By induction hypothesis, (2) implies $\mathcal{I}_i, C \wedge
      \alpha \dot{=} \tau^{\prime} \vDash \Phi_i$ for $\Phi_i = \llbracket
      \Gamma \vdash e_i : \tau_i \rrbracket$, $i = 1, 2, \tau_1 \assign
      \tau^{\prime} \rightarrow \tau, \tau_2 \assign \tau^{\prime}$.
      
      \item By (2) and nonemptiness of sorts, we have $C \vDash \exists \alpha
      . (C \wedge \alpha \dot{=} \tau^{\prime})$.
      
      \item By (3), the premise and because $C \vDash D$ implies $\exists
      \alpha .C \vDash \exists \alpha .D$, we have $\mathcal{I}_1
      \mathcal{I}_2, \exists \alpha . (C \wedge \alpha \dot{=} \tau^{\prime})
      \vDash \exists \alpha . (\Phi_1 \wedge \Phi_2)$.
      
      \item By (4) and (5), we have the goal $\mathcal{I}, C \vDash \llbracket
      \Gamma \vdash e_1 e_2 : \tau \rrbracket$ with $\mathcal{I}=\mathcal{I}_1
      \mathcal{I}_2$.
    \end{enumerate}
    \item Case {\tmname{LetRec}}. Let $\Gamma^{\prime} \assign \Gamma \{x
    \mapsto \sigma\}$.
    \begin{enumerate}
      \item {\tmname{LetRec}}'s premises are $C, \Gamma^{\prime} \vdash e_1 :
      \sigma$, which can only be derived by {\tmname{Gen}} from $C^{\prime}
      \wedge D, \Gamma^{\prime} \vdash e_1 : \beta$, where $\sigma = \forall
      \beta [\exists \bar{\alpha} .D] . \beta$ and $C = C^{\prime} \wedge
      \exists \beta \bar{\alpha} .D$; by induction hypothesis we get
      $\mathcal{I}_1, C^{\prime} \wedge D \vDash \Phi_1$ for $\Phi_1 =
      \llbracket \Gamma^{\prime} \vdash e_1 : \beta \rrbracket$;
      
      \item and $C, \Gamma^{\prime} \vdash e_2 : \tau$; by induction
      hypothesis we get $\mathcal{I}_2, C \vDash \Phi_2$ for $\Phi_2 =
      \llbracket \Gamma^{\prime} \vdash e_2 : \tau \rrbracket$.
      
      \item $\beta \bar{\alpha} \# \tmop{FV} (\Gamma, C^{\prime})$. W.l.o.g.,
      assume additionally that $\beta \bar{\alpha} \# \tmop{FV} (\tau)$.
      
      \item $\mathcal{I}_1, C^{\prime} \vDash \forall \beta . (\exists
      \bar{\alpha} .D) \Rightarrow \Phi_1$ iff $\mathcal{I}_1, C^{\prime}
      \vDash (\exists \bar{\alpha} .D) \Rightarrow \Phi_1$ iff $\mathcal{I}_1,
      C^{\prime} \vDash \forall \bar{\alpha} .D \Rightarrow \Phi_1$ iff
      $\mathcal{I}_1, C^{\prime} \vDash D \Rightarrow \Phi_1$ iff
      $\mathcal{I}_1, C^{\prime} \wedge D \vDash \Phi_1$, which is exactly
      (1).
      
      \item $\mathcal{I}_2, C^{\prime} \wedge \exists \beta \bar{\alpha} .D
      \vDash \forall \beta . (\exists \bar{\alpha} .D) \Rightarrow \Phi_1$
      follows from (5), $\mathcal{I}_2, C^{\prime} \wedge \exists \beta
      \bar{\alpha} .D \vDash \exists \beta . \exists \bar{\alpha} .D$, and
      $\mathcal{I}_2, C \vDash \Phi_2$ is exactly (2).
      
      \item From (4), (5) and the premise, $\mathcal{I}_1 \mathcal{I}_2, C
      \vDash (\forall \beta . (\exists \bar{\alpha} .D) \Rightarrow \Phi_1)
      \wedge (\exists \beta . \exists \bar{\alpha} .D) \wedge \Phi_2$.
      
      \item Let $\mathcal{I}=\mathcal{I}_1 \mathcal{I}_2 ; \chi \assign
      \exists \bar{\alpha} .D [\beta \assign \delta]$, where $\chi \#
      \tmop{PV} (\Gamma, \Phi_1, \Phi_2)$. (6) gives $\mathcal{I}, C \vDash
      (\forall \beta . \chi (\beta) \Rightarrow \Phi_1) \wedge (\exists \beta
      . \chi (\beta)) \wedge \Phi_2$, which is $\mathcal{I}, {\nobreak} C
      \vDash \left\llbracket \Gamma \vdash \text{{\tmstrong{let rec}}} x = e_1
      \text{{\tmstrong{in}}} e_2 : \tau \right\rrbracket$.
    \end{enumerate}
    \item Case {\tmname{Clause}}.
    \begin{enumerate}
      \item {\tmname{Clause}}'s premises are: $C \vdash p : \tau_1
      \longrightarrow \exists \bar{\beta} [D] \Gamma^{\prime}$,
      
      \item $C \wedge D \wedge_i \tau_{m_i} \leq \tau_{n_i}, \Gamma
      \Gamma^{\prime} \vdash e : \tau_2$,
      
      \item $C \wedge D, \Gamma \Gamma^{\prime} \vdash m_i : \tmop{Num}
      (\tau_{m_i})$ and $C \wedge D, \Gamma \Gamma^{\prime} \vdash n_i :
      \tmop{Num} (\tau_{n_i})$,
      
      \item and $\bar{\beta} \# \tmop{FV} (C, \Gamma, \tau_2)$.
      
      \item Assume w.l.o.g. that $\bar{\beta} \# \tmop{FV} (\tau_1)$.
      
      \item Let $\alpha_i^1 \alpha_i^2 \# \tmop{FV} (C, \tau_1, \tau_2,
      \overline{\tau_{m_i}, \tau_{n_i}})$.
      
      \item Let $\llbracket \vdash p \uparrow \tau_1 \rrbracket = \exists
      \overline{\beta^{\prime}} [D^{\prime}] \Gamma^{\prime\prime}$, where
      $\overline{\beta^{\prime}} \# \tmop{FV} (\Gamma, C, \tau_1, \tau_2,
      \bar{\beta})$.
      
      \item By lemma \ref{CompletePat}, (1) and (7) gives $C \vDash \llbracket
      \vdash p \downarrow \tau_1 \rrbracket$
      
      \item and $C \vDash \forall \overline{\beta^{\prime}} .D^{\prime}
      \Rightarrow \exists \bar{\beta} .D \wedge \Gamma^{\prime\prime} \dot{=}
      \Gamma^{\prime}$, which is equivalent to $C \wedge D^{\prime} \vDash
      \exists \bar{\beta} .D \wedge \Gamma^{\prime\prime} \dot{=}
      \Gamma^{\prime}$.
      
      \item Recall that $\left\llbracket \Gamma \vdash p
      \text{{\tmstrong{when}}} \wedge_i m_i \leqslant n_i .e : \tau_1
      \rightarrow \tau_2 \right\rrbracket = \exists \overline{\alpha^1_i
      \alpha^2_i} . \Phi$, for $\Phi = \\
      \llbracket \vdash p \downarrow \tau_1 \rrbracket \wedge \forall
      \bar{\beta}^{\prime} .D^{\prime} \Rightarrow \wedge_i \llbracket \Gamma
      \Gamma^{\prime\prime} \vdash m_i : \tmop{Num} (\alpha^1_i) \rrbracket
      \wedge_i \llbracket \Gamma \Gamma^{\prime\prime} \vdash n_i : \tmop{Num}
      (\alpha^2_i) \rrbracket \wedge (\wedge_i \alpha^1_i \leq \alpha^2_i
      \Rightarrow \llbracket \Gamma \Gamma^{\prime\prime} \vdash e : \tau_2
      \rrbracket)$.
      
      \item By lemma \ref{MonoEnvEq}, weakening and {\tmname{Equ}}, (3)
      implies $C \wedge D \wedge \Gamma^{\prime\prime} \dot{=} \Gamma^{\prime}
      \wedge \alpha^1_i \dot{=} \tau_{m_i}, \Gamma \Gamma^{\prime\prime}
      \vdash m_i : \tmop{Num} (\alpha^1_i)$ and $C \wedge D \wedge
      \Gamma^{\prime\prime} \dot{=} \Gamma^{\prime} \wedge \alpha^2_i \dot{=}
      \tau_{n_i}, \Gamma \Gamma^{\prime\prime} \vdash n_i : \tmop{Num}
      (\alpha^2_i)$.
      
      \item From (9) we have $C \wedge D^{\prime} \wedge \alpha^1_i \dot{=}
      \tau_{m_i}^{\prime} \vDash \exists \bar{\beta} .D \wedge
      \Gamma^{\prime\prime} \dot{=} \Gamma^{\prime} \wedge \alpha^1_i \dot{=}
      \tau_{m_i}$ and $C \wedge D^{\prime} \wedge \alpha^2_i \dot{=}
      \tau_{n_i}^{\prime} \vDash \exists \bar{\beta} .D \wedge
      \Gamma^{\prime\prime} \dot{=} \Gamma^{\prime} \wedge \alpha^2_i \dot{=}
      \tau_{n_i}$ for some $\tau_{m_i}^{\prime}, \tau_{n_i}^{\prime}$.
      
      \item By (11), {\tmname{Hide}} and (12), we get $C \wedge D^{\prime}
      \wedge \alpha^1_i \dot{=} \tau_{m_i}^{\prime}, \Gamma
      \Gamma^{\prime\prime} \vdash m_i : \tmop{Num} (\alpha^1_i)$ and $C
      \wedge D^{\prime} \wedge \alpha^2_i \dot{=} \tau_{n_i}^{\prime}, \Gamma
      \Gamma^{\prime\prime} \vdash n_i : \tmop{Num} (\alpha^2_i)$.
      
      \item By induction hypothesis applied to (13) we have $\mathcal{I}_i^1,
      C \wedge D^{\prime} \wedge \alpha^1_i \dot{=} \tau_{m_i}^{\prime} \vDash
      \llbracket \Gamma \Gamma^{\prime\prime} \vdash m_i : \tmop{Num}
      (\alpha^1_i) \rrbracket$ and $\mathcal{I}_i^2, C \wedge D^{\prime}
      \wedge \alpha^2_i \dot{=} \tau_{n_i}^{\prime} \vDash \llbracket \Gamma
      \Gamma^{\prime\prime} \vdash n_i : \tmop{Num} (\alpha^2_i) \rrbracket$.
      
      \item By lemma \ref{MonoEnvEq} and weakening, (2) implies $C \wedge D
      \wedge \Gamma^{\prime\prime} \dot{=} \Gamma^{\prime} \wedge_i \alpha^1_i
      \dot{=} \tau_{m_i} \wedge_i \alpha^2_i \dot{=} \tau_{n_i} \wedge_i
      \alpha_i^1 \leq \alpha_i^2, \Gamma \Gamma^{\prime\prime} \vdash e :
      \tau_2$.
      
      \item By (15), {\tmname{Hide}} and (12), we get $C \wedge D^{\prime}
      \wedge_i \alpha^1_i \dot{=} \tau_{m_i}^{\prime} \wedge_i \alpha^2_i
      \dot{=} \tau_{n_i}^{\prime} \wedge_i \alpha_i^1 \leq \alpha_i^2, \Gamma
      \Gamma^{\prime\prime} \vdash e : \tau_2$.
      
      \item By induction hypothesis, (16) gives $\mathcal{I}^3, C \wedge
      D^{\prime} \wedge_i \alpha^1_i \dot{=} \tau_{m_i}^{\prime} \wedge_i
      \alpha^2_i \dot{=} \tau_{n_i}^{\prime} \wedge_i \alpha_i^1 \leq
      \alpha_i^2 \vDash \llbracket \Gamma \Gamma^{\prime\prime} \vdash e :
      \tau_2 \rrbracket$.
      
      \item By (8), (10), (14) and (17), we get $\overline{\mathcal{I}_i^1}
      \overline{\mathcal{I}_i^2} \mathcal{I}^3, C \wedge_i \alpha^1_i \dot{=}
      \tau_{m_i}^{\prime} \wedge_i \alpha^2_i \dot{=} \tau_{n_i}^{\prime}
      \vDash \Phi$.
      
      \item By nonemptiness of the domain, from (18) we get the goal
      $\overline{\mathcal{I}_i^1} \overline{\mathcal{I}_i^2} \mathcal{I}^3, C
      \vDash \exists \overline{\alpha^1_i \alpha^2_i} . \Phi$.
    \end{enumerate}
    \item Case {\tmname{NegClause}}. The proof is nearly identical as above.
    \begin{enumerate}
      \item {\tmname{NegClause}}'s premises are: $C \vdash p : \tau_3
      \longrightarrow \exists \bar{\beta} [D] \Gamma^{\prime}$,
      
      \item $C \wedge D \wedge \tau_1 \dot{=} \tau_3 \wedge_i \tau_{m_i} \leq
      \tau_{n_i}, \Gamma \Gamma^{\prime} \vdash e : \tau_2$,
      
      \item $C \wedge D, \Gamma \Gamma^{\prime} \vdash m_i : \tmop{Num}
      (\tau_{m_i})$ and $C \wedge D, \Gamma \Gamma^{\prime} \vdash n_i :
      \tmop{Num} (\tau_{n_i})$,
      
      \item and $\bar{\beta} \# \tmop{FV} (C, \Gamma, \tau_2)$.
      
      \item Assume w.l.o.g. that $\bar{\beta} \# \tmop{FV} (\tau_3)$.
      
      \item Let $\alpha_3 \alpha_i^1 \alpha_i^2 \# \tmop{FV} (C, \tau_1,
      \tau_2, \tau_3, \overline{\tau_{m_i}, \tau_{n_i}})$.
      
      \item Let $\llbracket \vdash p \uparrow \tau_3 \rrbracket = \exists
      \overline{\beta^{\prime}} [D^{\prime}] \Gamma^{\prime\prime}$, where
      $\overline{\beta^{\prime}} \# \tmop{FV} (\Gamma, C, \tau_1, \tau_2,
      \tau_3, \bar{\beta})$.
      
      \item By lemma \ref{CompletePat}, (1) and (7) gives $C \vDash \llbracket
      \vdash p \downarrow \tau_3 \rrbracket$
      
      \item and $C \vDash \forall \overline{\beta^{\prime}} .D^{\prime}
      \Rightarrow \exists \bar{\beta} .D \wedge \Gamma^{\prime\prime} \dot{=}
      \Gamma^{\prime}$, which is equivalent to $C \wedge D^{\prime} \vDash
      \exists \bar{\beta} .D \wedge \Gamma^{\prime\prime} \dot{=}
      \Gamma^{\prime}$.
      
      \item Recall that $\left\llbracket \Gamma \vdash p
      \text{{\tmstrong{when}}} \wedge_i m_i \leqslant n_i .e : \tau_1
      \rightarrow \tau_2 \right\rrbracket = \exists \alpha_3
      \overline{\alpha^1_i \alpha^2_i} . \Phi$, for $\Phi = \\
      \llbracket \vdash p \downarrow \alpha_3 \rrbracket \wedge \forall
      \bar{\beta}^{\prime} .D^{\prime} \Rightarrow \wedge_i \llbracket \Gamma
      \Gamma^{\prime\prime} \vdash m_i : \tmop{Num} (\alpha^1_i) \rrbracket
      \wedge_i \llbracket \Gamma \Gamma^{\prime\prime} \vdash n_i : \tmop{Num}
      (\alpha^2_i) \rrbracket \wedge \\
      (\tau_1 \dot{=} \alpha_3 \wedge_i \alpha^1_i \leq \alpha^2_i
      \Rightarrow \llbracket \Gamma \Gamma^{\prime\prime} \vdash e : \tau_2
      \rrbracket)$.
      
      \item By lemma \ref{MonoEnvEq}, weakening and {\tmname{Equ}}, (3)
      implies $C \wedge D \wedge \Gamma^{\prime\prime} \dot{=} \Gamma^{\prime}
      \wedge \alpha^1_i \dot{=} \tau_{m_i}, \Gamma \Gamma^{\prime\prime}
      \vdash m_i : \tmop{Num} (\alpha^1_i)$ and $C \wedge D \wedge
      \Gamma^{\prime\prime} \dot{=} \Gamma^{\prime} \wedge \alpha^2_i \dot{=}
      \tau_{n_i}, \Gamma \Gamma^{\prime\prime} \vdash n_i : \tmop{Num}
      (\alpha^2_i)$.
      
      \item From (9) we have $C \wedge D^{\prime} \wedge \alpha^1_i \dot{=}
      \tau_{m_i}^{\prime} \vDash \exists \bar{\beta} .D \wedge
      \Gamma^{\prime\prime} \dot{=} \Gamma^{\prime} \wedge \alpha^1_i \dot{=}
      \tau_{m_i}$ and $C \wedge D^{\prime} \wedge \alpha^2_i \dot{=}
      \tau_{n_i}^{\prime} \vDash \exists \bar{\beta} .D \wedge
      \Gamma^{\prime\prime} \dot{=} \Gamma^{\prime} \wedge \alpha^2_i \dot{=}
      \tau_{n_i}$ for some $\tau_{m_i}^{\prime}, \tau_{n_i}^{\prime}$.
      
      \item By (11), {\tmname{Hide}} and (12), we get $C \wedge D^{\prime}
      \wedge \alpha^1_i \dot{=} \tau_{m_i}^{\prime}, \Gamma
      \Gamma^{\prime\prime} \vdash m_i : \tmop{Num} (\alpha^1_i)$ and $C
      \wedge D^{\prime} \wedge \alpha^2_i \dot{=} \tau_{n_i}^{\prime}, \Gamma
      \Gamma^{\prime\prime} \vdash n_i : \tmop{Num} (\alpha^2_i)$.
      
      \item By induction hypothesis applied to (13) we have $\mathcal{I}_i^1,
      C \wedge D^{\prime} \wedge \alpha^1_i \dot{=} \tau_{m_i}^{\prime} \vDash
      \llbracket \Gamma \Gamma^{\prime\prime} \vdash m_i : \tmop{Num}
      (\alpha^1_i) \rrbracket$ and $\mathcal{I}_i^2, C \wedge D^{\prime}
      \wedge \alpha^2_i \dot{=} \tau_{n_i}^{\prime} \vDash \llbracket \Gamma
      \Gamma^{\prime\prime} \vdash n_i : \tmop{Num} (\alpha^2_i) \rrbracket$.
      
      \item By lemma \ref{MonoEnvEq} and weakening, (2) implies $C \wedge D
      \wedge \Gamma^{\prime\prime} \dot{=} \Gamma^{\prime} \wedge \alpha_3
      \dot{=} \tau_3 \wedge_i \alpha^1_i \dot{=} \tau_{m_i} \wedge_i
      \alpha^2_i \dot{=} \tau_{n_i} \wedge \tau_1 \dot{=} \alpha_3 \wedge_i
      \alpha_i^1 \leq \alpha_i^2, \Gamma \Gamma^{\prime\prime} \vdash e :
      \tau_2$.
      
      \item By (15), {\tmname{Hide}} and (12), we get $C \wedge D^{\prime}
      \wedge \alpha_3 \dot{=} \tau_3 \wedge_i \alpha^1_i \dot{=}
      \tau_{m_i}^{\prime} \wedge_i \alpha^2_i \dot{=} \tau_{n_i}^{\prime}
      \wedge \tau_1 \dot{=} \alpha_3 \wedge_i \alpha_i^1 \leq \alpha_i^2,
      \Gamma \Gamma^{\prime\prime} \vdash e : \tau_2$.
      
      \item By induction hypothesis, (16) gives $\mathcal{I}^3, C \wedge
      D^{\prime} \wedge \alpha_3 \dot{=} \tau_3 \wedge_i \alpha^1_i \dot{=}
      \tau_{m_i}^{\prime} \wedge_i \alpha^2_i \dot{=} \tau_{n_i}^{\prime}
      \wedge \tau_1 \dot{=} \alpha_3 \wedge_i \alpha_i^1 \leq \alpha_i^2
      \vDash \llbracket \Gamma \Gamma^{\prime\prime} \vdash e : \tau_2
      \rrbracket$.
      
      \item By (8), (10), (14) and (17), we get $\overline{\mathcal{I}_i^1}
      \overline{\mathcal{I}_i^2} \mathcal{I}^3, C \wedge \alpha_3 \dot{=}
      \tau_3 \wedge_i \alpha^1_i \dot{=} \tau_{m_i}^{\prime} \wedge_i
      \alpha^2_i \dot{=} \tau_{n_i}^{\prime} \vDash \Phi$.
      
      \item By nonemptiness of the domains, from (18) we get the goal
      $\overline{\mathcal{I}_i^1} \overline{\mathcal{I}_i^2} \mathcal{I}^3, C
      \vDash \exists \alpha_3 \overline{\alpha^1_i \alpha^2_i} . \Phi$.
    \end{enumerate}
    \item Case {\tmname{FailClause}}. The proof is nearly identical as above.
    \begin{enumerate}
      \item {\tmname{FailClause}}'s premises are: $C \vdash p : \tau_3
      \longrightarrow \exists \bar{\beta} [D] \Gamma^{\prime}$,
      
      \item $C \wedge D \wedge \tau_1 \dot{=} \tau_3 \wedge_i \tau_{m_i} \leq
      \tau_{n_i}, \Gamma \Gamma^{\prime} \vdash s : \tmop{String}$,
      
      \item $C \wedge D, \Gamma \Gamma^{\prime} \vdash m_i : \tmop{Num}
      (\tau_{m_i})$ and $C \wedge D, \Gamma \Gamma^{\prime} \vdash n_i :
      \tmop{Num} (\tau_{n_i})$,
      
      \item and $\bar{\beta} \# \tmop{FV} (C, \Gamma)$.
      
      \item Assume w.l.o.g. that $\bar{\beta} \# \tmop{FV} (\tau_3)$.
      
      \item Let $\alpha_3 \alpha_i^1 \alpha_i^2 \# \tmop{FV} (C, \tau_1,
      \tau_3, \overline{\tau_{m_i}, \tau_{n_i}})$.
      
      \item Let $\llbracket \vdash p \uparrow \tau_3 \rrbracket = \exists
      \overline{\beta^{\prime}} [D^{\prime}] \Gamma^{\prime\prime}$, where
      $\overline{\beta^{\prime}} \# \tmop{FV} (\Gamma, C, \tau_1, \tau_3,
      \bar{\beta})$.
      
      \item By lemma \ref{CompletePat}, (1) and (7) gives $C \vDash \llbracket
      \vdash p \downarrow \tau_3 \rrbracket$
      
      \item and $C \vDash \forall \overline{\beta^{\prime}} .D^{\prime}
      \Rightarrow \exists \bar{\beta} .D \wedge \Gamma^{\prime\prime} \dot{=}
      \Gamma^{\prime}$, which is equivalent to $C \wedge D^{\prime} \vDash
      \exists \bar{\beta} .D \wedge \Gamma^{\prime\prime} \dot{=}
      \Gamma^{\prime}$.
      
      \item Recall that $\left\llbracket \Gamma \vdash p
      \text{{\tmstrong{when}}} \wedge_i m_i \leqslant n_i .e : \tau_1
      \rightarrow \tau_2 \right\rrbracket = \exists \alpha_3
      \overline{\alpha^1_i \alpha^2_i} . \Phi$, for $\Phi = \\
      \llbracket \vdash p \downarrow \alpha_3 \rrbracket \wedge \forall
      \bar{\beta}^{\prime} .D^{\prime} \Rightarrow \wedge_i \llbracket \Gamma
      \Gamma^{\prime\prime} \vdash m_i : \tmop{Num} (\alpha^1_i) \rrbracket
      \wedge_i \llbracket \Gamma \Gamma^{\prime\prime} \vdash n_i : \tmop{Num}
      (\alpha^2_i) \rrbracket \wedge \\
      (\tau_1 \dot{=} \alpha_3 \wedge_i \alpha^1_i \leq \alpha^2_i
      \Rightarrow \llbracket \Gamma \Gamma^{\prime\prime} \vdash e : \tau_2
      \rrbracket)$.
      
      \item By lemma \ref{MonoEnvEq}, weakening and {\tmname{Equ}}, (3)
      implies $C \wedge D \wedge \Gamma^{\prime\prime} \dot{=} \Gamma^{\prime}
      \wedge \alpha^1_i \dot{=} \tau_{m_i}, \Gamma \Gamma^{\prime\prime}
      \vdash m_i : \tmop{Num} (\alpha^1_i)$ and $C \wedge D \wedge
      \Gamma^{\prime\prime} \dot{=} \Gamma^{\prime} \wedge \alpha^2_i \dot{=}
      \tau_{n_i}, \Gamma \Gamma^{\prime\prime} \vdash n_i : \tmop{Num}
      (\alpha^2_i)$.
      
      \item From (9) we have $C \wedge D^{\prime} \wedge \alpha^1_i \dot{=}
      \tau_{m_i}^{\prime} \vDash \exists \bar{\beta} .D \wedge
      \Gamma^{\prime\prime} \dot{=} \Gamma^{\prime} \wedge \alpha^1_i \dot{=}
      \tau_{m_i}$ and $C \wedge D^{\prime} \wedge \alpha^2_i \dot{=}
      \tau_{n_i}^{\prime} \vDash \exists \bar{\beta} .D \wedge
      \Gamma^{\prime\prime} \dot{=} \Gamma^{\prime} \wedge \alpha^2_i \dot{=}
      \tau_{n_i}$ for some $\tau_{m_i}^{\prime}, \tau_{n_i}^{\prime}$.
      
      \item By (11), {\tmname{Hide}} and (12), we get $C \wedge D^{\prime}
      \wedge \alpha^1_i \dot{=} \tau_{m_i}^{\prime}, \Gamma
      \Gamma^{\prime\prime} \vdash m_i : \tmop{Num} (\alpha^1_i)$ and $C
      \wedge D^{\prime} \wedge \alpha^2_i \dot{=} \tau_{n_i}^{\prime}, \Gamma
      \Gamma^{\prime\prime} \vdash n_i : \tmop{Num} (\alpha^2_i)$.
      
      \item By induction hypothesis applied to (13) we have $\mathcal{I}_i^1,
      C \wedge D^{\prime} \wedge \alpha^1_i \dot{=} \tau_{m_i}^{\prime} \vDash
      \llbracket \Gamma \Gamma^{\prime\prime} \vdash m_i : \tmop{Num}
      (\alpha^1_i) \rrbracket$ and $\mathcal{I}_i^2, C \wedge D^{\prime}
      \wedge \alpha^2_i \dot{=} \tau_{n_i}^{\prime} \vDash \llbracket \Gamma
      \Gamma^{\prime\prime} \vdash n_i : \tmop{Num} (\alpha^2_i) \rrbracket$.
      
      \item By lemma \ref{MonoEnvEq} and weakening, (2) implies $C \wedge D
      \wedge \Gamma^{\prime\prime} \dot{=} \Gamma^{\prime} \wedge \alpha_3
      \dot{=} \tau_3 \wedge_i \alpha^1_i \dot{=} \tau_{m_i} \wedge_i
      \alpha^2_i \dot{=} \tau_{n_i} \wedge \tau_1 \dot{=} \alpha_3 \wedge_i
      \alpha_i^1 \leq \alpha_i^2, \Gamma \Gamma^{\prime\prime} \vdash s :
      \tmop{String}$.
      
      \item By (15), {\tmname{Hide}} and (12), we get $C \wedge D^{\prime}
      \wedge \alpha_3 \dot{=} \tau_3 \wedge_i \alpha^1_i \dot{=}
      \tau_{m_i}^{\prime} \wedge_i \alpha^2_i \dot{=} \tau_{n_i}^{\prime}
      \wedge \tau_1 \dot{=} \alpha_3 \wedge_i \alpha_i^1 \leq \alpha_i^2,
      \Gamma \Gamma^{\prime\prime} \vdash s : \tmop{String}$.
      
      \item By induction hypothesis, (16) gives $\mathcal{I}^3, C \wedge
      D^{\prime} \wedge \alpha_3 \dot{=} \tau_3 \wedge_i \alpha^1_i \dot{=}
      \tau_{m_i}^{\prime} \wedge_i \alpha^2_i \dot{=} \tau_{n_i}^{\prime}
      \wedge \tau_1 \dot{=} \alpha_3 \wedge_i \alpha_i^1 \leq \alpha_i^2
      \vDash \llbracket \Gamma \Gamma^{\prime\prime} \vdash s : \tmop{String}
      \rrbracket$.
      
      \item By (8), (10), (14) and (17), we get $\overline{\mathcal{I}_i^1}
      \overline{\mathcal{I}_i^2} \mathcal{I}^3, C \wedge \alpha_3 \dot{=}
      \tau_3 \wedge_i \alpha^1_i \dot{=} \tau_{m_i}^{\prime} \wedge_i
      \alpha^2_i \dot{=} \tau_{n_i}^{\prime} \vDash \Phi$.
      
      \item By nonemptiness of the domains, from (18) we get the goal
      $\overline{\mathcal{I}_i^1} \overline{\mathcal{I}_i^2} \mathcal{I}^3, C
      \vDash \exists \alpha_3 \overline{\alpha^1_i \alpha^2_i} . \Phi$.
    \end{enumerate}
    \item Case {\tmname{Equ}}.
    \begin{enumerate}
      \item {\tmname{Equ}}'s premises are $C, \Gamma \vdash e :
      \tau^{\prime}$, which by induction hypothesis gives $\mathcal{I}, C
      \vDash \llbracket \Gamma \vdash e : \tau^{\prime} \rrbracket$,
      
      \item and $C \vDash \tau^{\prime} \dot{=} \tau$.
      
      \item Let $\Phi_{\tau} \assign \llbracket \Gamma \vdash e : \tau
      \rrbracket$. Observe, that $\tau$ occurs in $\Phi_{\tau}$ only as a
      subterm in a side of equation: $\dot{=} \tau$, $\dot{=} \ldots
      \rightarrow \tau$, $\dot{=} (\ldots \rightarrow (\ldots \rightarrow
      \tau) \ldots)$. Therefore, $\tau^{\prime} \dot{=} \tau \vDash
      \Phi_{\tau^{\prime}} \Leftrightarrow \Phi_{\tau}$.
      
      \item (1), (2) and (3) imply that $\mathcal{I}, C \vDash \llbracket
      \Gamma \vdash e : \tau \rrbracket$.
    \end{enumerate}
    \item Case {\tmname{Hide}}.
    \begin{enumerate}
      \item {\tmname{Hide}}'s premises are $C, \Gamma \vdash e : \tau$, that
      by induction hypothesis gives $\mathcal{I}, C \vDash \llbracket \Gamma
      \vdash e : \tau \rrbracket$,
      
      \item and $\bar{\beta} \# \tmop{FV} (\Gamma, \tau)$.
      
      \item By (2), w.l.o.g. $\bar{\beta} \# \tmop{FV} (\llbracket \Gamma
      \vdash e : \tau \rrbracket)$.
      
      \item (1) implies that $\mathcal{I} \vDash \forall \bar{\beta} . (C
      \Rightarrow \Phi_1)$ which by (3) is equivalent to $\mathcal{I}, \exists
      \bar{\beta} .C \vDash \llbracket \Gamma \vdash e : \tau \rrbracket$.
    \end{enumerate}
    \item Case {\tmname{FElim}}. $\mathcal{I}, \tmmathbf{F} \vDash \Phi$ holds
    for any $\Phi$.
    
    \item Case {\tmname{DisjElim}}.
    \begin{enumerate}
      \item {\tmname{DisjElim}} premises are $C, \Gamma \vdash e : \tau$ and
      $D, \Gamma \vdash e : \tau$. Induction hypothesis gives $\mathcal{I}_1,
      C \vDash \llbracket \Gamma \vdash e : \tau \rrbracket$ and
      $\mathcal{I}_2, D \vDash \llbracket \Gamma \vdash e : \tau \rrbracket$
      for some interpretations of predicate variables $\mathcal{I}_1,
      \mathcal{I}_2$.
      
      \item Therefore, we have $\mathcal{I}, C \vee D \vDash \llbracket \Gamma
      \vdash e : \tau \rrbracket$, for both $\mathcal{I}=\mathcal{I}_1$ and
      $\mathcal{I}=\mathcal{I}_2$.
    \end{enumerate}
  \end{itemize}
\end{proof}

Proof of corollary \ref{InfComplCoro}.

\begin{proof}
  $C, \Gamma \vdash e : \forall \bar{\alpha} [D] . \tau$ can only be derived
  by the {\tmname{Gen}} rule, therefore we have $C^{\prime} \wedge D, \Gamma
  \vdash e : \tau$ for $\bar{\alpha} \# \tmop{FV} (\Gamma, C^{\prime})$ and $C
  = C^{\prime} \wedge \exists \bar{\alpha} .D$. By theorem \ref{InfComplExpr},
  there exists an interpretation $\mathcal{I}$ such that $\mathcal{I},
  C^{\prime} \wedge D \vDash \llbracket \Gamma \vdash e : \tau \rrbracket$.
  $\mathcal{I}, C^{\prime} \wedge D \vDash \llbracket \Gamma \vdash e : \tau
  \rrbracket$ iff $\mathcal{I} \vDash C^{\prime} \wedge D \Rightarrow
  \llbracket \Gamma \vdash e : \tau \rrbracket$ iff $\mathcal{I} \vDash
  \forall \bar{\alpha} .C^{\prime} \wedge D \Rightarrow \llbracket \Gamma
  \vdash e : \tau \rrbracket$ iff $\mathcal{I}, C^{\prime} \vDash \forall
  \bar{\alpha} .D \Rightarrow \llbracket \Gamma \vdash e : \tau \rrbracket$.
  Therefore $\mathcal{I}, C \vDash \forall \bar{\alpha} .D \Rightarrow
  \llbracket \Gamma \vdash e : \tau \rrbracket$.
\end{proof}

\subsection{Existential Types}\label{ExGADTsProofs}

Proof of theorem \ref{InfCorrComplEx}.

\begin{proof}
  By inspecting Table \ref{NormExpr}, note that $\lambda [K] e$ subexpressions
  are absent from $n (e)$. Thus $\mathcal{I}_e$ is empty in all cases other
  than {\tmname{ExIntro}}. We therefore shorten these cases by not mentioning
  $\mathcal{I}_e$ and $\Sigma$. Below we extend the inductive proofs with the
  cases for expressions introduced by, or rule applications of,
  {\tmname{ExIntro}}, {\tmname{LetIn}} and {\tmname{ExLetIn}}.
  \begin{itemize}
    \item Theorem \ref{InfCorrExpr} (Correctness) $\llbracket \Gamma, \Sigma_0
    \vdash e : \tau \rrbracket, \Gamma, \Sigma_0 \vdash e : \tau$. Case:
    $\mathcal{E} (e) \neq \varnothing$.
    \begin{enumerate}
      \item Induction hypothesis states $\llbracket \Gamma, \Sigma \vdash n
      (e) : \tau \rrbracket, \Gamma, \Sigma \vdash n (e) : \tau$.
      
      \item The goal follows by {\tmname{ExIntro}}.
    \end{enumerate}
    \item Theorem \ref{InfCorrExpr} (Correctness) Case: $e$ is
    $\text{{\tmstrong{let}}} p = e_1  \text{{\tmstrong{in}}} e_2$.
    \begin{enumerate}
      \item Induction hypothesis yields $\llbracket \Gamma \vdash K p.e_2 :
      \alpha_0 \rightarrow \tau \rrbracket, \Gamma \vdash K p.e_2 : \alpha_0
      \rightarrow \tau$, $\llbracket \Gamma \vdash p.e_2 : \alpha_0
      \rightarrow \tau \rrbracket, \Gamma \vdash p.e_2 : \alpha_0 \rightarrow
      \tau$ and $\llbracket \Gamma \vdash e_1 : \alpha_0 \rrbracket, \Gamma
      \vdash e_1 : \alpha_0$.
      
      \item By weakening, (1), {\tmname{Abs}} and {\tmname{App}}, we get
      $\llbracket \Gamma \vdash e_1 : \alpha_0 \rrbracket \wedge \llbracket
      \Gamma \vdash p.e_2 : \alpha_0 \rightarrow \tau \rrbracket \wedge
      \not{E} (\alpha_0), \Gamma \vdash \lambda (p.e_2) e_1 : \tau$.
      
      \item By {\tmname{ExLetIn}} we get $\llbracket \Gamma \vdash e_1 :
      \alpha_0 \rrbracket \wedge \llbracket \Gamma \vdash K p.e_2 : \alpha_0
      \rightarrow \tau \rrbracket, \Gamma \vdash \text{{\tmstrong{let}}} p =
      e_1  \text{{\tmstrong{in}}} e_2 : \tau$, and by {\tmname{LetIn}}:
      $\llbracket \Gamma \vdash e_1 : \alpha_0 \rrbracket \wedge \llbracket
      \Gamma \vdash p.e_2 : \alpha_0 \rightarrow \tau \rrbracket \wedge
      \not{E} (\alpha_0), \Gamma \vdash \text{{\tmstrong{let}}} p = e_1 
      \text{{\tmstrong{in}}} e_2 : \tau$.
      
      \item By (3) and {\tmname{DisjElim}} we get \ $\left( \llbracket \Gamma
      \vdash e_1 : \alpha_0 \rrbracket \wedge \llbracket \Gamma \vdash p.e_2 :
      \alpha_0 \rightarrow \tau \rrbracket \wedge \not{E} (\alpha_0) \right)
      \vee_{\mathcal{E}} (\llbracket \Gamma \vdash e_1 : \alpha_0 \rrbracket
      \wedge \llbracket \Gamma \vdash K p.e_2 : \alpha_0 \rightarrow \tau
      \rrbracket), \Gamma \vdash \text{{\tmstrong{let}}} p = e_1 
      \text{{\tmstrong{in}}} e_2 : {\nobreak} \tau$ for $\mathcal{E}= \{ K | K
      \colons \forall \overline{\alpha_K} \bar{\beta} [E] . \tau \rightarrow
      \varepsilon_K (\overline{\alpha_K}) \}$.
      
      \item By (4), weakening and {\tmname{Hide}}, we get the goal.
    \end{enumerate}
    \item Theorem \ref{InfCorrExpr} (Correctness) Case: $e$ is $e_1 e_2$.
    \begin{enumerate}
      \item Let $\alpha \# \tmop{FV} (\Gamma, \tau)$.
      
      \item By the induction hypothesis, we have $\llbracket \Gamma \vdash e_1
      : \alpha \rightarrow \tau \rrbracket, \Gamma \vdash e_1 : \alpha
      \rightarrow \tau$ and $\llbracket \Gamma \vdash e_2 : \alpha \rrbracket,
      \Gamma \vdash e_2 : \alpha$.
      
      \item By weakening and {\tmname{App}}, this yields $\llbracket \Gamma
      \vdash e_1 : \alpha \rightarrow \tau \rrbracket \wedge \left\llbracket
      \Gamma \vdash e_2 : {\nobreak} \alpha \right\rrbracket \wedge \not{E}
      (\alpha), \Gamma \vdash e_1 e_2 : {\nobreak} \tau$.
      
      \item By {\tmname{Hide}} using (1), $\llbracket \Gamma \vdash e_1 e_2 :
      \tau \rrbracket, \Gamma \vdash e_1 e_2 : \tau$.
    \end{enumerate}
    \item Theorem \ref{InfCorrExpr} (Correctness) Case: $e$ is $\lambda_K
    \bar{c}$.
    \begin{enumerate}
      \item Let $\alpha_0, \alpha_1 \# \tmop{FV} (\Gamma, \tau)$.
      
      \item By the induction hypothesis, we have $\llbracket \Gamma \vdash
      \lambda \bar{c} : \tau \rrbracket, \Gamma \vdash \lambda \bar{c} :
      \tau$.
      
      \item By weakening, {\tmname{Equ}}, {\tmname{Hide}} and
      {\tmname{ExAbs}}, we have $\llbracket \Gamma \vdash \lambda \bar{c} :
      \tau \rrbracket \wedge \tmop{RetType} \left( \tau, {\nobreak} \alpha_0
      \right) \wedge \exists \alpha_1 . \alpha_0 \dot{=} \varepsilon_K
      (\alpha_1), \Gamma \vdash \lambda_K \bar{c} : \tau$.
      
      \item By weakening, this yields the goal.
    \end{enumerate}
    \item Theorem \ref{InfComplExpr} (Completeness) Case {\tmname{ExIntro}}:
    premise $C, \Gamma, \Sigma^{\prime} \vdash n (e) : \tau$ for $\tmop{Dom}
    (\Sigma^{\prime}) \backslash \tmop{Dom} (\Sigma) =\mathcal{E} (e)$.
    \begin{enumerate}
      \item By induction hypothesis we have $\mathcal{I}_u, C \vDash
      \llbracket \Gamma, \Sigma^{\prime} \vdash n (e) : \tau \rrbracket$.
      
      \item Let $\Sigma_1 = \Sigma \overline{K \colons \forall \alpha_K
      \gamma_K [\chi_K (\gamma_K, \alpha_K)] . \gamma_K \rightarrow
      \varepsilon_K (\alpha_K)}$. The goal is $\mathcal{I}_u, C \vDash
      \mathcal{I}_e (\llbracket \Gamma, \Sigma_1 \vdash n (e) : \tau
      \rrbracket) [\overline{\varepsilon_K (\vec{\tau})} \assign
      \overline{\varepsilon_K (\bar{\tau})}]$.
      
      \item The goal follows by setting $\mathcal{I}_e = \Sigma^{\prime} /
      \Sigma$.
    \end{enumerate}
    \item Theorem \ref{InfComplExpr} (Completeness) Case {\tmname{App}}.
    \begin{enumerate}
      \item {\tmname{App}}'s premises are $C, \Gamma \vdash e_1 :
      \tau^{\prime} \rightarrow \tau$, $C, \Gamma \vdash e_2 : \tau^{\prime}$
      and $C \vDash \not{E} (\tau^{\prime})$.
      
      \item Pick w.l.o.g. $\alpha \nin \tmop{FV} (C, \tau^{\prime}, \Gamma,
      \tau)$. (1) implies $C \wedge \alpha \dot{=} \tau^{\prime} \vDash
      \not{E} (\alpha)$. By rule {\tmname{Equ}}, (1) implies $C \wedge \alpha
      \dot{=} \tau^{\prime}, \Gamma \vdash e_1 : \alpha \rightarrow \tau$ and
      $C \wedge \alpha \dot{=} \tau^{\prime}, \Gamma \vdash e_2 : \alpha$.
      
      \item By induction hypothesis, (2) implies $\mathcal{I}_i, C \wedge
      \alpha \dot{=} \tau^{\prime} \vDash \Phi_i$ for $\Phi_i = \llbracket
      \Gamma \vdash e_i : \tau_i \rrbracket$, $i = 1, 2, \tau_1 \assign
      \tau^{\prime} \rightarrow \tau, \tau_2 \assign \tau^{\prime}$.
      
      \item By (2) and nonemptiness of sorts, we have $C \vDash \exists \alpha
      . (C \wedge \alpha \dot{=} \tau^{\prime})$.
      
      \item By (2), (3), and because $C \vDash D$ implies $\exists \alpha .C
      \vDash \exists \alpha .D$, we have $\mathcal{I}_1 \mathcal{I}_2, \exists
      \alpha . (C \wedge \alpha \dot{=} \tau^{\prime}) \vDash \exists \alpha .
      (\Phi_1 \wedge \Phi_2 \wedge \not{E} (\alpha))$.
      
      \item By (4) and (5), we have the goal $\mathcal{I}, C \vDash \llbracket
      \Gamma \vdash e_1 e_2 : \tau \rrbracket$ with $\mathcal{I}=\mathcal{I}_1
      \mathcal{I}_2$.
    \end{enumerate}
    \item Theorem \ref{InfComplExpr} (Completeness) Case {\tmname{LetIn}}:
    premise $C, \Gamma \vdash \text{{\tmstrong{let}}} p = e_1 
    \text{{\tmstrong{in}}} e_2 : \tau$.
    \begin{enumerate}
      \item {\tmname{LetIn}}'s premise is: $C, \Gamma \vdash \lambda (p.e_2)
      e_1 : \tau$,
      
      \item derived by {\tmname{App}} and {\tmname{Abs}} from $C, \Gamma
      \vdash p.e_2 : \tau^{\prime} \rightarrow \tau$, $C, \Gamma \vdash e_1 :
      \tau^{\prime}$ and $C \vDash \not{E} (\tau^{\prime})$.
      
      \item Inductive hypothesis gives $\mathcal{I}_1, C \vDash \llbracket
      \Gamma \vdash p.e_2 : \tau^{\prime} \rightarrow \tau \rrbracket$ and
      $\mathcal{I}_2, C \vDash \llbracket \Gamma \vdash e_1 : \tau^{\prime}
      \rrbracket$.
      
      \item (1) and (3) imply $\mathcal{I}_1, C \vDash \llbracket \Gamma
      \vdash p.e_2 : \tau^{\prime} \rightarrow \tau \rrbracket \wedge \not{E}
      (\tau^{\prime}) \vee_{\mathcal{E}} \llbracket \Gamma \vdash K p.e_2 :
      \alpha_0 \rightarrow \tau \rrbracket$ as the first disjunct holds.
      
      \item As the premise $\tmop{PV} (C, \Gamma) = \varnothing$ gives
      disjoint domains for the $\mathcal{I}_i$, we have $\mathcal{I}, C \vDash
      \llbracket \Gamma \vdash e_1 : \tau^{\prime} \rrbracket \wedge \left(
      \llbracket \Gamma \vdash p.e_2 : \tau^{\prime} \rightarrow \tau
      \rrbracket \wedge \not{E} (\tau^{\prime}) \vee_{\mathcal{E}} \llbracket
      \Gamma \vdash K p.e_2 : \tau^{\prime} \rightarrow \tau \rrbracket
      \right)$ for $\mathcal{I}=\mathcal{I}_1 \mathcal{I}_2$.
      
      \item $\mathcal{I}, C \vDash \exists \alpha_0 . \llbracket \Gamma \vdash
      e_1 : \alpha_0 \rrbracket \wedge \left( \llbracket \Gamma \vdash p.e_2 :
      \alpha_0 \rightarrow \tau \rrbracket \wedge \not{E} (\alpha_0)
      \vee_{\mathcal{E}} \llbracket \Gamma \vdash K p.e_2 : \alpha_0
      \rightarrow \tau \rrbracket \right)$ by abstracting $\alpha_0 =
      \tau^{\prime}$.
    \end{enumerate}
    \item Theorem \ref{InfComplExpr} (Completeness) Case {\tmname{ExLetIn}}:\\
    premise $C, \Gamma \vdash \text{{\tmstrong{let}}} p = e_1 
    \text{{\tmstrong{in}}} e_2 : \tau$.
    \begin{enumerate}
      \item {\tmname{ExLetIn}}'s premises are: $C, \Gamma \vdash K p.e_2 :
      \tau^{\prime} \rightarrow \tau$ and $C, \Gamma \vdash e_1 :
      \tau^{\prime}$,
      
      \item Inductive hypothesis gives $\mathcal{I}_1, C \vDash \llbracket
      \Gamma \vdash K p.e_2 : \tau^{\prime} \rightarrow \tau \rrbracket$ and
      $\mathcal{I}_2, C \vDash \llbracket \Gamma \vdash e_1 : \tau^{\prime}
      \rrbracket$.
      
      \item (3) implies $\mathcal{I}_1, C \vDash \llbracket \Gamma \vdash
      p.e_2 : \tau^{\prime} \rightarrow \tau \rrbracket \wedge \not{E}
      (\tau^{\prime}) \vee_{\mathcal{E}} \llbracket \Gamma \vdash K p.e_2 :
      \tau^{\prime} \rightarrow \tau \rrbracket$ as one of the
      $\vee_{\mathcal{E}}$ disjuncts holds. The proof concludes as in the
      {\tmname{LetIn}} case.
    \end{enumerate}
    \item Theorem \ref{InfComplExpr} (Completeness) Case {\tmname{ExAbs}}:
    premise $C, \Gamma \vdash \lambda_K \bar{c} : \tau$.
    \begin{enumerate}
      \item {\tmname{ExAbs}}'s premises are: {\tmem{(a)}} $C, \Gamma \vdash
      \lambda \bar{c} : \tau$ and {\tmem{(b)}} $C \vDash \tmop{RetType} (\tau,
      \varepsilon_K (\bar{\tau}^{\prime}))$.
      
      \item Inductive hypothesis gives $\mathcal{I}_1, C \vDash \llbracket
      \Gamma \vdash \lambda \bar{c} : \tau \rrbracket$.
      
      \item Let $\alpha_0, \alpha_1 \# \tmop{FV} (\Gamma, \tau)$. (1b) and (2)
      give $\mathcal{I}_1, C \vDash \exists \alpha_0 . \llbracket \Gamma
      \vdash \lambda \bar{c} : \tau \rrbracket \wedge \tmop{RetType} \left(
      \tau, {\nobreak} \alpha_0 \right) \wedge \exists \alpha_1 . \alpha_0
      \dot{=} \varepsilon_K (\alpha_1) \wedge \alpha_1 \dot{=}
      \bar{\tau}^{\prime}$.
      
      \item Let $\mathcal{I}_0 = [\chi_K (\alpha) \assign \alpha \dot{=}
      \bar{\tau}^{\prime}]$. By (3), $\mathcal{I}_0 \mathcal{I}_1, C \vDash
      \exists \alpha_0 . \wedge \tmop{RetType} (\tau, \alpha_0) \wedge
      (\exists \alpha_1 . \alpha_0 \dot{=} \varepsilon_K (\alpha_1) \wedge
      \chi_K (\alpha_1)) \wedge \llbracket \Gamma \vdash \lambda \bar{c} :
      \tau \rrbracket$ which is the goal.
    \end{enumerate}
  \end{itemize}
\end{proof}

\subsection{Semantics by Reduction to $\tmop{HMG} (X)$}\label{SemanticProofs}

By induction on the structure of the derivation, we can show the following:

\begin{proposition}
  \label{PredVarSubstDeriv}Let $\mathcal{I}$ be an interpretation of predicate
  variables for formula $C$, as in Definition \ref{InterprPredVars}. If a
  typing judgment $C, \Gamma, \Sigma \vdash e : \tau$ is derivable in the type
  system $\tmop{MMG}_{\exists} (X)$, then $\mathcal{I} (C), \mathcal{I}
  (\Gamma), \mathcal{I} (\Sigma) \vdash e : \tau$ is derivable in
  $\tmop{MMG}_{\exists} (X)$.
\end{proposition}

Proof of Theorem \ref{SemRedHMG}.

\begin{proof}
  Let $\mathcal{I}$ be a subsitution of predicate variables such that
  $\mathcal{M}, \mathcal{I} \vDash \exists \tmop{FV} (\tau) . \llbracket
  \Gamma, \Sigma \vdash e : \tau \rrbracket$. By Theorems \ref{InfCorrExpr}
  and \ref{InfCorrComplEx}, there exists a derivation of $\llbracket \Gamma,
  \Sigma \vdash e : \tau \rrbracket, \Gamma, \Sigma \vdash e : \tau$. Without
  loss of generality, let {\tmname{ExIntro}} be the rule applied at the root
  of the derivation, and let $\Delta$ be the derivation of the premise
  $\llbracket \Gamma, \Sigma \vdash e : \tau \rrbracket, \Gamma,
  \Sigma^{\prime\prime} \vdash n (e) : {\nobreak} \tau$. Let $\Gamma^{\prime}
  =\mathcal{I} (\Gamma)$, $C =\mathcal{I} (\llbracket \Gamma, \Sigma \vdash e
  : \tau \rrbracket)$ and $\Sigma^{\prime} =\mathcal{I}
  (\Sigma^{\prime\prime})$. By Proposition \ref{PredVarSubstDeriv}, there
  exists a derivation $\Delta$ of $C, \Gamma^{\prime}, \Sigma^{\prime} \vdash
  e : {\nobreak} \tau$. We need to construct a derivation $\Delta^{\prime}$ of
  $C^{\prime}, \Gamma^{\prime} \vdash e^{\prime} : \tau$ under constructor
  environment $\Sigma^{\prime}$ in the $\tmop{HMG} (X)$ type system. The tag
  erasure of $e^{\prime}$ w.r.t. $\Sigma$ has to be computationally equivalent
  to the HMG-form of $e$ and $C^{\prime}$ has to be satisfiable. We will
  transform $\Delta$ into a derivation in $\tmop{HMG} (X)$ while preserving
  these conditions. We transform the resulting constraint, the resulting
  expression and the derivation simultaneously, as follows:
  \begin{itemize}
    \item For a derivation node applying rule {\tmname{App}}, we erase the $C
    \vDash \not{E} (\tau^{\prime})$ condition in the premise of the node, and
    we erase $\not{E} (\alpha)$ in the corresponding $\llbracket \Gamma \vdash
    e_1 e_2 : \tau \rrbracket$ part of the resulting constraint, where $C,
    \tau^{\prime}, \alpha$ etc. name the actual pieces of the derivation or
    constraint.
    
    \item For a derivation node {\tmname{Abs}}, with $\lambda (\overline{p_i
    .e_i})$ in conclusion, we erase derivation subtrees with
    {\tmname{NegClause}} or {\tmname{FailClause}} nodes at the root.
    Correspondingly, we replace the expression in the conclusion by $\lambda
    (\overline{(p_i .e_i^{\prime})}_{i : \neg \tmop{unreach} (e_i^{\prime})
    \vee \tmop{failure} (e_i)})$, which preserves computational equivalence;
    and we erase the $\llbracket \Gamma \vdash p_i .e_i : \tau_1 \rightarrow
    \tau_2 \rrbracket$ conjuncts of the resulting constraint for $i :
    \tmop{unreach} (e_i) \vee \tmop{failure} (e_i)$, which leads to a weaker
    formula and thus preserves satisfiability.
    
    \item We replace derivation nodes {\tmname{ExLetIn}} by applications of
    the {\tmname{App}} rule to the premises $C, \Gamma \vdash (K p.e_2) e_1 :
    \tau^{\prime} \rightarrow \tau$ and $C, \Gamma \vdash e_1 :
    \tau^{\prime}$, and replace the corresponding subexpressions
    $\text{{\tmstrong{let}}} p = e_1  \text{{\tmstrong{in}}} e_2$ of the
    resulting expression by $(K p.e_2) e_1$. We replace the $\left\llbracket
    \Gamma \vdash \text{{\tmstrong{let}}} p = e_1  \text{{\tmstrong{in}}} e_2
    : \tau \right\rrbracket$ part of the resulting constraint by $\exists
    \alpha_0 . \llbracket \Gamma \vdash e_1 : \alpha_0 \rrbracket \wedge
    \llbracket \Gamma \vdash K p.e_2 : \alpha_0 \rightarrow \tau \rrbracket$.
    To see that the satisfiability of the constraint is preserved, recall that
    $\left\llbracket \Gamma \vdash \text{{\tmstrong{let}}} p = e_1 
    \text{{\tmstrong{in}}} e_2 : \tau \right\rrbracket$ is a disjunction where
    only one disjunct is consistent with $\exists \bar{\alpha} . \tau^{\prime}
    \dot{=} \varepsilon_K (\bar{\alpha})$. The tag erasure of $(K p.e_2) e_1$
    w.r.t. the original constructor environment $\Sigma$ is computationally
    equivalent to $\text{{\tmstrong{let}}} p = e_1  \text{{\tmstrong{in}}}
    e_2$.
    
    \item We excise derivation nodes {\tmname{ExAbs}}, and remove the
    corresponding $\tmop{RetType}$ atom from $C$, preserving computational
    equivalence and satisfiability.
    
    \item For remaining nodes are easy to rearrange into applications of the
    corresponding $\tmop{HMG} (X)$ rules.
  \end{itemize}
\end{proof}

Proof of Corollary \ref{TypeSoundness}.

\begin{proof}
  The semantics of expression $e$ in $\tmop{MMG}_{\exists} (X)$ is given by
  the semantics of its HMG-form in $\tmop{HMG} (X)$. By well-typedness and
  closedness, we have $C, \varnothing, \Sigma \vdash e : \sigma$ is derivable
  for some type scheme $\sigma$, and $\mathcal{M} \vDash C$. By Theorems
  \ref{InfComplExpr} and \ref{InfCorrComplEx}, there exists an interpretation
  of predicate variables $\mathcal{I}$ such that $\mathcal{M}, \mathcal{I}
  \vDash C \Rightarrow \left\llbracket \varnothing, \Sigma \vdash e :
  {\nobreak} \sigma \right\rrbracket$, thus $\mathcal{M}, \mathcal{I} \vDash
  \left\llbracket \varnothing, \Sigma \vdash e : {\nobreak} \sigma
  \right\rrbracket$.
  
  By Theorem \ref{SemRedHMG}, we have an $\tmop{HMG} (X)$ environment
  $\Gamma^{\prime}$, constructor environment $\Sigma^{\prime}$ and an
  expression $e^{\prime}$ such that $C, \Gamma^{\prime} \vdash e^{\prime} :
  \tau$ is derivable in $\tmop{HMG} (X)$ for a valid $C$, and the tag erasure
  of $e^{\prime}$ w.r.t. $\Sigma$ is computationally equivalent to the
  HMG-form of $e$. By {\cite{simonet-pottier-hmg-toplas}} Theorem 3.31, see
  also Theorem \ref{HMGSoundness}, $e^{\prime}$ does not go wrong. Given the
  call-by-value semantics presented in {\cite{simonet-pottier-hmg-toplas}}, it
  is easy to see that the tag erasure of $e^{\prime}$ does not go wrong. By
  computational equivalence, HMG-form of $e$ does not go wrong.
\end{proof}

\section{Constraint Abduction}

\subsection{Formulating the Joint Constraint Abduction Problem}

\begin{proposition}
  \label{SolvedForm}{\tmem{Solved form property for terms}}. Let $\bar{\beta}
  = \tmop{Dom} (\tmmathbf{U}(C))$ be all variables occurring on the left-hand
  sides of solved form equations $\tmmathbf{U} (C)$, and $\alpha \in \tmop{FV}
  (\tmop{Image} (\tmmathbf{U}(C))$ be a variable occurring on the right-hand
  side. If $T \vDash \mathcal{Q}.C$ (equivalently $\mathcal{Q}.\tmmathbf{U}
  (C)$) holds, then $T \vDash \mathcal{Q}.C [\alpha \assign \tau]$
  (equivalently $T \vDash \mathcal{Q}.\tmmathbf{U} (C) [\alpha \assign
  \tau]$), for any $\tau$ such that for all $\alpha^{\prime} \in \tmop{FV}
  (\tau)$, $\alpha^{\prime} \leqslant_{\mathcal{Q}} \alpha$.
\end{proposition}

\subsection{Abduction Algorithm for The Combination of
Domains}\label{SortCombProof}

\begin{lemma}
  \label{SolvedFormProj}For any conjunction of atoms $D \in
  \mathcal{L}_{s_{\tmop{type}}}$, let $D^t$ be $D$ with all alien subterms
  $\bar{r}$ replaced with fresh variables $\overline{\alpha_r}$, $D^t \in
  \mathcal{L}_{\tmop{ty}}$. Let $\wedge_s D_s^t =\tmmathbf{U} (D^t)$ be a
  solved form with $D_s^t \in \mathcal{L}_s$. Observe that for any $\Psi$, for
  any $C \in \mathcal{L}_s$, if $\mathcal{M} \vDash D \wedge \Psi \Rightarrow
  C$, then
  \begin{enumerate}
    \item $\mathcal{M} \vDash D_s^t [\overline{\alpha_r} \assign \bar{r}]
    \wedge \Psi \Rightarrow C$ for $s \neq s_{\tmop{type}}$.
    
    \item For $s = s_{\tmop{type}}$, let $\wedge_s C_s^t =\tmmathbf{U} (C^t)$,
    where $C^t$ is $C$ with all alien subterms $\bar{r}^{\prime}$ replaced
    with fresh variables $\overline{\alpha_r}^{\prime}$. Then there exist a
    conjunction of equations $E$ over $\overline{\alpha_r}
    \overline{\alpha_r}^{\prime}$ such that $T \vDash D_{s_{\tmop{type}}}^t
    \wedge E \wedge \Psi \Rightarrow C_{s_{\tmop{type}}}^t$ and (for $s \neq
    s_{\tmop{type}}$) $\mathcal{M} \vDash D_s^t [\overline{\alpha_r} \assign
    \bar{r}] \wedge \Psi \wedge \overline{\alpha_r^s} \dot{=} \overline{r^s }
    \Rightarrow E_s$ where $E_s$ are $E \cap \mathcal{L}_s$.
  \end{enumerate}
\end{lemma}

The proof uses the {\tmem{type conservation}} and {\tmem{free generation}}
($s_{\tmop{type}}$ properties) and interpolation techniques. Below follows a
sketch of a proof of (2).

\begin{proof}
  Observe a proof of $D^t [\overline{\alpha_r} \assign \bar{r}] \wedge \Psi
  \Rightarrow C^t [\overline{\alpha_r}^{\prime} \assign \bar{r}^{\prime}]$
  (assuming a complete proof system for $\mathcal{M} \vDash$). It can be
  transformed into a proof of $D^t \wedge \Psi \wedge \overline{\alpha_r}
  \dot{=} \bar{r} \wedge \overline{\alpha_r}^{\prime} \dot{=} \bar{r}^{\prime}
  \Rightarrow C^t$. By $s_{\tmop{type}}$ properties we get
  $D^t_{s_{\tmop{type}}} \wedge \Psi \wedge E_0 \wedge \overline{\alpha_r}
  \dot{=} \bar{r} \wedge \overline{\alpha_r}^{\prime} \dot{=} \bar{r}^{\prime}
  \Rightarrow C^t_{s_{\tmop{type}}}$ where $E_0$ are equations among
  $\overline{\alpha_r}$. By interpolation we get $D_{s_{\tmop{type}}}^t \wedge
  E \wedge \Psi \Rightarrow C_{s_{\tmop{type}}}^t$ and $D^t \wedge \Psi \wedge
  \overline{\alpha_r} \dot{=} \bar{r} \wedge \overline{\alpha_r}^{\prime}
  \dot{=} \bar{r}^{\prime} \Rightarrow E$ for $E_0 \subseteq E$ and $E$ as
  required by the theorem, since $\overline{\alpha_r}$ are the only
  non-$s_{\tmop{type}}$ subterms of $C^t_{s_{\tmop{type}}}$.
\end{proof}

Proof of Theorem \ref{MultisortAbdThm}.

\begin{proof}
  We use the same notation as in Table \ref{MultisortAbd}. Note, that the
  JCAQP with a branch with unsatisfiable conclusion does not have an answer, \
  so we do not consider unsatisfiable conclusions.
  
  Observe that validity of $\mathcal{L}_{\tmop{ty}}$ formula in $T$ is
  equivalent to its validity in $\mathcal{M}$. We will therefore sometimes
  drop prefixes $\mathcal{M} \vDash$ and $T \vDash$ for readability.
  
  Let us check correctness, i.e. that $\exists \overline{\alpha_{\tmop{ans}}}
  .A_{\tmop{ans}} \in \tmop{Abd} (\mathcal{Q}, \bar{\beta}, \overline{D_i,
  C_i})$ are JCAQP\tmrsub{$\mathcal{M}$} answers.
  \begin{enumerate}
    \item We need to show: $\mathcal{M} \vDash \wedge_i (D_i \wedge
    A_{\tmop{ans}} \Rightarrow C_i)$,
    
    \item $\mathcal{M} \vDash \mathcal{Q}.A_{\tmop{ans}}
    [\bar{\alpha}_{\tmop{ans}} \bar{\beta} \assign \bar{t}]$ for some
    $\bar{t}$,
    
    \item $\mathcal{M} \vDash \wedge_i \exists \tmop{FV} (D_i \wedge
    A_{\tmop{ans}}) .D_i \wedge A_{\tmop{ans}}$,
    
    \item for all atoms $c \in A_{\tmop{ans}}$ such that $\mathcal{M} \nvDash
    \mathcal{Q}.c$ and $\tmop{FV} (c) \cap \bar{\beta} \neq \varnothing$, then
    for all $\beta_1 \in \tmop{FV} (c)$ such that $(\forall \beta_1) \in
    \mathcal{Q}$, there exists $\beta_2 \in \tmop{FV} (c) \cap \bar{\beta}$
    such that $\beta_1 \leqslant_{\mathcal{Q}} \beta_2$.
    
    \item We have: $\mathcal{M} \vDash \wedge_i (D_{i, s_{\tmop{type}}}^t
    \wedge A_T^j \Rightarrow C_{i, s_{\tmop{type}}}^t)$,
    
    \item $\mathcal{M} \vDash \mathcal{Q}.A_T^j \left[ \bar{\alpha}_j^T
    \overline{\overline{\alpha_r}} \bar{\beta} \assign \bar{t}_j^T \right]$
    for some $\bar{t}_j^T$,
    
    \item $\mathcal{M} \vDash \wedge_i \exists \tmop{FV} (D_{i,
    s_{\tmop{type}}}^t \wedge A_T^j) .D_{i, s_{\tmop{type}}}^t \wedge A_T^j$,
    
    \item for all atoms $\beta_2 \dot{=} t \in A_T^j$ such that $\beta_2 \in
    \overline{\overline{\alpha_r}} \bar{\beta}$, if $\beta_1 \in \tmop{FV}
    (c)$ such that $(\forall \beta_1) \in \mathcal{Q}$, then $\beta_1
    \leqslant_{\mathcal{Q}} \beta_2$.
    
    \item We have: $\mathcal{M} \vDash \wedge_i (D_i^s \wedge (D^t_{i, s}
    \wedge A^i_{p, j, s}) [\overline{\overline{\alpha_r}} \assign
    \overline{\bar{r}}] \wedge A_s^{k_j^s} \Rightarrow C_i^s \wedge (C_{i,
    s}^t \wedge A^i_{c, j, s}) [\overline{\overline{\alpha_r}} \assign
    \overline{\bar{r}}])$,
    
    \item $\mathcal{M} \vDash \mathcal{Q}.A_s^{k_j^s} \left[
    \overline{\alpha_s^{k_j^s}} \bar{\alpha}_r^j \bar{\beta} \assign
    \bar{t}_{k_j^s} \right]$ for some $\bar{t}_{k_j^s}$,
    
    \item $\mathcal{M} \vDash \wedge_i \exists \tmop{FV} (D_i^s \wedge
    (D^t_{i, s} \wedge A^i_{p, j, s}) [\overline{\overline{\alpha_r}} \assign
    \overline{\bar{r}}] \wedge A_s^{k_j^s}) .D_i^s \wedge (D^t_{i, s} \wedge
    A^i_{p, j, s}) [\overline{\overline{\alpha_r}} \assign \overline{\bar{r}}]
    \wedge A_s^{k_j^s}$,
    
    \item for all atoms $c \in A_s^{k_j^s}$ such that $\mathcal{M} \nvDash
    \mathcal{Q}.c$ and $\tmop{FV} (c) \cap \bar{\beta} \neq \varnothing$, then
    for all $\beta_1 \in \tmop{FV} (c)$ such that $(\forall \beta_1) \in
    \mathcal{Q}$, there exists $\beta_2 \in \tmop{FV} (c) \cap \bar{\beta}$
    such that $\beta_1 \leqslant_{\mathcal{Q}} \beta_2$.
    
    \item We have: $D_i \equiv \wedge_s D_i^s$ and $C_i \equiv \wedge_s
    C_i^s$; $D_i^t [\overline{\overline{\alpha_r}} \assign \overline{\bar{r}}]
    = D^{s_{\tmop{type}}}_i$, $C_i^t [\overline{\overline{\alpha_r}} \assign
    \overline{\bar{r}}] = C^{s_{\tmop{type}}}_i$; $A^i_{p j} = \{ x \dot{=} t
    \in \tmmathbf{U} (D_{i, s_{\tmop{type}}}^t \wedge A_T^{j\prime}) |x \in
    X_s, s \neq s_{\tmop{type}} \}$ and $A^i_{c j} = \{ x \dot{=} t \in
    \tmmathbf{U} (D_{i, s_{\tmop{type}}}^t \wedge C_{i, s_{\tmop{type}}}^t
    \wedge A_T^{j\prime}) |x \in X_s, s \neq s_{\tmop{type}} \}$, $A^i_{p / c,
    j} = \wedge_s A_{p / c, j, s}^i$, where $A_{p / c, j, s}^i \in
    \mathcal{L}_s$.
    
    \item $D_i \wedge \overline{r_i} \dot{=} \overline{\alpha_{r_i}}
    \Rightarrow \wedge_s D_i^s \wedge D_{i, s}^t$, $\wedge_s C_i^s \wedge
    C_{i, s}^t \wedge \overline{r_i} \dot{=} \overline{\alpha_{r_i}}
    \Rightarrow C_i$, by (13).
    
    \item We have: $A_{\tmop{ans}} = A_T^{j\prime} \wedge_{s \in
    \tmop{usorts}} A_s^{k_j^s}$ for some $j, \overline{k_j^s}$ and
    $\bar{\alpha}_{\tmop{ans}} = \overline{\alpha_j^T}^{\prime}
    \overline{\overline{\alpha_s^{k_j^s}}} \cap \tmop{FV} (A_T^{j\prime}
    \wedge_s A_s^{k_j^s})$, $\overline{\overline{\alpha_s^{k_j^s}}}$ is a
    concatenation of $\overline{\alpha_s^{k_j^s}}$ for $s \in \tmop{usorts}$.
    
    \item By type preservation and free generation, (5) we can have $D_{i,
    s_{\tmop{type}}}^t \wedge A_T^{j\prime} \Rightarrow C_{i,
    s_{\tmop{type}}}^{t\prime}$, where $C_{i, s_{\tmop{type}}}^{t\prime}$
    differs from $C_{i, s_{\tmop{type}}}^t$ only at alien subterm positions.
    Pick $C_{i, s_{\tmop{type}}}^{t\prime}$ which differs from $C_{i,
    s_{\tmop{type}}}^t$ on the least number of alien subterm positions.
    
    \item Consider $S = \{ x \dot{=} t \in \tmmathbf{U} (C_{i,
    s_{\tmop{type}}}^t \wedge C_{i, s_{\tmop{type}}}^{t\prime}) |x \in X_s, s
    \neq s_{\tmop{type}} \}$. Note that $S \wedge C_{i,
    s_{\tmop{type}}}^{t\prime} \Rightarrow C^t_{i, s_{\tmop{type}}}$.
    
    \item By (16) and separation of sorts, we have $\tmmathbf{U} (D_{i,
    s_{\tmop{type}}}^t \wedge C_{i, s_{\tmop{type}}}^t \wedge A_T^{j\prime})
    \backslash\mathcal{L}_{s_{\tmop{type}}} \Rightarrow \tmmathbf{U} (C_{i,
    s_{\tmop{type}}}^t \wedge C_{i, s_{\tmop{type}}}^{t\prime})
    \backslash\mathcal{L}_{s_{\tmop{type}}}$, i.e. $A^i_{c j} \Rightarrow S$.
    
    \item $D_i \wedge A_{\tmop{ans}} \wedge \overline{\bar{r}} \dot{=}
    \overline{\overline{\alpha_r}} \overset{(13, 14)}{\Longrightarrow}
    D_i^{s_{\tmop{type}}} \wedge_s D_i^s \wedge D_{i, s}^t
    [\overline{\overline{\alpha_r}} \assign \overline{\bar{r}}] \wedge
    A_T^{j\prime} \wedge_{s \neq s_{\tmop{type}}} A_s^{k_j^s} \wedge
    \overline{\bar{r}} \dot{=} \overline{\overline{\alpha_r}}$
    
    \item $\ldots \overset{(9, 13)}{\Longrightarrow} D_{i, s_{\tmop{type}}}^t
    \wedge A_T^{j\prime} \wedge_{s \neq s_{\tmop{type}}} C_i^s \wedge C_{i,
    s}^t [\overline{\alpha_{r_i}} \assign \overline{r_i}] \wedge A_{c, j, s}^i
    [\overline{\overline{\alpha_r}} \assign \overline{\bar{r}}] \wedge
    \overline{\bar{r}} \dot{=} \overline{\overline{\alpha_r}}$
    
    \item $\ldots \overset{(17 - 18)}{\Longrightarrow} C_{i,
    s_{\tmop{type}}}^t \wedge_{s \neq s_{\tmop{type}}} C_i^s \wedge C_{i, s}^t
    [\overline{\alpha_{r_i}} \assign \overline{r_i}] \wedge \overline{\bar{r}}
    \dot{=} \overline{\overline{\alpha_r}} \overset{(13)}{\Longrightarrow}
    C_i$.
    
    \item (19)-(21) and interpolation give the goal (1).
    
    \item From type preservation and free generation properties (as in
    Proposition \ref{SolvedForm}), by (6) we get $\mathcal{M} \vDash
    \mathcal{Q}.A_T^j \left[ \bar{\alpha}_j^T (\bar{\beta} \cap
    X_{s_{\tmop{type}}}) \assign \bar{t}_j^T ; (\bar{\beta} \backslash
    X_{s_{\tmop{type}}}) \assign \bar{t}_j^b ; \overline{\overline{\alpha_r}}
    \assign \bar{t}_j^r \right]$ for some $\bar{t}_j^T$, for all $\bar{t}_j^r,
    \bar{t}_j^b$.
    
    \item Since $\mathcal{L}_s$ are single-sorted for $s \neq
    s_{\tmop{type}}$, by (15), (23) and (10) we get the goal (2).
    
    \item Similarly, from (3), (7), type preservation and free generation
    properties, and because $\mathcal{L}_s$ are single-sorted for $s \neq
    s_{\tmop{type}}$, we get the goal (3).
    
    \item Consider $c \in A_{\tmop{ans}}$ such that $\mathcal{M} \nvDash
    \mathcal{Q}.c$ and $\tmop{FV} (c) \cap \bar{\beta} \neq \varnothing$, and
    a $\beta_1 \in \tmop{FV} (c)$ such that $(\forall \beta_1) \in
    \mathcal{Q}$.
    
    \item If $\beta_1 \in \bar{\beta}$, then $\beta_2 = \beta_1$ satisfies the
    goal (4). Consider the cases where $\beta_1 \nin \bar{\beta}$.
    
    \item Consider $\beta_1 \in X_{s_{\tmop{type}}}$. By (26) $\mathcal{M}
    \nvDash \mathcal{Q}.c$.
    
    \item By (6) -- or (24) -- $\mathcal{M} \vDash \mathcal{Q}.c
    [\bar{\alpha}_j^T \bar{\beta} \assign \bar{t}_j^T]$ for some
    $\bar{t}_j^T$.
    
    \item By (15), $c$ is in solved form $\beta_3 \dot{=} t$.
    
    \item By (28)-(30), $\beta_3 \in \bar{\beta}$. By (8) we have the goal (4)
    for case $\beta_1 \in X_{s_{\tmop{type}}}$.
    
    \item Consider $\beta_1 \in X_s$ for $s \neq s_{\tmop{type}}$. Because
    $\mathcal{L}_s$ are single-sorted for $s \neq s_{\tmop{type}}$, by (15) we
    have $c \in A_s^{k_j^s}$.
    
    \item From (12) we have the goal (4).
  \end{enumerate}
  Now we sketch a proof of completeness, i.e. that no JCAQP answer ``falls
  through''.
  \begin{enumerate}
    \item We need to show that for any JCAQP $\mathcal{Q}. \wedge_i (D_i
    \Rightarrow C_i)$ with parameters $\bar{\beta}$ answer $\exists
    \bar{\alpha} .A$, there is an $\exists \overline{\alpha_{\tmop{ans}}}
    .A_{\tmop{ans}} \in \tmop{Abd} (\mathcal{Q}, \bar{\beta}, \overline{D_i,
    C_i})$ and $\bar{t}$ such that $A \Rightarrow A_{\tmop{ans}}
    [\overline{\alpha_{\tmop{ans}}} \assign \bar{t}]$.
    
    \item Let $A \equiv \wedge_s A_s$ where $A_s$ has atoms of sort $s$ only.
    Let $A^{\prime}$ be a formula obtained from $A_{s_{\tmop{type}}}$ by
    substituting all alien subterms $\bar{r}^{\prime}$ of sorts other than
    $s_{\tmop{type}}$ by fresh variables $\overline{\alpha_r}^{\prime}$. Let
    $\wedge_s A^{\prime}_s =\tmmathbf{U} (A^{\prime})$ be solved form of
    $A^{\prime}$ with $A^{\prime}_s \in \mathcal{L}_s$.
    
    \item W.l.o.g., all $C_i, D_i$ are satisfiable. Observe, that $\vDash D_i
    \wedge A \Rightarrow C_i$ implies, by lemma \ref{SolvedFormProj}, that
    there is a conjunction of equations over variables
    $\overline{\alpha_{r_i}} \overline{\alpha_r}^{\prime}$:
    $A^{\prime\prime}_i \assign \wedge_{s \neq s_{\tmop{type}}} A_{i,
    s}^{\prime\prime}$, such that:
    \begin{enumerate}
      \item  $D_i \wedge A \wedge \overline{\alpha_{r_i}}
      \overline{\alpha_r}^{\prime} \dot{=} \overline{r_i} \bar{r}^{\prime}
      \Rightarrow A^{\prime\prime}_i$,
      
      \item $D^t_{i, s_{\tmop{type}}} \wedge A^{\prime}_{s_{\tmop{type}}}
      \wedge A^{\prime\prime}_i \Rightarrow C^t_{i, s_{\tmop{type}}}$,
      
      \item $D_i^s \wedge D^t_{i, s} [\overline{\alpha_{r_i}} \assign
      \overline{r_i}] \wedge A_s \wedge A^{\prime}_s
      [\overline{\alpha_r}^{\prime} \assign \bar{r}^{\prime}] \Rightarrow
      C_i^s \wedge C^t_{i, s} [\overline{\alpha_{r_i}} \assign
      \overline{r_i}]$ for $s \neq s_{\tmop{type}}$.
    \end{enumerate}
    \item (3b) and the assumption about $\tmop{Abd}_T$ give that for some
    $(\exists \overline{\alpha_j^T} .A_T^j) \in \tmop{Abd}_T \left(
    \mathcal{Q}, \bar{\beta}, \overline{D_{i, s_{\tmop{type}}}^t, C_{i,
    s_{\tmop{type}}}^t} \right)$, $A^{\prime}_{s_{\tmop{type}}} \wedge
    A^{\prime\prime}_i \Rightarrow A_T^j [\overline{\alpha_j^T} \assign
    \overline{t_T}]$ for some $\overline{t_T}$.
    
    \item (3a), (3c) and lemma \ref{SolvedFormProj} imply, by substituting
    free variables, $D_i^s \wedge D^t_{i, s} [\overline{\alpha_{r_i}} \assign
    \overline{r_i}] \wedge A_s \wedge A^{\prime}_s
    [\overline{\alpha_r}^{\prime} \assign \bar{r}^{\prime}] \Rightarrow C_i^s
    \wedge C^t_{i, s} [\overline{\alpha_{r_i}} \assign \overline{r_i}] \wedge
    A^{\prime\prime}_{i, s} [\overline{\overline{\alpha_r}}
    \overline{\alpha_r}^{\prime} \assign \overline{\bar{r}}
    \bar{r}^{\prime}]$.
    
    \item To use the assumption about $\tmop{Abd}_s$, we weaken (5). With a
    stronger premise, we also get a stronger conclusion: $D_i^s \wedge
    (D^t_{i, s} \wedge A^i_{p, j, s} \wedge A^i_{c, j, s})
    [\overline{\overline{\alpha_r}} \assign \overline{\bar{r}}] \wedge A_s
    \wedge A^{\prime}_s [\overline{\alpha_r}^{\prime} \assign
    \bar{r}^{\prime}] \Rightarrow C_i^s \wedge (C^t_{i, s} \wedge A^i_{c, j,
    s}) [\overline{\alpha_{r_i}} \assign \overline{r_i}] \wedge
    A^{\prime\prime}_{i, s} [\overline{\overline{\alpha_r}}
    \overline{\alpha_r}^{\prime} \assign \overline{\bar{r}}
    \bar{r}^{\prime}]$.
    
    \item (6) gives by the assumption about $\tmop{Abd}_s$: $(A^i_{c, j, s}
    \backslash A^i_{p, j, s}) [\overline{\overline{\alpha_r}} \assign
    \overline{\bar{r}}] \wedge A_s \wedge A^{\prime}_s
    [\overline{\alpha_r}^{\prime} \assign \bar{r}^{\prime}] \Rightarrow
    A_s^{k_j^s} \left[ \overline{\alpha_s^{k_j^s}} \assign
    \overline{t_s^{k_j^s}} \right]$ for some $k_j^s$ and
    $\overline{t_s^{k_j^s}}$.
    
    \item Check using (3b) and (4) that $A^{\prime}_{s_{\tmop{type}}} \wedge
    A^{\prime\prime}_i \wedge \overline{\alpha_{r_i}}
    \overline{\alpha_r}^{\prime} \dot{=} \overline{r_i} \bar{r}^{\prime}
    \Rightarrow \exists \bar{\alpha}_r^j .A_T^{j\prime} [\overline{\alpha_j^T}
    \assign \overline{t_T}] \wedge (A^i_{c, j, s} \backslash A^i_{p, j, s})
    [\overline{\overline{\alpha_r}} \assign \overline{\bar{r}}]$. The
    subtraction in $(A^i_{c, j, s} \backslash A^i_{p, j, s})$ allows us to
    drop $D^t_{i, s_{\tmop{type}}}$ from the premise in (3b).
    
    \item Select $\exists \overline{\alpha_{\tmop{ans}}} .A_{\tmop{ans}}
    \assign \exists \overline{\alpha_{\overline{k_j^s}}} .A_T^{j\prime}
    [\overline{\overline{\alpha_r}} \assign \overline{\bar{r}}] \wedge_s
    A_s^{k_j^s}$ for $\overline{\alpha_{\overline{k_j^s}}} \assign
    \overline{\alpha_j^T} \bar{\alpha}_r^j
    \overline{\overline{\alpha_s^{k_j^s}}} \cap \tmop{FV} (A_T^j
    [\overline{\alpha_r} \assign \bar{r}] \wedge_s A_s^{k_j^s})$. By (7) and
    (8), we have $A \wedge [\overline{\overline{\alpha_r}}
    \overline{\alpha_r}^{\prime} \dot{=} \overline{\bar{r}} \bar{r}^{\prime}]
    \Rightarrow A_{\tmop{ans}} \left[ \overline{\alpha_j^T} \bar{\alpha}_r^j
    \overline{\overline{\alpha_s^{k_j^s}}} \assign \overline{t_T} \bar{t}_r^j
    \overline{\overline{t_s^{k_j^s}}} \right]$ for some $\bar{t}_r^j$, which
    implies (1).
  \end{enumerate}
\end{proof}

\section{Constraint Generalization}\label{LUBProofs}

\begin{proposition}
  \label{SepProp}Let $D_s \in \mathcal{L}_s$ for all sorts $s$ such that
  $D_{s_{\tmop{type}}} =\tmmathbf{U} (D_{s_{\tmop{type}}})$, and
  $C_{s^{\prime}} \in \mathcal{L}_{s^{\prime}}$ for $s^{\prime} \neq
  s_{\tmop{type}}$. Then $\mathcal{M} \vDash (\wedge_s D_s) \Rightarrow
  C_{s^{\prime}}$ if and only if $\mathcal{M} \vDash D_{s^{\prime}}
  \Rightarrow C_{s^{\prime}}$.
\end{proposition}

\begin{proposition}
  {\tmem{(MGU property.)}} MGU: $T (F, X) \vDash C \Leftrightarrow
  \tmmathbf{U} (C)$.
\end{proposition}

\begin{proposition}
  \label{ImplSubst}Let $D$ be a conjunction of equations with $D =\tmmathbf{U}
  (D)$, let $A$ be a conjunction of equations and $\bar{\alpha}$ variables
  such that $\tmop{FV} (A) \subset \bar{\alpha} \cup \tmop{FV} (D)$. $T (F, X)
  \vDash D \Rightarrow \exists \bar{\alpha} .A$ if and only if there exists a
  substitution $S = [\bar{\alpha} \assign \overline{g_{\alpha}}]$ and a solved
  form $A^{\prime} \Leftrightarrow A$ with $A^{\prime} =\tmmathbf{U}
  (A^{\prime})$ such that for every $x \dot{=} t \in A^{\prime}$, either $S
  (x) = S (t)$, or $x \dot{=} S (t) \in D$.
\end{proposition}

Proof of Theorem \ref{CorComplLUBThm}.

\begin{proof}
  First, we show $\mathcal{M} \vDash \wedge_i (D_i \Rightarrow \exists
  \bar{\alpha} .A)$.
  \begin{enumerate}
    \item $\mathcal{M} \vDash D_i \wedge D_i^a \wedge D_i^u \Rightarrow c$,
    for each conjunct $c \in A_{s_{\tmop{type}}}$.
    \begin{enumerate}
      \item Case $c = x_j \dot{=} g_j$. By properties of MGU.
      
      \item Case $c \in D_i^g$ follows from $D_i \wedge D_i^a \Rightarrow
      D_{i, s_{\tmop{type}}}^t$.
      
      \item Case $c \in D_i^v$. The premises are: $x \dot{=} t_1 \in D_i^g, y
      \dot{=} t_2 \in D_i^g, \mathcal{M} \vDash D_i^g \Rightarrow t_1 \dot{=}
      t_2$. The goal follows from $D_i \wedge D_i^a \wedge D_i^u \Rightarrow
      D_i^g$, by transitivity. 
    \end{enumerate}
    \item By properties of $\tmop{LUB}_s$, we have $\vDash D_i^s \wedge
    D_i^{t, s} \wedge D_{i, s} \Rightarrow \exists \bar{\alpha}_s .A_s$ for $s
    \neq s_{\tmop{type}}$, and therefore $\vDash D_i \wedge D_i^a \wedge D_i^u
    \Rightarrow \exists \bar{\alpha}_s .A_s$.
    
    \item We have $\vDash D_i \Rightarrow \exists \overline{\alpha^j_i} .D_i
    \wedge D_i^a$, and by properties of most specific anti-unification,
    $\vDash D_i \Rightarrow \exists \overline{\bar{\alpha}_j} .D_i \wedge
    D_i^u$.
    
    \item Collecting the above points, we get $\vDash D_i \Rightarrow \exists
    \overline{\alpha^j_i} \overline{\bar{\alpha}_s} . \wedge_s A_s$.
  \end{enumerate}
  Now, we sketch a proof of: For every $\exists \bar{\alpha}_r .A_r$ such that
  $\mathcal{M} \vDash \wedge_i (D_i \Rightarrow \exists \bar{\alpha}_r .A_r)$,
  with variables renamed so that $\bar{\alpha} \# \tmop{FV} (A_r)$,
  $\mathcal{M} \vDash A \Rightarrow \exists \bar{\alpha}_r .A_r$.
  \begin{enumerate}
    \item The assumption is $\mathcal{M} \vDash \wedge_i (D_i \Rightarrow
    \exists \bar{\alpha}_r .A_r)$.
    
    \item By definition, $\mathcal{M} \vDash D_i \wedge D_i^a \Leftrightarrow
    \wedge_{s \neq s_{\tmop{type}}} D_i^s \wedge_s D_i^{t, s} \wedge_s D^t_{i,
    s}$.
    
    \item Let $\exists \bar{\alpha}_r .A_r \Leftrightarrow \exists
    \overline{\bar{\alpha}_{t, s}^r} . \wedge_s \exists \bar{\alpha}_{r, s}
    .D_{r, s}^a \wedge A_{r, s}$ where $\exists \bar{\alpha}_{r, s} .A_{r, s}
    \in \mathcal{L}_s$, $\bar{a}_{t, s}^r$ are all alien subterms of sort $s$
    in $A_r$, $\bar{\alpha}_{t, s}^r$ are fresh variables $\bar{\alpha}_{t,
    s}^r \# \tmop{FV} (A_r, \overline{D_i})$ and also all alien subterms in
    $A_{r, s_{\tmop{ty}}}$, $A_{r, s_{\tmop{type}}} =\tmmathbf{U} (A_{r,
    s_{\tmop{type}}})$, and $D_{r, s}^a = \bar{\alpha}_{t, s}^r \dot{=}
    \bar{a}_{t, s}^r$ for $s \neq s_{\tmop{type}}$, $D_{r, s_{\tmop{type}}}^a
    =\tmmathbf{T}$.
    
    \item By Proposition \ref{ImplSubst}, for each branch $i$ there is a
    substitution $\eta_i$ of $\overline{\bar{\alpha}_{t, s}^r}
    \bar{\alpha}_{r, s_{\tmop{type}}}$ that is an anti-unifier of
    $\overline{t_{i, j}^{\prime}}$, for all $x_j \nin
    \overline{\bar{\alpha}_{t, s}^r} \bar{\alpha}_{r, s_{\tmop{type}}}$ such
    that $x_j \dot{=} u_{r, j} \in A_{r, s_{\tmop{type}}}$, where $x_j,
    \overline{t_{i, j}} \in V$ and $\mathcal{M} \vDash D_i \Rightarrow
    t^{\prime}_{i, j} \dot{=} t_{i, j}$. $t^{\prime}_{i, j}$ can be made equal
    to $t_{i, j}$ up to a variable-variable substitution.
    \begin{itemize}
      \item I.e. a substitution $[\bar{\alpha} \assign \bar{\beta}]$ that can
      assign the same $\beta$ to several $\alpha$.
    \end{itemize}
    \item There exists a substitution $\rho$ which establishes $\mathcal{M}
    \vDash A_{s_{\tmop{type}}} \Rightarrow \exists \overline{\bar{\alpha}_{t,
    s}^r} . \exists \bar{\alpha}_{r, s_{\tmop{type}}} . \overline{x_j \dot{=}
    u_{r, j}}$: by variable-variable substitution, including alien subterm
    variables, according to $\mathcal{M} \vDash D_i \Rightarrow t^{\prime}_{i,
    j} \dot{=} t_{i, j}$, and by the factoring via most specific
    anti-unification.
    
    \item Let $\exists \bar{\alpha}_s .A_s = \tmop{LUB}_s \left(
    \overline{D_i^s \wedge \widetilde{D^a_i} (D^t_{i, s} \wedge D^u_{i, s})}
    \right)$ as in the definition of $\tmop{LUB} (\cdot)$.
    
    \item $\mathcal{M} \vDash D_i \Rightarrow \exists \bar{\alpha}_r .A_r$ by
    (1), $\mathcal{M} \vDash D_i \wedge D_i^a \wedge D^u_{i, s} \Rightarrow
    \exists \bar{\alpha}_r .A_r$ by weakening on the left, $\mathcal{M} \vDash
    D_i \wedge D_i^a \wedge D^u_{i, s} \Rightarrow \exists
    \overline{\bar{\alpha}_{t, s}^r} . \wedge_s \exists \bar{\alpha}_{r, s}
    .D_{r, s}^a \wedge A_{r, s}$ by (3), $\mathcal{M} \vDash D_i \wedge D_i^a
    \wedge D^u_{i, s} \Rightarrow \exists \bar{\alpha}_{r, s} .A_{r, s}$ by
    weakening on the right, $\mathcal{M} \vDash D_i^s \wedge D_i^{t, s} \wedge
    D^t_{i, s} \wedge D^u_{i, s} \Rightarrow \exists \bar{\alpha}_{r, s}
    .A_{r, s}$ by (2) and Proposition \ref{SepProp}.
    
    \item By (6), (7), interpolation and assumption about $\tmop{LUB}_s$,
    $\mathcal{M} \vDash A_s \Rightarrow \exists \bar{\alpha}_{r, s} .A_{r, s}$
    for every $s \neq s_{\tmop{type}}$.
    
    \item More specifically, we need $\mathcal{M} \vDash A_s \wedge
    A_{s_{\tmop{type}}} \Rightarrow \exists \bar{\alpha}_{r, s} .A_{r, s}
    \wedge \rho (\overline{\alpha^r_{t, s}}) \dot{=} \overline{a^r_{t, s}}$.
    It can be shown by extending the argument for (8) using the fact that
    $D^u_{i, s}$ relates $\rho (\overline{\alpha^r_{t, s}})$ introduced in (5)
    to the remaining constraints of branch $i$.
  \end{enumerate}
\end{proof}

\section{Solving for Predicate Variables}\label{PredVarProofs}

In this section, we provide proof sketches for correctness and a limited form
of completeness of the type inference implemented in {\tmname{InvarGenT}},
correctness with respect to the type system $\tmop{MMG}_{\exists} (X)$ and the
limited form of completeness wrt. $\tmop{MMG} (X)$.

\begin{lemma}
  \label{SplitProps}Let $(\mathcal{Q}^{\prime}, \bar{\alpha}_{\tmop{res}},
  A_{\tmop{res}}, \overline{\exists \bar{\alpha}^{\chi} .A_{\chi}}) \in
  \tmop{Split} (\mathcal{Q}, \bar{\alpha}, A, \overline{\bar{\beta}^{\chi}})$.
  Then:
  \begin{enumerate}
    \item $\mathcal{M} \vDash A_{\tmop{res}} \wedge_{\chi} A_{\chi}
    \Rightarrow A$.
    
    \item $\mathcal{M} \vDash A \Rightarrow \exists
    \overline{\bar{\alpha}^{\chi}_+} .A_{\tmop{res}} \wedge_{\chi} A_{\chi}$.
    
    \item $\mathcal{M} \vDash \mathcal{Q}^{\prime} .A_{\tmop{res}}
    [\bar{\alpha}_{\tmop{res}} \assign \bar{t}] \text{ \ for some } \bar{t}$.
    
    \item If $A$ only restricts $\bar{\beta}^{\chi}$ in ways that are
    expressible as $\mathcal{Q}_{< \beta_{\chi}} \exists \bar{\beta}^{\chi}
    \bar{\beta}^{\chi}_1 . \varphi$ for some $\bar{\beta}^{\chi}_1$, and
    $\mathcal{M} \vDash \mathcal{Q}.A [\bar{\alpha}
    \overline{\bar{\beta}^{\chi}} \assign \bar{t}]$ for some $\bar{t}$, then
    $\tmop{Split} (\mathcal{Q}, \bar{\alpha}, A,
    \overline{\bar{\beta}^{\chi}}) \neq \varnothing$.
    
    \item $\tmop{Split}$ preserves atomized form: if $\exists \bar{\alpha} .A$
    is in atomized form with respect to $\mathcal{Q}$ and parameters
    $\overline{\bar{\beta}^{\chi}}$, then $\exists \bar{\alpha}_{\tmop{res}}
    .A_{\tmop{res}}$ is in atomized form with respect to
    $\mathcal{Q}^{\prime}$ and parameters $\overline{\bar{\alpha}^{\chi}
    \bar{\beta}^{\chi}}$.
  \end{enumerate}
\end{lemma}

\begin{proof}
  Recall the definition of $\tmop{Split}$ and the notation used there.
  
  (1) and (2) follow from the observation that $A_{\chi}^L \wedge A^R_{\chi}
  \Rightarrow A_{\chi}^+$ and $A_{\chi}^+ \Rightarrow \exists
  \overline{\bar{\alpha}^{\chi}_+} .A_{\chi}^L \wedge A^R_{\chi}$.
  
  (3) follows from $\mathcal{M} \vDash \mathcal{Q}. (A \setminus \cup_{\chi}
  A_{\chi}^+) [\bar{\alpha}_{\tmop{res}} \assign \bar{t}]$ for some $\bar{t}$
  because $\cup_{\chi} \bar{\alpha}^{\chi} = \varnothing$ and $A_{\tmop{res}}
  = A \setminus \cup_{\chi} A_{\chi}$, in the last iteration.
  
  For (4) to hold the algorithm should at each recursion level be able to find
  $\overline{A_{\chi}^+}$ for which $\mathcal{M} \vDash \mathcal{Q}. \left( A
  \setminus {\nobreak} \cup_{\chi} A_{\chi}^+ \right) [\bar{\alpha} \assign
  \bar{t}] \text{ for some } \bar{t}$ and $\tmop{Strat} (A^+_{\chi},
  \bar{\beta}^{\chi})$ does not return $\bot$ for any $\chi$, where
  $\mathcal{Q}$ and $\bar{\beta}^{\chi}$ are either the initial arguments or
  the recursive call arguments. The algorithm terminates because
  $\bar{\alpha}$ always decreases in the recursive call.
  
  $A_{\chi}^+$ contains restrictions on the variables $\bar{\beta}^{\chi}$. If
  $\tmop{Strat} (A^+_{\chi}, \bar{\beta}^{\chi})$ returns $\bot$, then,
  because $A_{\chi}^+$ is in atomized form, $A_{\chi}^+$ restricts
  $\bar{\beta}^{\chi}$ in a way not expressible as $\mathcal{Q}_{<
  \beta_{\chi}} \exists \bar{\beta}^{\chi} \bar{\beta}^{\chi}_1 . \varphi$.
  
  The final remark follows from the minimality of $\overline{A_{\chi}^+}$
  w.r.t. inclusion.
\end{proof}

\begin{lemma}
  \label{OpCorrectLemma}Let $\Phi \in \mathcal{L}$, $S = \overline{\exists
  \bar{\alpha}^{\chi} .F_{\chi}}$ and $R = \overline{\exists
  \bar{\alpha}^{\chi_K} .F_{\chi_K}}$. Let $\tmop{NF} (RS (\Phi))
  =\mathcal{Q}. \wedge_i (D_i \Rightarrow C_i) =\mathcal{Q}. \Phi_{\tmop{PN}}
  \in \mathcal{L}^{}$.
  
  Let $(\exists \bar{\alpha}_{\tmop{res}}^{\prime} .F_{\tmop{res}}^{\prime},
  S^{\prime}, R^{\prime}) \in \Psi (k, \Phi, S, R)$ for any $k$. Let
  \[ \mathcal{Q}^{\prime} . \Phi_{\tmop{PN}}^{\prime} = \tmop{NF} (R^{\prime
     -} S^{\prime -} R^+ S^+ (\Phi)) \]
  Then $\mathcal{M} \vDash F_{\tmop{res}}^{\prime} \Rightarrow
  \Phi_{\tmop{PN}}^{^{\prime}}$ and $\mathcal{M} \vDash \mathcal{Q}^{\prime}
  .F^{\prime}_{\tmop{res}} [\bar{\alpha}_{\tmop{res}} \assign \bar{t}]$ for
  some $\bar{t}$.
\end{lemma}

\begin{proof}
  We use the notation from the definition of $\Psi$, with $S_k = S, R_k = R,
  \mathcal{Q}^{\prime} =\mathcal{Q}^{k + 1}$. We have
  \[ \exists \bar{\alpha} .A_0 \in \tmop{Abd} \left( \mathcal{Q}^{k\prime},
     \overline{\beta_{\chi} \overline{\beta^{}}^{\chi}},
     \overline{D_i^{k\prime}, C_i^{k\prime}} \right) \]
  and therefore {\tmem{(1)}} $\mathcal{M} \vDash \wedge_i (D_i^{k\prime}
  \wedge A_0 \Rightarrow C_i^{k\prime})$. Continuing the definition of $\Psi$,
  we have
  \[ \left( \mathcal{Q}^{k + 1}, \bar{\alpha}_{\tmop{res}}, A_{\tmop{res}},
     \overline{\exists \bar{\alpha}^{\beta_{\chi}} .A_{\beta_{\chi}}} \right)
     \in \tmop{Split} \left( \mathcal{Q}^{k\prime}, \bar{\alpha}, A,
     \overline{\beta_{\chi} \overline{\beta^{}}^{\chi}} \right) \]
  and $S^{\prime} (\chi) = \exists \bar{\beta}^{\chi, k} . \tmop{Simpl} \left(
  \exists \overline{\bar{\alpha}^{\beta_{\chi}}} .F_{\chi}^{\prime} \wedge
  A_{\beta_{\chi}} \left[ \overline{\beta_{\chi} \overline{\beta^{}}^{\chi}}
  \assign \overline{\delta \overline{\beta^{}}^{\chi, k}} \right] \right)$,
  $F_{\tmop{res}}^{\prime} = A_{\tmop{res}}$. Therefore by Lemma
  \ref{SplitProps} point 1, {\tmem{(2)}} $\mathcal{M} \vDash
  F_{\tmop{res}}^{\prime} \wedge_{\beta_{\chi}} A_{\beta_{\chi}} \Rightarrow
  {\nobreak} A$, and by Lemma \ref{SplitProps} point 3, the goal $\mathcal{M}
  \vDash \mathcal{Q}^{\prime} .F_{\tmop{res}}^{\prime}
  [\bar{\alpha}_{\tmop{res}} \assign \bar{t}] \text{ \ for some } \bar{t}$.
  
  It remains to show {\tmem{(3)}} $\mathcal{M} \vDash F_{\tmop{res}}^{\prime}
  \Rightarrow {\nobreak} \Phi_{\tmop{PN}}^{^{\prime}}$. Observe that
  {\tmem{(4)}} $\Phi^{\prime}_{\tmop{PN}} \equiv (\wedge_{\beta_{\chi}}
  A_{\beta_{\chi}} \Rightarrow \Phi_{\tmop{PN}} \wedge_{\chi_K} U_{\chi_K})$.
  (1) is $\mathcal{M} \vDash A_0 \Rightarrow \Phi_{\tmop{PN}} \wedge_{\chi_K}
  U_{\chi_K}$. But $F^{\prime}_{\tmop{res}} \wedge_{\chi} A_{\beta_{\chi}}
  \Rightarrow A_0$ by (2), so (1) and (4) give (3).
\end{proof}

Sketch of proof of Theorem \ref{ThCorrect}.

\begin{proof}
  Lemma \ref{OpCorrectLemma} gives the goal $\mathcal{M}, \mathcal{I} \vDash
  \Phi$ with \ $\mathcal{I} (\Phi) \equiv \mathcal{Q}^{\prime} .
  \Phi_{\tmop{PN}}^{\prime}$, because $S^{\prime} R^{\prime} = SR$ implies
  $\wedge_{\beta_{\chi}} A_{\beta_{\chi}} =\tmmathbf{T}$.
\end{proof}

\begin{axiom}
  (Interpolation property for $\mathcal{M}$.) We assume that for conjunctions
  of atoms $A$ and any quantifier-free formula $\Phi$, $\mathcal{M} \vDash A
  \Rightarrow \Phi$ implies that there is a conjunction of atoms $B$ with
  $\tmop{FV} (B) \subset \tmop{FV} (A) \cap \tmop{FV} (\Phi)$, $\mathcal{M}
  \vDash A \Rightarrow B$ and $\mathcal{M} \vDash B \Rightarrow \Phi$.
\end{axiom}

\begin{lemma}
  \label{OpComplLemma}Let $\tmop{NF} (\Phi [\overline{\chi (\tau)} \assign
  \overline{\exists \bar{\alpha}^{\chi} .F_{\chi} [\delta \assign \tau]}])
  =\mathcal{Q}. \Phi_N^{}$ and $\tmop{PV} (\Phi) = \tmop{PV}^1 (\Phi)$. Let
  \[ \mathcal{Q}^{\Delta} . \Phi_N^{\Delta} = \tmop{NF} (\Phi
     [\overline{\chi^- (\beta^{\chi})} \assign \overline{\exists
     \bar{\alpha}^{\chi} \bar{\alpha}^{\chi}_{\Delta} . (\Delta_{\chi} \wedge
     F_{\chi}) [\delta \assign \beta^{\chi}]} ; \overline{\chi^+ (\tau)}
     \assign \overline{\exists \bar{\alpha}^{\chi} .F_{\chi} [\delta \assign
     \tau]}]) \]
  for variables $\overline{\bar{\alpha}^{\chi}_{\Delta}}$ and conjunctions of
  atoms $\overline{\Delta_{\chi}} \in \mathcal{L}$ such that the
  JCAQP$_{\mathcal{M}}$ problem $\mathcal{Q}^{\Delta} . \Phi_N^{\Delta}$ has a
  solution. Assume that $\tmop{Abd}$ in the definition of $\Psi$ is a complete
  abduction algorithm returning an atomized form formula. Then there is
  $(F_{\tmop{res}}^{\prime}, {\nobreak} \overline{\exists \bar{\alpha}^{\chi
  \prime} .F_{\chi}^{\prime}}) \in \Psi (k, \Phi, \overline{\exists
  \bar{\alpha}^{\chi} .F_{\chi}})$ such that
  \[ \mathcal{M} \vDash \wedge_{\chi} \Delta_{\chi}^{\beta} \Rightarrow
     (\wedge_{\chi} (F^{\prime}_{\chi} \setminus F_{\chi}))
     [\overline{\bar{\alpha}^{\chi \prime}} \assign \bar{t}] \]
  for all $\chi$, for some $\bar{t}$.
\end{lemma}

\begin{proof}
  Let $\exists \bar{\alpha}^d_{\tmop{res}} .D_{\tmop{res}}$ be a solution to
  the JCAQP$_{\mathcal{M}}$ problem $\mathcal{Q}^{\Delta} . \Phi_N^{\Delta}$.
  We have $\mathcal{M} \vDash D_{\tmop{res}} \wedge_{\chi}
  \Delta_{\chi}^{\beta} \Rightarrow \Phi_N$. By the completeness assumption
  about $\tmop{Abd}$, we have $\mathcal{M} \vDash D_{\tmop{res}} \wedge_{\chi}
  \Delta_{\chi}^{\beta} \Rightarrow A [\bar{\alpha}
  \overline{\bar{\alpha}^{\chi}} \assign \bar{t}]$ for some $\bar{t}$. Since
  $D_{\tmop{res}}$ does not participate in restricting
  $\overline{\bar{\beta}^{\chi}}$ and $\Delta_{\chi}^{\beta}$ is limited to
  $\mathcal{Q}_{< \beta_{\chi}} \bar{\alpha}^{\chi}
  \bar{\alpha}^{\chi}_{\Delta}$ variables, by atomized form of $\exists
  \bar{\alpha} .A$ and interpolation, we get the expressiveness constraint
  needed for Lemma \ref{SplitProps} point 4, therefore $\tmop{Split}
  (\mathcal{Q}, \bar{\alpha}, A, \overline{\bar{\beta}^{\chi}}) \neq
  \varnothing$. Thus by Lemma \ref{SplitProps} point 2, $\mathcal{M} \vDash
  D_{\tmop{res}} \wedge_{\chi} \Delta_{\chi}^{\beta} \wedge
  \overline{\bar{\alpha}^{\chi \prime}} \dot{=} \bar{t} \Rightarrow \exists
  \overline{\bar{\alpha}^{\chi}_+} .A_{\tmop{res}} \wedge_{\chi} A_{\chi}$
  where by definition of $\Psi$, $A_{\chi} = F^{\prime}_{\chi} \setminus
  F_{\chi}$. Since the atomized form is preserved by $\tmop{Split}$ (Lemma
  \ref{SplitProps} point 5), we have the goal.
\end{proof}

Sketch of proof of Theorem \ref{InferCompleteness}.

\begin{proof}
  The thesis $\mathcal{M} \vDash \wedge_{\chi} F^s_{\chi} \Rightarrow
  (\wedge_{\chi} F^k_{\chi}) [\overline{\bar{\alpha}^{\chi, k}} \assign
  \bar{t}^k]$ is true for $k = 0$. Assume it holds for arbitrary $k$. We need
  to show $\mathcal{M} \vDash \wedge_{\chi} F^s_{\chi} [\delta \assign
  \beta_{\chi}] \Rightarrow (\wedge_{\chi} F^{k + 1}_{\chi} [\delta \assign
  \beta_{\chi}]) [\overline{\bar{\alpha}^{\chi, k + 1}} \assign \bar{t}^{k +
  1}]$ for some $(F_{\tmop{res}}^{k + 1}, \overline{\exists
  \bar{\alpha}^{\chi, k + 1} .F_{\chi}^{k + 1}})$ and a substitution of
  variables $\overline{\bar{\alpha}^{\chi, k + 1}}$.
  
  By the assumption that $\bar{\chi} = \tmop{PV} (\Phi) = \tmop{PV}^1 (\Phi)$,
  definition of $\Psi$ gives $F_{\chi}^{k + 1} = F_{\chi}^k \wedge A_{\chi}$.
  We will apply Lemma \ref{OpComplLemma} with $\Delta_{\chi} = F_{\chi}^s$ and
  $F_{\chi} = F^k_{\chi}$. By the inductive assumption, JCAQP$_{\mathcal{M}}$
  answer $\exists \bar{\alpha}^{\tmop{res}}_s .F_{\tmop{res}}^s$ to $\tmop{NF}
  (\Phi [\overline{\chi (\tau)} \assign \overline{\exists
  \bar{\alpha}^{\chi}_s .F_{\chi}^s [\delta \assign \tau]}])$ is also an
  answer to $\mathcal{Q}^{\Delta} . \Phi_N^{\Delta}$. By Lemma
  \ref{OpComplLemma}, we get the goal.
\end{proof}

\chapter{Algorithmic Details}

\section{Generating and Normalizing Formulas}\label{NormConstrn}

We inject the existential type and value constructors during parsing for
user-provided existential types, and during constraint generation for inferred
existential types, into the list of toplevel items. It facilitates exporting
inference results as OCaml source code.

Toplevel definitions are intended as boundaries for constraint solving. This
way the programmer can decompose functions that could be too complex for the
solver. A toplevel \tmverbatim{let rec} only binds a single identifier, while
\tmverbatim{let} binds variables in a pattern. To preserve the flexibility of
expression-level pattern matching, for \tmverbatim{let} -- unless it just
binds a value as discussed above -- we pack the constraints $\llbracket \Sigma
\vdash p \uparrow \alpha \rrbracket$ which the pattern makes available, into
existential types. Each pattern variable is a separate entry to the global
environment, therefore the connection between them is lost.

The \tmverbatim{let}$\ldots$\tmverbatim{in} syntax has two uses: binding
values of existential types means ``eliminating the quantification'' -- the
programmer has control over the scope of the existential constraint. The
second use is if the value is not of existential type, the constraint is
replaced by one that would be generated for a pattern matching branch. This
recovers the common use of the \tmverbatim{let}...\tmverbatim{in} syntax, with
exception of polymorphic \tmverbatim{let} cases, where \tmverbatim{let rec}
needs to be used.

We optimize the disjunctive constraint coming from \tmverbatim{let} bindings
as follows. Rather than inspecting $K$ in $\Sigma$ directly, let $\llbracket
\Sigma \vdash K x \uparrow \alpha_0 \rrbracket = \exists \bar{\beta} [D] \{ x
\mapsto \tau^{\prime} \}$. $\left\llbracket \Gamma \vdash K p.e_2 : {\nobreak}
\alpha_0 \rightarrow {\nobreak} \tau \right\rrbracket$ is equivalent to, or at
least implied by:
\[ \llbracket \Sigma \vdash K p \downarrow \alpha_0 \rrbracket \wedge \forall
   \bar{\beta} .D \Rightarrow \llbracket \Gamma \vdash p.e_2 : \tau^{\prime}
   \rightarrow \tau \rrbracket \]
Let us set $\tmmathbf{C}_0 = \not{E} (\alpha_0)$, $\tmmathbf{C}_K = \llbracket
\Sigma \vdash K p \downarrow \alpha_0 \rrbracket$, $\tmmathbf{D}_0 = \beta_0
\dot{=} \alpha_0$ and $\tmmathbf{D}_K = \beta_0 \dot{=} \tau^{\prime} \wedge
D$. Then,
\[ \exists \alpha_0 . \llbracket \Gamma \vdash e_1 : \alpha_0 \rrbracket
   \wedge \left( \llbracket \Gamma \vdash p.e_2 : \alpha_0 \rightarrow \tau
   \rrbracket \wedge \not{E} (\alpha_0) \vee_{\mathcal{E}} \llbracket \Gamma
   \vdash K p.e_2 : \alpha_0 \rightarrow \tau \rrbracket \right) \]
is equivalent to:
\[ \exists \alpha_0 . \llbracket \Gamma \vdash e_1 : \alpha_0 \rrbracket
   \wedge \vee_{i \in \{ 0 \} \cup \mathcal{E}} (\tmmathbf{C}_i \wedge \forall
   \bar{\beta} \beta_0 . (\tmmathbf{D}_i \Rightarrow \llbracket \Gamma \vdash
   p.e_2 : \beta_0 \rightarrow \tau \rrbracket)) \]
We use the same variable $\beta_0$ across disjuncts of the same disjunction,
so that $\llbracket \Gamma \vdash p.e_2 : \beta_0 \rightarrow \tau \rrbracket$
can be derived and simplified only once.

The second argument of the predicate variable $\chi_K (\gamma, \alpha)
\nosymbol$ provides an ``escape route'' for free variables, i.e. precondition
variables used in postcondition. In the implementation, we have user-defined
existential types with explicit constraints in addition to inferred
existential types. We expand the inferred existential types after they are
solved into the fuller format. In the inferred form, the result type has a
single parameter $\delta^{\prime}$, without loss of generality because the
actual parameters are passed as a tuple type. In the full format we recover
after inference, we extract the parameters $\delta^{\prime} \dot{=}
(\bar{\beta})$, the non-local variables of the existential type, and the
partially abstract type $\delta \dot{=} \tau$, and store them separately, i.e.
$\varepsilon_K (\bar{\beta}) = \forall \bar{\beta} \exists \bar{\alpha} [D] .
\tau$. The variables $\bar{\beta}$ are instantiated whenever the constructor
is used. For a toplevel \tmverbatim{let}, we form existential types after
solving the generated constraint, to have less intermediate variables in them.

Both during parsing and during inference, we inject new structure items to the
program, which capture the existential types. During printing existential
types in concrete syntax $\exists i : \bar{\beta} [\varphi] .t$ for an
occurrence $\varepsilon_K ((\bar{r}))$, the variables $\bar{\alpha}$ coming
from $\delta^{\prime} \dot{=} (\bar{\alpha}) \in \varphi$ are substituted-out
by $[\bar{\alpha} \assign \bar{r}]$.

For simplicity, only toplevel definitions accept type and invariant
annotations from the user. The constraints are modified according to the
$\llbracket \Gamma, \Sigma \vdash \tmop{ce} : \forall \bar{\alpha} [D] . \tau
\rrbracket$ rule. Where \tmverbatim{let rec}$\ldots$\tmverbatim{in} uses a
fresh variable $\beta$, a toplevel \tmverbatim{let rec} incorporates the type
from the annotation. The annotation is considered partial, $D$ becomes part of
the constraint generated for the recursive function but more constraints will
be added if needed. The polymorphism of $\forall \bar{\alpha}$ variables from
the annotation is preserved since they are universally quantified in the
generated constraint.

The constraints solver returns three components: the {\tmem{residue}}, which
implies the constraint when the predicate variables are instantiated, and the
solutions to unary and binary predicate variables. The residue and the
predicate variable solutions are separated into {\tmem{solved variables}}
part, which is a substitution, and remaining constraints (which are currently
limited to linear inequalities). To get a predicate variable solution we look
for the predicate variable identifier association and apply it to one or two
type variable identifiers, which will instantiate the parameters of the
predicate variable. We considered several ways to deal with multiple
solutions:
\begin{enumerate}
  \item report a failure to the user;
  
  \item ask the user for decision;
  
  \item silently pick one solution, if the wrong one is picked the subsequent
  program might fail;
  
  \item perform backtracking search for the first solution that satisfies the
  subsequent program.
\end{enumerate}
We use approach 3 as it is simplest to implement. Traditional type inference
workflow rules out approach 2, approach 4 is computationally too expensive. We
might use approach 1 in a future version of the system. Upon ``multiple
solutions'' failure -- or currently, when a wrong type or invariant is picked
-- the user can add \tmverbatim{assert} clauses (e.g. \tmverbatim{assert
false} stating that a program branch is impossible), and \tmverbatim{test}
clauses. The \tmverbatim{test} clauses are boolean expressions with
operational semantics of run-time tests: the test clauses are executed right
after the definition is executed, and run-time error is reported when a clause
returns \tmverbatim{false}. The constraints from test clauses are included in
the constraint for the toplevel definition, thus propagate more efficiently
than backtracking would. The \tmverbatim{assert} clauses are:
\tmverbatim{assert type e1 = e2} which translates as equality of types of
\tmverbatim{e1} and \tmverbatim{e2}, \tmverbatim{assert false} which
translates as \tmverbatim{CFalse}, and \tmverbatim{assert num e1 <= e2}, which
translates as inequality $n_1 \leqslant n_2$ assuming that \tmverbatim{e1} has
type \tmverbatim{Num n1} and \tmverbatim{e2} has type \tmverbatim{Num n2}.

We treat a chain of single branch functions with only \tmverbatim{assert
false} in the body of the last function specially. We put all information
about the type of the functions in the premise of the generated constraint.
Therefore the user can use them to exclude unintended types. See the example
\tmverbatim{equal\_assert.gadt}.

\subsection{Normalization}

We reduce the constraint to alternation-minimizing prenex-normal form, as in
the formalization. We return a variable comparison function. The branches we
return from normalization have unified conclusions, since we need to unify for
solving disjunctions anyway.

Releasing constraints \ from under disjunctions is done iteratively, somewhat
similar to how disjunction would be treated in constraint solvers. Releasing
the sub-constraints is essential for eliminating cases of further disjunction
constraints. When at the end more than one disjunct remains, we assume it is
the traditional {\tmname{LetIn}} rule and select its disjunct.

When one $\lambda [K]$ expression is a branch of another $\lambda [K]$
expression, the corresponding branch does not introduce a disjunction
constraint -- the case is settled syntactically to be the same existential
type.

\subsubsection{Implementation Details}

The unsolved constraints are particularly weak with regard to variables
constrained by predicate variables. We need to propagate which existential
type to select for result type of recursive functions, if any. Normalization
starts by flattening constraints into implications with conjunctions of atoms
as premises and conclusions, and disjunctions with disjuncts and additional
information. The additional information kept with a disjunct is the
conjunction of atoms that hold together with the disjunction. We try to
eliminate disjuncts by using unification to check for contradiction. If only
one disjunct is left, or we decide to pick {\tmname{LetIn}} anyway (when no
progress can be made otherwise), we return the disjunct. Otherwise we return
the filtered disjunction.

To help eliminate unintended disjuncts, we collect information about
existential return types of recursive definitions by:
\begin{enumerate}
  \item solving the conclusion of a branch together with additional
  conclusions, to know the return types of variables,
  
  \item registering existential return types for all variables in the
  substitution,
  
  \item registering existential return types for all $\tmop{RetType}$ first
  arguments and their substitution instances,
  
  \item traversing the premise and conclusion to find new variables that are
  types of recursive definitions,
  
  \item registering as ``type of recursive definition'' the return types for
  all variables in the substitution registered as types of recursive
  definitions,
  
  \item traversing all variables known to be types of recursive definitions,
  and registering existential type with recursive definition (i.e. unary
  predicate variable) if it has been learned,
  
  \item traversing all variables known to be types of recursive definitions
  again, and registering existential type of the recursive definition (if any)
  with the variable.
\end{enumerate}

\subsection{Simplification}

During normalization, we remove from a nested premise the atoms it is
conjoined with (as in ``modus ponens'').

After normalization, we simplify the constraints by removing redundant atoms.
We remove atoms that bind variables not occurring anywhere else in the
constraint, and in case of atoms not in premises, not universally quantified.
The simplification step is not currently proven correct and might need
refining. We merge implications with the same premise, unless one of them is
non-recursive and the other is recursive. We call an implication branch
recursive when an unary predicate variable $\chi$ (not a $\chi_K$) appears in
the conclusion {\tmem{or}} a binary predicate variable $\chi_K$ appears in the
premise.

\section{Abduction}

Our formal specification of abduction provides a scheme for combining sorts
that substitutes number sort subterms from type sort terms with variables, so
that a single-sort term abduction algorithm can be called. Since we implement
term abduction over the multisorted datatype \tmverbatim{typ}, we keep these
{\tmem{alien subterms}} in terms passed to term abduction.

\subsection{Abduction for Terms with Alien Subterms}\label{JCAAlienAlg}

Here we expand on the overview from Section \ref{JCAAlienMain}. The JCAQPAS
problem is more complex than simply substituting alien subterms with variables
and performing joint constraint abduction on resulting implications. The
ability to ``outsource'' constraints to the alien sorts enables more general
answers to the target sort, in our case the term algebra $T (F)$. Term
abduction will offer answers that cannot be extended to multisort answers.

One might mitigate the problem by preserving the joint abduction for terms
algorithm, and after a solution $\exists \bar{\alpha} .A$ is found,
``{\tmem{dissociating}}'' the alien subterms (including variables) in $A$ as
follows. We replace every alien subterm $n_s$ in $A$ (including variables,
even parameters) with a fresh variable $\alpha_s$, which results in
$A^{\prime}$ (in particular $A^{\prime} [\overline{\alpha_s} \assign
\overline{n_s}] = A$). Subsets $A^i_p \wedge A^i_c = A^i \subset
\overline{\alpha_s \dot{=} n_s}$ such that $\exists \bar{\alpha}
\overline{\alpha_s} .A^{\prime}, \overline{A^i_p, A^i_c}$ is a JCAQPAS answer
will be recovered automatically by a residuum-finding process after abduction
for terms ends. This process is needed regardless of the ``dissociation''
issue, to uncover the full content of numeric sort constraints.

To face efficiency of numerical abduction with many variables, we modify the
approach. On the first iteration of the main algorithm, we remove (purge)
alien subterms both from the branches and from the answer, but we do not
perform other-sort abduction at all. On the next iteration, we do not purge
alien subterms, neither from the branches nor from the answer, as we expect
the dissociation in the partial solutions (to predicate variables) from the
first step to be sufficient. Other-sort abduction algorithms now have less
work, because only a fraction of alien subterm variables $\alpha_s$ remain in
the partial solutions (see main algorithm in section \ref{MainAlgo}). They
also have more information to work with, present in the instatiation of
partial solutions. However, this optimization violates completeness guarantees
of the combination of sorts algorithm. To faciliate finding term abduction
solutions that hold under the quantifiers, we substitute-out other sort
variables, by variables more to the left in the quantifier, using equations
from the premise.

The dissociation interacts with the discard list mechanism. Since dissociation
introduces fresh variables, no answers with alien subterms would be
syntactically identical. When checking whether a partial answer should be
discarded, in case alien subterm dissociation is {\tmem{on}}, we ignore alien
sort subterms in the comparison.

\subsection{Joint Constraint Abduction}\label{JCAAlg}

Here we expand on the overview from Section \ref{JCAMain}. We further lose
generality by using a heuristic search scheme instead of testing all
combinations of simple abduction answers. In particular, our search scheme
returns from joint abduction for types with a single answer, which eliminates
deeper interaction between the sort of types and other sorts. Some amount of
interaction is provided by the validation procedure, which checks for
consistency of the partial answer, the premise and the conclusion of each
branch, including consistency for other sorts.

We accumulate simple abduction answers into the partial abduction answer, we
set aside branches that do not have any answer satisfiable with the partial
answer so far. After all branches have been tried and the partial answer is
not an empty conjunction (i.e. not $\top$), we retry the set-aside branches.
If during the retry, any of the set-aside branches fails, we add the partial
answer to discarded answers -- which are avoided during simple abduction --
and restart. Restart puts the set-aside branches to be tried first. If, when
left with set-aside branches only, the partial answer is an empty conjunction,
i.e. all the answer-contributing branches have been set aside, we fail --
return $\bot$ from the joint abduction. This does not peform complete
backtracking (no completeness guarantee), but is therefore quicker to report
unsolvable cases and does sufficient backtracking. After an answer working for
all branches has been found, we perform additional check, which encapsulates
negative constraints introduced by the \tmverbatim{assert false} construct. If
the check fails, we add the answer to discarded answers and repeat the search.

If a partial answer becomes as strong as one of the discarded answers inside
SCA, simple constraint abduction skips to find a different answer. The
discarded answers are initialized with a discard list passed from the main
algorithm.

To check validity of answers, we use a modified variant of unification under
quantifiers: unification with parameters, where the parameters do not interact
with the quantifiers and thus can be freely used and eliminated. Note that to
compute conjunction of the candidate answer with a premise, unification does
not check for validity under quantifiers.

Because it would be difficult to track other sort constraints while updating
the partial answer, we discard numeric sort constraints in simple abduction
algorithm, and recover them after the final answer for terms (i.e. for the
type sort) is found.

Searching for an abduction answer can fail in only one way: we have set aside
all the branches that could contribute to the answer. It is difficult to
pin-point the culprit. We remember which branch caused the restart when the
number of set-aside branches was the smallest. The conclusion of that branch
can be used to construct the error report.

\subsection{Simple Constraint Abduction for Terms}\label{SCATermsAlg}

Here we expand on the overview from Section \ref{SCATermsMain}. Our initial
implementation of simple constraint abduction for terms follows
{\cite{AbductionSolvMaher}} p. 13. The mentioned algorithm only gives
{\tmem{fully maximal answers}} which is loss of generality w.r.t. our
requirements. To solve $D \Rightarrow C$ the algorithm starts with
$\tmmathbf{U} (D \wedge C)$ and iteratively replaces subterms by fresh
variables $\alpha \in \bar{\alpha}$ for a final solution $\exists \bar{\alpha}
.A$. As our primary approach to mitigate some of the limitations of fully
maximal answers, we start from $\tmmathbf{U}_{} (\tilde{A} (D \wedge C))$,
where $\exists \bar{\alpha} .A$ is the solution to previous problems solved by
the joint abduction algorithm, and $\tilde{A} (\cdummy)$ is the corresponding
substitution. Moreover, motivated by examples from Chuan-kai Lin
{\cite{PracticalGADTsInfer}}, we intruduce variable-variable equations
$\beta_1 \dot{=} \beta_2$, for $\beta_1 \beta_2 \subset \bar{\beta}$, not
implied by $\tilde{A} (D \wedge C)$, as additional candidate answer atoms.
During abduction $\tmop{Abd} (\mathcal{Q}, \bar{\beta}, \overline{D_i, C_i})$,
we ensure that the (partial as well as final) answer $\exists \bar{\alpha} .A$
satisfies $\vDash \mathcal{Q}.A [\bar{\alpha} \bar{\beta} \assign \bar{t}]$
for some $\bar{t}$. We achieve this by normalizing the answer using
parameterized unification under quantifiers $\tmmathbf{U}_{\bar{\alpha}
\bar{\beta}} (\mathcal{Q}.A)$. $\bar{\beta}$ are the parameters of the
invariants.

In fact, when performing unification, we check more than
$\tmmathbf{U}_{\bar{\alpha} \bar{\beta}} (\mathcal{Q}.A)$ requires. We also
ensure that the use of parameters will not cause problems in the
$\tmop{Split}$ phase of the main algorithm. To this effect, we forbid
substitution of a variable $\beta_1$ from $\bar{\beta}$ with a term containing
a universally quantified variable that is not in $\bar{\beta}$ and to the
right of $\beta_1$ in $\mathcal{Q}$. Also, we forbid substitution of a
variable $\beta_1$ from $\bar{\beta}^{\chi}$ with a term containing a variable
$\beta_2 \in \bar{\beta}^{\chi^{\prime}}$ for $\chi \neq \chi^{\prime}$.

In implementing {\cite{AbductionSolvMaher}} p. 13, we follow a top-down
approach where bigger subterms are abstracted first -- replaced by a fresh
variable, \ together with an arbitrary selection of other occurrences of the
subterm. If dropping the candidate atom maintains $T (F) \vDash A \wedge D
\Rightarrow C$, we proceed to neighboring subterm or next equation. Otherwise,
we try all of: replacing the subterm by the fresh variable; proceeding to
subterms of the subterm; preserving the subterm; replacing the subterm by
variables corresponding to earlier occurrences of the subterm. This results in
a single, branching pass over all subterms considered. Finally, we clean-up
the solution by eliminating fresh variables when possible (i.e.
substituting-out equations $x \dot{=} \alpha$ for variable $x$ and fresh
variable $\alpha$).

Although there could be an infinite number of abduction answers, there is
always a finite number of {\tmem{fully maximal}} answers, or more generally, a
finite number of equivalence classes of formulas strictly stronger than a
given conjunction of equations in the domain $T (F)$. We use a search scheme
that tests as soon as possible. The simple abduction algorithm takes a partial
solution -- a conjunction of candidate solutions for some other branches --
and checks if the solution being generated is satisfiable together with the
candidate partial solution. The algorithm also takes several indicators to let
it select the expected answer:
\begin{itemize}
  \item a number that determines how many correct solutions to skip;
  
  \item a validation procedure that checks whether the partial answer meets a
  condition, in joint abduction the condition is consistency with premise and
  conclusion of each branch;
  
  \item the parameters and candidates for parameters of the invariants,
  $\bar{\beta}$, updated as we add new atoms to the partial answer;
  existential variables that are not to the left of parameters and are
  connected to parameters become parameters; we process atoms containing
  parameters first;
  
  \item the quantifier $\mathcal{Q}$ (source \tmverbatim{q}) so that the
  partial answer $\exists \bar{\alpha} .A$ (source \tmverbatim{vs,ans}) can be
  checked for validity with parameters: $\vDash \mathcal{Q}.A [\bar{\alpha}
  \assign \bar{t}]$ for some $\bar{t}$;
  
  \item a discard list of partial answers to avoid (a tabu list) -- implements
  backtracking, with answers from abductions raising ``fallback'' going there.
\end{itemize}
Since an atom can be mistakenly discarded when some variable could be
considered an invariant parameter but is not at the time, we process atoms
incident with candidates for invariant parameters first. A variable becomes a
candidate for a parameter if there is a parameter that depends on it. That is,
we process atoms $x \dot{=} t$ such that $x \in \bar{\beta}$ first, and if
equation $x \dot{=} t$ is added to the partial solution, we add to
$\bar{\beta}$ existential variables in $t$. Note that $x \dot{=} t$ can stand
for either $x \assign t$, or $y \assign x$ for $t = y$. For a universally
quantified variable $x \nin \bar{\beta}$, $x \assign t$ will not be part of
the answer as it does not hold under the quantifier.

To simplify the search in presence of a quantifier prefix, we preprocess the
initial candidate by eliminating universally quantified variables:
\begin{eqnarray*}
  S & = & [\overline{t_u} \assign \overline{t^{\prime}_u}] \text{ for }
  \tmop{FV} (t_u) \cap \overline{\beta_u} \neq \varnothing, \overline{\forall
  \beta_u} \subset \mathcal{Q} \text{ such that } \mathcal{M} \vDash D
  \Rightarrow \dot{S},\\
  S^{\prime} & = & [\bar{u} \assign \overline{t^{\prime}_u}] \text{ for }
  \bar{u} \subset \overline{\beta_u}, \overline{\forall \beta_u} \subset
  \mathcal{Q} \text{ such that } \mathcal{M} \vDash D \wedge C \Rightarrow
  \dot{S}^{\prime},\\
  \tmop{Rev}_{\forall} (\mathcal{Q}, \bar{\beta}, D, C) & = & \{ c^{\prime} |
  c = x \dot{=} t \in C, \text{if } x = S^{\prime} (t) \text{ then }
  c^{\prime} = S (c) \text{ else } c^{\prime} = SS^{\prime} (c) \}
\end{eqnarray*}
Note that $S$ above is a substitution of subterms rather than of variables. To
move further beyond fully maximal answers, we incorporate candidates $\beta_1
\dot{=} \beta_2$ for which the following conditions hold: $\beta_1 \beta_2
\subset \bar{\beta}$, $\beta_1 \assign t_1 \in \tmmathbf{U}_{} (\tilde{A} (D
\wedge C))$, $\beta_2 \assign t_2 \in \tmmathbf{U}_{} (\tilde{A} (D \wedge
C))$ and $t_1 \dot{=} t_2$ is satisfiable. We also need to include the unifier
of $t_1 \dot{=} t_2$ among the candidates, since otherwise the equation
$\beta_1 \dot{=} \beta_2$ would not suffice to imply that of the atoms
$\beta_1 \dot{=} t_1$, $\beta_2 \dot{=} t_2$ which belongs to the conclusion
$C$. The {\tmem{full candidates}} $\tmmathbf{U}_{} (\tilde{A} (D \wedge C))$
and the {\tmem{guess candidates}} $\overline{\beta_1 \assign \beta_2 ;
\tmmathbf{U} (t_1 \dot{=} t_2)}$ are kept apart, the guess candidates are
guessed before the full candidates. By default, we additionally limit
consideration to atoms $\beta_1 \dot{=} t_1$, $\beta_2 \dot{=} t_2$ where
$t_1, t_2$ are not themselves variables.

To recapitulate, the implementation is:
\begin{itemize}
  \item If there are no more candidates to add to the partial solution: check
  for repeated answers, skipping, and discarded answers.
  
  \item If there are no more guessed candidates, pick the next full candidate
  atom $\tmop{FV} (c) \cap \bar{\beta} \neq \varnothing$ if any, reordering
  the candidates until one is found. Otherwise, pick the first guess candidate
  atom without reordering.
  
  \item The are 6 mutually exclusive choices through which the algorithm
  loops.
  \begin{enumerate}
    \item Try to drop the atom (if the partial answer plus remaining
    candidates can still be completed to a correct answer).
    
    \item Replace the current subterm of the atom with a fresh parameter,
    adding the subterm to replacements; if at the root of the atom, check
    connected and validate before proceeding to remaining candidates.
    
    \item Step into subterms of the current subterm, if any, and if at the
    sort of types.
    
    \item Keep the current part of the atom unchanged; if at the root of the
    atom, check connected and validate before proceeding to remaining
    candidates.
    
    \item Replace the current subterm with a parameter introduced for an
    earlier occurrence; branch over all matching parameters; if at the root of
    the atom, check connected and validate before proceeding to remaining
    candidates.
    
    \item Keep a variant of the original atom, but with constants
    substituted-out by variable-constant equations from the premise.
    Redundant, not available when the option \tmverbatim{-more\_general} is
    on.
  \end{enumerate}
  \item Each iteration is a backtracking point, the choices are tried in turn,
  depending on options selected.
  
  \item Default ordering of choices is 1, 6, 2, 4, 3, 5 -- pushing 4 up
  minimizes the amount of branching in 5.
  \begin{itemize}
    \item There is an option \tmverbatim{-more\_general}, which reorders the
    choices to: 1, 6, 4, 2, 3, 5; however the option is not exposed in the
    interface because the cost of this reordering is prohibitive.
    
    \item An option \tmverbatim{-richer\_answers} reorders the choices to: 6,
    1, 2, 4, 3, 5; it does not increase computational cost but sometimes leads
    to answers that are not most general.
    
    \item If choice 6 would lead to more negative constraints contradicted
    than choice 1, we pick choice 6 first for a particular candidate atom.
    
    \item An option \tmverbatim{-prefer\_guess} reorders choice 6 prior to
    choice 1, but only for guess candidates.
  \end{itemize}
  \item Form initial candidates $\tmop{Rev}_{\forall} (\mathcal{Q},
  \bar{\beta}, \tmmathbf{U} (D \wedge A_p), \tmmathbf{U} (A_p \wedge D \wedge
  C))$.
  
  \item Form the substitution of subterms for choice-6 counterparts of initial
  candidate atoms. For $\alpha_1 \dot{=} \tau, \ldots, \alpha_n \dot{=} \tau
  \in \tmmathbf{U} (D \wedge A_p)$, form the substitution of subterms
  $\alpha_1 \assign \alpha_i, \ldots, \alpha_n \assign \alpha_i, \tau \assign
  \alpha_i$ (excluding $\alpha_i \assign \alpha_i$) where $\alpha_i$ is the
  most upstream existential variable (or parameter) and $\tau$ is a constant.
  Note that analogous transformation for a universally quantified variable as
  right-hand-sides is performed by $\tmop{Rev}_{\forall}$.
  \begin{itemize}
    \item Since for efficiency reasons we do not always remove alien subterms,
    we need to mitigate the problem of alien subterm variables causing
    violation of the quantifier prefix. To this effect, we include the premise
    equations from other sorts in the process generating the initial
    candidates and choice 6 candidates, but not as candidates. Not to lose
    generality of answer, we only keep a renaming substitution, in particular
    we try to eliminate universal variables.
  \end{itemize}
  \item Sort the initial candidates by decreasing size, because shorter answer
  atoms are more valuable and dropping a candidate from participating in an
  answer is the first choice.
  \begin{itemize}
    \item There is an argument in favor of sorting by increasing size: so that
    the replacements of step 2 are formed at a root position before they are
    used in step 5 -- instead of forming a replacement at a subterm, and using
    it in step 5 at a root.
    
    \item If ordering in increasing size turns out to be necessary, a
    workaround should be introduced to favor answers that, if possible, do not
    have parameters $\bar{\beta}$ as left-hand-sides.
  \end{itemize}
\end{itemize}
The above ordering of choices ensures that more general answers are tried
first. Moreover:
\begin{itemize}
  \item choice 1 could be dropped as it is equivalent to choice 2 applied on
  the root term;
  
  \item choices 4 and 5 could be reordered but having choice 4 as soon as
  possible is important for efficiency.
\end{itemize}
We perform a two-layer iterative deepening (when without the
\tmverbatim{-more\_general} option): in the first run we only try choices 1
and 6. It is an imperfect optimization since the running time gets longer
whenever choices 2-5 are needed.

\subsubsection{Heuristic for Better Answers to Invariants}

We implement an optional heuristic in forming the candidates proposed by
choice 6. It may lead to better invariants when multiple maximally general
types are possible, but also it may lead to getting the most general type
without the need for backtracking across iterations of the main algorithm,
which unfortunately often takes very long.

We look at the types of substitutions for the variables that are invariant
parameters, in the partial answer, and try to form the initial candidates for
choice 6 so that the return type variables cover the most of argument types
variables, for each term substituted for an invariant parameter. We select
from the candidates equations between any variable, or only non-argument-type
variable, and a $\tmop{FV} (\text{argument types}) \backslash \tmop{FV}
(\text{return type})$ variable -- we turn the equation so that the latter is
the RHS. We locate the equations among the candidates that have an invariant
parameter variable or a $\tmop{FV} (\text{return type}) \backslash \tmop{FV}
(\text{argument types})$ variable as LHS. We apply the substitution to the RHS
of these equations; if the LHS is an invariant parameter, we use the
substitution based on equations where one side is a non-argument-type
variable. We preserve the order of equations in the candidate list.

\subsection{Simple Constraint Abduction for Linear
Arithmetics}\label{SCANumAlg}

Here we expand on the overview from Section \ref{SCANumMain}. For checking
validity or satisfiability, we use {\tmem{Fourier-Motzkin elimination}}. To
avoid complexities we only handle the rational number domain. To extend the
algorithm to integers, {\tmem{Omega-test}} procedure as presented in
{\cite{ArithQuantElim}} needs to be adapted. The major operations are:
\begin{itemize}
  \item {\tmem{Elimination}} of a variable takes an equation and selects a
  variable that is not upstream, i.e. to the left in $\mathcal{Q}$ regarding
  alternations, of any other variable of the equation, and substitutes-out
  this variable from the rest of the constraint. The solved form contains an
  equation for this variable.
  
  \item {\tmem{Projection}} of a variable takes a variable $x$ that is not
  upstream of any other variable in the unsolved part of the constraint, and
  reduces all inequalities containing $x$ to the form $x \dot{\leqslant} a$ or
  $b \dot{\leqslant} x$, depending on whether the coefficient of $x$ is
  positive or negative. For each such pair of inequalities: if $b = a$, we add
  $x \dot{=} a$ to implicit equalities; otherwise, we add the inequality $b
  \dot{\leqslant} a$ to the unsolved part of the constraint.
\end{itemize}
We use elimination to solve all equations before we proceed to inequalities.
The starting point of our algorithm is {\cite{ArithQuantElim}} section
{\tmem{4.2 Online Fourier-Motzkin Elimination for Reals}}. We add detection of
implicit equalities, and more online treatment of equations, introducing known
inequalities on eliminated variables to the projection process. When implicit
equalities have been found, we iterate the process to normalize them as well.

Our abduction algorithm follows a familiar incrementally-generate-and-test
scheme as in term abduction. There are two new ideas, which can also be
applied to abduction in other domains. We build a lazy list of possible
transformations with linear combinations involving equations $a$ implied by $D
\wedge C$. We pair each inequality $c$ in $C$ with all inequalities $d$
implied by $D$ which share a variable with $c$ and try out the abduction
answers to $d \Rightarrow c$ as contributions to the partial abduction answer
to $D \Rightarrow C$.

To simplify the search in presence of a quantifier prefix, we preprocess the
initial candidate by eliminating universally quantified variables:
\begin{eqnarray*}
  S & = & [\overline{\beta_u} \assign \overline{t_u}] \text{ for }
  \overline{\forall \beta_u} \subset \mathcal{Q} \text{ such that }
  \mathcal{M} \vDash D \Rightarrow \dot{S},\\
  \tmop{Rev}_{\forall} (\mathcal{Q}, \bar{\beta}, D, C) & = & \{ c^{\prime} |
  c \in C, \text{if } \mathcal{M} \vDash \mathcal{Q}.c [\bar{\beta} \assign
  \bar{t}] \text{ for some } \bar{t} \text{ then } c^{\prime} = c \text{ else
  } c^{\prime} = S (c) \}
\end{eqnarray*}
Before accepting a new atom into the partial answer, we check that it would
not violate the quantifier conditions from the {\tmem{split}} phase of the
main algorithm, and that the partial answer is satisfiable with all
implication branches of the joint abduction problem. As part of the quantifier
conditions, we ensure that the escaping parameters are upward in the
constraint before prenexization, i.e. are parameters of the type of a parent
definition (containing a given definition in its body) rather than a parallel
definition. The check of satisfiability we call {\tmem{validation}}. For
domains other than the term domain, validation also involves instantiating
use-sites of recursive definitions with parts of the partial answer, split in
a simplified way.

Abduction algorithm:
\begin{enumerate}
  \item Let $C^{=\prime} = \widetilde{A_i} (C^=)$, resp. $C^{\leqslant \prime}
  = \widetilde{A_i} (C^{\leqslant})$ where $C^=$, resp. $C^{\leqslant}$ are
  the equations, resp. inequalities in $C$ and $\widetilde{A_i}$ is the
  substitution according to equations in $A_i$. Let $D^{\prime} =
  \widetilde{A_i} \left( D \wedge {\nobreak} A_i \right)$ and $D^{=\prime}$ be
  the equations in $D^{\prime}$, i.e. substituted equations and implicit
  equalities. Let $D^{\leqslant \prime}$ be a solved form of inequalities.
  \begin{enumerate}
    \item Let $C^=_0 = \tmop{Rev}_{\forall} (\mathcal{Q}, \bar{\beta},
    D^{=\prime}, C^{=\prime})$ and $C^{\leqslant}_0 = \tmop{Rev}_{\forall}
    (\mathcal{Q}, \bar{\beta}, D^{=\prime}, C^{\leqslant \prime})$.
  \end{enumerate}
  \item Prepare the initial transformations from atoms $a \in D^{=\prime}$:
  \begin{enumerate}
    \item Add combinations $k^s a + b$ for $k = - n \ldots n, s = - 1, 1$ to
    the stack of transformations to be tried for atoms $b$.
    
    \item The final transformations have the form: $b \mapsto b + \Sigma_{a
    \in D} k_a^{s_a} a$.
  \end{enumerate}
  \item Modify $C^{=\prime}$ to promote answers with variables rather than
  constants, as in term abduction: For $\alpha_1 \dot{=} \tau, \ldots,
  \alpha_n \dot{=} \tau \in C^{=\prime}$, form the substitution of subterms
  $\alpha_1 \assign \alpha_i, \ldots, \alpha_n \assign \alpha_i, \tau \assign
  \alpha_i$ (excluding $\alpha_i \assign \alpha_i$) where $\alpha_i$ is the
  most upstream existential variable (or parameter) and $\tau$ is a constant.
  
  \item Start from $\tmop{Acc} \assign \{ \}$ and $C_0 \assign C^{=\prime}
  \wedge C^{\leqslant \prime}$. Handle atoms $a$ in $C_0 = a C_0^{\prime}$,
  equations first.
  
  \item Let $B = A_i \wedge D \wedge C_0^{\prime} \wedge \tmop{Acc}$.
  
  \item If $a$ is a tautology ($0 \dot{=} 0$ or $c \leq 0$ for $c \leqslant
  0$) or $B \Rightarrow C$, repeat with $C_0 \assign C_0^{\prime}$.
  Corresponds to choice 1 of term abduction.
  
  \item If $B \nRightarrow C$, for a candidate $a^{\prime}$ generated for $a$,
  starting with $a$, which passes validation against other branches in a joint
  problem: $\tmop{Acc} \assign \tmop{Acc} \cup \{ a^{\prime} \}$, or fail if
  all $a^{\prime}$ fail.
  \begin{enumerate}
    \item If $a$ is an inequality, let $a_0$ be an abduction answer to $d
    \Rightarrow \tmop{Rev}_{\forall} (\mathcal{Q}, \bar{\beta}, D^{=\prime},
    \widetilde{\tmop{Acc}^=} (a))$, where $d$ belongs to the solved form
    inequalities implied by $D^{\prime}$. Otherwise, let $a_0 =
    \widetilde{\tmop{Acc}^=} (a)$.
    
    \item Sort the candidates $a_0$ in order of decreasing value, described
    below. Let $a^{\prime}$ be $a_0$ with some $D^{=\prime}$-derived
    transformation applied.
    \begin{itemize}
      \item We increase the value of $a_0$ for each variable that is bound by
      $a_0$ on the side that it is unbound in $B$. We decrease the value of
      $a_0$ for:
      \begin{itemize}
        \item its size, proportionally to the sum of nominators and
        denominators,
        
        \item containing a constant term,
        
        \item containing parameters from multiple invariants,
        
        \item introducing an implicit equality,
        
        \item binding a variable by a constant on the side where it is already
        bounded by $B$,
        
        \item binding a variable by a constant on the right (i.e. upper
        bound),
        
        \item optionally, being less general than some other candidate, modulo
        $B$,
        
        \item we include the value of inequalities implied by the candidate
        together with the partial answer (but not the partial answer alone)
        and not holding when parameters are universally quantified (see
        {\tmem{atomization}}).
      \end{itemize}
      The score determines the order in which atoms are tried.
      
      \item Currently, we do not perform transformations when $a_0$ is an
      inequality, for simplicity and speed at cost of missing some answers.
      However, we eliminate the universal variables, i.e. we use
      $\tmop{Rev}_{\forall}$ above. If it proves necessary, we will also try
      the transformations for the inequalities. For example, by trying all
      candidates $a_0$ before proceeding to the next transformation.
    \end{itemize}
    \item If $A_i \wedge (\tmop{Acc} \cup \{ a^{\prime} \})$ (resp. $A_i
    \wedge (\tmop{Acc} \cup \{ a^{\prime\prime} \})$) does not pass validation
    for all $a^{\prime}$, backtrack.
    
    \item If $A_i \wedge (\tmop{Acc} \cup \{ a^{\prime} \})$ (resp. $A_i
    \wedge (\tmop{Acc} \cup \{ a^{\prime\prime} \})$) passes validation,
    repeat from step 5 with $C_0 \assign C_0^{\prime}, \tmop{Acc} \assign
    \tmop{Acc} \cup \{ a^{\prime} \}$ (resp. $\tmop{Acc} \assign \tmop{Acc}
    \cup \{ a^{\prime\prime} \}$).
  \end{enumerate}
  \item The answers are $A_{i + 1} = A_i \wedge \tmop{Acc}$.
\end{enumerate}
We precompute the tranformation variants to try out. The parameter $n$ is set
by option \tmverbatim{-{\nobreak}num\_abduction\_rotations} and defaults to a
small value (currently 3).

To check whether $B \Rightarrow C$, we check for each $c \in C$:
\begin{itemize}
  \item if $c = x \dot{=} y$, that $A (x) = A (y)$, where $A (\cdummy)$ is the
  substitution corresponding to equations and implicit equalities in $A$;
  
  \item if $c = x \dot{\leqslant} y$, that $B \wedge y \dot{<} x$ is not
  satisfiable.
\end{itemize}
We use the \tmverbatim{nums} library for exact precision rationals.

To find the abduction answers to $d \Rightarrow c$, pick a common variable
$\alpha \in \tmop{FV} (d) \cap \tmop{FV} (c)$ or the constant $\alpha = 1$. We
have four possibilities:
\begin{enumerate}
  \item $d \Leftrightarrow \alpha \leqslant d_{\alpha}$ and $c \Leftrightarrow
  \alpha \leqslant c_{\alpha}$: the abduction answers are $c$ and $d_{\alpha}
  \leqslant c_{\alpha}$,
  
  \item $d \Leftrightarrow \alpha \leqslant d_{\alpha}$ and $c \Leftrightarrow
  c_{\alpha} \leqslant \alpha$: the abduction answer is only $c$,
  
  \item $d \Leftrightarrow d_{\alpha} \leqslant \alpha$ and $c \Leftrightarrow
  \alpha \leqslant c_{\alpha}$: the abduction answer is only $c$,
  
  \item $d \Leftrightarrow d_{\alpha} \leqslant \alpha$ and $c \Leftrightarrow
  c_{\alpha} \leqslant \alpha$: the abduction answers are $c$ and $c_{\alpha}
  \leqslant d_{\alpha}$.
\end{enumerate}
Thanks to cases (1) and (4) above, the abduction algorithm can find some
answers which are not fully maximal. The joint constraint abduction algorithm
can help in some of the remaining cases where fully maximal abduction is
insufficient for some implications, by solving simpler implications first.

We provide an optional optimization: we do not pass, in the first call to
numerical abduction (the second iteration of the main algorithm), branches
that contain unary predicate variables in the conclusion, i.e. we only use the
``non-recursive'' branches. Other optimizations that we use are: iterative
deepening on the constant $n$ used to generate $k^s$ factors. We also
constrain the algorithm by filtering out transformations that contain ``too
many'' variables, which can lead to missing answers if the setting
\tmverbatim{-num\_prune\_at} -- ``too many'' -- is too low. Similarly to term
abduction, we count the number of steps of the loop and fail if more than the
option \tmverbatim{-num\_abduction\_timeout} steps have been taken.

\section{Constraint Generalization}\label{ConstrnGenAlg}

Here we expand on the exposition from Section \ref{ConstrnGenMain}.
{\tmem{Constraint generalization}} answers are the maximally specific
conjunctions of atoms that are implied by each of a given set of conjunction
of atoms. In case of term equations the constraint generalization algorithm is
based on the {\tmem{anti-unification}} algorithm. In case of linear arithmetic
inequalities, constraint generalization is exactly finding the convex hull of
a set of possibly unbounded polyhedra. We employ our unification algorithm to
separate sorts. Since as a result we do not introduce variables for
{\tmem{alien subterms}}, we include the variables introduced by
anti-unification in constraints sent to constraint generalization for their
respective sorts.

The adjusted algorithm looks as follows:
\begin{enumerate}
  \item Let $\wedge_s D_{i, s} \equiv \tmmathbf{U} (D_i)$ where $D_{i, s}$ is
  of sort $s$, be the result of our sort-separating unification.
  
  \item For the sort $s_{\tmop{ty}}$:
  \begin{enumerate}
    \item Let $V = \{ x_j, \overline{t_{i, j}} | \forall i \exists t_{i, j}
    .x_j \dot{=} t_{i, j} \in D_{i, s_{\tmop{ty}}} \}$.
    
    \item Let $G = \{ \bar{\alpha}_j, u_j, \overline{\theta_{i, j}} |
    \theta_{i, j} = [\bar{\alpha}_j \assign \bar{g}_j^i], \theta_{i, j} (u_j)
    = t_{i, j} \}$ be the most specific anti-unifiers of $\overline{t_{i, j}}$
    for each $j$.
    
    \item Let $D_i^u = \wedge_j \bar{\alpha}_j \dot{=} \bar{g}_j^i$ and $D_i^g
    = D_{i, s_{\tmop{ty}}} \wedge D_i^u$.
    
    \item Let $D^v_i = \{ x \dot{=} y | x \dot{=} t_1 \in D_i^g, y \dot{=} t_2
    \in D_i^g, D_i^g \vDash t_1 \dot{=} t_2 \}$.
    
    \item Let $A_{s_{\tmop{ty}}} = \wedge_j x_j \dot{=} u_j \wedge \bigcap_i
    (D_i^g \wedge D_i^v)$ (where conjunctions are treated as sets of conjuncts
    and equations are ordered so that only one of $a \dot{=} b, b \dot{=} a$
    appears anywhere), and $\bar{\alpha}_{s_{\tmop{ty}}} =
    \overline{\bar{\alpha}_j}$.
    
    \item Let $\wedge_s D^u_{i, s} \equiv D^u_i$ for $D^u_{i, s}$ of sort $s$.
  \end{enumerate}
  \item For sorts $s \neq s_{\tmop{ty}}$, let $\exists \bar{\alpha}_s .A_s =
  \tmop{LUB}_s (\overline{D_i^s \wedge D^u_{i, s}})$.
  
  \item The answer is $\exists \overline{\overline{\alpha^j_i}}
  \overline{\bar{\alpha}_s} . \wedge_s A_s$.
\end{enumerate}
We simplify the result by substituting-out redundant answer variables.

We follow the anti-unification algorithm provided in {\cite{AntiUnifAlg}},
fig. 2.

\subsection{Extended Convex Hull}\label{NumConstrnGenAlg}

{\cite{ConvexHull}} provides a polynomial-time algorithm to find the
half-space represented convex hull of closed polytopes. It can be generalized
to unbounded polytopes -- conjunctions of linear inequalities. Our
implementation is inspired by this algorithm but very much simpler, at cost of
losing the optimality requirement.

First we find among the given inequalities those which are also the faces of
resulting convex hull. The negation of such inequality is not satisfiable in
conjunction with any of the polytopes -- any of the given sets of
inequalities. Next we iterate over {\tmem{ridges}} touching the selected
faces: pairs of the selected face and another face from the same polytope. We
rotate one face towards the other: we compute a convex combination of the two
faces of a ridge. We add to the result those half-spaces whose complements lie
outside of the convex hull (i.e. negation of the inequality is unsatisfiable
in conjunction with every polytope). For a given ridge, we add at most one
face, the one which is farthest away from the already selected face, i.e. the
coefficient of the selected face in the convex combination is smallest. We
check a small number of rotations, where the algorithm from
{\cite{ConvexHull}} would solve a linear programming problem to find the
rotation which exactly touches another one of the polytopes.

When all variables of an equation $a \dot{=} b$ appear in all branches $D_i$,
we can turn the equation $a \dot{=} b$ into pair of inequalities $a \leqslant
b \wedge b \leqslant a$. We eliminate all equations and implicit equalities
which contain a variable not shared by all $D_i$, by substituting out such
variables. We pass the resulting inequalities to the convex hull algorithm.
Separately, we compute the equations common to all branches, because the
convex hull algorithm is not guaranteed to recover them.

\subsection{Issues in Inferring Postconditions}\label{NumConv}

Although finding recursive function invariants -- predicate variables solved
by abduction -- could theoretically fail to converge for both the type sort
and the numerical sort constraints, neither problem was observed. Finding
existential type constraints can only fail to converge for numerical sort,
because solutions are expected to decrease in strength. But such diverging
numerical constraints are commonplace. The main algorithm starts by performing
constraint generalization only on implication branches corresponding to
non-recursive cases, i.e. without binary predicate variables in premise (or
unary predicate variables in conclusion). This generates a stronger constraint
than the correct one. Subsequent iterations include all branches in constraint
generalization, weakening the constraints, and so still weaker constraints are
fed to constraint generalization in each following step. To ensure convergence
of the numerical part, starting from some step of the main loop, we compare
consecutive solutions and extrapolate the trend. Currently we simply intersect
the sets of atoms, but first we expand equations into pairs of inequalities.

We ``lift'' variables escaping the scope of a postcondition by renaming them
to fresh variables, which are added to the answer variables of generalization.
We apply this renaming ater anti-unification, prior to performing
generalization in other domains. We pass to other domains as parameters to
preserve only non-escaping variables. The non-escaping variables are these
which have in their scope a variable $\alpha_K$, passed as argument to the
generalization algorithm.

Constraint generalization limited to non-recursive branches, the initial
iteration of postcondition inference, will often generate constraints that
contradict other branches. For another iteration to go through, the partial
solutions need to be consistent. Therefore we filter the constraints using the
same validation mechanism as in abduction. We add atoms to a constraint
greedily, but to favor relevant atoms, we do the filtering while computing the
connected component of constraint generalization result. See the details of
the main algorithm in section \ref{MainAlgoBody}.

While reading section \ref{MainAlgoBody}, you will notice that postconditions
are not subjected to stratification. This is because the type system does not
support nested existential types.

In the simplification step at the end of constraint generalization, we try to
preserve alien variables that are parameters rather than substituting them by
constants. A parameter can only equal a constant if not all branches have been
considered for constraint generalization. The parameter is both as informative
as the constant, and less likely to contradict other branches.

\subsection{Abductive Constraint Generalization}\label{AbdConstrnGenAlg}

Here we expand on the overview from Section \ref{AbdConstrnGenMain}.

{\tmem{Global variables}} here are the variables shared by all disjuncts, i.e.
$\cap_i \tmop{FV} (D_i)$, remaining variables are {\tmem{non-global}}. Recall
that for numerical constraint generalization, we either substitute-out a
non-global variable in a branch if it appears in an equation, or we drop the
inequalities it appers in if it is not part of any equation. Non-global
variables can also pose problems for the term sort, by forcing constraint
generalization answers to be too general. When inferring the type for a
function, which has a branch that does not use one of arguments of the
function, the existential type inferred would hide the corresponding
information in the result, even if the remaining branches assume the argument
has a single concrete type. We would like the corresponding non-global
variable to resolve to the concrete type suggested by other branches of the
resulting constraint.

We extend the notion of constraint generalization: substitution $U$ and solved
form $\exists \bar{\alpha} .A$ is an answer to {\tmem{abductive constraint
generalization}} problem $\overline{D_i}$ given a quantifier prefix
$\mathcal{Q}$ when:
\begin{enumerate}
  \item $(\forall i) \vDash U (D_i) \Rightarrow \exists \bar{\alpha}
  \backslash \tmop{FV} (U) .A$;
  
  \item If $\alpha \in \tmop{Dom} (U)$, then $(\exists \alpha) \in
  \mathcal{Q}$ -- variables substituted by $U$ are existentially quantified;
  
  \item $(\forall i) \vDash \forall (\tmop{Dom} (U)) \exists (\tmop{FV} (D_i)
  \backslash \tmop{Dom} (U)) .D_i$.
\end{enumerate}
Our generalized anti-unification algorithm is in Table
\ref{AntiUnifAlg}.\begin{table}[t]
  \begin{eqnarray*}
    \tmop{au}_{U, G} (t ; \ldots ; t) & = & \varnothing, t, U, G\\
    \tmop{au}_{U, G} (f (\bar{t}^1) ; \ldots ; f (\bar{t}^n)) & = &
    \bar{\alpha}, f (\bar{g}), U^{\prime}, G^{\prime}\\
    \text{where \ } \bar{\alpha}, \bar{g}, U^{\prime}, G^{\prime} & = &
    \tmop{aun}_{U, G} (\bar{t}^1 ; \ldots ; \bar{t}^n)\\
    \tmop{au}_{U, G} (t_1 ; \ldots ; t_n) & = & \varnothing, \alpha, U, G\\
    \text{when} &  & ([t_1 ; \ldots ; t_n] \mapsto \alpha) \in G\\
    \tmop{au}_{U, G} (\ldots ; \beta_i ; \ldots ; f (\bar{t}^j) ; \ldots
    \text{ as } \bar{t}) & = & \bar{\alpha} \bar{\alpha}^{\prime}, g,
    U^{\prime\prime}, G^{\prime}\\
    \text{where \ } \bar{\alpha}^{\prime}, g, U^{\prime\prime}, G^{\prime} & =
    & \tmop{au}_{U^{\prime}, G} (\bar{t} [\beta_i \assign f (\bar{\alpha})])\\
    U^{\prime} & = & U [\beta_i \assign f (\bar{\alpha})] \wedge \beta_i
    \dot{=} f (\bar{\alpha})\\
    \text{when} &  & (\exists \beta_i) \in \mathcal{Q} \vee \beta_i \in
    \bar{\beta} \text{, treat } \bar{\alpha} \text{ as quantified with }
    \exists \beta_i\\
    \tmop{au}_{U, G} (\ldots ; \beta_i ; \ldots ; \beta_j ; \ldots \text{ as }
    \bar{t}) & = & \bar{\alpha}^{\prime}, g, U^{\prime\prime}, G^{\prime}\\
    \text{where \ } \bar{\alpha}^{\prime}, g, U^{\prime\prime}, G^{\prime} & =
    & \tmop{au}_{U^{\prime}, G} (\bar{t} [\beta_i \assign \beta_j])\\
    U^{\prime} & = & U [\beta_i \assign \beta_j] \wedge \beta_i \dot{=}
    \beta_j\\
    \text{when} &  & \exists \beta_i \in \mathcal{Q}, \beta_j
    \leqslant_{\mathcal{Q}} \beta_i\\
    \tmop{au}_{U, G} (t_1 ; \ldots ; t_n) & = & \alpha, \alpha, U, ([t_1 ;
    \ldots ; t_n] \mapsto \alpha) G\\
    \text{otherwise,} &  & \text{where } \alpha \# \tmop{FV} (t_1 ; \ldots ;
    t_n, U, G)\\
    \tmop{aun}_{U, G} (\varnothing) & = & \varnothing, \varnothing, U, G\\
    \tmop{aun}_{U, G} (\varnothing ; \ldots) & = & \varnothing, \varnothing,
    U, G\\
    \tmop{aun}_{U, G} (t^1_1, \bar{t}^1 ; \ldots ; t^n_1, \bar{t}^n) & = &
    \bar{\alpha} \bar{\alpha}^{\prime}, g \bar{g}^{\prime}, U^{\prime\prime},
    G^{\prime\prime}\\
    \text{where \ } \bar{\alpha}, g, U^{\prime}, G^{\prime} & = &
    \tmop{au}_{U, G} (t_1^1 ; \ldots ; t_1^n)\\
    \text{and \ } \bar{\alpha}^{\prime}, \bar{g}^{\prime}, U^{\prime\prime},
    G^{\prime\prime} & = & \tmop{aun}_{U^{\prime}, G^{\prime}} (\bar{t}^1 ;
    \ldots ; \bar{t}^n)
  \end{eqnarray*}
  \label{AntiUnifAlg}
  \caption{Abductive anti-unification algorithm}
\end{table} The notational shorthand $\ldots ; \beta_i ; \ldots ; f
(\bar{t}^j) ; \ldots$ represents the case where all terms are either
existential variables or start with a function symbol $f$. Similarly, $\ldots
; \beta_i ; \ldots ; \beta_j ; \ldots$ represents the case when there is a
variable $\beta_j$ such that all terms are either $\beta_j$ or are existential
variables to the right of $\beta_j$ in the quantifier.

The sort-integrating algorithm essentially does not change:
\begin{enumerate}
  \item Let $\wedge_s D_{i, s} \equiv \tmmathbf{U} (D_i)$ where $D_{i, s}$ is
  of sort $s$, be the result of our sort-separating unification.
  
  \item For the sort $s_{\tmop{ty}}$:
  \begin{enumerate}
    \item Let $V = \{ x_j, \overline{t_{i, j}} | \forall i \exists t_{i, j}
    .x_j \dot{=} t_{i, j} \in D_{i, s_{\tmop{ty}}} \}$.
    
    \item Let $G = \{ \bar{\alpha}_j, g_j, u_j, \overline{\theta_{i, j}} |
    \theta_{i, j} = [\bar{\alpha}_j \assign \bar{g}_j^i], \theta_{i, j} (g_j)
    = (\wedge_{k \leqslant j} u_k) (t_{i, j}) \}$ be the most specific
    anti-unifiers of $\overline{_{} (\wedge_{k \leqslant j} u_k) (t_{i, j})}$
    for each $j$, where $\wedge_j u_j$ is the resulting $U_{}$.
    
    \item Let $D_i^u = \wedge_j \bar{\alpha}_j \dot{=} \bar{g}_j^i$ and $D_i^g
    = D_{i, s_{\tmop{ty}}} \wedge D_i^u$.
    
    \item Let $D^v_i = \{ x \dot{=} y | x \dot{=} t_1 \in D_i^g, y \dot{=} t_2
    \in D_i^g, D_i^g \vDash t_1 \dot{=} t_2 \}$.
    
    \item Let $A_{s_{\tmop{ty}}} = \wedge_j x_j \dot{=} g_j \wedge \bigcap_i
    (D_i^g \wedge D_i^v)$ (where conjunctions are treated as sets of conjuncts
    and equations are ordered so that only one of $a \dot{=} b, b \dot{=} a$
    appears anywhere), and $\bar{\alpha}_{s_{\tmop{ty}}} =
    \overline{\bar{\alpha}_j}$.
    
    \item Let $\wedge_s D^u_{i, s} \equiv D^u_i$ for $D^u_{i, s}$ of sort $s$.
  \end{enumerate}
  \item For sorts $s \neq s_{\tmop{ty}}$, let $\exists \bar{\alpha}_s .A_s =
  \tmop{DisjElim}_s (\overline{D_i^s \wedge D^u_{i, s}})$.
  
  \item The answer is substitution $U = \wedge_j u_j$ and solved form $\exists
  \overline{\overline{\alpha^j_i}} \overline{\bar{\alpha}_s} . \wedge_s A_s$.
\end{enumerate}
The task of constraint generalization is to find postconditions. Usually, it
is beneficial to make the postcondition, i.e. the existential type that is the
return type of a function, more specific, at the expense of making the overal
type of the function less general. To this effect, constraint generalization,
just as abduction takes invariant parameters $\bar{\beta}$. We replace
conditions $\exists \beta_i \in \mathcal{Q}$ above by $\beta_i \in \bar{\beta}
\vee (\exists \beta_i) \in \mathcal{Q}$. Recall that the right-hand-side (RHS)
variable $\beta_j$ can in general be universally quantified: $\forall \beta_j
\in \mathcal{Q}$. We exclude universal non-parameter RHS when a parameter is
present: if for any $\beta_i$, $\beta_i \in \bar{\beta}$, then for all
$\beta_i$ including RHS, $\beta_i \in \bar{\beta} \vee \exists \beta_i \in
\mathcal{Q}$. Note that having weaker postconditions also results in correct
types, just not the intended ones. In rare cases a weaker postcondition but a
more general invariant can be beneficial. To this effect, the option
\tmverbatim{-{\nobreak}more\_existential} turns off generating the
substitution entries when the RHS is a variable, i.e. the case $\tmop{au}_{U,
G} (\ldots ; \beta_i ; \ldots ; \beta_j ; \ldots \text{ as } \bar{t})$ is
skipped.

Due to greater flexibility of the numerical domain, abductive extension of
numerical constraint generalization does not seem necessary and is turned off
by default. It could take a similar form, we experiment with the following
heuristic. If atoms specific to a disjunct (i.e. not shared by all disjuncts)
do not contain a variable, do not include the disjunct when considering
inclusion of inequalities containing the variable in the constraint
generalization answer.

\section{Incorporating Negative Constraints}\label{NegElimAlg}

Here we expand on the overview from Section \ref{NegElimMain}. We call a
{\tmem{negative constraint}} an implication $D \Rightarrow \tmmathbf{F}$ in
the normalized constraint $\mathcal{Q}. \wedge_i (D_i \Rightarrow C_i)$, and
we call $D$ the {\tmem{negated constraint}}. Such constraints are generated
for pattern matching branches whose right-hand-side is the expression
\tmverbatim{assert false}. A generic approach to account for negative
constraints is as follows. We check whether the solution found by multisort
abduction contradicts the negated constraints. If some negated constraint is
satisfiable, we discard the answer and ``fall back'' to try finding another
answer. Unfortunately, this approach is insufficient when the answer sought
for is not among the maximally general answers to the remaining,
{\tmem{positive part}} of the constraint. Therefore, we introduce
{\tmem{negation elimination}}.

For the numerical sort, our implementation of negation elimination can produce
too strong constraints when the numerical domain is intended to contain
non-integer rational numbers, and can be turned off by the
\tmverbatim{-no\_int\_negation} option. It can also produce too weak
constraints, because it produces at most one atom for a given negated
constraint. If the abduction answer for terms does already contradict a
negated constraint $D$, we are done. Otherwise, let $D = c_1 \wedge \ldots
\wedge c_n$. For numerical sort atoms $c_i$, either drop them from
consideration or convert their negation $\neg c_i$ into $d_i$ or $d_{i_1} \vee
d_{i_2}$ as follows. The conversion assumes that the numerical domain is
integers. We convert an inequality $w \leqslant 0$, e.g. $\frac{1}{3} x -
\frac{1}{2} y - 2 \leqslant 0$, to $- kw + 1 \leqslant 0$, e.g. $3 y - 2 x +
13 \leqslant 0$, where $k$ is the common denominator so that $kw$ has only
integer numbers. Note that $\neg (w \leqslant 0) \Leftrightarrow w > 0
\Leftrightarrow - w < 0 \Leftrightarrow - kw < 0 \Leftarrow - kw + 1 \leqslant
0$. Similarly, we convert an equation $w \dot{=} 0$ to inequalities $- kw + 1
\leqslant 0 \vee kw + 1 \leqslant 0$. Note that $\neg (w \geqslant 0)
\Leftrightarrow w < 0 \Leftrightarrow kw < 0 \Leftarrow kw + 1 \leqslant 0$.
In both cases the implications $\Leftarrow$ would be equivalences if the
numerical domain was integers rather than rational numbers. At present, we
ignore {\tmem{opti}} atoms. The disjunct $d_i$ is a conjunction of
inequalities if $c_i$ is a {\tmem{subopti}} atom.

Assuming that each negative constraint points to a single atomic fact, we try
to find one disjunct $d_i$, resp. $d_{i_1}$ or $d_{i_2}$, corresponding to
$\neg c_i$, discarding those disjuncts that contradict any implication branch.
Specifically, let $\mathcal{Q}. \wedge_i (D_i \Rightarrow C_i)$ be the
constraint we solve, and let $\exists \bar{\alpha} .A$ be a term abduction
answer for $\mathcal{Q}. \wedge_{i : C_i \neq \tmmathbf{F}} (D_i \Rightarrow
C_i)$. We search for $i$ such that for all $k$ with $C_k \neq \tmmathbf{F}$
and $D_k$ satisfiable, $d_i \wedge A \wedge D_k \wedge C_k$ is satisfiable. We
provide a function $\tmop{NegElim} (\neg D, \overline{B_i}) = d_{i_0}$, where
$d_{i_0}$ is the biggest, syntactically, such atom, and $B_i = A \wedge D_i
\wedge C_i$.

Unfortunately, for the sort of terms we do not have well-defined
representation of dis-equalities not using negations. The generic approach of
relying on backtracking to find another abduction answer is more useful for
terms than for the numeric sort. It falls short, however, when negation was
intended to prevent the answer from being too general. Ideally, we would
introduce disequation atoms $\tau \dot{\neq} \tau$ and follow the scheme we
use for the numerical sort. For now, we only cover a very specific use of
negation, to discriminate among type-level ``enumeration''. We limit negation
elimination to considering atoms of the form $\beta \dot{=} \varepsilon_1$,
and contradict them by introducing atoms $\beta \dot{=} \varepsilon_2$, for
types $\varepsilon_1 \neq \varepsilon_2$ without parameters which we call
{\tmem{phantom enumerations}}. The variables $\beta$ are limited to the answer
variables generated in the previous iteration of the main algorithm. The
nullary datatype constructor $\varepsilon_2$ is picked so that the atom is
valid, using the same validation procedure as the one passed to the abduction
algorithm. The heuristic defines phantom enumerations as nullary phantom types
that do not share datype parameter position (in GADT constructor definitions)
with non-enumeration types. When the equations derived for different negated
constraints involve a common variable as the left-hand-side, we select a
common right-hand-side. In the end, we only introduce the negation elimination
result to the answer when a single disjunct remains.

Since the {\tmem{discard}} (or taboo) list used by backtracking is based on
complete answers, it is preferable to perform negation elimination prior to
abduction. Otherwise, the search might fall into a loop where abduction keeps
returning the same answer, since it is more general than the discarded answer
incorporating negation elimination result.

\section{{\tmem{opti}} and {\tmem{subopti}}: {\tmem{minimum}} and
{\tmem{maximum}} Relations in \tmverbatim{num}}\label{OptiAtoms}

We extend the numerical domain with relations {\tmem{opti}} and
{\tmem{subopti}} defined below. Operations {\tmem{min}} and {\tmem{max}} can
then be defined using it. Let $k, v, w$ be any linear combinations. Note that
the relations are introduced to smuggle in a limited form of disjunction into
the solved forms.
\begin{eqnarray*}
  \tmop{opti} (v, w) & = & v \leqslant 0 \wedge w \leqslant 0 \wedge (v
  \dot{=} 0 \vee w \dot{=} 0)\\
  k \dot{=} \min (v, w) & = & \tmop{opti} (k - v, k - w)\\
  k \dot{=} \max (v, w) & = & \tmop{opti} (v - k, w - k)\\
  \tmop{subopti} (v, w) & = & v \leqslant 0 \vee w \leqslant 0\\
  k \leqslant \max (v, w) & = & \tmop{subopti} (k - v, k - w)\\
  \min (v, w) \leqslant k & = & \tmop{subopti} (v - k, w - k)
\end{eqnarray*}
In particular, $\tmop{opti} (v, w) \equiv \max (v, w) \dot{=} 0$ and
$\tmop{subopti} (v, w) \equiv \min (v, w) \leqslant 0$. We call an
{\tmem{opti}} or {\tmem{subopti}} atom {\tmem{directed}} when there is a
variable $n$ that appears in $v$ and $w$ with the same sign. We do not
prohibit undirected {\tmem{opti}} or {\tmem{subopti}} atoms, but we do not
introduce them, to avoid bloat.

For simplicity, we do not support {\tmem{min}} and {\tmem{max}} as subterms in
concrete syntax. Instead, we parse atoms of the form
\tmverbatim{k=min(}$\ldots$\tmverbatim{,}$\ldots$\tmverbatim{)}, resp.
\tmverbatim{k=max(}$\ldots$\tmverbatim{,}$\ldots$\tmverbatim{)} into the
corresponding {\tmem{opti}} atoms, where \tmverbatim{k} is any numerical term.
Similarly for {\tmem{subopti}}. We also print directed {\tmem{opti}} and
{\tmem{subopti}} atoms using the syntax with $\min$ and $\max$expressions. Not
to pollute the syntax with a new keyword, we use concrete syntax
\tmverbatim{min\textbar max(}$\ldots$\tmverbatim{,}$\ldots$\tmverbatim{)} for
parsing arbitrary, and printing non-directed, {\tmem{opti}} atoms, and
\tmverbatim{min\textbar\textbar
max(}$\ldots$\tmverbatim{,}$\ldots$\tmverbatim{)} for {\tmem{subopti}}
respectively.

If need arises, in a future version, we can extend {\tmem{opti}} to a larger
arity $N$.

\subsection{Normalization, Validity and Implication Checking}

In the function that produces solved forms of numerical constraints, we treat
{\tmem{opti}} clauses in an efficient but incomplete manner, doing a single
step of constraint solving. We include the {\tmem{opti}} terms in processed
inequalities. After equations have been solved, we apply the substitution to
the {\tmem{opti}} and {\tmem{subopti}} disjunctions. When one of the
{\tmem{opti}} resp. {\tmem{subopti}} disjunct terms becomes contradictory or
the disjunct terms become equal, we include the other in implicit equalities,
resp. in inequalities to solve. When one of the {\tmem{opti}} or
{\tmem{subopti}} terms becomes tautological, we drop the disjunction. We
iterate calls to the solver function to propagate implicit equalities.

We do not perform case splitting on {\tmem{opti}} and {\tmem{subopti}}
disjunctions, therefore some contradictions may be undetected. However,
abduction and constraint generalization currently perform upfront case
splitting on {\tmem{opti}} and {\tmem{subopti}} disjunctions, sometimes
leading to splits that a smarter solver would avoid.

\subsection{Abduction}

We eliminate {\tmem{opti}} and {\tmem{subopti}} in premises by expanding the
definition and converting the branch into two branches, e.g. $D \wedge (v
\dot{=} 0 \vee w \dot{=} 0) \Rightarrow C$ into $(D \wedge v \dot{=} 0
\Rightarrow C) \wedge (D \wedge w \dot{=} 0 \Rightarrow C)$. Recall that an
{\tmem{opti}} atom also implies inequalities $v \leqslant \wedge w \leqslant
0$ assumed to be in $D$ above. This is one form of {\tmem{case splitting}}: we
consider cases $v \dot{=} 0$ and $w \dot{=} 0$, resp. $v \leqslant 0$ and $w
\leqslant 0$, separately. We do not eliminate {\tmem{opti}} and
{\tmem{subopti}} in conclusions. Rather, we consider whether to keep or drop
it in the answer, like with other candidate atoms. The transformations apply
to an {\tmem{opti}} atom by applying to both its arguments.

Generating a new {\tmem{opti}} atom for inclusion in an answer means finding a
pair of equations such that the following conditions hold. Each equation,
together with remaining atoms of an answer but without the remaining equation
selected, is a correct answer to a simple abduction problem. The equations
selected share a variable and are oriented so that the variable appears with
the same sign in them. The resulting {\tmem{opti}} atom passes the validation
test for joint constraint abduction. We may implement generating new
{\tmem{opti}} atoms for abduction answers in a future version, when need
arises. Currently, we only generate new {\tmem{opti}} and {\tmem{subopti}}
atoms for postconditions, i.e. during constraint generalization.

\subsection{Constraint Generalization}

We eliminate {\tmem{opti}} and {\tmem{subopti}} atoms prior to finding the
extended convex hull of $\overline{D_i}$ by expanding the definition and
converting the disjunction $\vee_i D_i$ to disjunctive normal form. This is
another form of case splitting.

In addition to finding the extended convex hull, we need to discover
{\tmem{opti}} relations that are implied by $\vee_i D_i$. We select these
faces of the convex hull which also appear as an equation in some disjuncts.
Out of these faces, we find all minimal covers of size 2, i.e. pairs of faces
such that in each disjunct, either one or the other linear combination appears
as an equation. We only keep pairs of faces that share a same-sign variable.
For the purposes of detecting {\tmem{opti}} relations, we need to perform
transitive closure of the extended convex hull equations and inequalities,
because the redundant inequalities might be required to find a cover.

Finding {\tmem{subopti}} atoms is similar. We find all minimal covers of size
2, i.e. pairs of inequalities such that one or the other appears in each
disjunct. We only keep pairs of inequalities that share a same-sign variable.

We provide a function for the numerical domain to remove {\tmem{opti}} atoms
of the form $k \dot{=} \min (c, v)$, $k \dot{=} \min (v, c)$, $k \dot{=} \max
(c, v)$ or $k \dot{=} \min (v, c)$ for a constant $c$, similarly for
{\tmem{subopti}} atoms, while in initial iterations where constraint
generalization is only performed for non-recursive branches.

We need to further extend the notion of (abductive) constraint generalization,
to achieve the results required for postcondition inference. For constraint
branches $\overline{D_i \Rightarrow C_i}$, we need not only the disjuncts
$\overline{D_i \wedge C_i}$, but also the premises $\overline{D_i}$. We keep
{\tmem{opti}} and {\tmem{subopti}} atoms $\tmop{opti} (v, w), \tmop{subopti}
(v, w)$ such that either both $v \leqslant w$ and $w \leqslant v$ are
{\tmem{satisfiable with}} all {\tmem{implication branches}}, or neither is. An
atom $c$ is satisfiable with implication $D_i \Rightarrow C_i$ here, when
either $c \wedge D_i$ is not satisfiable, or $c \wedge D_i \wedge C_i$ is
satisfiable. The underlying idea is that since {\tmem{opti}} and
{\tmem{subopti}} atoms express a disjunction, resp. $v \leqslant 0 \wedge w
\leqslant 0 \wedge (v \dot{=} 0 \vee w \dot{=} 0)$ and $v \leqslant 0 \vee w
\leqslant 0$, they are meaningful when both cases of the disjunction can
obtain under some circumstances. When only one of the atoms, say $v \leqslant
w$, is satisfiable with all implication branches, we make an abductive guess
that $v \leqslant w$.

\section{Solving for Predicate Variables}\label{PredVarAlg}

Here we discuss the implementation in {\tmname{InvarGenT}} of ideas in the
overview from Section \ref{PredVarOverview} and in the formal discussion from
Section \ref{PredVarMain}. As we decided to provide the first solution to
abduction problems, we accordingly simplify the task of solving for predicate
variables. Instead of a tree of solutions being refined, we have a single
sequence which we unfold until reaching fixpoint or contradiction. Another
choice point besides abduction in the original algorithm is the selection of
invariants that leaves a consistent subset of atoms as residuum. Here we also
select the first solution found. We introduce a form of backtracking,
described in section \ref{Details}.

\subsection{Solving for Predicates in Premises}

The core of our inference algorithm consists of distributing atoms of an
abduction answer among the predicate variables. The negative occurrences of
predicate variables, when instantiated with the updated solution, enrich the
premises so that the next round of abduction leads to a smaller answer (in
number of atoms).

Let us discuss the algorithm for $\tmop{Split} (\mathcal{Q}, \bar{\alpha}, A,
\overline{\bar{\beta}^{\chi}}, \overline{A_{\chi}^0})$. Note that due to
existential types predicates, we actually compute $\tmop{Split} \left(
\mathcal{Q}, \bar{\alpha}, A, \overline{\bar{\beta}^{\beta_{\chi}}},
\overline{A_{\beta_{\chi}}^0} \right)$, i.e. we index by $\beta_{\chi}$ (which
can be multiple for a single $\chi$) rather than $\chi$. We retain the
notation indexing by $\chi$ as it better conveys the intent. We do not pass
quantifiers around to reflect the source code: the helper function
\tmverbatim{loop avs ans sol} of function \tmverbatim{split} corresponds to
$\tmop{Split} \left( \bar{\alpha}, A, \overline{A_{\beta_{\chi}}^0} \right)$.
\begin{eqnarray*}
  \alpha \prec \beta & \equiv & \alpha <_{\mathcal{Q}} \beta \vee (\alpha
  \leqslant_{\mathcal{Q}} \beta \wedge \beta \nless_{\mathcal{Q}} \alpha
  \wedge \alpha \in \overline{\bar{\beta}^{\chi}} \wedge \beta \nin
  \overline{\bar{\beta}^{\chi}})\\
  A_{\alpha \beta} & = & \{ \beta \dot{=} \alpha \in A | \beta \in
  \overline{\bar{\beta}^{\chi}} \wedge (\exists \alpha) \in \mathcal{Q} \wedge
  \beta \prec \alpha \}\\
  A_0 & = & A\backslash A_{\alpha \beta}\\
  A^1_{\chi} & = & \left\{ c \in A_0 \middle| \forall \alpha \in \tmop{FV} (c)
  . (\exists \alpha) \in \mathcal{Q} \vee \\
  \text{ \ \ \ } \alpha <_{\mathcal{Q}} \beta_{\chi} \wedge \alpha \nin
  \tmop{PrimCV} (c) \vee \alpha \in \overline{\bar{\beta}^{\chi}} \wedge
  \alpha \in \tmop{PrimCV} (c) \right\}\\
  A^2_{\chi} & = & \tmop{Atomized} (\overline{\bar{\beta}^{\chi}},
  A_{\chi}^1)\\
  A^3_{\chi} & = & A^2_{\chi} \backslash \cup_{\chi^{\prime}}
  A^1_{\chi^{\prime}}\\
  \text{if} &  & \mathcal{M} \nvDash \mathcal{Q}. (A \setminus \cup_{\chi}
  A^2_{\chi}) [\bar{\alpha} \assign \bar{t}] \text{ \ for all } \bar{t}\\
  \text{then return} &  & \bot\\
  \text{for all } \overline{A_{\chi}^+} \text{ min. w.r.t. } \subset \text{
  s.t.} &  & \wedge_{\chi} (A_{\chi}^+ \subset A^2_{\chi}) \wedge \mathcal{M}
  \vDash \mathcal{Q}. (\cup_{\chi} A_{\chi}^+ \Rightarrow A) [\bar{\alpha}
  \assign \bar{t}] \text{ \ for some } \bar{t} :\\
  \text{if} &  & \tmop{Strat} (A^+_{\chi}, \bar{\beta}^{\chi}) \text{ \
  returns } \bot \text{ for some } \chi\\
  \text{then return} &  & \bot\\
  \text{else \ } \bar{\alpha}_+^{\chi}, A_{\chi}^L, A^R_{\chi} & = &
  \tmop{Strat} (A^+_{\chi}, \bar{\beta}^{\chi})\\
  A_{\chi} & = & A_{\chi}^0 \cup A_{\chi}^L\\
  \bar{\alpha}^{\chi}_0 & = & \bar{\alpha} \cap \tmop{FV} (A_{\chi})\\
  \bar{\alpha}^{\chi} & = & \left( \bar{\alpha}^{\chi}_0 \setminus
  \bigcup_{\chi^{\prime} <_{\mathcal{Q}} \chi} \bar{\alpha}^{\chi^{\prime}}_0
  \right) \bar{\alpha}_+^{\chi}\\
  A_+ & = & \cup_{\chi} A_{\chi}^R\\
  A_{\tmop{res}} & = & A_+ \cup \widetilde{A_+} (A \setminus \cup_{\chi}
  A_{\chi}^+)\\
  \text{if} &  & \cup_{\chi} \bar{\alpha}^{\chi} \neq \varnothing \vee
  \cup_{\chi} A^3_{\chi} \neq \varnothing \text{ \ then}\\
  \mathcal{Q}^{\prime}, A_{\tmop{res}}^{\prime}, \overline{\exists
  \bar{\alpha}^{\prime \chi} .A_{\chi}^{\prime}} & \in & \tmop{Split} \left(
  \mathcal{Q} [\overline{\forall \bar{\beta}^{\chi}} \assign \overline{\forall
  (\bar{\beta}^{\chi} \cup \bar{\alpha}^{\chi})}], \bar{\alpha} \setminus
  \cup_{\chi} \bar{\alpha}^{\chi}, A_{\tmop{res}} \wedge_{\chi} A^3_{\chi}, \\
  \text{ \ \ \ \ \ \ \ } \overline{\bar{\beta}^{\chi} \cup
  \bar{\alpha}^{\chi}}, \overline{A_{\chi}} \right)\\
  \text{return} &  & \mathcal{Q}^{\prime}, A_{\alpha \beta} \wedge
  A_{\tmop{res}}^{\prime}, \overline{\exists \bar{\alpha}^{\chi}
  \bar{\alpha}^{\prime \chi} .A_{\chi}^{\prime}}\\
  \text{else return} &  & \mathcal{Q} \exists (\bar{\alpha} \setminus
  \cup_{\chi} \bar{\alpha}^{\chi}), A_{\alpha \beta} \wedge A_{\tmop{res}},
  \overline{\exists \bar{\alpha}^{\chi} .A_{\chi}}
\end{eqnarray*}
where $\tmop{Strat} (A, \bar{\beta}^{\chi})$ is computed as follows: for every
$c \in A$, and for every $\beta_2 \in \tmop{FV} (c)$ such that $\beta_1
<_{\mathcal{Q}} \beta_2$ for $\beta_1 \in \bar{\beta}^{\chi}$, if $\beta_2$ is
universally quantified in $\mathcal{Q}$ and $\beta_2 \nin
\overline{\bar{\beta}^{\chi}}$, then return $\bot$; otherwise, introduce a
fresh variable $\alpha_f$, replace $c \assign c [\beta_2 \assign \alpha_f]$,
add $\beta_2 \dot{=} \alpha_f$ to $A^R_{\chi}$ and $\alpha_f$ to
$\bar{\alpha}_+^{\chi}$, after replacing all such $\beta_2$ add the resulting
$c$ to $A_{\chi}^L$. $\tmop{PrimCV} (c)$ stands for {\tmem{primary constrained
variables}} of an atom $c$. These are defined as the substituted variables in
solved forms for term constraints, and in the case of a numerical atom, the
shared variable of a directed opti or subopti atom, i.e. the variable $v$ in
$v \leqslant \max (\ldots)$, $\min (\ldots) \leqslant v$, $v = \max (\ldots)$,
$v = \min (\ldots)$.

Description of the algorithm in more detail:
\begin{enumerate}
  \item $\begin{array}{lll}
    \alpha \prec \beta & \equiv & \alpha <_{\mathcal{Q}} \beta \vee (\alpha
    \leqslant_{\mathcal{Q}} \beta \wedge \beta \nless_{\mathcal{Q}} \alpha
    \wedge \alpha \in \overline{\bar{\beta}^{\chi}} \wedge \beta \nin
    \overline{\bar{\beta}^{\chi}})
  \end{array}$ The variables $\overline{\bar{\beta}^{\chi}}$ are the invariant
  parameters of the solution from the previous round. We need to keep them
  apart from other variables even when they're not separated by quantifier
  alternation.
  
  \item $\begin{array}{lll}
    A_0 & = & A\backslash A_{\alpha \beta}
  \end{array}$ where$\begin{array}{lll}
    A_{\alpha \beta} & = & \{ \beta \dot{=} \alpha \in A | \beta \in
    \overline{\bar{\beta}^{\chi}} \wedge (\exists \alpha) \in \mathcal{Q}
    \wedge \beta \prec \alpha \}
  \end{array}$ Discard the ``scaffolding'' information, in particular the
  $A_+$ equations introduced by $\tmop{Strat}$ in an earlier iteration.
  
  \item $\begin{array}{lll}
    A^1_{\chi} & = & \{ c \in A_0 | \forall \alpha \in \tmop{FV} (c)
    \backslash \bar{\beta}^{\chi} . (\exists \alpha) \in \mathcal{Q} \vee
    \alpha <_{\mathcal{Q}} \beta_{\chi} \}
  \end{array}$ Gather atoms pertaining to predicate $\chi$, which should not
  remain in the residuum.
  
  \item $\begin{array}{lll}
    A^2_{\chi} & = & \tmop{Atomized} (\overline{\bar{\beta}^{\chi}},
    A_{\chi}^1)
  \end{array}$ An atomized form of $A_{\chi}^1$ wrt.
  $\overline{\bar{\beta}^{\chi}}$: a conjunction of atoms equivalent to
  $A_{\chi}^1$ containing some of the implications of $A_{\chi}^1$, see the
  formal exposition of InvarGenT. Actually, rather than computing the atomized
  form upfront, we generate its contribution to $A_{\chi}^+$ while performing
  Fourier-Motzkin elimination to check $\mathcal{M} \vDash \mathcal{Q}. \left(
  \cup_{\chi} A_{\chi}^+ \Rightarrow {\nobreak} A \right)$.
  
  \item $\begin{array}{lll}
    A^3_{\chi} & = & A^2_{\chi} \backslash \cup_{\chi^{\prime}}
    A^1_{\chi^{\prime}}
  \end{array}$ We prune atoms coming from atomization that are already present
  in the invariants. Actually, in the implementation we prune by redundancy
  with the invariants assigned so far.
  
  \item $\begin{array}{l}
    \text{if } \mathcal{M} \nvDash \mathcal{Q}. (A \setminus \cup_{\chi}
    A^1_{\chi}) [\bar{\alpha} \assign \bar{t}] \text{ \ for all } \bar{t}
    \text{ then return } \bot
  \end{array}$: Failed solution attempt. A common example is when the use site
  of recursive definition, resp. the existential type introduction site, is
  not in scope of a defining site of recursive definition, resp. an
  existential type elimination site, and has too strong requirements. FIXME:
  
  \item $\text{for all } \overline{A_{\chi}^+} \text{ min. w.r.t. } \subset
  \text{ s.t.} \wedge_{\chi} (A_{\chi}^+ \subset A_{\chi}^2) \wedge
  \mathcal{M} \vDash \mathcal{Q}. (\cup_{\chi} A_{\chi}^+ \Rightarrow A)
  [\bar{\alpha} \assign \bar{t}] \text{ \ for some } \bar{t} :$ Select
  invariants such that the residuum $A \setminus \cup_{\chi} A_{\chi}^+$ is
  consistent. The final residuum $A_{\tmop{res}}$ represents the global
  constraints, the solution for global type variables. The solutions
  $A_{\chi}^+$ represent the invariants, the solution for invariant type
  parameters.
  
  \item $\begin{array}{lll}
    \text{if} &  & \tmop{Strat} (A^+_{\chi}, \bar{\beta}^{\chi}) \text{ \
    returns } \bot \text{ for some } \chi
  \end{array} \text{then return } \bot$ In the implementation, we address
  stratification issues already during abduction.
  
  \item $\bar{\alpha}_+^{\chi}, A_{\chi}^L, A^R_{\chi} = \tmop{Strat}
  (A^+_{\chi}, \bar{\beta}^{\chi})$ is computed as follows: for every $c \in
  A^+_{\chi}$, and for every $\beta_2 \in \tmop{FV} (c)$ such that $\beta_1
  <_{\mathcal{Q}} \beta_2$ for $\beta_1 \in \bar{\beta}^{\chi}$, if $\beta_2$
  is universally quantified in $\mathcal{Q}$, then return $\bot$; otherwise,
  introduce a fresh variable $\alpha_f$, replace $c \assign c [\beta_2 \assign
  \alpha_f]$, add $\beta_2 \dot{=} \alpha_f$ to $A^R_{\chi}$ and $\alpha_f$ to
  $\bar{\alpha}_+^{\chi}$, after replacing all such $\beta_2$ add the
  resulting $c$ to $A_{\chi}^L$.
  \begin{itemize}
    \item We will add $\bar{\alpha}_+^{\chi}$ to $\bar{\beta}^{\chi}$.
  \end{itemize}
  \item $\begin{array}{lll}
    A_{\chi} & = & A_{\chi}^0 \cup A_{\chi}^L
  \end{array}$ is the updated solution formula for $\chi$, where $A_{\chi}^0$
  is the solution from previous round.
  
  \item $\begin{array}{lll}
    \bar{\alpha}^{\chi}_0 & = & \bar{\alpha} \cap \tmop{FV} (A_{\chi})
  \end{array}$ are the additional solution parameters coming from variables
  generated by abduction.
  
  \item $\begin{array}{lll}
    \bar{\alpha}^{\chi} & = & \left( \bar{\alpha}^{\chi}_0 \setminus
    \bigcup_{\chi^{\prime} <_{\mathcal{Q}} \chi}
    \bar{\alpha}^{\chi^{\prime}}_0 \right) \bar{\alpha}_+^{\chi}
  \end{array}$ The final solution parameters also include the variables
  generated by $\tmop{Strat}$.
  
  \item $\begin{array}{lll}
    A_+ & = & \cup_{\chi} A_{\chi}^R
  \end{array}$ and $\begin{array}{lll}
    A_{\tmop{res}} & = & A_+ \cup \widetilde{A_+} (A \setminus \cup_{\chi}
    A_{\chi}^+)
  \end{array}$ is the resulting global constraint, where $\widetilde{A_+}$ is
  the substitution corresponding to $A_+$.
  
  \item $\begin{array}{lll}
    \text{if} &  & \cup_{\chi} \bar{\alpha}^{\chi} \neq \varnothing \vee
    \cup_{\chi} A^3_{\chi} \neq \varnothing \text{ \ then}
  \end{array}$ -- If new parameters or new atoms have been introduced, we need
  to redistribute the remaining atoms to make $\mathcal{Q}^{\prime}
  .A_{\tmop{res}}$ valid again.
  
  \item $\begin{array}{lll}
    \mathcal{Q}^{\prime}, A_{\tmop{res}}^{\prime}, \overline{\exists
    \bar{\alpha}^{\prime \chi} .A_{\chi}^{\prime}} & \in & \tmop{Split}
    (\mathcal{Q} [\overline{\forall \bar{\beta}^{\chi}} \assign
    \overline{\forall (\bar{\beta}^{\chi} \cup \bar{\alpha}^{\chi})}],
    \bar{\alpha} \setminus \cup_{\chi} \bar{\alpha}^{\chi}, A_{\tmop{res}}
    \wedge_{\chi} A^3, \overline{\bar{\beta}^{\chi} \cup \bar{\alpha}^{\chi}},
    \overline{A_{\chi}})
  \end{array}$ Recursive call includes $\bar{\alpha}^{\chi}_+$ in
  $\bar{\beta}^{\chi}$ so that, among other things, $A_+$ are redistributed
  into $A_{\chi}$.
  
  \item $\begin{array}{lll}
    \text{return} &  & \mathcal{Q}^{\prime}, A_{\tmop{res}}^{\prime},
    \overline{\exists \bar{\alpha}^{\chi} \bar{\alpha}^{\prime \chi}
    .A_{\chi}^{\prime}}
  \end{array}$ We do not add $\forall \bar{\alpha}$ in front of
  $\mathcal{Q}^{\prime}$ because it already includes these variables.
  $\overline{\bar{\alpha}_+^{\chi} \overline{\alpha_{}}^{\chi}_+^{\prime}}$
  lists all variables introduced by $\tmop{Strat}$.
  
  \item $\begin{array}{lll}
    \text{else return} &  & \mathcal{Q} \exists (\bar{\alpha} \setminus
    \cup_{\chi} \bar{\alpha}^{\chi}), A_{\tmop{res}}, \overline{\exists
    \bar{\alpha}^{\chi} .A_{\chi}}
  \end{array}$ Note that $\bar{\alpha} \setminus \cup_{\chi}
  \bar{\alpha}^{\chi}$ does not contain the current $\bar{\beta}^{\chi}$,
  because $\bar{\alpha}$ does not contain it initially and the recursive call
  maintains that: $\bar{\alpha} \assign \bar{\alpha} \setminus \cup_{\chi}
  \bar{\alpha}^{\chi}, \bar{\beta}^{\chi} \assign \bar{\beta}^{\chi}
  \bar{\alpha}^{\chi}$.
\end{enumerate}
Finally we define $\tmop{Split} (\bar{\alpha}, A) \assign \tmop{Split}
(\bar{\alpha}, A, \bar{\tmmathbf{T}})$. The complete algorithm for solving
predicate variables is presented in the next section.

\subsection{Solving for Existential Types Predicates and Main
Algorithm}\label{MainAlgoBody}

The general scheme is that we perform constraint generalization on branches
with positive occurrences of existential type predicate variables on each
round. Since the branches are substituted with the solution from previous
round, constraint generalization will automatically preserve monotonicity. We
retain existential types predicate variables in the later abduction step.

What differentiates existential type predicate variables from recursive
definition predicate variables is that the former are not local to context,
while the latter are treated as local to context. It means that free variables
need to be bound in the former while they are retained in the latter. However,
it is just a technical difference because

In the algorithm we operate on all predicate variables, and perform additional
operations on existential type predicate variables; there are no operations
pertaining only to recursive definition predicate variables. The equation
numbers $(N)$ below are matched by comment numbers \tmverbatim{(* N *)} in the
source code. Let $\chi_{\varepsilon K}$ be the invariant which will constrain
the existential type $\varepsilon_K$, or $\top$. When $\chi_{\varepsilon K} =
\top$, then we define $\bar{\beta}^{\chi_{\varepsilon K}, k} = \varnothing$.
Step $k$ of the loop of the final algorithm:
\begin{eqnarray}
  \overline{\exists \bar{\beta}^{\chi, k} .F_{\chi}} & = & S_k \nonumber\\
  &  & \text{In iteration 2, remove non-term-sort atoms} \nonumber\\
  &  & \text{containing outer-scope parameters from } S_k . \nonumber\\
  D_K^{\alpha} \Rightarrow C_K^{\alpha} \in R_k^- S_k (\Phi) & = & \text{all
  such that } \chi_K (\alpha, \alpha_{\alpha}^K) \in C_K^{\alpha}, \\
  &  & \overline{C^{\alpha}_j} = \{ C | D \Rightarrow C \in S_k (\Phi) \wedge
  D \subseteq D^{\alpha}_K \} \nonumber\\
  \alpha_K & : & \chi_K (\alpha_K) \text{ is a positive atom in } \Phi
  \nonumber\\
  U_{\chi_K}, \exists \bar{\alpha}^{\chi_K}_g .G_{\chi_K}, \overline{B_d^K} &
  = & \tmop{DisjElim} \left( \alpha_K, \overline{\delta \dot{=} \alpha \wedge
  D^{\alpha}_K \wedge_j C^{\alpha}_j}_{\alpha \in \overline{\alpha_3^{i, K}}}
  \right) \\
  \vec{\tau}_{\varepsilon_K} & = & \overrightarrow{\{ \alpha \in \tmop{FV}
  (G_{\chi_K}) \backslash \delta \delta^{\prime} \bar{\alpha}^{\chi_K}_g |
  (\exists \alpha) \in \mathcal{Q} \vee \alpha \in
  \overline{\bar{\beta}^{\chi}} \backslash \bar{\beta}^{\chi_K} \}}
  \nonumber\\
  \Xi (\exists \bar{\alpha}^{\chi_K}_g .G_{\chi_K}) & = & \exists \tmop{FV}
  (\vec{\tau}_{\varepsilon_K}, G_{\chi_K}) \backslash \delta \delta^{\prime} .
  \delta^{\prime} \dot{=} \vec{\tau}_{\varepsilon_K} \wedge G_{\chi_K}
  \nonumber\\
  (\exists \bar{\alpha}^{\chi_K} .F_{\chi_K}), \rho & = & H (R_k (\chi_K), \Xi
  (\exists \bar{\alpha}_g^{\chi_K} .G_{\chi_K}))  \label{Rg}\\
  R_g (\chi_K) & = & \exists \bar{\alpha}^{\chi_K} .F_{\chi_K} \nonumber\\
  P_g (\chi_K (\delta, \delta^{\prime})) = P_g (\chi_K (\delta^{\prime})) & =
  & \delta^{\prime} \dot{=} \vec{\tau}_{\varepsilon_K} \nonumber\\
  F_{\chi}^{\prime} & = & F_{\chi} [\varepsilon_K (\vec{\tau}_{\tmop{old}})
  \assign \varepsilon_K (\vec{\tau}_{\varepsilon_K}) [\tmop{recover}
  (\vec{\tau}_{\varepsilon_K}, \vec{\tau}_{\tmop{old}})]] \nonumber\\
  S_k^{\prime} & = & \overline{\exists \bar{\beta}^{\chi, k} \{ \alpha \in
  \tmop{FV} (F^{\prime}_{\chi}) | \beta_{\chi} <_{\mathcal{Q}} \alpha \}
  .F^{\prime}_{\chi}^{}}  \label{CaptureEscParams}\\
  \mathcal{Q}^{\prime} . \wedge_i (D_i \Rightarrow C_i) \wedge_j (D_j^-
  \Rightarrow \tmmathbf{F}) & = & R_g^- P_g^+ S_k^{\prime} (\Phi
  \wedge_{\chi_K} U_{\chi_K}) \nonumber\\
  \exists \bar{\alpha} .A_0 & = & \tmop{Abd} \left( \mathcal{Q}^{\prime},
  \bar{\beta} = \overline{\beta_{\chi} \overline{\beta^{}}^{\chi}},
  \overline{D_i, C_i} \right) \\
  A & = & A_0 \wedge_j \tmop{NegElim} (\neg \tmop{Simpl} (\tmop{FV} (D_j^-)
  \backslash \overline{\bar{\beta}^{\chi}} .D_j^-), A_0, \overline{D_i, C_i})
  \nonumber\\
  &  & \text{In later iterations, check negative constraints.} \\
  \left( \mathcal{Q}^{k + 1}, A_{\tmop{res}}, \overline{\exists
  \bar{\alpha}^{\beta_{\chi}} .A_{\beta_{\chi}}} \right) & = & \tmop{Split}
  \left( \mathcal{Q}^{\prime}, \bar{\alpha}, A, \overline{\beta_{\chi}
  \overline{\beta^{}}^{\chi}} \right) \nonumber\\
  \vec{\tau}_{\varepsilon_K}^{\prime} & = & \overrightarrow{\tmop{FV}
  (\widetilde{A_{\tmop{res}}} (\vec{\tau}_{\varepsilon_K}))} \nonumber\\
  R_{k + 1} (\chi_K) & = & \exists \bar{\beta}^{\chi_K, k}
  \overline{\bar{\alpha}^{\beta_{\chi_K}}} \bar{\alpha}^{\chi_K} \backslash
  \tmop{FV} (\vec{\tau}_{\varepsilon_K}^{\prime}) . \delta^{\prime} \dot{=}
  \vec{\tau}_{\varepsilon_K}^{\prime} \wedge \widetilde{A_{\tmop{res}}}
  (F_{\chi_K} \backslash \delta^{\prime} \dot{=} \ldots) \\
  \text{ \ \ \ } \wedge_{\chi_K} A_{\beta_{\chi_K}} \left[
  \overline{\beta_{\chi_K} \overline{\beta^{}}^{\beta_{\chi_K}}} \assign
  \overline{\delta \overline{\beta^{}}^{\chi_K, k}} \right] \\
  A_K^d & = & \left\{ c \in A_{\beta_{\chi_K}} \left[ \overline{\beta_{\chi_K}
  \overline{\beta^{}}^{\beta_{\chi_K}}} \assign \overline{\delta
  \overline{\beta^{}}^{\chi_K, k}} \right] \middle| \\
  \text{ \ \ \ \ \ } \tmop{FV} (c) \subseteq \tmop{FV} (\rho (B_d^K))
  \right\} \nonumber\\
  \exists \bar{\alpha}^{\prime}_K .A^{\prime}_K & = & \tmop{Abd} \left(
  \mathcal{Q}^{\prime}, \overline{\beta_{\chi} \overline{\beta^{}}^{\chi}}
  \backslash \overline{\beta_{\chi_K} \overline{\beta^{}}^{\chi_K}},
  \overline{\rho (B_d^K), A_K^d} \right) \\
  \left( \mathcal{Q}^{k + 1}, \varnothing, \overline{\exists
  \bar{\alpha}^{\beta_{\chi}\prime} .A_{\beta_{\chi}}^{\prime}} \right) & = &
  \tmop{Split} \left( \mathcal{Q}^{k + 1}, \overline{\bar{\alpha}^{\prime}_K},
  \wedge_K A^{\prime}_K, \overline{\beta_{\chi} \overline{\beta^{}}^{\chi}}
  \backslash \overline{\beta_{\chi_K} \overline{\beta^{}}^{\chi_K}} \right)
  \nonumber\\
  S_{k + 1} (\chi) & = & \exists \bar{\beta}^{\chi, k} . \tmop{Simpl} \left(
  \exists \overline{\bar{\alpha}^{\beta_{\chi}}} .F_{\chi}^{\prime} \\
  \text{ \ \ \ \ \ \ \ \ } \wedge A_{\beta_{\chi}} \left[
  \overline{\beta_{\chi} \overline{\beta^{}}^{\chi}} \assign \overline{\delta
  \overline{\beta^{}}^{\chi, k}} \right] \wedge A_{\beta_{\chi}}^{\prime}
  \right)  \label{Skp1}\\
  \text{if} &  & (\forall \chi) S_{k + 1} (\chi) \subseteq S_k (\chi), \\
  &  & (\forall \chi_K) R_{k + 1} (\chi_K) = R_k (\chi_K), \nonumber\\
  &  & (\forall \beta_{\chi_K}) A_{\beta_{\chi_K}} =\tmmathbf{T}, \nonumber\\
  &  & k > 1 \nonumber\\
  \text{then return} &  & A_{\tmop{res}}, S_{k + 1}, R_{k + 1} \nonumber\\
  \text{repeat} &  & k \assign k + 1 \nonumber
\end{eqnarray}
Note that $\tmop{Split}$ returns $\overline{\exists
\bar{\alpha}^{\beta_{\chi}} .A_{\beta_{\chi}}}$ rather than $\overline{\exists
\bar{\alpha}^{\chi} .A_{\chi}}$. This is because in case of existential type
predicate variables $\chi_K$, there can be multiple negative position
occurrences $\chi_K (\beta_{\chi_K}, \cdummy)$ with different $\beta_{\chi_K}$
when the corresponding value is used in multiple $\text{{\tmstrong{let}}}
\ldots \text{{\tmstrong{in}}}$ expressions. The variant of the algorithm to
achieve completeness would compute all answers for variants of $\tmop{Abd}$
and $\tmop{Split}$ algorithms that return multiple answers. Unary predicate
variables $\chi (\beta_{\chi})$ can also have multiple negative occurrences in
the normalized form, but always with the same argument $\beta_{\chi}$. The
substitution $\left[ \overline{\beta_{\chi_K}
\overline{\beta^{}}^{\beta_{\chi_K}}} \assign \overline{\delta
\overline{\beta^{}}^{\chi_K, k}} \right]$ replaces the instance parameters
introduced in $\chi_K (\beta_{\chi_K}, \cdummy)$ by the formal parameters used
in $R_k (\chi_K)$.

Even for a fixed $K$, the same branch $D_K^{\alpha} \Rightarrow C_K^{\alpha}$
can contribute to multiple disjuncts, with different $\alpha = \alpha_3^{i,
K}$. Substitution $R_g^-$ substitutes only negative occurrences of $\chi_K$,
i.e. it affects only the premises. Substitution $P_g^+$ substitutes only
positive occurrences of $\chi_K$, i.e. it affects only the conclusions.
$P_g^+$ ensures that the parameter instances of the postcondition are
connected with the argument variables. As a matter of implementation, the
substitution $\tmop{recover} (\vec{\tau}_{\varepsilon_K},
\vec{\tau}_{\tmop{old}})$ is actually derived from the way the variables are
freshened in step \ref{CaptureEscParams} above. Its effect is currently
implemented as a substitution over the intermediate formula $F_{\chi}
[\varepsilon_K (\vec{\tau}_{\tmop{old}}) \assign \varepsilon_K
(\vec{\tau}_{\varepsilon_K})]$.

We start with $S_0 \assign \bar{\tmmathbf{T}}$ and $R_0 \assign
\bar{\tmmathbf{T}}$. $S_k$ grow in strength by definition. The constraint
generalization parts $G_{\chi_K}$ of $R_1$ and $R_2$ are computed from
non-recusrive branches only. Starting from $R_2$, $R_k$ are expected to
decrease in strength, but monotonicity is not guaranteed because of
contributions from abduction: mostly in form of $A_{\beta_{\chi_K}}$, but also
from stronger premises due to $S_k$. We remove non-term-sort atoms containing
containing outer-scope parameters from $S_2$, to enable considering a
different solution when new facts about outer-scope parameters propagate.

$\tmop{Connected} (\alpha, G)$ is the connected component of hypergraph $G$
containing node $\alpha$, where nodes are variables $\tmop{FV} (G)$ and
hyperedges are atoms $c \in G$. $\tmop{Connected}_0 (\delta, (\bar{\alpha},
G))$ additionally removes atoms $c : \tmop{FV} (c) \# \delta \bar{\alpha}$. In
initial iterations, when the branches $D_K^{\alpha} \Rightarrow C_K^{\alpha}$
are selected from non-recursive branches only, we include a connected atom
only if it is satisfiable in all branches. $H (R_k, R_{k + 1})$ is a
convergence improving heuristic, with $H (R_k, R_{k + 1}) = \rho (R_g)$ for
early iterations and ``roughly'' $H (R_k, R_g) = R_k \cap \rho (R_g)$ later.
The substitution $H$, also computed by the function $H$, is a renaming that
replaces the variables introduced by constraint generalization in the current
iteration, by the corresponding variables introduced by constraint
generalization in the previous iteration.

The use sites of definitions with postconditions can introduce requirements
$A_{\beta_{\chi_K}}$ on postconditions. The inferred postconditions can only
meet those requirements which are guaranteed by all defining cases. We use
abduction again to strengthen the preconditions, so that the postconditions
meet the requirements. Abduction is applied directly to the constraint
branches preprocessed by constraint generalization, so that the required
postconditions will follow in the next iteration.

The condition $(\forall \beta_{\chi_K}) A_{\beta_{\chi_K}} =\tmmathbf{T}$ is
required for correctness of returned solution. Fixpoint of postconditions $R_k
(\chi_K)$ is not a sufficient stopping condition, because it does not imply
$A_{\beta_{\chi_K}} =\tmmathbf{T}$, the same $A_{\beta_{\chi_K}} \neq
\tmmathbf{T}$ may be introduced in consecutive iterations. This is not the
case for invariants, where $S_k (\chi)$ and $S_{k + 1} (\chi)$ differ by the
portion $A_{\beta_{\chi}}$ of abduction answer.

We introduced the \tmverbatim{assert false} construct into the programming
language to indicate that a branch of code should not be reached. Type
inference generates for it the logical connective $\tmmathbf{F}$ (falsehood).
We partition the implication branches $D_i, C_i$ into $\{ D_i, C_i |
\tmmathbf{F} \nin C_i \}$ which are fed to the algorithm and $\{ D_j^- \} = \{
D_i | \tmmathbf{F} \in C_i \}$. After the main algorithm ends we check that
for each $D_j^-$, \ $S_k (D_i)$ fails. Optionally, but by default, we perform
the check in each iteration, starting from the third iteration, i.e. $k > 1$.
Turning this option {\tmem{on}} gives a way to express negative constraints
for the term domain. The option should be turned {\tmem{off}} when a single
iteration (plus fallback backtracking described below) is insufficient to
solve for the invariants. Moreover, for domains with negation elimination,
i.e. the numerical domain, we incorporate negative constraints as described in
section \ref{NegElim}. The corresponding computation is $A_0 \wedge_j
\tmop{NegElim} (\neg \tmop{Simpl} (\tmop{FV} (D_j^-) \backslash
\overline{\bar{\beta}^{\chi}} .D_j^-), \overline{D_i, C_i})$ in the
specification of the main algorithm. The simplification reduces the number of
atoms in the formula and keeps those that are most relevant to finding
invariants and postconditions. For convenience, negation elimination is called
from the function that computes the abduction answer.

We implement backtracking using a tabu-search-like discard list. When
abduction raises an exception: for example contradiction arises in the
branches $S_k (\Phi)$ passed to it, or it cannot find an answer and raises
\tmverbatim{Suspect} with information on potential problem, we fall-back to
step $k - 1$. Similarly, with checking for negative constraints {\tmem{on}},
when the check of negative branches $D_i \in \Phi_{\tmmathbf{F}}$, \
$\mathcal{M} \nvDash \exists \tmop{FV} (S_k (D_i)) .S_k (D_i)$ fails. In step
$k - 1$, we maintain a discard list of different answers found not to work in
this step and previous steps: initially empty, after fall-back we add there
the latest partial solution. We redo step $k - 1$ starting from $S_{k - 1}
(\Phi)$. Infinite loop is avoided because answers already attempted are
discarded. When step $k - 1$ cannot find a new partial solution, we fall back
to step $k - 2$, etc. We store discard lists for distinct sorts separately and
we only add to the discard list of the sort that caused fallback.
Unfortunately this means a slight loss of completeness, as an error in another
sort might be due to bad choice in the sort of types. The loss is ``slight''
because of the dissociation step described previously. Moreover, the sort
information from checking negative branches is likewise only approximate.

\subsection{Stages of Iteration}

Changes in the algorithm between iterations were mentioned above but not
clearly exposed. We divide iterations into multiple stages, demarcated by
parameters $j_k$ we now describe.
\begin{enumerate}
  \item $j_0$ is when inferring any postconditions starts; we use the initial
  iteration to infer the shape of types using abduction.
  
  \item $j_1$ is when inferring numerical postconditions starts.
  
  \item $j_2$ is when using all branches of the constraint in inference
  starts; before $j_2$, we only use the branches of the constraint that do not
  contain requirements derived from recursive calls. If we use all branches
  too early, we suffer from missing information.
  
  \item $j_3$ is when second-phase abduction, to enrich preconditions based on
  postconditions, starts.
  
  \item $j_4$ is when any forms of guessing in constraint generation should
  end; we have not yet implemented any such guessing mechanisms.
  
  \item $j_5$ is when convergence of postconditions is enforced; from this
  step onward, a solution to postconditions in an iteration is truncated to
  contain only atoms matching atoms from the previous iteration.
  
  \item $j_6$ is the last iteration; if the fixpoint is not reached, we signal
  an error {\tmem{``Answers do not converge''}}.
\end{enumerate}
Default settings are $[j_0 ; j_1 ; j_2 ; j_3 ; j_4 ; j_5 ; j_6] = [1 ; 1 ; 2 ;
3 ; 4 ; 8 ; 11]$ where the initial iteration has index 0. In a single
iteration, constraint generalization precedes abduction.

When existential types are used, the expected number of iterations is $k = 5$
(six iterations), because the last iteration needs to verify that the
last-but-one iteration has found the correct solution. The minimal number of
iterations is $k = 2$ (three iterations), so that all branches are considered.

\subsection{Implementation Details}\label{Details}

We represent $\vec{\alpha}$ as a tuple type rather than as a function type. We
modify the quantifier $\mathcal{Q}$ imperatively, because it mostly grows
monotonically, and when variables are dropped they do not conflict with fresh
variables generated later.

The code that selects $\wedge_{\chi} (A_{\chi}^+ \subset A_{\chi}^1) \wedge
\mathcal{M} \vDash \mathcal{Q}.A \setminus \cup_{\chi} A_{\chi}^+$ is an
incremental validity checker. It starts with $A \setminus \cup_{\chi}
A_{\chi}^1$ and tries to add as many atoms $c \in \cup_{\chi} A_{\chi}^1$ as
possible to what in effect becomes $A_{\tmop{res}}$.

We count the number of iterations of the main loop, a fallback decreases the
iteration number to the previous value. The main loop decides whether
multisort abduction should dissociate alien subterms -- in the first iteration
of the loop -- or should perform abductions for other sorts -- in subsequent
iteration. See discussion in subsection \ref{AlienSubterms}. In the first two
iterations, we remove branches that contain unary predicate variables in the
conclusion (or binary predicate variables in the premise, keeping only
non-recursive branches), as discussed at the end of subsection \ref{SCAlinear}
and beginning of subsection \ref{NumConv}. As discussed in subsection
\ref{NumConv}, starting from iteration $k_3$, we enforce convergence on
solutions for binary predicate variables.

Computing abduction is the ``axis'' of the main loop. If anything fails, the
previous abduction answer is the culprit. We add the previous answer to the
discard list and retry, without incrementing the iteration number. If
abduction and splitting succeeds, we reset the discard list and increment the
iteration number. We use recursion for backtracking, instead of making
\tmverbatim{loop} tail-recursive.

\section{Generating OCaml Source and Interface Code}

We have a single basis from which to generate all generated output files:
\tmverbatim{.gadti}, \tmverbatim{.ml} -- \tmverbatim{annot\_item}. It contains
a superset of information in \tmverbatim{struct\_item}: type scheme
annotations on introduced names, and source code annotated with type schemes
at recursive definition nodes. We use \tmverbatim{type a.} syntax instead of
\tmverbatim{'a.} syntax because the former supports inference for GADTs in
OCaml. A benefit of the nicer \tmverbatim{type a.} syntax is that nested type
schemes can have free variables, which will be correctly captured by the outer
type scheme. For completeness we sometimes need to annotate all
\tmverbatim{function} nodes with types. To avoid clutter, we start by only
annotating \tmverbatim{let rec} nodes, and in case \tmverbatim{ocamlc -c}
fails on generated code, we re-annotate by putting type schemes on
\tmverbatim{let rec} nodes and types on \tmverbatim{function} nodes. If need
arises, \tmverbatim{let-in} node annotations can also be introduced in this
fallback -- we provide an option to annotate the definitions on
\tmverbatim{let-in} nodes. Type annotations are optional because they
introduce a slight burden on the solver -- the corresponding variables cannot
be removed by the initial simplification of the constraints.
\tmverbatim{let-in} node annotations are more burdensome than
\tmverbatim{function} node annotations.

In the signature declarations for existential types, we replace existential
identifiers with regular indentifiers of constructors, to get informative
output for printing the various result files. We print constraint formulas and
alien subterms in the original InvarGenT syntax, commented out.

The types \tmverbatim{Int}, \tmverbatim{Num}, \tmverbatim{Bool} and
\tmverbatim{String} should be considered built-in. \tmverbatim{Int},
\tmverbatim{Bool} and \tmverbatim{String} follow the general scheme of
exporting a datatype constructor with the same name, only lower-case. However,
numerals \tmverbatim{0}, \tmverbatim{1}, ... are always type-checked as
\tmverbatim{Num 0}, \tmverbatim{Num 1}... A parameter \tmverbatim{-num\_is}
decides the type alias definition added in the generated code:
\tmverbatim{-num\_is bar} adds \tmverbatim{type num = bar} in front of an
\tmverbatim{.ml} file, by default \tmverbatim{int}. Numerals are exported as
integers passed to a \tmverbatim{bar\_of\_int} function. The variant
\tmverbatim{-num\_is\_mod} exports numerals by passing to a
\tmverbatim{Bar.of\_int} function. Special treatment for \tmverbatim{Bool}
amounts to exporting \tmverbatim{True} and \tmverbatim{False} as
\tmverbatim{true} and \tmverbatim{false}, unlike other constants. In addition,
pattern matching \tmverbatim{match}$\ldots$\tmverbatim{ with True
->}$\ldots$\tmverbatim{ \textbar  False ->}$\ldots$, i.e. the corresponding
beta-redex, is exported as the \tmverbatim{if}$\ldots$\tmverbatim{
then}$\ldots$\tmverbatim{ else}$\ldots$ expression.

In declarations which have concrete syntax starting with the word
\tmverbatim{external}, we provide names assumed to have given type scheme. The
syntax has two variants, differing in the way the declaration is exported. It
can be either an \tmverbatim{external} declaration in OCaml, which is the
binding mechanism of the {\tmem{foreign function interface}}. But in the
InvarGenT form \tmverbatim{external let}, the declaration provides an OCaml
definition, which is exported as the toplevel \tmverbatim{let} definition of
OCaml. It has the benefit that the OCaml compiler will verify this definition,
since InvarGenT calls \tmverbatim{ocamlc -c} to verify the exported
code.{\newpage}

\chapter{Source Code of Examples}\label{ExamplesCode}

Subsection titles start with ``Function'' for examples from Tables
\ref{Experiments} and \ref{ExperimentsNum}, and ``Program'' for examples from
Table \ref{ExperimentsBounds}.

\section{Incompleteness Example for \tmverbatim{OutsideIn}: Function
\tmverbatim{rx}}

Example from {\cite{OutsideIn}}, described on page 50. Source file
{\tmem{non\_outsidein\_rx.gadt}}:
\begin{tmcode}
datatype R : type
datacons RBool : R Bool
external let fortytwo : Int = "42"

let rx = function RBool -> fortytwo
\end{tmcode}
Value items from {\tmem{non\_outsidein\_rx.gadti.target}}:
\begin{tmcode}
val rx : a. R a  Int
\end{tmcode}

\section{Pointwise Examples}

These are the longer examples from {\cite{PointwiseGADTs}} within the scope of
algorithm $\mathcal{P}$.

\subsection{Function \tmverbatim{rotate}}

Example from {\cite{PracticalGADTsInfer}}, described on pages 193-194. Source
file {\tmem{pointwise\_rbtree\_rotate.gadt}}:
\begin{tmcode}
datatype Z
datatype S : type

datatype RoB : type * type
datatype Black
datatype Red
datacons Leaf : RoB (Black, Z)
datacons RNode : a. RoB (Black, a) * Int * RoB (Black, a)  RoB (Red, a)
datacons BNode :
  a, b, c. RoB (a, c) * Int * RoB (b, c)  RoB (Black, S c)

datatype Dir
datacons LeftD : Dir
datacons RightD : Dir

let rotate = fun dir1 p_e sib dir2 g_e uncle -> function
| RNode (x, e, y) ->
  (match dir1, dir2 with
  | RightD, RightD -> BNode (RNode (x, e, y), p_e, RNode (sib, g_e, uncle))
  | RightD, LeftD -> BNode (RNode (uncle, g_e, x), e, RNode (y, p_e, sib))
  | LeftD, RightD -> BNode (RNode (sib, p_e, x), e, RNode (y, g_e, uncle))
  | LeftD, LeftD -> BNode (RNode (uncle, g_e, sib), p_e, RNode (x, e, y)))
| BNode (_, _, _) -> assert false
\end{tmcode}
Value items from {\tmem{pointwise\_rbtree\_rotate.gadti.target}}:
\begin{tmcode}
val rotate :
   a.
   Dir  Int  RoB (Black, a)  Dir  Int  RoB (Black, a) 
     RoB (Red, a)  RoB (Black, S a)
\end{tmcode}

\subsection{Function \tmverbatim{zip2}: $N$-way \tmverbatim{zip\_with}}

Example from {\cite{PracticalGADTsInfer}}, provided on page 206. The
definition of \tmverbatim{z2} is $\eta$-expanded here to fit the call-by-value
semantics of {\tmname{InvarGenT}}. Source file {\tmem{pointwise\_zip2.gadt}}:
\begin{tmcode}
datatype List : type
datacons N : a. List a
datacons C : a. a * List a  List a

datatype Zip2 : type * type
datacons Zero2 : a. Zip2 (a, List a)
datacons Succ2 :
  a, b, c. Zip2 (a, b)  Zip2 ((c  a), (List c  b))

let zip2 =
  let rec zipZ = function
    | Zero2 -> N
    | Succ2 n -> fun _ -> zipZ n in
  let rec zipS = fun f r -> function
    | Zero2 -> C (f, r)
    | Succ2 n -> function
      | N -> zipZ n
      | C (z, zs) -> zipS (f z) (r zs) n in
  let rec z2 = fun n f -> zipS f (z2 n f) n in
  z2
\end{tmcode}
Value items from {\tmem{pointwise\_zip2.gadti.target}}:
\begin{tmcode}
val zip2 : a, b. Zip2 (b, a)  b  a
\end{tmcode}

\subsection{Function \tmverbatim{rotl}}

Example from {\cite{PracticalGADTsInfer}}, described on pages 198-199. Source
file {\tmem{pointwise\_avl\_rotl.gadt}}:
\begin{tmcode}
datatype Z
datatype S : type
datatype AVL : type
datacons Tip : AVL Z
datacons LNode : a. AVL a * Int * AVL (S a)  AVL (S (S a))
datacons SNode : a. AVL a * Int * AVL a  AVL (S a)
datacons MNode : a. AVL (S a) * Int * AVL a  AVL (S (S a))

datatype Choice : type * type
datacons L : a, b. a  Choice (a, b)
datacons R : a, b. b  Choice (a, b)

let rotl = fun u v -> function
  | Tip -> assert false
  | SNode (a, x, b) -> R (MNode (LNode (u, v, a), x, b))
  | LNode (a, x, b) -> L (SNode (SNode (u, v, a), x, b))
  | MNode (k, y, c) ->
    (match k with
      | SNode (a, x, b) -> L (SNode (SNode (u, v, a), x, SNode (b, y, c)))
      | LNode (a, x, b) -> L (SNode (MNode (u, v, a), x, SNode (b, y, c)))
      | MNode (a, x, b) -> L (SNode (SNode (u, v, a), x, LNode (b, y, c))))
\end{tmcode}
Value items from {\tmem{pointwise\_avl\_rotl.gadti.target}}:
\begin{tmcode}
val rotl :
   a.
   AVL a  Int  AVL (S (S a)) 
     Choice (AVL (S (S a)), AVL (S (S (S a))))
\end{tmcode}

\subsection{Function \tmverbatim{ins}}

Example from {\cite{PracticalGADTsInfer}}, provided on page 209. Source file
{\tmem{pointwise\_avl\_ins.gadt}}:
\begin{tmcode}
datatype Z
datatype S : type
datatype AVL : type
datacons Tip : AVL Z
datacons LNode : a. AVL a * Int * AVL (S a)  AVL (S (S a))
datacons SNode : a. AVL a * Int * AVL a  AVL (S a)
datacons MNode : a. AVL (S a) * Int * AVL a  AVL (S (S a))

datatype Choice : type * type
datacons L : a, b. a  Choice (a, b)
datacons R : a, b. b  Choice (a, b)

datatype LinOrder
datacons LT : LinOrder
datacons EQ : LinOrder
datacons GT : LinOrder
external let compare : a. a  a  LinOrder =
  "fun x y -> let c=Pervasives.compare x y in
              if c<0 then LT else if c=0 then EQ else GT"

external rotl :
  a. AVL a  Int  AVL (S (S a)) 
    Choice (AVL (S (S a)), AVL (S (S (S a))))

external rotr :
  a. AVL (S (S a))  Int  AVL a 
    Choice (AVL (S (S a)), AVL (S (S (S a))))

let rec ins = fun i -> function
  | Tip -> R (SNode (Tip, i, Tip))
  | SNode (a, x, b) as t ->
    (match compare i x with
      | EQ -> L t
      | LT ->
        (match ins i a with
          | L a -> L (SNode (a, x, b))
          | R a -> R (MNode (a, x, b)))
      | GT ->
        (match ins i b with
          | L b -> L (SNode (a, x, b))
          | R b -> R (LNode (a, x, b))))
  | LNode (a, x, b) as t ->
    (match compare i x with
      | EQ -> L t
      | LT ->
        (match ins i a with
          | L a -> L (LNode (a, x, b))
          | R a -> L (SNode (a, x, b)))
      | GT ->
        (match ins i b with
          | L b -> L (LNode (a, x, b))
          | R b -> rotl a x b))
  | MNode (a, x, b) as t ->
    (match compare i x with
      | EQ -> L t
      | LT ->
        (match ins i a with
          | L a -> L (MNode (a, x, b))
          | R a -> rotr a x b)
      | GT ->
        (match ins i b with
          | L b -> L (MNode (a, x, b))
          | R b -> L (SNode (a, x, b))))
\end{tmcode}
Value items from {\tmem{pointwise\_avl\_ins.gadti.target}}:
\begin{tmcode}
val ins : a. Int  AVL a  Choice (AVL a, AVL (S a))
\end{tmcode}

\subsection{Function \tmverbatim{extract}}

Example from {\cite{PracticalGADTsInfer}}, provided on page 207. Source file
{\tmem{pointwise\_extract.gadt}}:
\begin{tmcode}
datatype List : type
datacons N : a. List a
datacons C : a. a * List a  List a

datatype Nd
datatype Fk : type * type

datatype Tree : type * type
datacons End : a. a  Tree (Nd, a)
datacons Fork : a, b, c. Tree (a, c) * Tree (b, c)  Tree (Fk (a, b), c)

datatype Path : type
datacons Here : Path Nd
datacons ForkL : a, b. Path a  Path (Fk (a, b))
datacons ForkR : a, b. Path b  Path (Fk (a, b))

external append : a. List a  List a  List a
external map : a, b. (a  b)  List a  List b

let rec find f = function
  | End m -> if f m then C (Here, N) else N
  | Fork (x, y) -> append (map (fun y -> ForkL y) (find f x))
                          (map (fun y -> ForkR y) (find f y))

let rec extract = fun p t ->
  match p with
    | Here -> (match t with End m -> m | Fork (_, _) -> assert false)
    | ForkL p1 -> (match t with Fork (x, y) -> extract p1 x)
    | ForkR p1 -> (match t with Fork (x, y) -> extract p1 y)
\end{tmcode}
Value items from {\tmem{pointwise\_extract.gadti.target}}:
\begin{tmcode}
val find : a, b. (a  Bool)  Tree (b, a)  List (Path b)

val extract : a, b. Path b  Tree (b, a)  a
\end{tmcode}

\subsection{Function \tmverbatim{run\_state}}

Example from {\cite{PracticalGADTsInfer}}, provided on page 208. Source file
{\tmem{pointwise\_run\_state.gadt}}:
\begin{tmcode}
datatype State : type * type
datacons Bind :
  a, b, s. State (s, a) * (a  State (s, b))  State (s, b)
datacons Return : a, s. a  State (s, a)
datacons Get : s. State (s, s)
datacons Put : s. s  State (s, ())

let rec run_state = fun s -> function
  | Return a -> (s, a)
  | Get -> (s, s)
  | Put u -> (u, ())
  | Bind (m, k) ->
    let s1, a1 = run_state s m in
    run_state s1 (k a1)
\end{tmcode}
Value items from {\tmem{pointwise\_run\_state.gadti.target}}:
\begin{tmcode}
val run_state : a, b. a  State (a, b)  (a, b)
\end{tmcode}

\section{Non-Pointwise Examples}

These are the examples from {\cite{PointwiseGADTs}} falling outside the scope
of algorithm $\mathcal{P}$.

\subsection{Function \tmverbatim{joint}}

Example from {\cite{PracticalGADTsInfer}}, described on page 86. Source file
{\tmem{non\_pointwise\_split.gadt}}:
\begin{tmcode}
datatype Split : type * type
datacons Whole : Split (Int, Int)
datacons Parts : a, b. Split ((Int, a), (b, Bool))
external let seven : Int = "7"
external let three : Int = "3"

let joint = function
  | Whole -> seven
  | Parts -> three, True
\end{tmcode}
Value items from {\tmem{non\_pointwise\_split.gadti.target}}:
\begin{tmcode}
val joint : a. Split (a, a)  a
\end{tmcode}

\subsection{Function \tmverbatim{rotr}}

Example from {\cite{PracticalGADTsInfer}}, described on page 88. Source file
{\tmem{non\_pointwise\_avl.gadt}}:
\begin{tmcode}
(** Normally we would use [num], but this is a stress test for [type]. *)
datatype Z
datatype S : type
datatype Balance : type * type * type
datacons Less : a. Balance (a, S a, S a)
datacons Same : a. Balance (a, a, a)
datacons More : a. Balance (S a, a, S a)
datatype AVL : type
datacons Leaf : AVL Z
datacons Node :
  a, b, c. Balance (a, b, c) * AVL a * Int * AVL b  AVL (S c)

datatype Choice : type * type
datacons Left : a, b. a  Choice (a, b)
datacons Right : a, b. b  Choice (a, b)

let rotr = fun z d -> function
  | Leaf -> assert false
  | Node (Less, a, x, Leaf) -> assert false
  | Node (Same, a, x, (Node (_,_,_,_) as b)) ->
    Right (Node (Less, a, x, Node (More, b, z, d)))
  | Node (More, a, x, (Node (_,_,_,_) as b)) ->
    Left (Node (Same, a, x, Node (Same, b, z, d)))
  | Node (Less, a, x, Node (Same, b, y, c)) ->
    Left (Node (Same, Node (Same, a, x, b), y, Node (Same, c, z, d)))
  | Node (Less, a, x, Node (Less, b, y, c)) ->
    Left (Node (Same, Node (More, a, x, b), y, Node (Same, c, z, d)))
  | Node (Less, a, x, Node (More, b, y, c)) ->
    Left (Node (Same, Node (Same, a, x, b), y, Node (Less, c, z, d)))
\end{tmcode}
Value items from {\tmem{non\_pointwise\_avl.gadti.target}}:
\begin{tmcode}
val rotr :
   a.
   Int  AVL a  AVL (S (S a)) 
     Choice (AVL (S (S a)), AVL (S (S (S a))))
\end{tmcode}

\subsection{Function \tmverbatim{delmin}}

Example from {\cite{PracticalGADTsInfer}}, described on pages 212-213. Source
file {\tmem{non\_pointwise\_avl\_delmin.gadt}}:
\begin{tmcode}
(** Normally we would use [num], but this is a stress test for [type]. *)
datatype Z
datatype S : type
(** This datatype definition is modified to make type inference for
    rotr, rotl, ins functions pointwise. *)
datatype AVL : type
datacons Tip : AVL Z
datacons LNode : a. AVL a * Int * AVL (S a)  AVL (S (S a))
datacons SNode : a. AVL a * Int * AVL a  AVL (S a)
datacons MNode : a. AVL (S a) * Int * AVL a  AVL (S (S a))

datatype Choice : type * type
datacons L : a, b. a  Choice (a, b)
datacons R : a, b. b  Choice (a, b)

datatype Zero : type
datacons IsZ : Zero Z
datacons NotZ : a. Zero (S a)

external rotl :
  a. AVL a  Int  AVL (S (S a))  Choice (AVL (S (S a)),
                                               AVL (S (S (S a))))

external rotr :
  a. AVL (S (S a))  Int  AVL a  Choice (AVL (S (S a)),
                                               AVL (S (S (S a))))

let empty = function
  | Tip -> IsZ
  | LNode (_, _, _) -> NotZ
  | SNode (_, _, _) -> NotZ
  | MNode (_, _, _) -> NotZ

let rec delmin = function
  | LNode (a, x, b) ->
    (match empty a with
      | IsZ -> x, L b
      | NotZ ->
        (match delmin a with
          y, k ->
          (match k with
            | L a -> y, rotl a x b
            | R a -> y, R (LNode (a, x, b)))))
  | SNode (a, x, b) ->
    (match empty a with
      | IsZ -> x, L b
      | NotZ ->
        (match delmin a with
          y, k ->
          (match k with
            | L a -> y, R (LNode (a, x, b))
            | R a -> y, R (SNode (a, x, b)))))
  | MNode (a, x, b) ->
    (match delmin a with
      y, k ->
      (match k with
        | L a -> y, L (SNode (a, x, b))
        | R a -> y, R (MNode (a, x, b))))
\end{tmcode}
Value items from {\tmem{non\_pointwise\_avl\_delmin.gadti.target}}:
\begin{tmcode}
val empty : a. AVL a  Zero a

val delmin : a. AVL (S a)  (Int, Choice (AVL a, AVL (S a)))
\end{tmcode}

\subsection{Function \tmverbatim{fd\_comp}}

Example from {\cite{PracticalGADTsInfer}}, described on pages 213-215. The
only difference is that the original has \tmverbatim{\textbar  FDI -> fd2} as
the fourth line of \tmverbatim{fd\_comp} instead of our expanded
\tmverbatim{\textbar  FDI -> (match fd2 with \textbar  FDI -> fd2 \textbar 
FDC \_ -> fd2 \textbar  FDG \_ -> fd2)}. Source file
{\tmem{non\_pointwise\_fd\_comp.gadt}}:
\begin{tmcode}
datatype FunDesc : type * type
datacons FDI : a. FunDesc (a, a)
datacons FDC : a, b. b  FunDesc (a, b)
datacons FDG : a, b. (a  b)  FunDesc (a, b)

external fd_fun : a, b. FunDesc (a, b)  a  b

let fd_comp = fun fd1 fd2 ->
  let o = fun f g x -> f (g x) in
  match fd1 with
    | FDI -> (match fd2 with | FDI -> fd2 | FDC _ -> fd2 | FDG _ -> fd2)
    | FDC b -> 
      (match fd2 with
        | FDI -> fd1
        | FDC c -> FDC (fd_fun fd2 b)
        | FDG g -> FDC (fd_fun fd2 b))
    | FDG f ->
      (match fd2 with
        | FDI -> fd1
        | FDC c -> FDC c
        | FDG g -> FDG (o (fd_fun fd2) f))
\end{tmcode}
Compare also the example {\tmem{non\_pointwise\_fd\_comp2.gadt}}:
\begin{tmcode}
datatype FunDesc : type * type
datacons FDI : a. FunDesc (a, a)
datacons FDC : a, b. b  FunDesc (a, b)
datacons FDG : a, b. (a  b)  FunDesc (a, b)

external fd_fun : a, b. FunDesc (a, b)  a  b

let o = fun f g x -> f (g x)

let fd_comp =
  function
    | FDI -> (function FDI -> FDI | FDC c -> FDC c | FDG g -> FDG g)
    | FDC b as fd1 ->
      (function
        | FDI -> fd1
        | FDC c as fd2 -> FDC (fd_fun fd2 b)
        | FDG g as fd2 -> FDC (fd_fun fd2 b))
    | FDG f as fd1 ->
      (function
        | FDI -> fd1
        | FDC c -> FDC c
        | FDG g as fd2 -> FDG (o (fd_fun fd2) f))
\end{tmcode}
Value items from {\tmem{non\_pointwise\_fd\_comp2.gadti.target}}:
\begin{tmcode}
val o : a, b, c. (b  a)  (c  b)  c  a

val fd_comp :
   a, b, c. FunDesc (b, c)  FunDesc (c, a)  FunDesc (b, a)
\end{tmcode}

\subsection{Function \tmverbatim{zip1}: $N$-way \tmverbatim{zip\_with}}

Example from {\cite{PracticalGADTsInfer}}, described on pages 216-217. The
local definition of \tmverbatim{apply} is a \tmverbatim{let rec} here because
it needs to be polymorphic. Source file {\tmem{non\_pointwise\_zip1.gadt}}:
\begin{tmcode}
datatype List : type
datacons N : a. List a
datacons C : a. a * List a  List a

datatype Zip1 : type * type
datacons Zero1 : a. Zip1 (List a, List a)
datacons Succ1 :
  a, b, c. Zip1 (List a, b)  Zip1 (List (c  a), (List c  b))

external zip_with : a, b, c. (a  b  c)  List a  List b  List c

external repeat : a. a  List a

let zip1 =
  let rec apply = fun f x -> f x in
  let rec z1 = fun fs -> function
    | Zero1 -> fs
    | Succ1 n2 -> fun xs -> z1 (zip_with apply fs xs) n2 in
  fun n f -> z1 (repeat f) n
\end{tmcode}
Value items from {\tmem{non\_pointwise\_zip1.gadt}}:
\begin{tmcode}
val zip1 : a, b. Zip1 (List b, a)  b  a
\end{tmcode}

\subsection{Function \tmverbatim{leq}}

Example from {\cite{PracticalGADTsInfer}}, described on pages 217-219. Source
file {\tmem{non\_pointwise\_leq.gadt}}, requires option
\tmverbatim{-prefer\_guess} for inference:
\begin{tmcode}
datatype Z
datatype S : type

datatype Nat : type
datacons Zn : Nat Z
datacons Sn : a. Nat a  Nat (S a)

datatype NatLeq : type * type
datacons LeZ : a. NatLeq (Z, a)
datacons LeS : a, b. NatLeq (a, b)   NatLeq (S a, S b)

let rec leq = function
  | Zn -> LeZ
  | Sn n -> LeS (leq n)
\end{tmcode}
Value items from {\tmem{non\_pointwise\_leq.gadti.target}}:
\begin{tmcode}
val leq : a. Nat a  NatLeq (a, a)
\end{tmcode}

\subsection{Function \tmverbatim{run\_state}}

Example from {\cite{PracticalGADTsInfer}}, described on pages 219-220. Source
file {\tmem{non\_pointwise\_run\_state.gadt}}:
\begin{tmcode}
datatype State : type * type
datacons Bind :
  a, b, s. State (s, a) * (a  State (s, b))  State (s, b)
datacons Return : a, s. a  State (s, a)
datacons Get : s. State (s, s)
datacons Put : s. s  State (s, ())

let rec run_state = fun s -> function
  | Return a -> (s, a)
  | Get -> (s, s)
  | Put u -> (u, ())
  | Bind (m, k) ->
    match m with
      | Return a -> run_state s (k a)
      | Get -> run_state s (k s)
      | Put u -> run_state u (k ())
      | Bind (n, j) -> run_state s (Bind (n, fun x -> Bind (j x, k)))
\end{tmcode}
Value items from {\tmem{non\_pointwise\_run\_state.gadti.target}}:
\begin{tmcode}
val run_state : a, b. a  State (a, b)  (a, b)
\end{tmcode}

\section{Run-time Type Representations}

\subsection{Function \tmverbatim{eval}}

Source file {\tmem{eval.gadt}}:
\begin{tmcode}
datatype Term : type

external let plus : Int  Int  Int = "(+)"
external let is_zero : Int  Bool = "(=) 0"

datacons Lit : Int  Term Int
datacons Plus : Term Int * Term Int  Term Int
datacons IsZero : Term Int  Term Bool
datacons If : a. Term Bool * Term a * Term a  Term a
datacons Pair : a, b. Term a * Term b  Term (a, b)
datacons Fst : a, b. Term (a, b)  Term a
datacons Snd : a, b. Term (a, b)  Term b

let rec eval = function
  | Lit i -> i
  | IsZero x -> is_zero (eval x)
  | Plus (x, y) -> plus (eval x) (eval y)
  | If (b, t, e) -> (match eval b with True -> eval t | False -> eval e)
  | Pair (x, y) -> eval x, eval y
  | Fst p -> (match eval p with x, y -> x)
  | Snd p -> (match eval p with x, y -> y)
\end{tmcode}
Value items from {\tmem{eval.gadti.target}}:
\begin{tmcode}
val eval : a. Term a  a
\end{tmcode}
Exported OCaml source file {\tmem{eval.ml.target}}:
\begin{tmcode}
type num = int
type _ term =
  | Lit : int -> int term
  | Plus : int term * int term -> int term
  | IsZero : int term -> bool term
  | If : (*'a.*)bool term * 'a term * 'a term -> 'a term
  | Pair : (*'a, 'b.*)'a term * 'b term -> (('a * 'b)) term
  | Fst : (*'a, 'b.*)(('a * 'b)) term -> 'a term
  | Snd : (*'a, 'b.*)(('a * 'b)) term -> 'b term
let plus : (int -> int -> int) = (+)
let is_zero : (int -> bool) = (=) 0
let rec eval : type a . (a term -> a) =
  ((function Lit i -> i | IsZero x -> is_zero (eval x)
    | Plus (x, y) -> plus (eval x) (eval y)
    | If (b, t, e) -> (if eval b then eval t else eval e)
    | Pair (x, y) -> (eval x, eval y)
    | Fst p -> let ((x, y): (a * _)) = eval p in x
    | Snd p -> let ((x, y): (_ * a)) = eval p in y): a term -> a)
\end{tmcode}

\subsection{Function \tmverbatim{equal}}

Source file {\tmem{equal\_assert.gadt}}:
\begin{tmcode}
datatype Ty : type
datatype List : type
datacons Zero : Int
datacons Nil : a. List a
datacons TInt : Ty Int
datacons TPair : a, b. Ty a * Ty b  Ty (a, b)
datacons TList : a. Ty a  Ty (List a)
external let eq_int : Int  Int  Bool = "(=)"
external let b_and : Bool  Bool  Bool = "(&&)"
external let b_not : Bool  Bool = "not"
external forall2 : a, b. (a  b  Bool)  List a  List b  Bool

let rec equal = function
  | TInt, TInt -> fun x y -> eq_int x y
  | TPair (t1, t2), TPair (u1, u2) ->  
    (fun (x1, x2) (y1, y2) ->
        b_and (equal (t1, u1) x1 y1)
              (equal (t2, u2) x2 y2))
  | TList t, TList u -> forall2 (equal (t, u))
  | _ -> fun _ _ -> False
  | TInt, TList l -> (function Nil -> assert false)
  | TList l, TInt -> (fun _ -> function Nil -> assert false)
\end{tmcode}
Source file {\tmem{equal\_test.gadt}}:
\begin{tmcode}
datatype Ty : type
datatype List : type
datacons Nil : a. List a
datacons TInt : Ty Int
datacons TPair : a, b. Ty a * Ty b  Ty (a, b)
datacons TList : a. Ty a  Ty (List a)
external let zero : Int = "0"
external let eq_int : Int  Int  Bool = "(=)"
external let b_and : Bool  Bool  Bool = "(&&)"
external let b_not : Bool  Bool = "not"
external forall2 : a, b. (a  b  Bool)  List a  List b  Bool

let rec equal = function
  | TInt, TInt -> fun x y -> eq_int x y
  | TPair (t1, t2), TPair (u1, u2) ->  
    (fun (x1, x2) (y1, y2) ->
        b_and (equal (t1, u1) x1 y1)
              (equal (t2, u2) x2 y2))
  | TList t, TList u -> forall2 (equal (t, u))
  | _ -> fun _ _ -> False
test b_not (equal (TInt, TList TInt) zero Nil)
\end{tmcode}
Value items from {\tmem{equal\_test.gadti.target}}:
\begin{tmcode}
val equal : a, b. (Ty a, Ty b)  a  b  Bool
\end{tmcode}

\section{Lists with Length}

\subsection{Function \tmverbatim{head}}

Source file {\tmem{list\_head.gadt}}:
\begin{tmcode}
datatype List : type * num
datacons LNil : a. List(a, 0)
datacons LCons : n, a [0n]. a * List(a, n)  List(a, n+1)

let head = function
  | LCons (x, _) -> x
  | LNil -> assert false
\end{tmcode}
Value items from {\tmem{list\_head.gadti.target}}:
\begin{tmcode}
val head : n, a[1  n]. List (a, n)  a
\end{tmcode}

\subsection{Function \tmverbatim{append}}

This example is not featured in the {\tmem{examples}} directory. The natural
solution would be to propagate one of the arguments when the other list is
empty: \tmverbatim{function LNil -> (fun l -{\nobreak}> l) \textbar}$\ldots$
This currently leads to the problem of insufficient information about
\tmverbatim{l}. We can expand all cases of both arguments:
\begin{tmcode}
datatype Elem
datatype List : num
datacons LNil : List 0
datacons LCons : n [0n]. Elem * List n  List (n+1)

let rec append =
  function LNil -> (function LNil -> LNil | LCons (_,_) as l -> l)
    | LCons (x, xs) ->
      (function LNil -> LCons (x, append xs LNil)
        | LCons (_,_) as l -> LCons (x, append xs l))
\end{tmcode}
Inferred type:
\begin{tmcode}
val append : n, k. List (a, k)  List (a, n)  List (n + k)
\end{tmcode}
We can also use assertions:
\begin{tmcode}
let rec append =
  function
    | LNil ->
      (function l when (length l + 1) <= 0 -> assert false | l -> l)
    | LCons (x, xs) ->
      (function l when (length l + 1) <= 0 -> assert false
      | l -> LCons (x, append xs l))
\end{tmcode}
Inferred type:
\begin{tmcode}
val append : n, k[0n  0n + k]. List(a, k)  List(a, n)  List(n + k)
\end{tmcode}

\subsection{Function \tmverbatim{flatten\_pairs}}

Source file {\tmem{flatten\_pairs.gadt}}:
\begin{tmcode}
datatype List : type * num
datacons LNil : a. List(a, 0)
datacons LCons : n, a [0n]. a * List(a, n)  List(a, n+1)

let rec flatten_pairs =
  function LNil -> LNil
    | LCons ((x, y), l) ->
      LCons (x, LCons (y, flatten_pairs l))
\end{tmcode}
Value items from {\tmem{flatten\_pairs.gadti.target}}:
\begin{tmcode}
val flatten_pairs : n, a. List ((a, a), n)  List (a, 2 n)
\end{tmcode}

\subsection{Function \tmverbatim{filter}}

Source file {\tmem{filter.gadt}}:
\begin{tmcode}
datatype List : type * num
datacons LNil : a. List(a, 0)
datacons LCons : n, a [0n]. a * List(a, n)  List(a, n+1)

let rec filter = fun f ->
  efunction LNil -> LNil
    | LCons (x, xs) ->
      ematch f x with
        | True ->
          let ys = filter f xs in
          LCons (x, ys)
        | False ->
          filter f xs
\end{tmcode}
Value items from {\tmem{filter.gadti.target}}:
\begin{tmcode}
val filter :
   n, a.
   (a  Bool)  List (a, n)  k[0  k  k  n].List (a, k)
\end{tmcode}

\subsection{Function \tmverbatim{zip}}

Source file {\tmem{zip.gadt}}:
\begin{tmcode}
datatype List : type * num
datacons LNil : a. List(a, 0)
datacons LCons : n, a [0n]. a * List(a, n)  List(a, n+1)

let rec zip =
  efunction
    | LNil, LNil -> LNil
    | LNil, LCons (_, _) -> LNil
    | LCons (_, _), LNil -> LNil
    | LCons (x, xs), LCons (y, ys) ->
      let zs = zip (xs, ys) in
      LCons ((x, y), zs)
\end{tmcode}
The inferred type is:
\begin{tmcode}
val n, k, a, b. (List (a, n), List (b, k)) 
                  i[i=min (k, n)].List ((a, b), i)
\end{tmcode}

\section{Binary Numbers}

\subsection{Function \tmverbatim{plus}}

Source file {\tmem{binary\_plus.gadt}}:
\begin{tmcode}
datatype Binary : num
datatype Carry : num

datacons Zero : Binary 0
datacons PZero : n [0n]. Binary n  Binary(2 n)
datacons POne : n [0n]. Binary n  Binary(2 n + 1)

datacons CZero : Carry 0
datacons COne : Carry 1

let rec plus =
  function CZero ->
    (function
      | Zero ->
        (function Zero -> Zero
          | PZero _ as b -> b
          | POne _ as b -> b)
      | PZero a1 as a ->
        (function Zero -> a
          | PZero b1 -> PZero (plus CZero a1 b1)
          | POne b1 -> POne (plus CZero a1 b1))
      | POne a1 as a ->
        (function Zero -> a
          | PZero b1 -> POne (plus CZero a1 b1)
          | POne b1 -> PZero (plus COne a1 b1)))
    | COne ->
    (function Zero ->
        (function Zero -> POne(Zero)
          | PZero b1 -> POne b1
          | POne b1 -> PZero (plus COne Zero b1))
      | PZero a1 ->
        (function Zero -> POne a1
          | PZero b1 -> POne (plus CZero a1 b1)
          | POne b1 -> PZero (plus COne a1 b1))
      | POne a1 ->
        (function Zero -> PZero (plus COne a1 Zero)
          | PZero b1 -> PZero (plus COne a1 b1)
          | POne b1 -> POne (plus COne a1 b1)))
\end{tmcode}
Value items from {\tmem{binary\_plus.gadti.target}}:
\begin{tmcode}
val plus :
   i, k, n. Carry i  Binary k  Binary n  Binary (n + k + i)
\end{tmcode}

\subsection{Function \tmverbatim{increment}}

This example is not featured in the {\tmem{examples}} directory.
\begin{tmcode}
datatype Binary : num

datacons Zero : Binary 0
datacons PZero : n [0n]. Binary n  Binary(2 n)
datacons POne : n [0n]. Binary n  Binary(2 n + 1)

let rec increment =
  function Zero -> POne Zero
    | PZero a1 -> POne a1
    | POne a1 -> PZero (increment a1)
\end{tmcode}
Inferred type:
\begin{tmcode}
val increment : n. Binary n  Binary (n + 1)
\end{tmcode}

\subsection{Function \tmverbatim{bitwise\_or}}

For brevity, the function is named \tmverbatim{ub} rather than
\tmverbatim{bitwise\_or} in the code example. Source file
{\tmem{binary\_upper\_bound.gadt}}:
\begin{tmcode}
datatype Binary : num
datacons Zero : Binary 0
datacons PZero : n [0n]. Binary n  Binary(2 n)
datacons POne : n [0n]. Binary n  Binary(2 n + 1)

let rec ub = efunction
  | Zero ->
      (efunction Zero -> Zero
        | PZero b1 as b -> b
        | POne b1 as b -> b)
  | PZero a1 as a ->
      (efunction Zero -> a
        | PZero b1 ->
          let r = ub a1 b1 in
          PZero r
        | POne b1 ->
          let r = ub a1 b1 in
          POne r)
  | POne a1 as a ->
      (efunction Zero -> a
        | PZero b1 ->
          let r = ub a1 b1 in
          POne r
        | POne b1 ->
          let r = ub a1 b1 in
          POne r)
\end{tmcode}
Value items from
\begin{tmcode}
val plus :
   i, k, n. Carry i  Binary k  Binary n  Binary (n + k + i)
\end{tmcode}

\section{AVL Trees}

Source file {\tmem{avl\_tree.gadt}}:
\begin{tmcode}
(** We follow the AVL tree algorithm from OCaml Standard Library, where 
    the branches of a node are allowed to differ in height by at most 2. *)

datatype Avl : type * num
datacons Empty : a. Avl (a, 0)
datacons Node :
  a,k,m,n [k=max(m,n)  0m  0n  nm+2  mn+2].
     Avl (a, m) * a * Avl (a, n) * Num (k+1)  Avl (a, k+1)

datatype LinOrder
datacons LT : LinOrder
datacons EQ : LinOrder
datacons GT : LinOrder
external let compare : a. a  a  LinOrder =
  "fun x y -> let c=Pervasives.compare x y in
              if c<0 then LT else if c=0 then EQ else GT"

let height = function
  | Empty -> 0
  | Node (_, _, _, k) -> k

let create = fun l x r ->
  eif height l <= height r then Node (l, x, r, height r + 1)
  else Node (l, x, r, height l + 1)

let singleton x = Node (Empty, x, Empty, 1)

let rotr = fun l x r ->
    ematch l with
    | Empty -> assert false
    | Node (ll, lx, lr, _) ->
      eif height lr <= height ll then
        let r' = create lr x r in
        create ll lx r'
      else
        ematch lr with
        | Empty -> assert false
        | Node (lrl, lrx, lrr, _) ->
          let l' = create ll lx lrl in
          let r' = create lrr x r in
          create l' lrx r'

let rotl = fun l x r ->
    ematch r with
    | Empty -> assert false
    | Node (rl, rx, rr, _) ->
      eif height rl <= height rr then
        let l' = create l x rl in
        create l' rx rr
      else
        ematch rl with
        | Empty -> assert false
        | Node (rll, rlx, rlr, _) ->
          let l' = create l x rll in
          let r' = create rlr rx rr in
          create l' rlx r'

let rec add x = efunction
  | Empty -> Node (Empty, x, Empty, 1)
  | Node (l, y, r, h) ->
    ematch compare x y, height l, height r with
    | EQ, _, _ -> Node (l, x, r, h)
    | LT, hl, hr ->
      let l' = add x l in
      eif height l' <= hr+2 then create l' y r else rotr l' y r
    | GT, hl, hr ->
      let r' = add x r in
      eif height r' <= hl+2 then create l y r' else rotl l y r'

let rec mem x = function
  | Empty -> False
  | Node (l, y, r, _) ->
    match compare x y with
    | LT -> mem x l
    | EQ -> True
    | GT -> mem x r

let rec min_binding = function
  | Empty -> assert false
  | Node (Empty, x, r, _) -> x
  | Node ((Node (_,_,_,_) as l), x, r, _) -> min_binding l

let rec remove_min_binding = efunction
  | Empty -> assert false
  | Node (Empty, x, r, _) -> r
  | Node ((Node (_,_,_,_) as l), x, r, _) ->
    let l' = remove_min_binding l in
    eif height r <= height l' + 2 then create l' x r
    else rotl l' x r

let merge = efunction
  | Empty, Empty -> Empty
  | Empty, (Node (_,_,_,_) as t) -> t
  | (Node (_,_,_,_) as t), Empty -> t
  | (Node (_,_,_,_) as t1), (Node (_,_,_,_) as t2) ->
    let x = min_binding t2 in
    let t2' = remove_min_binding t2 in
    eif height t1 <= height t2' + 2 then create t1 x t2'
    else rotr t1 x t2'

let rec remove = fun x -> efunction
  | Empty -> Empty
  | Node (l, y, r, h) ->
    ematch compare x y with
    | EQ -> merge (l, r)
    | LT ->
      let l' = remove x l in
      eif height r <= height l' + 2 then create l' y r
      else rotl l' y r
    | GT ->
      let r' = remove x r in
      eif height l <= height r' + 2 then create l y r'
      else rotr l y r'
\end{tmcode}
Value items from {\tmem{avl\_tree.gadti.target}}:
\begin{tmcode}
val height : n, a. Avl (a, n)  Num n

val create :
   k, n, a[0  n  0  k  n  k + 2  k  n + 2].
   Avl (a, k)  a  Avl (a, n)  i[i=max (k + 1, n + 1)].Avl (a, i)

val singleton : a. a  Avl (a, 1)

val rotr :
   k, n, a[0  n  n + 2  k  k  n + 3].
   Avl (a, k)  a  Avl (a, n)  n[k  n 
     n  k + 1].Avl (a, n)

val rotl :
   k, n, a[0  k  n  k + 3  k + 2  n].
   Avl (a, k)  a  Avl (a, n)  k[k  n + 1 
     n  k].Avl (a, k)

val add :
   n, a.
   a  Avl (a, n)  k[1  k  n  k  k  n + 1].Avl (a, k)

val mem : n, a. a  Avl (a, n)  Bool

val min_binding : n, a[1  n]. Avl (a, n)  a

val remove_min_binding :
   n, a[1  n].
   Avl (a, n)  k[n  k + 1  k  n  k + 2  2 n].Avl (a, k)

val merge :
   k, n, a[n  k + 2  k  n + 2].
   (Avl (a, n), Avl (a, k))  i[n  i  k  i  i  n + k 
     imax (k + 1, n + 1)].Avl (a, i)

val remove :
   n, a.
   a  Avl (a, n)  k[n  k + 1  0  k  k  n].Avl (a, k)
\end{tmcode}

\section{Arrays and Matrices}

For clarity, we present the array and matrix operations prelude separately
here. Arrays are polymorphic, matrices only store floating point numbers.
\begin{tmcode}
datatype Array : type * num
external let array_make :
  n, a [0n]. Num n  a  Array (a, n) = "fun a b -> Array.make a b"
external let array_get :
  n, k, a [0k  k+1n]. Array (a, n)  Num k  a = "fun a b -> Array.get a b"
external let array_set :
  n, k, a [0k  k+1n]. Array (a, n)  Num k  a  () =
  "fun a b c -> Array.set a b c"
external let array_length :
  n, a. Array (a, n)  Num n = "fun a -> Array.length a"

external type Matrix : num * num =
  "(int, Bigarray.int_elt, Bigarray.c_layout) Bigarray.Array2.t"
external let matrix_make :
  n, k [0n  0k]. Num n  Num k  Matrix (n, k) =
  "fun a b -> Bigarray.Array2.create Bigarray.int Bigarray.c_layout a b"
external let matrix_get :
  n, k, i, j [0i  i+1n  0j  j+1k].
   Matrix (n, k)  Num i  Num j  Int = "Bigarray.Array2.get"
external let matrix_set :
  n, k, i, j [0i  i+1n  0j  j+1k].
   Matrix (n, k)  Num i  Num j  Int  () = "Bigarray.Array2.set"
external let matrix_dim1 :
  n, k [0n  0k]. Matrix (n, k)  Num n = "Bigarray.Array2.dim1"
external let matrix_dim2 :
  n, k [0n  0k]. Matrix (n, k)  Num k = "Bigarray.Array2.dim2"
\end{tmcode}

\subsection{Program \tmverbatim{dotprod}}

Source file {\tmem{liquid\_dotprod.gadt}}:
\begin{tmcode}
external let add : Int  Int  Int = "fun n k -> n + k"
external let prod : Int  Int  Int = "fun n k -> n * k"
external let int0 : Int = "0"

let dotprod v1 v2 =
  let rec loop n sum i =
    if n <= i then sum else
        loop n (add (prod (array_get v1 i) (array_get v2 i)) sum)
          (i + 1) in
  loop (array_length v1) int0 0
\end{tmcode}
Value items from {\tmem{liquid\_dotprod.gadti.target}}:
\begin{tmcode}
val dotprod : k, n[k  n]. Array (Int, k)  Array (Int, n)  Int
\end{tmcode}

\subsection{Program \tmverbatim{bcopy}}

Source file {\tmem{liquid\_bcopy.gadt}}:
\begin{tmcode}
let rec bcopy_aux src des = function
  | i, m when m <= i -> ()
  | i, m when i+1 <= m ->
      let n = array_get src i in
      array_set des i n;
      let j = i + 1 in
      bcopy_aux src des (j, m)

let bcopy src des =
  let sz = array_length src in
  bcopy_aux src des (0, sz)
\end{tmcode}
Value items from {\tmem{liquid\_bcopy.gadti.target}}:
\begin{tmcode}
val bcopy_aux :
   i, j, k, n, a[0  n  k  i  k  j].
   Array (a, j)  Array (a, i)  (Num n, Num k)  ()

val bcopy : k, n, a[k  n]. Array (a, k)  Array (a, n)  ()
\end{tmcode}

\subsection{Program \tmverbatim{bsearch}}

Source file {\tmem{liquid\_bsearch\_harder.gadt}}:
\begin{tmcode}
datatype LinOrder
datacons LE : LinOrder
datacons GT : LinOrder
datacons EQ : LinOrder

external let compare : a. a  a  LinOrder =
  "fun a b -> let c = Pervasives.compare a b in
              if c < 0 then LE else if c > 0 then GT else EQ"
external let equal : a. a  a  Bool = "fun a b -> a = b"

external let div2 : n. Num (2 n)  Num n = "fun x -> x / 2"

datatype Answer : type
datacons NotFound : a. Answer a
datacons Found : a. a  Answer a

let bsearch key vec =
  let rec look lo hi =
    if lo <= hi then
        let m = div2 (hi + lo) in
        let x = array_get vec m in
        match compare key x with
          | LE -> look lo (m + (-1))
          | GT -> look (m + 1) hi
          | EQ -> if equal key x then Found x else NotFound
    else NotFound in
  look 0 (array_length vec + (-1))
\end{tmcode}
Value items from {\tmem{liquid\_bsearch\_harder.gadti.target}}:
\begin{tmcode}
val bsearch : n, a[0  n]. a  Array (a, n)  Answer a
\end{tmcode}
\subsection{Function \tmverbatim{bsearch2}}\label{ExBsearch2}

Source file {\tmem{liquid\_bsearch2\_harder3.gadt}}:
\begin{tmcode}
datatype LinOrder
datacons LE : LinOrder
datacons GT : LinOrder
datacons EQ : LinOrder

external let compare : a. a  a  LinOrder =
  "fun a b -> let c = Pervasives.compare a b in
              if c < 0 then LE else if c > 0 then GT else EQ"
external let equal : a. a  a  Bool = "fun a b -> a = b"

external let div2 : n. Num (2 n)  Num n = "fun x -> x / 2"

let bsearch key = efunction vec ->
  let rec look key vec lo hi =
    assert num -1 <= hi;
    eif lo <= hi then
        let m = div2 (hi + lo) in
        let x = array_get vec m in
        ematch compare key x with
          | LE -> look key vec lo (m + (-1))
          | GT -> look key vec (m + 1) hi
          | EQ -> eif equal key x then m else -1
    else -1 in
  look key vec 0 (array_length vec + (-1))
\end{tmcode}
Value items from {\tmem{liquid\_bsearch2\_harder3.gadti.target}}:
\begin{tmcode}
val bsearch :
   n, a[0  n].
   a  Array (a, n)  k[0  k + 1  k + 1  n].Num k
\end{tmcode}
Exported OCaml source {\tmem{liquid\_bsearch2\_harder3.ml.target}}:
\begin{tmcode}
type num = int
(** type _ array = built-in *)
let array_make : (*0  n*) ((* n *) num -> 'a -> ('a (* n *)) array) =
  fun a b -> Array.make a b
let array_get :
  (*0  k  k + 1  n*) (('a (* n *)) array -> (* k *) num -> 'a) =
  fun a b -> Array.get a b
let array_set :
  (*0  k  k + 1  n*)
  (('a (* n *)) array -> (* k *) num -> 'a -> unit) =
  fun a b c -> Array.set a b c
let array_length : (*0  n*) (('a (* n *)) array -> (* n *) num) =
  fun a -> Array.length a
type linOrder =
  | LE : linOrder
  | GT : linOrder
  | EQ : linOrder
let compare : ('a -> 'a -> linOrder) =
  fun a b -> let c = Pervasives.compare a b in
              if c  0 then GT else EQ
let equal : ('a -> 'a -> bool) = fun a b -> a = b
let div2 : ((* 2 n *) num -> (* n *) num) = fun x -> x / 2
type ex4 =
  | Ex4 : (*'k, 'n[0  k + 1  k + 1  n].*)(* k *) num ->
    (* n *) ex4
type ex3 =
  | Ex3 : (*'k, 'n[0  k + 1  k  n].*)(* k *) num -> (* n *) ex3
let bsearch
  : type (*n*) a (*0  n*). (a -> (a (* n *)) array -> (* n *) ex4) =
  (fun key vec ->
    let Ex3 xcase =
    let rec look :
    type (*i*) (*j*) (*k*) b (*0  k + 1  0  i 
    k + 1  j*). (b -> (b (* j *)) array -> (* i *) num -> (* k *) num ->
                     (* k *) ex3)
    =
    (fun key vec lo hi ->
      (if lo <= hi then
      let m = div2 (hi + lo) in
      let x = array_get vec m in
      (((match (compare key x: linOrder) with
      LE -> let Ex3 xcase = look key vec lo (m + -1) in Ex3 xcase
        | GT -> let Ex3 xcase = look key vec (m + 1) hi in Ex3 xcase
        | EQ ->
            (if equal key x then let xcase = m in Ex3 xcase else
            let xcase = -1 in Ex3 xcase))) :
      (* k *) ex3) else let xcase = -1 in Ex3 xcase)) in
    look key vec 0 (array_length vec + -1) in Ex4 xcase)
\end{tmcode}

\subsection{Program \tmverbatim{queen}}

Source file {\tmem{liquid\_queen.gadt}}:
\begin{tmcode}
external let n2i : n. Num n  Int = "fun i -> i"
external let equal : a. a  a  Bool = "fun x y -> x = y"
external let leq : a. a  a  Bool = "fun x y -> x <= y"

external let print : String  () = "print_string"
external let string_make : Int  String  String =
  "fun n s -> String.make n s.[0]"
external let abs : Int  Int = "fun i -> if i < 0 then ~-i else i"
external let minus : Int  Int  Int = "(-)"
external let plus : Int  Int  Int = "(+)"

let queens size =
  let board = array_make size (n2i 0) in
  let print_row pos =
    print (string_make (minus pos (n2i 1)) "."); print "Q";
    print (string_make (minus (n2i size) pos) ".") in
  let print_queens () =
    let rec aux row =
      if size <= row then ()
      else (print_row (array_get board row); aux (row + 1)) in
    aux 0 in
  let rec solved j =
    let q2j = array_get board j in
    let rec aux i =
      if i + 1 <= j then
        let q2i = array_get board i in
        let qdiff = abs (minus q2j q2i) in
        if equal q2i q2j then False
        else if equal qdiff (n2i (j - i)) then False
        else aux (i + 1)
      else True in
    aux 0 in
  let rec loop row =
    let next = plus (array_get board row) (n2i 1) in
    if leq (n2i (size + 1)) next then (
      array_set board row (n2i 0);
      if row <= 0 then () else loop (row - 1))
    else (
      array_set board row next;
      if solved row then
        if size <= row + 1 then (print_queens (); loop row)
        else loop (row + 1)
      else loop row) in
  loop 0
\end{tmcode}
Value items from {\tmem{liquid\_queen.gadti.target}}:
\begin{tmcode}
val queens : n[1  n]. Num n  ()
\end{tmcode}

\subsection{Function \tmverbatim{swap\_interval}}

Source file {\tmem{liquid\_vecswap.gadt}}:
\begin{tmcode}
let swap_interval arr start middle n =
  let rec item i = array_get arr i in
  let rec swap i j =
    let tmpj = item j in
    let tmpi = item i in
    array_set arr i tmpj; array_set arr j tmpi in
  let rec vecswap i j n =
    if n <= 0 then () else (
      swap i j; vecswap (i + 1) (j + 1) (n - 1)) in
  vecswap start middle n
\end{tmcode}
Value items from {\tmem{liquid\_vecswap.gadti.target}}:
\begin{tmcode}
val swap_interval :
   i, j, k, n, a[1  j  0  n  0  k  i + k  j 
   i + n  j]. Array (a, j)  Num n  Num k  Num i  ()
\end{tmcode}

\subsection{Program \tmverbatim{isort}}

Source file {\tmem{liquid\_isort\_harder.gadt}}:
\begin{tmcode}
datatype LinOrder
datacons LE : LinOrder
datacons GT : LinOrder
datacons EQ : LinOrder

external let compare : a. a  a  LinOrder =
  "fun a b -> let c = Pervasives.compare a b in
              if c < 0 then LE else if c > 0 then GT else EQ"
external let equal : a. a  a  Bool = "fun a b -> a = b"

let rec insertSort arr start n =
  let rec item i = array_get arr i in
  let rec swap i j =
    let tmpj = item j in
    let tmpi = item i in
    array_set arr i tmpj; array_set arr j tmpi in
  let rec outer i =
    if start + n <= i then ()
    else
      let rec inner j =
        if j <= start then outer (i + 1)
        else if equal (compare (item j) (item (j - 1))) LE
        then (swap j (j - 1); inner (j - 1))
        else outer (i + 1) in
      inner i in
  outer (start + 1)
\end{tmcode}
Value items from {\tmem{liquid\_isort\_harder.gadti.target}}:
\begin{tmcode}
val insertSort :
   i, k, n, a[0  k  1  i  k + n  i].
   Array (a, i)  Num k  Num n  ()
\end{tmcode}

\subsection{Program \tmverbatim{tower}}

Source file {\tmem{liquid\_tower.gadt}}:
\begin{tmcode}
external let n2i : n. Num n  Int = "fun i -> i"
external let equal : a. a  a  Bool = "fun x y -> x = y"
external let leq : a. a  a  Bool = "fun x y -> x <= y"

external let print : String  () = "print_string"
external let string_make : Int  String  String =
  "fun n s -> String.make n s.[0]"
external let string_of_int : Int  String = "string_of_int"

external let abs : Int  Int = "fun i -> if i < 0 then ~-i else i"
external let minus : Int  Int  Int = "(-)"
external let plus : Int  Int  Int = "(+)"
external let int0 : Int = "0"

let play sz =
  let leftPost = array_make sz int0 in
  let middlePost = array_make sz int0 in
  let rightPost = array_make sz int0 in

  let initialize post =
    let rec init_rec i =
      if i + 1 <= sz - 1 then (
        array_set post i (n2i (i+1));
        init_rec (i+1))
      else () in
    init_rec 0 in

  let showpiece n =
    let rec r_rec i =
      if leq (plus i (n2i 2)) n then (
        print " "; r_rec (plus i (n2i 1)))
      else () in
    let rec r2_rec j =
      if leq (plus j (n2i 1)) (n2i sz)
      then (print "#"; r2_rec (plus j (n2i 1)))
      else () in
    r_rec (n2i 1);
    r2_rec (plus n (n2i 1)) in

  let showposts () =
    let rec show_rec i =
      if i + 1 <= sz - 1 then (
        showpiece (array_get leftPost i); print " ";
        showpiece (array_get middlePost i); print " ";
        showpiece (array_get rightPost i); print "\textbackslash n";
        show_rec (i+1))
      else () in
    show_rec 0; print "\textbackslash n" in

  initialize leftPost;
  let rec move n source s post p post' p' =
    if n <= 1 then (
      array_set post (p - 1) (array_get source s);
      array_set source s int0; showposts ())
    else (
      move (n - 1) source s post' p' post p;
      array_set post (p - 1) (array_get source (s + n - 1));
      array_set source (s + n - 1) int0;
      showposts ();
      move (n - 1) post' ((p' - n) + 1) post (p - 1) source (s + n)) in

  showposts ();
  move sz leftPost 0 rightPost sz middlePost sz
\end{tmcode}
Value items from {\tmem{liquid\_tower.gadti.target}}:
\begin{tmcode}
val play : n[1  n]. Num n  ()
\end{tmcode}

\subsection{Program \tmverbatim{matmult}}

Source file {\tmem{liquid\_matmult.gadt}}:
\begin{tmcode}
external let n2i : n. Num n  Int = "fun i -> i"
external let equal : a. a  a  Bool = "fun x y -> x = y"
external let leq : a. a  a  Bool = "fun x y -> x <= y"

external let abs : Int  Int = "fun i -> if i < 0 then ~-i else i"
external let minus : Int  Int  Int = "(-)"
external let plus : Int  Int  Int = "(+)"
external let mult : Int  Int  Int = "( * )"
external let int0 : Int = "0"

let fillar arr2 fill =
  let d1 = matrix_dim1 arr2 in
  let d2 = matrix_dim2 arr2 in
  let rec loop i =
    if i + 1 <= d1 
    then
      let rec loopi j =
        if j + 1 <= d2 then (
          matrix_set arr2 i j (fill ());
          loopi (j + 1))
        else () in
      loopi 0
    else loop (i + 1) in
  loop 0

let matmul a b =
  let p = matrix_dim1 a in
  let q = matrix_dim2 a in
  let r = matrix_dim2 b in

  let cdata = matrix_make p r in
  let callback0 () = int0 in
  fillar cdata callback0;

  let rec loop1 i =
    if i + 1 <= p then (
      let rec loop2 j =
        if j + 1 <= r then (
          let rec loop3 k sum =
            if k + 1 <= q then (
              let sum_p =
                plus sum (mult (matrix_get a i k) (matrix_get b k j)) in
              loop3 (k + 1) sum_p)
            else sum in
          let l3 = loop3 0 int0 in
          matrix_set cdata i j l3;
          loop2 (j + 1))
        else () in
      loop2 0;
      loop1 (i + 1))
    else () in
  loop1 0;
  cdata
\end{tmcode}
Value items from {\tmem{liquid\_matmult.gadti.target}}:
\begin{tmcode}
val fillar :
   k, n[0  n  0  k]. Matrix (n, k)  (()  Int)  ()

val matmul :
   i, j, k, n[0  n  0  k  0  j  j  i].
   Matrix (n, j)  Matrix (i, k)  Matrix (n, k)
\end{tmcode}

\subsection{Program \tmverbatim{heapsort}}

Source file {\tmem{liquid\_heapsort.gadt}}:
\begin{tmcode}
external let div2 : n. Num (2 n)  Num n = "fun x -> x / 2"

external let n2i : n. Num n  Int = "fun i -> i"
external let equal : a. a  a  Bool = "fun x y -> x = y"
external let leq : a. a  a  Bool = "fun x y -> x <= y"
external let less : a. a  a  Bool = "fun x y -> x < y"

external let print : String  () = "print_string"
external let string_make : Int  String  String =
  "fun n s -> String.make n s.[0]"
external let string_of_int : Int  String = "string_of_int"

external let abs : Int  Int = "fun i -> if i < 0 then ~-i else i"
external let minus : Int  Int  Int = "(-)"
external let plus : Int  Int  Int = "(+)"
external let int0 : Int = "0"

let rec heapify size data i =
  let left = 2 * i + 1 in
  let right = 2 * i + 2 in
  let largest1 =
    eif left + 1 <= size then
      eif less (array_get data i) (array_get data left)
      then left else i
    else i in
  let largest2 =
    eif right + 1 <= size then
      eif less (array_get data largest1) (array_get data right) then right
      else largest1
    else largest1 in
  if i + 1 <= largest2 then
    let temp = array_get data i in
    let temp2 = array_get data largest2 in
    array_set data i temp2;
    array_set data largest2 temp;
    heapify size data largest2
  else ()
(* We do not get the constraint [i + 1 
   the last [else] branch is entered. *)

let buildheap size data =
  let rec loop i =
    if 0 <= i then (
      heapify size data i;
      loop (i - 1))
    else () in
  loop (div2 size - 1)

let heapsort maxx size data =
  buildheap size data;
  let rec loop i =
    if 1 <= i then
      let temp = array_get data i in
      array_set data i (array_get data 0);
      array_set data 0 temp;
      heapify i data 0;
      loop (i - 1)
    else () in
  loop (maxx - 1)

let print_array data i j =
  let rec loop k =
    if k + 1 <= j then (
      print (array_get data k);
      loop (k + 1))
    else () in
  loop i
\end{tmcode}
Value items from {\tmem{liquid\_heapsort.gadti.target}}:
\begin{tmcode}
val heapify :
   i, k, n, a[0  n  i  k].
   Num i  Array (a, k)  Num n  ()

val buildheap : k, n, a[k  n]. Num k  Array (a, n)  ()

val heapsort :
   i, k, n, a[k  n  i  k].
   Num i  Num k  Array (a, n)  ()

val print_array :
   i, k, n[k  i  0  n].
   Array (String, i)  Num n  Num k  ()
\end{tmcode}

\subsection{Program \tmverbatim{simplex}}

Source file {\tmem{liquid\_simplex.gadt}}:
\begin{tmcode}
external let n2f : n. Num n  Float = "float_of_int"
external let equal : a. a  a  Bool = "fun x y -> x = y"
external let leq : a. a  a  Bool = "fun x y -> x <= y"
external let less : a. a  a  Bool = "fun x y -> x < y"

external let minus : Float  Float  Float = "(-.)"
external let plus : Float  Float  Float = "(+.)"
external let mult : Float  Float  Float = "( *. )"
external let div : Float  Float  Float = "( /. )"
external let fl0 : Float = "0.0"

(* step 1 *)

let rec is_neg_aux a j =
  if j + 2 <= matrix_dim2 a then
    if less (matrix_get a 0 j) fl0 then True
    else is_neg_aux a (j+1)
  else False

let is_neg a = is_neg_aux a 1

(* step 2 *)

let rec unb1 a i j =
  let rec unb2 a i j =
    if i + 1 <= matrix_dim1 a then
      if less (matrix_get a i j) fl0
      then unb2 a (i+1) j
      else unb1 a 0 (j+1)
    else True in

  if j + 2 <= matrix_dim2 a then
    if less (matrix_get a 0 j) fl0
    then unb2 a (i+1) j
    else unb1 a 0 (j+1)
  else False

(* step 6 *)

let rec norm_aux a i c j =
  if j + 1 <= matrix_dim2 a then (
    matrix_set a i j (div (matrix_get a i j) c);
    norm_aux a i c (j+1))
  else ()

let rec norm a i j =
  let c = matrix_get a i j in
  norm_aux a i c 1

let rec row_op_aux1 a i i' c j =
  if j + 1 <= matrix_dim2 a then
    matrix_set a i' j
      (minus (matrix_get a i' j)
        (mult (matrix_get a i j) c));
    row_op_aux1 a i i' c (j+1)
  else ()

let rec row_op_aux2 a i i' j =
  row_op_aux1 a i i' (matrix_get a i' j) 1

let rec row_op_aux3 a i j i' =
  if i' + 1 <= matrix_dim1 a then
    if i' <= i && i <= i' then (
      row_op_aux2 a i i' j;
      row_op_aux3 a i j (i'+1))
    else row_op_aux3 a i j (i'+1)
  else ()

let rec row_op a i j =
  norm a i j;
  row_op_aux3 a i j 0

let rec simplex a =

  (* step 3 *)

  let rec enter_var j c j' =
    eif j' + 2 <= matrix_dim2 a then
      let c' = matrix_get a 0 j' in
      eif less c' c then enter_var j' c' (j'+1)
      else enter_var j c (j'+1)
    else j in

  (* step 4 *)

  let rec depart_var j i r i' =
    eif i' + 1 <= matrix_dim1 a then
      let c' = matrix_get a i' j in
      eif less fl0 c' then
        let r' = div (matrix_get a i' (matrix_dim2 a + (-1))) c' in
        eif less r' r then depart_var j i' r' (i'+1)
        else depart_var j i r (i'+1)
      else depart_var j i r (i'+1)
    else i in

  (* step 5 *)

  let rec init_ratio j i =
    eif i + 1 <= matrix_dim1 a then
      let c = matrix_get a i j in
      eif less fl0 c then i, div (matrix_get a i (matrix_dim2 a + (-1))) c
      else init_ratio j (i+1)
    else runtime_failure "init_ratio: no variable found" in

  (* step 7 *)

  if is_neg a then
    if unb1 a 0 1 then () (* assert false *)
    else
      let j = enter_var 1 (matrix_get a 0 1) 2 in
      let i, r = init_ratio j 1 in
      let i = depart_var j i r (i+1) in
      row_op a i j;
      simplex a
  else ()
\end{tmcode}
The file {\tmem{liquid\_simplex\_harder.gadt}} has the above steps in order as
part of a single toplevel definition. Value items from
{\tmem{liquid\_simplex.gadti.target}}:
\begin{tmcode}
val is_neg_aux :
   i, k, n[0  n  1  k  0  i].
   Matrix (k, i)  Num n  Bool

val is_neg : k, n[1  n  0  k]. Matrix (n, k)  Bool

val unb1 :
   i, j, k, n[0  n  0  k + 1  1  i  0  j].
   Matrix (i, j)  Num k  Num n  Bool

val norm_aux :
   i, j, k, n[0  n  0  k  k + 1  i  0  j].
   Matrix (i, j)  Num k  Float  Num n  ()

val norm :
   i, j, k, n[0  n  0  k  k + 1  i  n + 1  j].
   Matrix (i, j)  Num k  Num n  ()

val row_op_aux1 :
   i, j, k, m, n[0  n  0  k  0  i  i + 1  j 
   k + 1  j  0  m].
   Matrix (j, m)  Num i  Num k  Float  Num n  ()

val row_op_aux2 :
   i, j, k, m, n[0  n  0  k  0  i  i + 1  j 
   k + 1  j  n + 1  m].
   Matrix (j, m)  Num i  Num k  Num n  ()

val row_op_aux3 :
   i, j, k, m, n[0  k  0  i  0  j  k + 1  m].
   Matrix (j, m)  Num i  Num k  Num n  ()

val row_op :
   i, j, k, n[0  n  0  k  k + 1  i  n + 1  j].
   Matrix (i, j)  Num k  Num n  ()

val simplex : k, n[1  n  2  k]. Matrix (n, k)  ()
\end{tmcode}
Exported OCaml source {\tmem{liquid\_simplex.ml.target}}:
\begin{tmcode}
type num = int
type matrix =
  (float, Bigarray.float64_elt, Bigarray.c_layout) Bigarray.Array2.t
let matrix_make :
  (*0  n  0  k*) ((* n *) num -> (* k *) num -> (* n, k *) matrix) =
  fun a b -> Bigarray.Array2.create Bigarray.float64 Bigarray.c_layout a b
let matrix_get :
  (*0  i  i + 1  n  0  j  j + 1  k*)
  ((* n, k *) matrix -> (* i *) num -> (* j *) num -> float) =
  Bigarray.Array2.get
let matrix_set :
  (*0  i  i + 1  n  0  j  j + 1  k*)
  ((* n, k *) matrix -> (* i *) num -> (* j *) num -> float -> unit) =
  Bigarray.Array2.set
let matrix_dim1 :
  (*0  n  0  k*) ((* n, k *) matrix -> (* n *) num) =
  Bigarray.Array2.dim1
let matrix_dim2 :
  (*0  n  0  k*) ((* n, k *) matrix -> (* k *) num) =
  Bigarray.Array2.dim2
let n2f : ((* n *) num -> float) = float_of_int
let equal : ('a -> 'a -> bool) = fun x y -> x = y
let leq : ('a -> 'a -> bool) = fun x y -> x 
let less : ('a -> 'a -> bool) = fun x y -> x 
let minus : (float -> float -> float) = (-.)
let plus : (float -> float -> float) = (+.)
let mult : (float -> float -> float) = ( *. )
let div : (float -> float -> float) = ( /. )
let fl0 : float = 0.0
let rec is_neg_aux :
  (*type i k n [0  n  1  k 
    0  i].*) ((* k, i *) matrix -> (* n *) num -> bool) =
  (fun a j ->
    (if j + 2 <= matrix_dim2 a then
    (if less (matrix_get a 0 j) fl0 then true else is_neg_aux a (j + 1)) else
    false))
let is_neg
  : (*type k n [1  n  0  k].*) ((* n, k *) matrix -> bool) =
  (fun a -> is_neg_aux a 1)
let rec unb1 :
  (*type i j k n [0  n  0  k + 1  1  i 
    0  j].*) ((* i, j *) matrix -> (* k *) num -> (* n *) num -> bool) =
  (fun a i j ->
    let rec unb2 :
    (*type k1 l m n1 [0  m  0  l  0  n1 
      m + 1  k1].*) ((* n1, k1 *) matrix -> (* l *) num -> (* m *) num ->
                         bool)
    =
    (fun a i j ->
      (if i + 1 <= matrix_dim1 a then
      (if less (matrix_get a i j) fl0 then unb2 a (i + 1) j else
      unb1 a 0 (j + 1)) else true)) in
    (if j + 2 <= matrix_dim2 a then
    (if less (matrix_get a 0 j) fl0 then unb2 a (i + 1) j else
    unb1 a 0 (j + 1)) else false))
let rec norm_aux :
  (*type i j k n [0  n  0  k  k + 1  i 
    0  j].*) ((* i, j *) matrix -> (* k *) num -> float -> (* n *) num ->
                  unit) =
  (fun a i c j ->
    (if j + 1 <= matrix_dim2 a then
    (matrix_set a i j (div (matrix_get a i j) c) ; norm_aux a i c (j + 1))
    else ()))
let rec norm :
  (*type i j k n [0  n  0  k  k + 1  i 
    n + 1  j].*) ((* i, j *) matrix -> (* k *) num -> (* n *) num -> unit) =
  (fun a i j -> let c = matrix_get a i j in norm_aux a i c 1)
let rec row_op_aux1 :
  (*type i j k m n [0  n  0  k  0  i  i + 1  j 
    k + 1  j 
    0  m].*) ((* j, m *) matrix -> (* i *) num -> (* k *) num -> float ->
                  (* n *) num -> unit) =
  (fun a i i' c j ->
    (if j + 1 <= matrix_dim2 a then
    (matrix_set a i' j
       (minus (matrix_get a i' j) (mult (matrix_get a i j) c))
    ; row_op_aux1 a i i' c (j + 1)) else ()))
let rec row_op_aux2 :
  (*type i j k m n [0  n  0  k  0  i  i + 1  j 
    k + 1  j 
    n + 1  m].*) ((* j, m *) matrix -> (* i *) num -> (* k *) num ->
                      (* n *) num -> unit) =
  (fun a i i' j -> row_op_aux1 a i i' (matrix_get a i' j) 1)
let rec row_op_aux3 :
  (*type i j k m n [0  k  0  i  0  j 
    k + 1  m].*) ((* j, m *) matrix -> (* i *) num -> (* k *) num ->
                      (* n *) num -> unit) =
  (fun a i j i' ->
    (if i' + 1 <= matrix_dim1 a then
    (if i <= i'&& i' <= i then
    (row_op_aux2 a i i' j ; row_op_aux3 a i j (i' + 1)) else
    row_op_aux3 a i j (i' + 1)) else ()))
let rec row_op :
  (*type i j k n [0  n  0  k  k + 1  i 
    n + 1  j].*) ((* i, j *) matrix -> (* k *) num -> (* n *) num -> unit) =
  (fun a i j -> (norm a i j ; row_op_aux3 a i j 0))
type ex7 =
  | Ex7 : (*'i, 'k, 'n[k  i  i + 1  n].*)((* i *) num * float) ->
    (* n, k *) ex7
type ex5 =
  | Ex5 : (*'i, 'k, 'n[0  k  kmax (i, n + -1)].*)(* k *) num ->
    (* n, i *) ex5
type ex2 =
  | Ex2 : (*'i, 'k, 'n[0  k  kmax (i, n + -1)].*)(* k *) num ->
    (* n, i *) ex2
let rec simplex :
  (*type k n [1  n  2  k].*) ((* n, k *) matrix -> unit) =
  (fun a ->
    let rec enter_var :
    (*type i j [0  j 
      0  i].*) ((* i *) num -> float -> (* j *) num -> (* k, i *) ex2)
    =
    (fun j c j' ->
      (if j' + 2 <= matrix_dim2 a then
      let c' = matrix_get a 0 j' in
      (if less c' c then
      let Ex2 xcase = enter_var j' c' (j' + 1) in Ex2 xcase else
      let Ex2 xcase = enter_var j c (j' + 1) in Ex2 xcase) else
      let xcase = j in Ex2 xcase)) in
    let rec depart_var :
    (*type l m n1 [0  l  0  n1  n1 + 1  k 
      0  m].*) ((* n1 *) num -> (* m *) num -> float -> (* l *) num ->
                    (* n, m *) ex5)
    =
    (fun j i r i' ->
      (if i' + 1 <= matrix_dim1 a then
      let c' = matrix_get a i' j in
      (if less fl0 c' then
      let r' = div (matrix_get a i' (matrix_dim2 a + -1)) c' in
      (if less r' r then
      let Ex5 xcase = depart_var j i' r' (i' + 1) in Ex5 xcase else
      let Ex5 xcase = depart_var j i r (i' + 1) in Ex5 xcase) else
      let Ex5 xcase = depart_var j i r (i' + 1) in Ex5 xcase) else
      let xcase = i in Ex5 xcase)) in
    let rec init_ratio :
    (*type i1 k1 [0  k1  0  i1 
      i1 + 1  k].*) ((* i1 *) num -> (* k1 *) num -> (* n, k1 *) ex7)
    =
    (fun j i ->
      (if i + 1 <= matrix_dim1 a then
        let c = matrix_get a i j in
        (if less fl0 c then
        let xcase = (i, div (matrix_get a i (matrix_dim2 a + -1)) c) in
        Ex7 xcase else let Ex7 xcase = init_ratio j (i + 1) in Ex7 xcase)
      else
        (failwith "init_ratio: no variable found"))) in
    (if is_neg a then
    (if unb1 a 0 1 then () else
      let Ex2 j = enter_var 1 (matrix_get a 0 1) 2 in
      let Ex7 (i, r) = init_ratio j 1 in
      let Ex5 i = depart_var j i r (i + 1) in (row_op a i j ; simplex a))
    else
      ()))
\end{tmcode}

\subsection{Program \tmverbatim{gauss}}

Source file {\tmem{liquid\_gauss.gadt}}:
\begin{tmcode}
external let n2f : n. Num n  Float = "float_of_int"
external let equal : a. a  a  Bool = "fun x y -> x = y"
external let leq : a. a  a  Bool = "fun x y -> x <= y"
external let less : a. a  a  Bool = "fun x y -> x < y"

external let minus : Float  Float  Float = "(-.)"
external let plus : Float  Float  Float = "(+.)"
external let mult : Float  Float  Float = "( *. )"
external let div : Float  Float  Float = "( /. )"
external let fabs : Float  Float = "abs_float"
external let fl0 : Float = "0.0"
external let fl1 : Float = "1.0"

let getRow data i =
  let stride = matrix_dim2 data in
  let rowData = array_make stride fl0 in
  let rec extract j =
    if j + 1 <= stride then (
      array_set rowData j (matrix_get data i j);
      (* lukstafi: the call below missing in the original source? *)
      extract (j + 1))
    else () in
  extract 0;
  rowData

let putRow data i row =
  let stride = array_length row in
  let rec put j =
    if j + 1 <= stride then (
      matrix_set data i j (array_get row j);
      (* lukstafi: the call below missing in the original source? *)
      put (j + 1))
    else () in
  put 0

let rowSwap data i j =
  let temp = getRow data i in
  putRow data i (getRow data j);
  putRow data j temp

let norm r n i =
  let x = array_get r i in
  array_set r i fl1;
  let rec loop k =
    if k + 1 <= n then (
      array_set r k (div (array_get r k) x);
      loop (k+1))
    else () in
  loop (i+1)

let rowElim r s n i =
  let x = array_get s i in
  array_set s i fl0;
  let rec loop k =
    if k + 1 <= n then (
      array_set s k (minus (array_get s k) (mult x (array_get r k)));
      loop (k+1))
    else () in
  loop (i+1)

let gauss data =
  let n = matrix_dim1 data in
  let m = matrix_dim2 data in

  let rec rowMax i j x mx =
    eif j + 1 <= n then
      let y = fabs (matrix_get data j i) in
      eif (less x y) then rowMax i (j+1) y j
      else rowMax i (j+1) x mx
    else mx in

  let rec loop1 i =
    if i + 1 <= n then
      let x = fabs (matrix_get data i i) in
      let mx = rowMax i (i+1) x i in
      norm (getRow data mx) m i;
      rowSwap data i mx;
      let rec loop2 j =
        if j + 1 <= n then (
          rowElim (getRow data i) (getRow data j) m i;
          loop2 (j+1))
        else () in
      loop2 (i+1);
      loop1 (i+1)
    else () in
  loop1 0
\end{tmcode}
Value items from {\tmem{liquid\_gauss.gadti.target}}:
\begin{tmcode}
val getRow :
   i, k, n[0  n  0  k  k + 1  i].
   Matrix (i, n)  Num k  Array (Float, n)

val putRow :
   i, j, k, n[0  n  0  k  k + 1  i  n  j].
   Matrix (i, j)  Num k  Array (Float, n)  ()

val rowSwap :
   i, j, k, n[0  n  0  k  k + 1  i  n + 1  i 
   0  j]. Matrix (i, j)  Num k  Num n  ()

val norm :
   i, k, n[0  n  n + 1  i  k  i].
   Array (Float, i)  Num k  Num n  ()

val rowElim :
   i, j, k, n[0  n  n + 1  i  k  i  k  j].
   Array (Float, j)  Array (Float, i)  Num k  Num n  ()

val gauss : k, n[0  n  1  k  n  k]. Matrix (n, k)  ()
\end{tmcode}
The last toplevel definition from {\tmem{liquid\_gauss2.gadt}}, requires
\tmverbatim{-prefer\_bound\_to\_local}:
\begin{tmcode}
let gauss data =
  let n = matrix_dim1 data in

  let rec rowMax i j x mx =
    eif j + 1 
      let y = fabs (matrix_get data j i) in
      eif (less x y) then rowMax i (j+1) y j
      else rowMax i (j+1) x mx
    else mx in

  let rec loop1 i =
    if i + 1 
      let x = fabs (matrix_get data i i) in
      let mx = rowMax i (i+1) x i in
      norm (getRow data mx) (n+1) i;
      rowSwap data i mx;
      let rec loop2 j =
        if j + 1 
          rowElim (getRow data i) (getRow data j) (n+1) i;
          loop2 (j+1))
        else () in
      loop2 (i+1);
      loop1 (i+1)
    else () in
  loop1 0
\end{tmcode}
The inferred types do not differ from the {\tmem{liquid\_gauss.gadt}} case.

The last toplevel definition from
{\tmem{liquid\_gauss\_harder\_asserted.gadt}}:
\begin{tmcode}
let gauss data =
  let n = matrix_dim1 data in

  let rec rowMax = efunction i ->
    let x = fabs (matrix_get data i i) in
    let rec loop j x mx =
      assert num mx + 1 
      eif j + 1 
        let y = fabs (matrix_get data j i) in
        eif (less x y) then loop (j+1) y j
        else loop (j+1) x mx
      else mx in
    loop (i+1) x i in

  let rec loop1 i =
    if i + 1 
      let mx = rowMax i in
      norm (getRow data mx) (n+1) i;
      rowSwap data i mx;
      let rec loop2 j =
        if j + 1 
          rowElim (getRow data i) (getRow data j) (n+1) i;
          loop2 (j+1))
        else () in
      loop2 (i+1);
      loop1 (i+1)
    else () in
  loop1 0
\end{tmcode}
The inferred types do not differ from the {\tmem{liquid\_gauss.gadt}} case.

\subsection{Program \tmverbatim{fft}}\label{FFTExample}

Source file {\tmem{liquid\_fft\_full.gadt}}:
\begin{tmcode}
external let n2i : n. Num n  Int = "fun i -> i"
external let div2 : n. Num (2 n)  Num n = "fun i -> i / 2"
external let div4 : n. Num (4 n)  Num n = "fun i -> i / 4"
external let n2f : n. Num n  Float = "float_of_int"
external let equal : a. a  a  Bool = "fun x y -> x = y"
external let leq : a. a  a  Bool = "fun x y -> x <= y"
external let less : a. a  a  Bool = "fun x y -> x < y"
external let ignore : a. a  () = "ignore"

external let abs : Float  Float = "abs_float"
external let cos : Float  Float = "cos"
external let sin : Float  Float = "sin"
external let neg : Float  Float = "(~-.)"
external let minus : Float  Float  Float = "(-.)"
external let plus : Float  Float  Float = "(+.)"
external let mult : Float  Float  Float = "( *. )"
external let div : Float  Float  Float = "( /. )"
external let fl0 : Float = "0.0"
external let fl05 : Float = "0.5"
external let fl1 : Float = "1.0"
external let fl2 : Float = "2.0"
external let fl3 : Float = "3.0"
external let fl4 : Float = "4.0"
external let pi : Float = "4.0 *. atan 1.0"
external let two_pi : Float = "8.0 *. atan 1.0"

datatype Bounded : num * num
datacons Index : i, k, n[n  i  i  k].Num i  Bounded (n, k)

let ffor s d body =
  let rec loop i =
    if i 
  loop s

external let ref : a. a  Ref a = "fun a -> ref a"
external let asgn : a. Ref a  a  () = "fun a b -> a := b"
external let deref : a. Ref a  a = "(!)"

let fft px py = (* n must be a power of 2! *)
  let n = array_length px + (-1) in
  let rec loop n2 n4 =
    if n2 
    let e = div two_pi (n2f n2) in
    let e3 = mult fl3 e in
    let a = ref fl0 in
    let a3 = ref fl0 in
    let rec forbod j' =
    match j' with Index j ->
    (*for j = 1 to n4 do*)
      let cc1 = cos (deref a) in
      let ss1 = sin (deref a) in
      let cc3 = cos (deref a3) in
      let ss3 = sin (deref a3) in
      asgn a (plus (deref a) e);
      asgn a3 (plus (deref a3) e3);
      let rec loop1 i0 i1 i2 i3 id =
        if n + 1 
        let g_px_i0 = array_get px i0 in
        let g_px_i2 = array_get px i2 in
        let r1 = minus g_px_i0 g_px_i2 in
        let r1' = plus g_px_i0 g_px_i2 in
        array_set px i0 r1';
        let g_px_i1 = array_get px i1 in
        let g_px_i3 = array_get px i3 in
        let r2 = minus g_px_i1 g_px_i3 in
        let r2' = plus g_px_i1 g_px_i3 in
        array_set px i1 r2';
        let g_py_i0 = array_get py i0 in
        let g_py_i2 = array_get py i2 in
        let s1 = minus g_py_i0 g_py_i2 in
        let s1' = plus g_py_i0 g_py_i2 in
        array_set py i0 s1';
        let g_py_i1 = array_get py i1 in
        let g_py_i3 = array_get py i3 in
        let s2 = minus g_py_i1 g_py_i3 in
        let s2' = plus g_py_i1 g_py_i3 in
        array_set py i1 s2';
        let s3 = minus r1 s2 in 
        let r1 = plus r1 s2 in
        let s2 = minus r2 s1 in
        let r2 = plus r2 s1 in
        array_set px i2 (minus (mult r1 cc1) (mult s2 ss1));
        array_set py i2 (minus (mult (neg s2) cc1) (mult r1 ss1));
        array_set px i3 (plus (mult s3 cc3) (mult r2 ss3));
        array_set py i3 (minus (mult r2 cc3) (mult s3 ss3));
        loop1 (i0 + id) (i1 + id) (i2 + id) (i3 + id) id in
      let rec loop2 is id =
        if n 
        else (
          let i1 = is + n4 in
          let i2 = i1 + n4 in
          let i3 = i2 + n4 in
          loop1 is i1 i2 i3 id;
          loop2 (2 * id + j - n2) (4 * id)) in
      loop2 j (2 * n2) in
    ffor 1 n4 forbod; 
    loop (div2 n2) (div2 n4) in
  loop n (div4 n);

  let rec loop1 i0 i1 id =
    if n + 1 
    let r1 = array_get px i0 in
    array_set px i0 (plus r1 (array_get px i1));
    array_set px i1 (minus r1 (array_get px i1));
    let r1 = array_get py i0 in
    array_set py i0 (plus r1 (array_get py i1));
    array_set py i1 (minus r1 (array_get py i1));
    loop1 (i0 + id) (i1 + id) id in
  let rec loop2 is id =
    if n 
      loop1 is (is + 1) id;
      loop2 (2 * id + (-1)) (4 * id)) in
  loop2 1 4;

  let rec loop1 j k =
    eif j 
    else loop1 (j - k) (div2 k) in

  let rec loop2 i j =
    if n 
      if j 
        let xt = array_get px j in
        array_set px j (array_get px i); array_set px i (xt);
        let xt = array_get py j in
        array_set py j (array_get py i); array_set py i (xt));
      let j' = loop1 j (div2 n) in
      loop2 (i + 1) j') in
  loop2 1 1; n

let ffttest np =
  let enp = n2f np in
  let n2 = div2 np in
  let npm = n2 - 1 in
  let pxr = array_make (np+1) fl0 in
  let pxi = array_make (np+1) fl0 in
  let t = div pi enp in
  array_set pxr 1 (mult (minus enp fl1) fl05);
  array_set pxi 1 (fl0);
  array_set pxr (n2+1) (neg (mult fl1 fl05));
  array_set pxi (n2+1) fl0;
  let rec forbod i =
    if i 
      let j = np - i in
      array_set pxr (i+1) (neg (mult fl1 fl05));
      array_set pxr (j+1) (neg (mult fl1 fl05));
      let z = mult t (n2f i) in
      let y = mult (div (cos z) (sin z)) fl05 in
      array_set pxi (i+1) (neg y); array_set pxi (j+1) (y)
    else () in
  forbod 1;
  ignore (fft pxr pxi);
  (* lukstafi: kr and ki are placeholders? *)
  let rec loop i zr zi kr ki =
    if np 
      let a = abs (minus (array_get pxr (i+1)) (n2f i)) in
      let b = less zr a in
      let zr = if b then a else zr in
      let kr = eif b then i else kr in
      let a = abs (array_get pxi (i+1)) in
      let b = less zi a in
      let zi = if b then a else zi in
      let ki = eif b then i else ki in
      loop (i+1) zr zi kr ki in
  let zr, zi = loop 0 fl0 fl0 0 0 in
  let zm = if less (abs zr) (abs zi) then zi else zr in
  (*in print_float zm; print_newline ()*) zm

let rec loop_np i np =
  if 17 
  ( ignore (ffttest np); loop_np (i + 1) (2 * np) )

let doit _ = loop_np 4 16
\end{tmcode}
Value items from {\tmem{liquid\_fft\_full.gadti.target}}:
\begin{tmcode}
val ffor :
   i, j, k, n[j  i  k  n].
   Num n  Num j  (Bounded (k, i)  ())  ()

val fft :
   k, n[1  n  n + 1  k].
   Array (Float, n + 1)  Array (Float, k)  Num n

val ffttest : n[2  n]. Num n  Float

val loop_np : k, n[2  n]. Num k  Num n  ()

val doit : a. a  ()
\end{tmcode}

\chapter{{\tmname{InvarGenT}}: Manual}

\section{Introduction}

Type systems are an established natural deduction-style means to reason about
programs. Dependent types can represent arbitrarily complex properties as they
use the same language for both types and programs, the type of value returned
by a function can itself be a function of the argument. Generalized Algebraic
Data Types bring some of that expressivity to type systems that deal with
data-types. Type systems with GADTs introduce the ability to reason about
return type by case analysis of the input value, while keeping the benefits of
a simple semantics of types, for example deciding equality between types can
be very simple. Existential types hide some information conveyed in a type,
usually when that information cannot be reconstructed in the type system. A
part of the type will often fail to be expressible in the simple language of
types, when the dependence on the input to the program is complex. GADTs
express existential types by using local type variables for the hidden parts
of the type encapsulated in a GADT.

The {\tmname{InvarGenT}} type system for GADTs differs from more pragmatic
approaches in mainstream functional languages in that we do not require any
type annotations on expressions, even on recursive functions. The
implementation also includes linear equations and inequalities over rational
numbers in the language of types, with the possibility to introduce more
domains in the future.

\section{Tutorial}

The concrete syntax of {\tmname{InvarGenT}} is similar to that of OCaml.
However, it does not currently cover records, the module system, objects, and
polymorphic variant types. It supports higher-order functions, algebraic
data-types including built-in tuple types, and linear pattern matching. It
supports conjunctive patterns using the \tmverbatim{as} keyword, but it
currently does not support disjunctive patterns. It currently has limited
support for guarded patterns: after \tmverbatim{when}, only inequality
\tmverbatim{<=} between values of the \tmverbatim{Num} type are allowed.

The sort of a type variable is identified by the first letter of the variable.
\tmverbatim{a}, \tmverbatim{b}, \tmverbatim{c}, \tmverbatim{r},
\tmverbatim{s}, \tmverbatim{t}, \tmverbatim{a1},... are in the sort of terms
called \tmverbatim{type}, i.e. ``types proper''. \tmverbatim{i},
\tmverbatim{j}, \tmverbatim{k}, \tmverbatim{l}, \tmverbatim{m},
\tmverbatim{n}, \tmverbatim{i1},... are in the sort of linear arithmetics over
rational numbers called \tmverbatim{num}. Remaining letters are reserved for
sorts that may be added in the future. Value constructors (like in OCaml) and
type constructors (unlike in OCaml) have the same syntax: capitalized name
followed by a tuple of arguments. They are introduced by \tmverbatim{datatype}
and \tmverbatim{datacons} respectively. The \tmverbatim{datatype} declaration
might be misleading in that it only lists the sorts of the arguments of the
type, the resulting sort is always \tmverbatim{type}. Values assumed into the
environment are introduced by \tmverbatim{external}. There is a built-in type
corresponding to declaration \tmverbatim{datatype Num : num} and definitions
of numeric constants \tmverbatim{newcons 0 : Num 0 newcons 1 : Num 1}... The
programmer can use \tmverbatim{external} declarations to give the semantics of
choice to the \tmverbatim{Num} data-type. The type with additional support as
\tmverbatim{Num} is the integers.

When solving negative constraints, arising from \tmverbatim{assert false}
clauses, we assume that the intended domain of the sort \tmverbatim{num} is
integers. This is a workaround to the lack of strict inequality in the sort
\tmverbatim{num}. We do not make the whole sort \tmverbatim{num} an integer
domain because it would complicate the algorithms.

In examples here we use Unicode characters. For ASCII equivalents, take a
quick look at the tables in the following section.

We start with a simple example, a function that can compute a value from a
representation of an expression -- a ready to use value whether it be
\tmverbatim{Int} or \tmverbatim{Bool}. Prior to the introduction of GADT
types, we could only implement a function \tmverbatim{eval : $\forall$a. Term
a $\rightarrow$ Value} where, using OCaml syntax, \tmverbatim{type value = Int
of int \textbar  Bool of bool}.
\begin{tmcode}
datatype Term : type
external let plus : Int  Int  Int = "(+)"
external let is_zero : Int  Bool = "(=) 0"
datacons Lit : Int  Term Int
datacons Plus : Term Int * Term Int  Term Int
datacons IsZero : Term Int  Term Bool
datacons If : a. Term Bool * Term a * Term a  Term a

let rec eval = function
  | Lit i -> i
  | IsZero x -> is_zero (eval x)
  | Plus (x, y) -> plus (eval x) (eval y)
  | If (b, t, e) -> if eval b then eval t else eval e
\end{tmcode}
Let us look at the corresponding generated, also called {\tmem{exported}},
OCaml source code:
\begin{tmcode}
type _ term =
  | Lit : int -> int term
  | Plus : int term * int term -> int term
  | IsZero : int term -> bool term
  | If : (*'a.*)bool term * 'a term * 'a term -> 'a term
  
let plus : (int -> int -> int) = (+)
let is_zero : (int -> bool) = (=) 0
let rec eval : type a . (a term -> a) =
  (function Lit i -> i | IsZero x -> is_zero (eval x)
    | Plus (x, y) -> plus (eval x) (eval y)
    | If (b, t, e) -> (if eval b then eval t else eval e))
\end{tmcode}
The \tmverbatim{Int}, \tmverbatim{Num} and \tmverbatim{Bool} types are
built-in. \tmverbatim{Int} and \tmverbatim{Bool} follow the general scheme of
exporting a datatype constructor with the same name, only lower-case. However,
numerals \tmverbatim{0}, \tmverbatim{1}, ... are always type-checked as
\tmverbatim{Num 0}, \tmverbatim{Num 1}... \tmverbatim{Num} can also be
exported as a type other than \tmverbatim{int}, and then numerals are exported
via an injection function (ending with) \tmverbatim{of\_int}.

The syntax \tmverbatim{external let} allows us to name an OCaml library
function or give an OCaml definition which we opt-out from translating to
{\tmname{InvarGenT}}. Such a definition will be verified against the rest of
the program when {\tmname{InvarGenT}} calls \tmverbatim{ocamlc -c} to verify
the exported code. Another variant of \tmverbatim{external} (omitting the
\tmverbatim{let} keyword) exports a value using \tmverbatim{external} in OCaml
code, which is OCaml source declaration of the foreign function interface of
OCaml. When we are not interested in linking and running the exported code, we
can omit the part starting with the \tmverbatim{=} sign. The exported code
will reuse the name in the FFI definition: \tmverbatim{external f :
}...\tmverbatim{ = "f"}.

The type inferred for the above example is \tmverbatim{eval : $\forall$a. Term
a$\rightarrow$a}. GADTs make it possible to reveal that \tmverbatim{IsZero x}
is a \tmverbatim{Term Bool} and therefore the result of \tmverbatim{eval}
should in its case be \tmverbatim{Bool}, \tmverbatim{Plus (x, y)} is a
\tmverbatim{Term Num} and the result of \tmverbatim{eval} should in its case
be \tmverbatim{Num}, etc. The
\tmverbatim{if/eif}$\ldots$\tmverbatim{then}$\ldots$\tmverbatim{else}$\ldots$
syntax is a syntactic sugar for
\tmverbatim{match$/$ematch}$\ldots$\tmverbatim{with True
->}$\ldots$\tmverbatim{ \textbar  False ->}$\ldots$, and any such expressions
are exported using \tmverbatim{if} expressions.

\tmverbatim{equal} is a function comparing values provided representation of
their types:
\begin{tmcode}
datatype Ty : type
datatype Int
datatype List : type
datacons Zero : Int
datacons Nil : a. List a
datacons TInt : Ty Int
datacons TPair : a, b. Ty a * Ty b  Ty (a, b)
datacons TList : a. Ty a  Ty (List a)
datatype Boolean
datacons True : Boolean
datacons False : Boolean
external eq_int : Int  Int  Bool
external b_and : Bool  Bool  Bool
external b_not : Bool  Bool
external forall2 : a, b. (a  b  Bool)  List a  List b  Bool

let rec equal = function
  | TInt, TInt -> fun x y -> eq_int x y
  | TPair (t1, t2), TPair (u1, u2) ->  
    (fun (x1, x2) (y1, y2) ->
        b_and (equal (t1, u1) x1 y1)
              (equal (t2, u2) x2 y2))
  | TList t, TList u -> forall2 (equal (t, u))
  | _ -> fun _ _ -> False
\end{tmcode}
{\tmname{InvarGenT}} returns an unexpected type: \tmverbatim{equal$:
\forall$a,b.(Ty a, Ty b)$\rightarrow$a$\rightarrow$a$\rightarrow$Bool}, one of
four maximally general types of \tmverbatim{equal} as defined above. The other
maximally general ``wrong'' types are \tmverbatim{$\forall$a,b.(Ty a, Ty
b)$\rightarrow$b$\rightarrow$b$\rightarrow$Bool} and
\tmverbatim{$\forall$a,b.(Ty a, Ty
b)$\rightarrow$b$\rightarrow$a$\rightarrow$Bool}. This illustrates that
unrestricted type systems with GADTs lack principal typing property.

{\tmname{InvarGenT}} commits to a type of a toplevel definition before
proceeding to the next one, so sometimes we need to provide more information
in the program. Besides type annotations, there are three means to enrich the
generated constraints: \tmverbatim{assert false} syntax for providing negative
constraints, \tmverbatim{assert type e1 = e2;}$\ldots$ and \tmverbatim{assert
num e1 <= e2;}$\ldots$ for positive constraints, and \tmverbatim{test} syntax
for including constraints of use cases with constraint of a toplevel
definition. To ensure only one maximally general type for \tmverbatim{equal},
we use \tmverbatim{assert false} and \tmverbatim{test}. We can either add the
\tmverbatim{assert false} clauses:
\begin{tmcode}
  | TInt, TList l -> (function Nil -> assert false)
  | TList l, TInt -> (fun _ -> function Nil -> assert false)
\end{tmcode}
The first assertion excludes independence of the first encoded type and the
second argument. The second assertion excludes independence of the second
encoded type and the third argument. Or we can add the \tmverbatim{test}
clause:
\begin{tmcode}
test b_not (equal (TInt, TList TInt) Zero Nil)
\end{tmcode}
The test ensures that arguments of distinct types can be~given.
{\tmname{InvarGenT}} returns the expected type \tmverbatim{equal$:
\forall$a,b.(Ty a, Ty b)$\rightarrow$a$\rightarrow$b$\rightarrow$Bool}.

Now we demonstrate numerical invariants:
\begin{tmcode}
datatype Binary : num
datatype Carry : num
datacons Zero : Binary 0
datacons PZero : n[0n]. Binary n  Binary(2 n)
datacons POne : n[0n]. Binary n  Binary(2 n + 1)
datacons CZero : Carry 0
datacons COne : Carry 1

let rec plus =
  function CZero ->
    (function Zero -> (fun b -> b)
      | PZero a1 as a ->
        (function Zero -> a
          | PZero b1 -> PZero (plus CZero a1 b1)
          | POne b1 -> POne (plus CZero a1 b1))
      | POne a1 as a ->
        (function Zero -> a
          | PZero b1 -> POne (plus CZero a1 b1)
          | POne b1 -> PZero (plus COne a1 b1)))
    | COne ->
    (function Zero ->
        (function Zero -> POne(Zero)
          | PZero b1 -> POne b1
          | POne b1 -> PZero (plus COne Zero b1))
      | PZero a1 as a ->
        (function Zero -> POne a1
          | PZero b1 -> POne (plus CZero a1 b1)
          | POne b1 -> PZero (plus COne a1 b1))
      | POne a1 as a ->
        (function Zero -> PZero (plus COne a1 Zero)
          | PZero b1 -> PZero (plus COne a1 b1)
          | POne b1 -> POne (plus COne a1 b1)))
\end{tmcode}
We get \tmverbatim{plus$: \forall$i,k,n.Carry i$\rightarrow$Binary
k$\rightarrow$Binary n$\rightarrow$Binary (n + k + i)}.

We can introduce existential types directly in type declarations. To have an
existential type inferred, we have to use \tmverbatim{efunction},
\tmverbatim{ematch} or \tmverbatim{eif} expressions, which differ from
\tmverbatim{function}, \tmverbatim{match}, \tmverbatim{eif} respectively in
that the (return) type is an existential type. To use a value of an
existential type, we have to bind it with a \tmverbatim{let}..\tmverbatim{in}
expression. Otherwise, the existential type will not be unpacked. An
existential type will be automatically unpacked before being ``repackaged'' as
another existential type. In the following artificial example, we abstract
away the particular resulting location.
\begin{tmcode}
datatype Room
datatype Yard
datatype Village
datatype Castle : type
datatype Place : type
datacons Room : Room  Castle Room
datacons Yard : Yard  Castle Yard
datacons CastleRoom : Room  Place Room
datacons CastleYard : Yard  Place Yard
datacons Village : Village  Place Village

external wander : a. Place a  b. Place b

let rec find_castle = efunction
  | CastleRoom x -> Room x
  | CastleYard x -> Yard x
  | Village _ as x ->
    let y = wander x in
    find_castle y
\end{tmcode}
We get \tmverbatim{find\_castle$: \forall$a. Place a$\rightarrow$ $\exists$b.
Castle b}. Next consider a slightly less artificial, toy example of computer
hardware configuration. It illustrates many aspects of existential types in
{\tmname{InvarGenT}}. We introduce functions \tmverbatim{config\_mem\_board}
and \tmverbatim{config\_gpu} as external, i.e., ones whose definition is not
type-checked by {\tmname{InvarGenT}}. Their types illustrate that existential
types can be used in type annotations. Types \tmverbatim{Slow} and
\tmverbatim{Fast}, although declared as data-types, are phantom types, i.e.
are not inhabited and convey information as parameters of other types.
\begin{tmcode}
datatype Slow  datatype Fast  datatype Budget
datacons Small : Budget  datacons Medium : Budget  datacons Large : Budget
datatype Memory : type  datacons Best_mem : Memory Fast
datatype Motherboard : type  datacons Best_board : Motherboard Fast
external config_mem_board : Budget  a. (Memory a, Motherboard a)
datatype CPU : type
datacons FastCPU : CPU Fast  datacons SlowCPU : CPU Slow
datatype GPU : type
datacons FastGPU : GPU Fast  datacons SlowGPU : GPU Slow
external config_gpu : Budget  a. GPU a
datatype PC : type * type * type * type
datacons PC :
  a,b,c,r. CPU a * GPU b * Memory c * Motherboard r  PC (a,b,c,r)
datatype Usecase  datacons Gaming : Usecase
datacons Scientific : Usecase  datacons Office : Usecase

let budget_to_cpu = efunction
  | Small -> SlowCPU | Medium -> FastCPU | Large -> FastCPU
let usecase_to_gpu budget = efunction
  | Gaming -> FastGPU | Scientific -> FastGPU
  | Office -> config_gpu budget
let rec configure = efunction
  | Small, Gaming -> configure (Small, Office)
  | Large, Gaming -> PC (FastCPU, FastGPU, Best_mem, Best_board)
  | budget, usecase ->
    let mem, board = config_mem_board budget in
    let cpu = budget_to_cpu budget in
    let gpu = usecase_to_gpu budget usecase in
    PC (cpu, gpu, mem, board)
\end{tmcode}
{\tmname{InvarGenT}} infers the following types:
\begin{tmcode}
budget_to_cpu : Size  a.CPU a
usecase_to_gpu : Usecase  a.GPU a
configure : (Size, Usecase)  a, b, c.PC (a, b, c, c)
\end{tmcode}
The definition of \tmverbatim{configure} illustrates explicit elimination of
existential types by \tmverbatim{let}..\tmverbatim{in} definitions: by design,
inlining of \tmverbatim{cpu} or \tmverbatim{gpu} definitions would make the
function not typeable. The call to \tmverbatim{config\_gpu} and the recursive
call to \tmverbatim{configure} illustrate implicit elimination of existential
types in return positions.

A more practical existential type example:
\begin{tmcode}
datatype Bool
datacons True : Bool
datacons False : Bool
datatype List : type * num
datacons LNil : a. List(a, 0)
datacons LCons : n,a[0n]. a * List(a, n)  List(a, n+1)

let rec filter f =
  efunction LNil -> LNil
    | LCons (x, xs) ->
      eif f x then
          let ys = filter f xs in
          LCons (x, ys)
      else filter f xs
\end{tmcode}
We get \tmverbatim{filter$: \forall$n, a.(a$\rightarrow$Bool)$\rightarrow$List
(a, n)$\rightarrow$ $\exists$k[0$\leqslant$k $\wedge$ k$\leqslant$n].List (a,
k)}. Note that we need to use both \tmverbatim{efunction} and \tmverbatim{eif}
above, since every use of \tmverbatim{function}, \tmverbatim{match} or
\tmverbatim{if} will force the types of its branches to be equal. In
particular, for lists with length the resulting length would have to be the
same in each branch. If the constraint cannot be met, as for
\tmverbatim{filter} with either \tmverbatim{function} or \tmverbatim{if}, the
code will not type-check.

A more complex example that computes bitwise {\tmem{or}} -- \tmverbatim{ub}
stands for ``upper bound'':
\begin{tmcode}
datatype Binary : num
datacons Zero : Binary 0
datacons PZero : n [0n]. Binary n  Binary(2 n)
datacons POne : n [0n]. Binary n  Binary(2 n + 1)

let rec ub = efunction
  | Zero ->
      (efunction Zero -> Zero
        | PZero b1 as b -> b
        | POne b1 as b -> b)
  | PZero a1 as a ->
      (efunction Zero -> a
        | PZero b1 ->
          let r = ub a1 b1 in
          PZero r
        | POne b1 ->
          let r = ub a1 b1 in
          POne r)
  | POne a1 as a ->
      (efunction Zero -> a
        | PZero b1 ->
          let r = ub a1 b1 in
          POne r
        | POne b1 ->
          let r = ub a1 b1 in
          POne r)
\end{tmcode}
\tmverbatim{ub:$\forall$k,n.Binary k$\rightarrow$Binary n$\rightarrow
\exists$:i[0$\leqslant$n $\wedge$ 0$\leqslant$k $\wedge$ n$\leqslant$i
$\wedge$ k$\leqslant$i $\wedge$ i$\leqslant$n+k].Binary i}.

Why cannot we shorten the above code by converting the initial cases to
\tmverbatim{Zero -> (efunction b -> b)}? Without pattern matching, we do not
make the contribution of \tmverbatim{Binary n} available. Knowing
\tmverbatim{n$=$i} and not knowing \tmverbatim{0$\leqslant$n}, for the case
\tmverbatim{k$=$0}, we get: \tmverbatim{ub:$\forall$k,n.Binary
k$\rightarrow$Binary n$\rightarrow
\exists$i[0$\leqslant$k$\wedge$n$\leqslant$i$\wedge$i$\leqslant$n+k].Binary
i}. \tmverbatim{n$\leq {\nobreak}$i} follows from \tmverbatim{n$=$i},
\tmverbatim{i$\leqslant$n+k} follows from \tmverbatim{n$=$i} and
\tmverbatim{0$\leqslant$k}, but \tmverbatim{k$\leqslant$i} cannot be inferred
from \tmverbatim{k$=$0} and \tmverbatim{n$=$i} without knowing that
\tmverbatim{0$\leqslant$n}.

Besides displaying types of toplevel definitions, {\tmname{InvarGenT}} can
also export an OCaml source file with all the required GADT definitions and
type annotations.

\section{Syntax}

Below we present, using examples, the syntax of {\tmname{InvarGenT}}: the
mathematical notation, the concrete syntax in ASCII and the concrete syntax
using Unicode.

\begin{tabular}{|l|l|l|l|}
  \hline
  type variable: types & $\alpha, \beta, \gamma, \tau$ & \tmverbatim{a},
  \tmverbatim{b}, \tmverbatim{c}, \tmverbatim{r}, \tmverbatim{s},
  \tmverbatim{t}, \tmverbatim{a1},... & \\
  \hline
  type variable: nums & $k, m, n$ & \tmverbatim{i}, \tmverbatim{j},
  \tmverbatim{k}, \tmverbatim{l}, \tmverbatim{m}, \tmverbatim{n},
  \tmverbatim{i1},... & \\
  \hline
  type var. with coef. & $\frac{1}{3} n$ & \tmverbatim{1/3 n} & \\
  \hline
  type constructor & $\tmop{List}$ & \tmverbatim{List} & \\
  \hline
  number (type) & $7$ & \tmverbatim{7} & \\
  \hline
  numerical sum (type) & $m + n$ & \tmverbatim{m+n} & \\
  \hline
  existential type & $\exists k, n [k \leqslant n] . \tau$ & \tmverbatim{ex k,
  n [k<=n].t} & \tmverbatim{$\exists$k,n[k$\leqslant$n].t}\\
  \hline
  type sort & $s_{\tmop{ty}}$ & \tmverbatim{type} & \\
  \hline
  number sort & $s_R$ & \tmverbatim{num} & \\
  \hline
  function type & $\tau_1 \rightarrow \tau_2$ & \tmverbatim{t1 -> t2} &
  \tmverbatim{t1 $\rightarrow$\tmverbatim{ t2}}\\
  \hline
  equation & $a \dot{=} b$ & \tmverbatim{a = b} & \\
  \hline
  inequation & $k \leqslant n$ & \tmverbatim{k <= n} & \tmverbatim{k
  $\leqslant$ n}\\
  \hline
  conjunction & $\varphi_1 \wedge \varphi_2$ & \tmverbatim{a=b \&\& b=a} &
  \tmverbatim{a=b $\wedge$ b=a}\\
  \hline
\end{tabular}

For the syntax of expressions, we discourage non-ASCII symbols. Below $e, e_i$
stand for any expression, $p, p_i$ stand for any pattern, $x$ stands for any
lower-case identifier and $K$ for an upper-case identifier. $K_T$ stands for
\tmverbatim{True}, $K_F$ for \tmverbatim{False}, and $K_u$ for
\tmverbatim{()}.

\begin{tabular}{|l|l|l|}
  \hline
  named value & $x$ & \tmverbatim{x} \ --lower-case identifier\\
  \hline
  numeral (expr.) & $7$ & \tmverbatim{7}\\
  \hline
  constructor & $K$ & \tmverbatim{K} \ --upper-case identifier\\
  \hline
  application & $e_1 e_2$ & \tmverbatim{e1 e2}\\
  \hline
  non-br. function & $\lambda (p_1 . \lambda (p_2 .e))$ & \tmverbatim{fun
  (p1,p2) p3 -> e}\\
  \hline
  branching function & $\lambda (p_1 .e_1 \ldots p_n .e_n)$ &
  \tmverbatim{function p1->e1 \textbar }...\tmverbatim{ \textbar  pn->en}\\
  \hline
  pattern match & $\lambda (p_1 .e_1 \ldots p_n .e_n) e$ & \tmverbatim{match e
  with p1->e1 \textbar }...\tmverbatim{}\\
  \hline
  if-then-else clause & $\lambda (K_T .e_1, K_F .e_2) e$ & \tmverbatim{if e
  then e1 else e2}\\
  \hline
  if-then-else condition & $\lambda \left( \_ \text{{\tmstrong{when}}} m
  \leqslant n.e_1, \ldots \right) K_u$ & \tmverbatim{if m <= n then e1 else
  e2}\\
  \hline
  postcond. function & $\lambda [K] (p_1 .e_1 \ldots p_n .e_n)$ &
  \tmverbatim{efunction p1->e1 \textbar }...\\
  \hline
  postcond. match & $\lambda [K] (p_1 .e_1 \ldots p_n .e_n) e$ &
  \tmverbatim{ematch e with p1->e1 \textbar }...\\
  \hline
  eif-then-else clause & $\lambda [K] (K_T .e_1, K_F .e_2) e$ &
  \tmverbatim{eif e then e1 else e2}\\
  \hline
  eif-then-else condition & $\lambda [K] \left( \_ \text{{\tmstrong{when}}} m
  \leqslant n.e_1, \ldots \right) K_u$ & \tmverbatim{eif m <= n then e1 else
  e2}\\
  \hline
  rec. definition & $\text{{\tmstrong{let rec}}} x = e_1 
  \text{{\tmstrong{in}}} e_2$ & \tmverbatim{let rec x = e1 in e2}\\
  \hline
  definition & $\text{{\tmstrong{let}}} p = e_1  \text{{\tmstrong{in}}} e_2$ &
  \tmverbatim{let p1,p2 = e1 in e2}\\
  \hline
  asserting dead br. & $\text{{\tmstrong{assert false}}}$ & \tmverbatim{assert
  false}\\
  \hline
  runtime failure & $\tmmathbf{\tmop{runtime} \tmop{failure}} s$ &
  \tmverbatim{runtime\_failure s}\\
  \hline
  assert equal types & $\text{{\tmstrong{assert type}}} \tau_{e_1} \dot{=}
  \tau_{e_2} ; e_3$ & \tmverbatim{assert type e1 = e2; e3}\\
  \hline
  assert inequality & $\text{{\tmstrong{assert num}}} e_1 \leqslant e_2 ; e_3$
  & \tmverbatim{assert num e1 <= e2; e3}\\
  \hline
\end{tabular}

A built-in fail at runtime with the given text message is only needed for
introducing existential types: a user-defined equivalent of
\tmverbatim{runtime\_failure} would introduce a spurious branch for
generalization.

Toplevel expressions (corresponding to structure items in OCaml) introduce
types, type and value constructors, global variables with given type \
(external names) or inferred type (definitions).

\begin{tabular}{|l|l|}
  \hline
  type constructor & \tmverbatim{datatype List : type * num}\\
  \hline
  value constructor & \tmverbatim{datacons Cons : all n a. a * List(a,n) -->
  List(a,n+1)}\\
  \hline
  & \tmverbatim{datacons Cons : $\forall$n,a. a * List(a,n) $\longrightarrow$
  List(a,n+1)}\\
  \hline
  declaration & \tmverbatim{external foo : $\forall$n,a. List(a,n)$\rightarrow
  \exists$k[k<=n].List(a,k)="c\_foo"}\\
  \hline
  & \tmverbatim{external filter : $\forall$n,a. List(a,n)$\rightarrow
  \exists$k[k<=n].List(a,k)}\\
  \hline
  let-declaration & \tmverbatim{external let mult : $\forall$n,m. Num
  n$\rightarrow$Num m$\rightarrow \exists$k.Num k = "( * )"}\\
  \hline
  rec. definition & \tmverbatim{let rec f =}...\\
  \hline
  non-rec. definition & \tmverbatim{let a, b =}...\\
  \hline
  definition with test & \tmverbatim{let rec f =}...\tmverbatim{ test e1;
  }...\tmverbatim{; en}\\
  \hline
\end{tabular}

Toplevel non-recursive \tmverbatim{let} definitions are polymorphic as an
exception. In expressions, \tmverbatim{let}$\ldots$\tmverbatim{in} definitions
are monomorphic, one should use the \tmverbatim{let
rec}$\ldots$\tmverbatim{in} syntax to get a polymorphic
\tmverbatim{let}-binding.

Tests list expressions of type \tmverbatim{Bool} that at runtime have to
evaluate to \tmverbatim{True}. Type inference is affected by the constraints
generated to typecheck the expressions.

There are variants of the if-then-else clause syntax supporting
$\tmmathbf{\tmop{when}}$ conditions:
\begin{itemize}
  \item \tmverbatim{if m1 <= n1 \&\& m2 <= n2 \&\&}$\ldots$\tmverbatim{ then
  e1 else e2} is $\lambda \left( \_ \text{{\tmstrong{when}}} \wedge_i m_i
  \leqslant n_i .e_1, \_.e_2 \right) K_u$,
  
  \item \tmverbatim{if m <= n then e1 else e2} is $\lambda \left( \_
  \text{{\tmstrong{when}}} m \leqslant n.e_1, \_ \text{{\tmstrong{when}}} n +
  1 \leqslant m.e_2 \right) K_u$ if integer mode is on (as in default
  setting),
  
  \item similarly for the \tmverbatim{eif} variants.
\end{itemize}
We add the standard syntactic sugar for function definitions:
\begin{itemize}
  \item \tmverbatim{let $p_1$ $p_2$}$\ldots$\tmverbatim{ $p_n$ = $e_1$ in
  $e_2$} expands to \tmverbatim{let $p_1$ = fun $p_2$}$\ldots$\tmverbatim{
  $p_n$ -> $e_1$ in $e_2$}
  
  \item \tmverbatim{let rec $l_1$ $p_2$}$\ldots$\tmverbatim{ $p_n$ = $e_1$ in
  $e_2$} expands to \tmverbatim{let rec $l_1$ = fun $p_2$}$\ldots$\tmverbatim{
  $p_n$ -> $e_1$ in $e_2$}
  
  \item toplevel \tmverbatim{let} and \tmverbatim{let rec} definitions expand
  correspondingly.
\end{itemize}
For simplicity of theory and implementation, mutual non-nested recursion and
or-patterns are not provided. For mutual recursion, nest one recursive
definition inside another.

Like in OCaml, types of arguments in declarations of constructors are
separated by asterisks. However, the type constructor for tuples is
represented by commas, like in Haskell but unlike in OCaml.

At any place between lexemes, regular comments encapsulated in
\tmverbatim{(*}$\ldots$\tmverbatim{*)} can occur. They are ignored during
lexing. In front of all toplevel definitions and declarations, e.g. before a
\tmverbatim{datatype}, \tmverbatim{datacons}, \tmverbatim{external},
\tmverbatim{let rec} or \tmverbatim{let}, and in front of \tmverbatim{let
rec}$\ldots$\tmverbatim{in} and \tmverbatim{let}$\ldots$\tmverbatim{in} nodes
in expressions, documentation comments \tmverbatim{(**}$\ldots$\tmverbatim{*)}
can be put. Documentation comments at other places are syntax errors.
Documentation comments are preserved both in generated interface files and in
exported source code files.

\section{Solver Parameters and CLI}

The default settings of {\tmname{InvarGenT}} parameters should be sufficient
for most cases. For example, after downloading {\tmname{InvarGenT}} source
code and changing current directory to \tmverbatim{invargent}, we can enter,
assuming a Unix-like shell:
\begin{tmcode}
$ make main
$ ./invargent examples/binary_upper_bound.gadt
\end{tmcode}
To get the inferred types printed on standard output, use the
\tmverbatim{-inform} option:
\begin{tmcode}
$ ./invargent -inform examples/avl_tree.gadt
\end{tmcode}
Below we demonstrate what happens with insufficiently high parameter setting.
Consider this example:
\begin{tmcode}
$ ./invargent -inform examples/flatten_septs.gadt
File "examples/flatten_septs.gadt", line 8, characters 6-104:
No answer in type: term abduction failed

Perhaps increase the -term_abduction_timeout parameter.
Perhaps increase the -term_abduction_fail parameter.
$ ./invargent -inform -term_abduction_timeout 3000 examples/flatten_septs.gadt
File "examples/flatten_septs.gadt", line 3, characters 26-261:
No answer in num: numerical abduction failed

Perhaps do not pass the -no_dead_code flag.
Perhaps increase the -num_abduction_timeout parameter.
Perhaps increase the -num_abduction_fail parameter.
$ ./invargent -inform -term_abduction_timeout 3000 \textbackslash 
  -num_abduction_rotations 4 examples/flatten_septs.gadt
val flatten_septs :
  n, a. List ((a, a, a, a, a, a, a), n)  List (a, 7 n)
InvarGenT: Generated file examples/flatten_septs.gadti
InvarGenT: Generated file examples/flatten_septs.ml
InvarGenT: Command "ocamlc -w -25 -c examples/flatten_septs.ml" exited with code 0
\end{tmcode}
The \tmverbatim{Perhaps increase} suggestions are generated only when the
corresponding limit has actually been exceeded. Remember however that the
limits will often be exceeded for erroneus programs which should not
type-check. Moreover, as illustrated above, other settings might be the
culprit.

To understand the intent of the solver parameters, we need a rough ``birds-eye
view'' understanding of how {\tmname{InvarGenT}} works. The invariants and
postconditions that we solve for are logical formulas and can be ordered by
strength. Least Upper Bounds (LUBs) and Greatest Lower Bounds (GLBs)
computations are traditional tools used for solving recursive equations over
an ordered structure. In case of implicational constraints that are generated
for type inference with GADTs, constraint abduction is a form of LUB
computation. {\tmem{Constraint generalization}} is our term for computing the
GLB wrt. strength for formulas that are conjunctions of atoms. We want the
invariants of recursive definitions -- i.e. the types of recursive functions
and formulas constraining their type variables -- to be as weak as possible,
to make the use of the corresponding definitions as easy as possible. The
weaker the invariant, the more general the type of definition. Therefore the
use of LUB, constraint abduction. For postconditions -- i.e. the existential
types of results computed by \tmverbatim{efunction} expressions and formulas
constraining their type variables -- we want the strongest possible solutions,
because stronger postcondition provides more information at use sites of a
definition. Therefore we use GLB, constraint generalization, but only if
existential types have been introduced by \tmverbatim{efunction} or
\tmverbatim{ematch}.

Below we discuss all of the {\tmname{InvarGenT}} options. We use the technical
term {\tmem{terms}} to mean type shapes, types without the concern for the
sort of numbers (or other sorts to come in the future).
\begin{description}
  \item[\tmverbatim{-inform}] Print type schemes of toplevel definitions as
  they are inferred.
  
  \item[\tmverbatim{-time}] Print the time it took to infer type schemes of
  toplevel definitions.
  
  \item[\tmverbatim{-no\_sig}] Do not generate the \tmverbatim{.gadti} file.
  
  \item[\tmverbatim{-no\_ml}] Do not generate the \tmverbatim{.ml} file.
  
  \item[\tmverbatim{-overwrite\_ml}] Overwrite the \tmverbatim{.ml} file if it
  already exists.
  
  \item[\tmverbatim{-ml\_file}] Generate the exported OCaml file under the
  provided name.
  
  \item[\tmverbatim{-no\_verif}] Do not call \tmverbatim{ocamlc -c} on the
  generated \tmverbatim{.ml} file.
  
  \item[\tmverbatim{-num\_is}] The exported type for which \tmverbatim{Num} is
  an alias (default \tmverbatim{int}). If \tmverbatim{-num\_is bar} for
  \tmverbatim{bar} different than \tmverbatim{int}, numerals are exported as
  integers passed to a \tmverbatim{bar\_of\_int} function. The variant
  \tmverbatim{-num\_is\_mod} exports numerals by passing to a
  \tmverbatim{Bar.of\_int} function.
  
  \item[\tmverbatim{-full\_annot}] Annotate the \tmverbatim{function} and
  \tmverbatim{let}..\tmverbatim{in} nodes in generated OCaml code. This
  increases the burden on inference a bit because the variables associated
  with the nodes cannot be eliminated from the constraint during initial
  simplification.
  
  \item[\tmverbatim{-keep\_assert\_false}] Keep \tmverbatim{assert false}
  clauses in exported code. When faced with multiple maximally general types
  of a function, we sometimes want to prevent some interpretations by
  asserting that a combination of arguments is not possible. These arguments
  will not be compatible with the type inferred, causing exported code to fail
  to typecheck. Sometimes we indicate unreachable cases just for
  documentation. If the type is tight this will cause exported code to fail to
  typecheck too. This option keeps pattern matching branches with
  \tmverbatim{assert false} in their bodies in exported code nevertheless.
  
  \item[\tmverbatim{-allow\_dead\_code}] Allow more programs with dead code
  than would otherwise pass.
  
  \item[\tmverbatim{-force\_no\_dead\_code}] Reject all programs with dead
  code (may misclassify programs using {\tmem{min}} or {\tmem{max}} atoms).
  Unreachable pattern matching branches lead to unsatisfiable premises of the
  type inference constraint, which we detect. However, sometimes multiple
  implications in the simplified form of the constraint can correspond to the
  same path through the program, in particular when solving constraints with
  {\tmem{min}} and {\tmem{max}} clauses. Dead code due to datatype mismatch,
  i.e. patterns unreachable without resort to numerical constraints, is
  detected even without using this option.
  
  \item[\tmverbatim{-term\_abduction\_timeout}] Limit on term simple abduction
  steps (default 700). Simple abduction works with a single implication
  branch, which roughly corresponds to a single branch -- an execution path --
  of the program.
  
  \item[\tmverbatim{-term\_abduction\_fail}] Limit on backtracking steps in
  term joint abduction (default 4). Joint abduction combines results for all
  branches of the constraints.
  
  \item[\tmverbatim{-no\_alien\_prem}] Do not include alien (e.g. numerical)
  premise information in term abduction.
  
  \item[\tmverbatim{-early\_num\_abduction}] Include recursive branches in
  numerical abduction from the start. By default, in the second iteration of
  solving constraints, which is the first iteration that numerical abduction
  is performed, we only pass non-recursive branches to numerical abduction.
  This makes it faster but less likely to find the correct solution.
  
  \item[\tmverbatim{-convergence\_step}] The iteration at which to start
  truncating postconditions by only keeping atoms present in the previous
  iteration, to force convergence (default 8).
  
  \item[\tmverbatim{-early\_postcond\_abd}] Include postconditions from
  recursive calls in abduction from the start. We do not derive requirements
  put on postconditions by recursive calls on first iteration. The
  requirements may turn smaller after some derived invariants are included in
  the premises. This option turns off the special treatment of postconditions
  on first iteration.
  
  \item[\tmverbatim{-num\_abduction\_rotations}] Numerical abduction:
  coefficients from $\pm 1 / N$ to $\pm N$ (default 3). Numerical abduction
  answers are built, roughly speaking, by adding premise equations of a branch
  with conclusion of a branch to get an equation or inequality that does not
  conflict with other branches, but is equivalent to the conclusion
  equation/inequality. This parameter decides what range of coefficients is
  tried. If the highest coefficient in correct answer is greater, abduction
  might fail. However, it often succeeds because of other mechanisms used by
  the abduction algorithm.
  
  \item[\tmverbatim{-num\_prune\_at}] Keep less than $N$ elements in abduction
  sums (default 6). By elements here we mean distinct variables -- lack of
  constant multipliers in concrete syntax of types is just a syntactic
  shortcoming.
  
  \item[\tmverbatim{-num\_abduction\_timeout}] Limit on numerical simple
  abduction steps (default 1000).
  
  \item[\tmverbatim{-num\_abduction\_fail}] Limit on backtracking steps in
  numerical joint abduction (default 10).
  
  \item[\tmverbatim{-affine\_penalty}] How much to penalize an abduction
  candidate inequality for containing a constant term (default 4). Too small a
  value may lead to divergence, e.g. in some examples abduction will pick an
  answer $a + 1$, which in the following step will force an answer $a + 2$,
  then $a + 3$, etc.
  
  \item[\tmverbatim{-reward\_constrn}] How much to reward introducing a
  constraint on so-far unconstrained varialbe, or penalize if negative
  (default 2).
  
  \item[\tmverbatim{-complexity\_penalty}] How much to penalize an abduction
  candidate inequality for complexity of its coefficients; the coefficient of
  either the linear or power scaling of the coefficients (default 2.5).
  
  \item[\tmverbatim{-abd\_lin\_thres\_scaling}] Scale the complexity cost of
  coefficients linearly with a jump of the given height after coefficient 1
  (default 2.0).
  
  \item[{\tmstrong{\tmverbatim{-abd\_pow\_scaling}}}] Scale the complexity
  cost of coefficients according to the given power.
  
  \item[\tmverbatim{-prefer\_bound\_to\_local}] Prefer a bound coming from
  outer scope, to inequality between two local parameters. In numerical
  abduction heuristic, such bounds are usually doubly penalized: for having a
  constant, and non-locality of parameters.
  
  \item[\tmverbatim{-prefer\_bound\_to\_outer}] Prefer a bound coming from
  outer scope, to inequality between two outer scope parameters. Outer-scope
  constraints sometimes lead to a solution not general enough.
  
  \item[\tmverbatim{-only\_off\_by\_1}] Limit the effect of
  \tmverbatim{-prefer\_bound\_to\_local} and
  \tmverbatim{-{\nobreak}prefer\_bound\_to\_outer} to inequalities with a
  constant 1. This corresponds to an upper bound of an index into a
  zero-indexed array/matrix/etc.
  
  \item[\tmverbatim{-same\_with\_assertions}] Do not treat definitions with
  positive assertions (\tmverbatim{assert num}, \tmverbatim{assert type})
  specially. The special treatment is currently equivalent to passing
  \tmverbatim{-{\nobreak}reward\_constrn -1} and
  \tmverbatim{-prefer\_bound\_to\_local}.
  
  \item[\tmverbatim{-concl\_abd\_penalty}] Penalize abductive guess when the
  supporting argument comes from the partial answer, instead of from the
  current premise (default 4). Guesses involving the partial answer are less
  secure, for example they depend on the order in which the constraint to
  explain is being processed.
  
  \item[\tmverbatim{-more\_general\_num}] Filter out less general abduction
  candidate atoms (does not guarantee overall more general answers). The
  filtering is currently not performed by default to save on computational
  cost.
  
  \item[\tmverbatim{-no\_num\_abduction}] Turn off numerical abduction; will
  not ensure correctness. Numerical abduction uses a brute-force algorithm and
  will fail to work in reasonable time for complex constraints. However,
  including the effects of \tmverbatim{assert false} clauses, and inference of
  postconditions, do not rely on numerical abduction. If the numerical
  invariant of a typeable (i.e. correct) function follows from
  \tmverbatim{assert false} facts alone, a call with
  \tmverbatim{-no\_num\_abduction} may still find the correct invariant and
  postcondition.
  
  \item[\tmverbatim{-if\_else\_no\_when}] Do not add \tmverbatim{when} clause
  to the \tmverbatim{else} branch of an \tmverbatim{if} expression with a
  single inequality as condition. Expressions \tmverbatim{if}, resp.
  \tmverbatim{eif}, with a single inequality as the condition are expanded
  into expressions \tmverbatim{match}, resp. \tmverbatim{ematch}, with
  \tmverbatim{when} conditions on both the \tmverbatim{True} branch and the
  \tmverbatim{False} branch. I.e. \tmverbatim{if m <= n then e1 else e2} is
  expanded into \tmverbatim{match () with \_ when m <= n -> e1 \textbar \_
  when n+1 <= m -> e2}. Passing \tmverbatim{-if\_else\_no\_when} will result
  in expansion \tmverbatim{match () with \_ when m <= n -> e1 \textbar \_ ->
  e2}. The same effect can be achieved for a particular expression by
  artificially incresing the number of inequalities: \tmverbatim{if m <= n
  \&\& m <= n then e1 else e2}.
  
  \item[\tmverbatim{-weaker\_pruning}] Do not assume integers as the numerical
  domain when pruning redundant atoms.
  
  \item[\tmverbatim{-stronger\_pruning}] Prune atoms that force a numerical
  variable to a single value under certain conditions; exclusive with
  \tmverbatim{-weaker\_pruning}.
  
  \item[\tmverbatim{-postcond\_rotations}] In postconditions, check
  coefficients from $1 / N$ (default 3). Numerical constraint generalization
  is performed by approximately finding the convex hull of the polytopes
  corresponding to disjuncts. A step in an exact algorithm involves rotating a
  side along a ridge -- an intersection with another side -- until the side
  touches yet another side. We approximate by trying out a couple of
  rotations: convex combinations of the inequalities defining the sides. This
  parameter decides how many rotations to try.
  
  \item[\tmverbatim{-postcond\_opti\_limit}] Limit the number of atoms $x =
  \min (a, b)$, $x = \max (a, b)$ in (intermediate and final) postconditions
  (default 4). Unfortunately, inference time is exponential in the number of
  atoms of this form. The final postconditions usually have few of these
  atoms, but a greater number is sometimes needed in the intermediate steps of
  the main loop.
  
  \item[\tmverbatim{-postcond\_subopti\_limit}] Limit the number of atoms
  $\min (a, b) \leqslant x$, $x \leqslant \max (a, b)$ in (intermediate and
  final) postconditions (default 4). Unfortunately, inference time is
  exponential in the number of atoms of this form. The final postconditions
  usually have few of these atoms, but a greater number is sometimes needed in
  the intermediate steps of the main loop.
  
  \item[\tmverbatim{-iterations\_timeout}] Limit on main algorithm iterations
  (default 6). Answers found in an iteration of the main algorithm are
  propagated to use sites in the next iteration. However, for about four
  initial iterations, each iteration turns on additional processing which
  makes better sense with the results from the previous iteration propagated.
  At least three iterations will always be performed.
  
  \item[\tmverbatim{-richer\_answers}] Keep some equations in term abduction
  answers even if redundant. Try keeping an initial guess out of a list of
  candidate equations before trying to drop the equation from consideration.
  We use fully maximal abduction for single branches, which cannot find
  answers not implied by premise and conclusion of a branch. But we seed it
  with partial answer to branches considered so far. Sometimes an atom is
  required to solve another branch although it is redundant in given branch.
  \tmverbatim{-richer\_answers} does not increase computational cost but
  sometimes leads to answers that are not most general. This can always be
  fixed by adding a \tmverbatim{test} clause to the definition which uses a
  type conflicting with the too specific type.
  
  \item[\tmverbatim{-prefer\_guess}] Try to guess equality-between-parameters
  before considering other possibilities. Implied by
  \tmverbatim{-richer\_answers} but less invasive.
  
  \item[\tmverbatim{-more\_existential}] More general invariant at expense of
  more existential postcondition. To avoid too abstract postconditions,
  constraint generalization can infer additional constraints over invariant
  parameters. In rare cases a weaker postcondition but a more general
  invariant can be beneficial.
  
  \item[\tmverbatim{-show\_extypes}] Show datatypes encoding existential
  types, and their identifiers with uses of existential types. The type system
  in {\tmname{InvarGenT}} encodes existential types as GADT types, but this
  representation is hidden from the user. Using \tmverbatim{-show\_extypes}
  exposes the representation as follows. The encodings are exported in
  \tmverbatim{.gadti} files as regular datatypes named \tmverbatim{exN}, and
  existential types are printed using syntax \tmverbatim{$\exists$N:}$\ldots$
  instead of $\exists \ldots$, where \tmverbatim{N} is the identifier of an
  existential type.
  
  \item[\tmverbatim{-passing\_ineq\_trs}] Include inequalities in conclusion
  when solving numerical abduction. This setting leads to more inequalities
  being tried for addition in numeric abduction answer.
  
  \item[\tmverbatim{-not\_annotating\_fun}] Do not keep information for
  annotating \tmverbatim{function} nodes. This may allow eliminating more
  variables during initial constraint simplification.
  
  \item[\tmverbatim{-annotating\_letin}] Keep information for annotating
  \tmverbatim{let}..\tmverbatim{in} nodes. Will be set automatically anyway
  when \tmverbatim{-full\_annot} is passed.
  
  \item[\tmverbatim{-let\_in\_fallback}] Annotate
  \tmverbatim{let}..\tmverbatim{in} nodes in fallback mode of \tmverbatim{.ml}
  generation. When verifying the resulting \tmverbatim{.ml} file fails, a
  retry is made with \tmverbatim{function} nodes annotated. This option
  additionally annotates \tmverbatim{let}..\tmverbatim{in} nodes with types in
  the regenerated \tmverbatim{.ml} file.
\end{description}
Let us have a look at tests from the \tmverbatim{examples} direcotory that
need a non-default parameter setting. The program
\tmverbatim{examples/flatten\_septs.gadt} has already been shown above. It is
the only example that needs the \tmverbatim{-term\_abduction\_timeout} and
\tmverbatim{-num\_abduction\_rotations} settings. The need for
\tmverbatim{-num\_abduction\_rotations} comes from having an equation with
large coefficient in the answer. The need for
\tmverbatim{-{\nobreak}term\_abduction\_timeout} comes from having bigger type
shapes to handle during term, i.e. type shape, abduction.
\begin{tmcode}
$ ./invargent -inform examples/non_pointwise_leq.gadt
File "examples/non_pointwise_leq.gadt", line 12, characters 14-60:
No answer in type: Answers do not converge

Perhaps increase the -iterations_timeout parameter or try one of the options: -more_existential, -prefer_guess, -prefer_bound_to_local.
$ ./invargent -inform -prefer_guess examples/non_pointwise_leq.gadt
val leq : a. Nat a  NatLeq (a, a)
InvarGenT: Generated file examples/non_pointwise_leq.gadti
InvarGenT: Generated file examples/non_pointwise_leq.ml
InvarGenT: Command "ocamlc -w -25 -c examples/non_pointwise_leq.ml" exited with code 0
\end{tmcode}
Other examples that need the \tmverbatim{-prefer\_guess} option:
\tmverbatim{non\_pointwise\_zip1\_simpler.gadt},
\tmverbatim{non\_pointwise\_zip1\_simpler2.gadt},
\tmverbatim{non\_pointwise\_zip1\_modified.gadt}. On the other hand,
\tmverbatim{non\_pointwise\_zip1.gadt} is inferred with default settings.

The response from the system does not always include an option which would
make the inference succeed.
\begin{tmcode}
$ ./invargent -inform examples/liquid_simplex_step_3a.gadt
File "examples/liquid_simplex_step_3a.gadt", line 7, characters 49-1651:
No answer in type: Answers do not converge

Perhaps do not pass the -no_dead_code flag.
Perhaps increase the -iterations_timeout parameter or try one of the options: -more_existential, -prefer_guess, -prefer_bound_to_local.
Perhaps some definition is used with requirements on
its inferred postcondition not warranted by the definition.
$ ./invargent -inform -prefer_bound_to_local -only_off_by_1 \textbackslash 
examples/liquid_simplex_step_3a.gadt
val main_step3_test : k, n[1  n  3  k]. Matrix (n, k)  Float
InvarGenT: Generated file examples/liquid_simplex_step_3a.gadti
InvarGenT: Generated file examples/liquid_simplex_step_3a.ml
File "examples/liquid_simplex_step_3a.ml", line 43, characters 8-9:
Warning 26: unused variable m.
InvarGenT: Command "ocamlc -w -25 -c examples/liquid_simplex_step_3a.ml" exited with code 0
\end{tmcode}
The other examples that need the \tmverbatim{-prefer\_bound\_to\_local}
option, but not the \tmverbatim{-{\nobreak}only\_off\_by\_1} option:
\tmverbatim{liquid\_simplex\_step\_6a\_2.gadt},
\tmverbatim{liquid\_tower\_harder.gadt}, \tmverbatim{liquid\_gauss2.gadt}.
\begin{tmcode}
$ ./invargent -inform examples/pointwise_zip2_harder.gadt
val zip2 : a, b. Zip2 (a, b)  a  b
InvarGenT: Generated file examples/pointwise_zip2_harder.gadti
InvarGenT: Generated file examples/pointwise_zip2_harder.ml
File "examples/pointwise_zip2_harder.ml", line 19, characters 21-32:
Error: This kind of expression is not allowed as right-hand side of ‘let rec'
InvarGenT: Command "ocamlc -w -25 -c examples/pointwise_zip2_harder.ml" exited with code 2
InvarGenT: Regenerated file examples/pointwise_zip2_harder.ml
File "examples/pointwise_zip2_harder.ml", line 21, characters 21-32:
Error: This kind of expression is not allowed as right-hand side of ‘let rec'
InvarGenT: Command "ocamlc -w -25 -c examples/pointwise_zip2_harder.ml" exited with code 2
$ ./invargent -inform -no_ml examples/pointwise_zip2_harder.gadt
val zip2 : a, b. Zip2 (a, b)  a  b
InvarGenT: Generated file examples/pointwise_zip2_harder.gadti
\end{tmcode}
The example \tmverbatim{pointwise\_zip2\_harder.gadt} is not compatible with
the pass-by-value semantics. We can avoid the complaint of the OCaml compiler
by passing either the \tmverbatim{-no\_ml} flag or the \tmverbatim{-no\_verif}
flag. More interestingly, we can notice that the file
\tmverbatim{pointwise\_zip2\_harder.ml} is generated twice. This happens
because {\tmname{InvarGenT}}, noticing the failure, generates an OCaml source
with more type information, as if the \tmverbatim{-full\_annot} option was
used.

The examples \tmverbatim{liquid\_fft\_simpler.gadt} and
\tmverbatim{liquid\_fft\_full\_asserted.gadt} contain assertions, but are
nearly as hard as \tmverbatim{liquid\_fft.gadt},
\tmverbatim{liquid\_fft\_full.gadt} respectively. They need the option
\tmverbatim{-same\_with\_assertions} to not switch to settings tuned for cases
where assertions capture the harder aspects of the invariants to infer.

Unfortunately, inference fails for some examples regardless of parameters
setting. We discuss them in the next section.

\section{Limitations of Current {\tmname{InvarGenT}} Inference}

Type inference for the type system underlying {\tmname{InvarGenT}} is
undecidable. In some cases, the failure to infer a type is not at all
problematic. Consider this example due to Chuan-kai Lin:
\begin{tmcode}
datatype EquLR : type * type * type
datacons EquL : a, b. EquLR (a, a, b)
datacons EquR : a, b. EquLR (a, b, b)
datatype Box : type
datacons Cons : a. a  Box a
external let eq : a. a  a  Bool = "(=)"

let vary = fun e y ->
  match e with
  | EquL, EquL -> eq y "c"
  | EquR, EquR -> Cons (match y with True -> 5 | False -> 7)
\end{tmcode}
Although \tmverbatim{vary} has multiple types, it is a contrived example
unlikely to have an intended type. However, not all cases of failure to infer
a type for a correct program are due to contrived examples. The problems are
not insurmountable theoretically. The algorithms used in the inference can
incorporate heuristics for special cases, and can be modified to do a more
exhaustive search.

The example \tmverbatim{pointwise\_head.gadt} fails because of the limitations
of the \tmverbatim{type} sort in representing disequalities.
\begin{tmcode}
datatype Z
datatype S : type
datatype List : type * num
datacons LNil : a. List(a, Z)
datacons LCons : a, b. a * List(a, b)  List(a, S b)

let head = function
  | LCons (x, _) -> x
  | LNil -> assert false
\end{tmcode}
If we omit the \tmverbatim{LNil} branch, we get the technically correct but
inadequate type \tmverbatim{$\forall$a, b. List(a, b) $\rightarrow$ a},
because the type system does not guarantee exhaustiveness of the pattern
matching. The intended type is \tmverbatim{$\forall$a, b. List(a, S b)
$\rightarrow$ a}.

The example \tmverbatim{non\_pointwise\_fd\_comp\_harder.gadt} is inferred an
insufficiently general type \tmverbatim{$\forall$a, b. FunDesc (b, b)
$\rightarrow$ FunDesc (b, a) $\rightarrow$ FunDesc (b, a)}.
\begin{tmcode}
datatype FunDesc : type * type
datacons FDI : a. FunDesc (a, a)
datacons FDC : a, b. b  FunDesc (a, b)
datacons FDG : a, b. (a  b)  FunDesc (a, b)
external fd_fun : a, b. FunDesc (a, b)  a  b

let fd_comp fd1 fd2 =
  let o f g x = f (g x) in
  match fd1 with
    | FDI -> fd2
    | FDC b -> 
      (match fd2 with
        | FDI -> fd1
        | FDC c -> FDC (fd_fun fd2 b)
        | FDG g -> FDC (fd_fun fd2 b))
    | FDG f ->
      (match fd2 with
        | FDI -> fd1
        | FDC c -> FDC c
        | FDG g -> FDG (o (fd_fun fd2) f))
\end{tmcode}
This happens because the second argument \tmverbatim{fd2} is not expanded when
\tmverbatim{fd1} is equal to \tmverbatim{FDI}. Type inference cannot carry out
the different reasoning steps leading to the more general type.

In the example \tmverbatim{liquid\_bsearch2\_harder4.gadt} it turns out to be
too hard to infer the full postcondition.
\begin{tmcode}
datatype Array : type * num
external let array_make :
  n, a [0n]. Num n  a  Array (a, n) = "fun a b -> Array.make a b"
external let array_get :
n, k, a [0k  k+1n]. Array (a, n)  Num k  a =
  "fun a b -> Array.get a b"
external let array_length :
  n, a [0n]. Array (a, n)  Num n = "fun a -> Array.length a"
datatype LinOrder
datacons LE : LinOrder
datacons GT : LinOrder
datacons EQ : LinOrder
external let compare : a. a  a  LinOrder =
  "fun a b -> let c = Pervasives.compare a b in
              if c  0 then GT else EQ"
external let equal : a. a  a  Bool = "fun a b -> a = b"
external let div2 : n. Num (2 n)  Num n = "fun x -> x / 2"

let bsearch key vec =
  let rec look key vec lo hi =
    eif lo 
        let m = div2 (hi + lo) in
        let x = array_get vec m in
        ematch compare key x with
          | LE -> look key vec lo (m + (-1))
          | GT -> look key vec (m + 1) hi
          | EQ -> eif equal key x then m else -1
    else -1 in
  look key vec 0 (array_length vec + (-1))
\end{tmcode}
We get the result type \tmverbatim{$\exists$n[0 $\leqslant$ n + 1].Num n}
instead of \tmverbatim{$\exists$k[k $\leqslant$ n $\wedge$ 0 $\leqslant$ k +
1].Num k}. The inference of the intended type succeeds after we introduce an
appropriate assertion, e.g. \tmverbatim{assert num -1 }. Alternatively, we
could include a use case for \tmverbatim{bsearch} where the full postcondition
is required.

Examples \tmverbatim{liquid\_simplex\_step\_3.gadt},
\tmverbatim{liquid\_simplex\_step\_4.gadt} and
\tmverbatim{liquid\_gauss\_rowMax.gadt} result in uninformative, empty
postconditions, because to tell more would require inspecting the behavior of
the respective function across recursive calls. To save space, we list just
the function definition from \tmverbatim{liquid\_simplex\_step\_3.gadt}:
\begin{tmcode}
let rec enter_var arr2 n j c j' =
  eif j' + 2 
    let c' = matrix_get arr2 0 j' in
    eif less c' c then enter_var arr2 n j' c' (j'+1)
    else enter_var arr2 n j c (j'+1)
  else j
\end{tmcode}
Fortunately, if the function is used in the same toplevel definition in which
it is defined, use-site requirements facilitate the inference of the intended
postcondition.

The example \tmverbatim{liquid\_gauss\_harder.gadt} poses too big a challenge
for {\tmname{InvarGenT}}. To get it pass the inference, we streamline one of
the nested definitions, to not introduce another, unnecessary level of
nesting. This gives the example \tmverbatim{liquid\_gauss2.gadt}, which needs
to be run with the option \tmverbatim{-prefer\_bound\_to\_local}.
Additionally, we can relax the constraint on the processed portion of the
matrix, coming from the restriction on the matrix size intended in the
original source of the \tmverbatim{liquid\_gauss\_harder.gadt} example. In
\tmverbatim{liquid\_gauss.gadt}, the whole matrix is processed and the
inferred type is most general, under the default settings -- no need to pass
any options to {\tmname{InvarGenT}}. The reason
\tmverbatim{liquid\_gauss\_harder.gadt} is too difficult for
{\tmname{InvarGenT}} is that the nesting interferes with the propagation of
use-site constraints to the postcondition of the nested definition (the
\tmverbatim{loop} inside \tmverbatim{rowMax}). Inference works for
\tmverbatim{liquid\_gauss\_harder\_asserted.gadt}, because the assertion
provides the required information to infer the \tmverbatim{rowMax} invariants
directly.

\end{document}
